\documentclass[11pt,a4paper]{article}
\usepackage[T1]{fontenc}
\usepackage[utf8]{inputenc}
\usepackage{lmodern}
\ifdefined\pdfmapfile
  \pdfmapfile{=lm.map}
  \pdfmapfile{+cm.map}
  \pdfmapfile{+symbols.map}
  \pdfmapfile{+euler.map}
  \pdfmapfile{+rsfs.map}
\fi
\usepackage[a4paper,margin=27mm]{geometry}
\usepackage{amsmath,amssymb,amsthm,mathtools,mathrsfs}
\usepackage{booktabs,tabularx,longtable,array,enumitem}
\usepackage{microtype,xcolor,etoolbox}
\usepackage{hyperref}
\usepackage{xurl}
\hypersetup{colorlinks=true,linkcolor=blue!45!black,citecolor=blue!45!black,
urlcolor=blue!45!black,pdfauthor={Non-Archimedean Bedford--Taylor research record},
pdftitle={Closed Bands, Energy Duals, and Nonlinear Source Gaps}}
\usepackage[capitalise,noabbrev]{cleveref}
\numberwithin{equation}{section}
\newtheorem{theorem}{Theorem}[section]
\newtheorem{lemma}[theorem]{Lemma}
\newtheorem{proposition}[theorem]{Proposition}
\newtheorem{corollary}[theorem]{Corollary}
\theoremstyle{definition}
\newtheorem{definition}[theorem]{Definition}
\newtheorem{example}[theorem]{Example}
\theoremstyle{remark}
\newtheorem{remark}[theorem]{Remark}
\newcommand{\R}{\mathbf R}
\newcommand{\Q}{\mathbf Q}
\newcommand{\Z}{\mathbf Z}
\newcommand{\PSH}{\operatorname{PSH}}
\newcommand{\supp}{\operatorname{supp}}
\newcommand{\dd}{d'd''}
\newcommand{\an}{\mathrm{an}}
\setlist{topsep=4pt,itemsep=2pt,parsep=0pt}
\setlength{\parindent}{1.3em}
\setlength{\parskip}{2pt}
\setlength{\emergencystretch}{2em}
\makeatletter
\patchcmd{\l@section}{1.5em}{2.1em}{}{}
\renewcommand*\l@subsection{\@dottedtocline{2}{1.5em}{3.6em}}
\makeatother
\title{Closed Bands, Energy Duals, and Nonlinear Source Gaps\\[5pt]
\large Twelfth research record in non-Archimedean Bedford--Taylor theory}
\author{Continuation of the exact authorized R8--R11 research chain}
\date{12 November 2026 (UTC)}
\begin{document}
\maketitle
\begin{abstract}
We work over the original finite extension of $\mathbf Q_p$, with its
original absolute value and the full original CLD test sheaf. The
background remains an arbitrary symmetric middle-positive current
closed separately under $d'$ and $d''$.
Parameter-annihilated currents admit closed Borel bands and a
disintegration into positive closed currents on the original ambient
analytic space. This gives mixed Bedford--Taylor products for bounded
Borel families whose specified slices are original classical psh,
with weighted joint decreasing continuity.
On the frozen strict-box first-energy domain, we construct a Banach
space of potential differences, an actual Borel order-energy Banach
lattice, and continuous all-degree original current multipliers.
The nonlinear curvature image is closed in an explicit signed dual
norm. Its variational normal decomposition, obstacle contact defect,
and Perron--variational energy gap distinguish a maximizer from a
source solution.
A single geometry-only test potential detects the complete source
normal of literal strong-continuous energy data. A full-degree
radial construction exhibits an exact one-sided solvability threshold
under arbitrarily small strictly positive original smooth density
perturbations of a bounded source.
An exact positive current represents nonlinear Jensen defects.
A single geometry-dependent actual finite-mass potential detects
all such weak-cluster defects through one weighted mass.
Closed moment lifts give source obstructions in every sub-top
nonlinear degree and outer-nullity of proper moment fibers.
Examples distinguish strong-continuous positive duals from literal
Radon data and exclude naive enlargements of the relative potential
cone. Closed-band capacitary completions and an explicit parameter
finite-energy module provide all-original products, comparison,
compatible source solutions, and an attained energy-normalized
capacity with full local Borel regularity in their stated scopes.
A curvature-depth filtration sharpens conormal zero detection,
while the unrestricted analytic PS/PL relation quotient and the
remaining general capacity and appropriate Dirichlet gates are
retained explicitly.
\end{abstract}
\newpage
{\small\tableofcontents}
\newpage

% BEGIN REVIEWED MODULE research_r12/scope_and_main_results.tex
\section{Scope, frozen inputs, and main results}
\label{sec:r12-scope}

\subsection{The fixed analytic problem}

Let $K/\mathbf Q_p$ be the original finite extension, equipped with its
original absolute value. Throughout, $X$ is good, Hausdorff, boundaryless,
quasi-smooth and pure of dimension $n$. Opens are ordinary analytic
opens. A background of degree $(p,p)$ means an original symmetric
CLD middle-positive current $S$ satisfying
\[
                         d'S=d''S=0.
 \tag{M1}
\]
Strong and weak positivity are not substituted for this middle cone.
Every current identity is tested against all original compactly
supported smooth complementary CLD forms, after all actual tropical
moment refinements. Ordinary current equality is never silently
promoted to equality on all GK delta tests.

The bounded arbitrary-background BT calculus of R4 and the exact
R11 energy and relative boundary results are
\texttt{INTERNAL\_INPUT}. The internal papers R, B and A2 remain
authorized inputs in their printed scopes. We have not re-audited
their proofs. In particular A2 geometric adelic nefness is not an
open item in this record.

The immediate frozen baseline is the complete R11 archive, with SHA256
\[
 \text{\small\nolinkurl{a7fef6ca478f79dde711f18c20e05db8d740963cfac8c8cc8281d5c6be10d7bf}}.
\]
Its 264 manifest entries and ZIP CRC were verified. The unchanged
archive is included in \path{previous/}; it contains the exact
R10--R3 chain, including the user's original R8 baseline. Earlier
proofs are used through exact named statements, not rewritten as
new results.

\subsection{The energy domain and the precise completion}

The global energy results have the following domain hypothesis.
There is one compact original affinoid $B$ and a compact closed
Laurent box
\[
 K_0=\{|f_a|\leq R_a,\ |g_b|\geq r_b\}
                      \Subset\operatorname{Int}_X(B),
 \qquad
 V=\{|f_a|<R_a,\ |g_b|>r_b\}\ne\varnothing,
 \tag{M2}
\]
with all thresholds positive and $S$ defined on an ordinary
neighborhood of $K_0$. Put $q=n-p\geq1$ and $d=q+1$.
A bounded continuous negative box barrier $\rho$ with compact
negative sublevels and finite top mass is part of the proved
R11 construction (TC3)--(TC4).

For Borel sets $E\subset V$, the raw and outer capacities are
\[
 \begin{split}
 C_S(E;V)&=\sup_{0\leq h\leq1,\ h\ {\rm original\ classical\ psh}}
                       ((Dh)^q\wedge S)(E),\\
 C_S^*(E;V)&=\inf_{\substack{G\supset E\\G\ {\rm ordinary\ open}}}
                                                     C_S(G;V).
 \end{split}
 \tag{M3}
\]
The supremum is over the full original class on the indicated
whole domain. No raw/outer equality on arbitrary sets is assumed.

Let $\mathcal E^1_{S,\partial}(V)$ be exactly the actual class of
R11: negative original classical potentials with the compact
negative-sublevel boundary condition, the proved canonical
truncation no-loss condition, and finite first energy.
Its full strong completion
\[
                       \mathscr C=
                       \widehat{\mathcal E}^{1}_{\rm str}(V)
 \tag{M4}
\]
is a distinct space of actual Borel classes modulo local
ordinary-open outer capacity zero. Its boundary condition is
strong approximation by the actual class, not a newly asserted
ordinary boundary value. Its members need not have ambient
classical representatives or a common actual lower barrier.

Write $\widehat N_k$ for its original mixed currents and
$\widehat\mu_u=\widehat N_q(u^{[q]})$. Its global mixed energy,
energy root, and strong discrepancy are
\[
 \begin{split}
 B(u_0,\ldots,u_q)&=\int_V(-u_0)\widehat N_q(u_1,\ldots,u_q),\\
 e(u)&=B(u^{[d]}),\qquad r(u)=e(u)^{1/d},\\
 I(u,v)&=e(u)+e(v)-B(u,v^{[q]})-B(v,u^{[q]}).
 \end{split}
 \tag{M5}
\]
The exact frozen inputs used below are symmetry and the constant-one
mixed energy inequality; positive integrated energy matrices; strong
completeness and all-original weighted current continuity; finite
maxima and monotone closure; the full relative Perron theorem; and
the original test decomposition into two actual finite-mass
potentials. Their sources are R11 (GE7)--(GE18), (EV1)--(EV8),
(IC10)--(IC22), (CC1)--(CC11), (CI1)--(CI8), (SE1)--(SE8), and
(HC14)--(HC15). They are retained with (M2). This record does not
turn them into theorems on arbitrary unrestricted domains.

\subsection{Closed bands provide a genuinely original conditional calculus}

The first mechanism does not require (M2). Let $s_i=\log|f_i|$
be unit-logarithm parameters on an ordinary open $U$. A positive
closed current $R$ is parameter-annihilated if
\[
                   d's_i\wedge R=d''s_i\wedge R=0
                   \quad\text{for every }i.
 \tag{M6}
\]
For every arbitrary original $S$, the derived current
$\Omega_s\wedge S$, where
$\Omega_s=\bigwedge_i d's_i\wedge d''s_i$, has this property
in the permitted degree. The background $S$ itself is unchanged.

\begin{theorem}[Closed bands and bounded parameter-family products]
\label{thm:r12-main-bands}
Let $R$ satisfy (M6), be symmetric middle-positive and separately
closed, and have degree $(P,P)$. Then every Borel parameter band
$\mathbf1_{\{s\in E\}}R$ is again an original positive separately
closed current. Ordinary-locally, there is a finite Radon measure
$\nu$ and a measurable family $R_a$ of original currents on the
same ambient $U$ such that
\[
 R=\int R_a\,d\nu(a),\qquad
 \mathbf1_{\{s\in E\}}R=\int_E R_a\,d\nu(a).
 \tag{M7}
\]
Outside one common null set, all $R_a$ are positive, separately
closed, parameter-annihilated, and carried by $s=a$ on every
original coefficient. They retain the entire original test sheaf.

Let $Y$ be an allowed original $K$-analytic space and let
$\pi:U\to Y$ be an original analytic map. For each $i$, let
$H_i(a,y)$ be jointly Borel, uniformly bounded on compact
parameter/reference products, with every actual slice $H_i(a,\cdot)$
locally bounded original classical psh. Put
$h_i(x)=H_i(s(x),\pi(x))$. For $P+k\leq n$, the intrinsic recursion
\[
 D\bigl(h_1D(h_2\cdots D(h_kR)\cdots)\bigr)
       =\int T_{R_a,k}(H_1(a,\pi(\cdot)),\ldots,H_k(a,\pi(\cdot)))
                                                        \,d\nu(a)
 \tag{M8}
\]
is positive, separately closed and mixed symmetric. It obeys local
relative CLN bounds, Borel parameter linearity, and full weighted
joint decreasing continuity when every specified family
$H_{i,m}$, including any scalar-weight slot, decreases on the
whole parameter/reference product to an admissible slice family,
with the printed compact slice bounds.
\end{theorem}

\noindent\textbf{Status: \texttt{NEW\_PROVED} (restatement).}
The complete proofs are in
Sections~\ref{sec:r12-closed-parameter-bands},
\ref{sec:r12-original-current-disintegration}, and
\ref{sec:r12-bounded-borel-family-bt}.
Conditional currents remain on the original $K$-space. No
residue-field BT theory, geometric inverse pushforward, or
product-test spanning assertion is used.

\subsection{Energy duals separate variational and source questions}

Under (M2), form $\mathbb D_S=\mathscr C-\mathscr C$ with
\[
 \|h\|_{\rm ed}
     =\inf_{h=a-b}\bigl(r(a)+r(b)\bigr),\qquad
 J(u)(h)=-\int_V h\,d\widehat\mu_u.
 \tag{M9}
\]
The first infimum is over $a,b\in\mathscr C$.
The strong-continuous dual consists of the norm-dual functionals
continuous on every strongly convergent cone row.

\begin{theorem}[A closed source image and an exact source gap]
\label{thm:r12-main-energy-source}
The normed space $(\mathbb D_S,\|\cdot\|_{\rm ed})$ is an
actual Borel Banach space and $\|u\|_{\rm ed}=r(u)$ on the cone.
The image $J(\mathscr C)$ is norm closed in $\mathbb D_S^*$, and
$\|J(u)-J(v)\|_*$ is a complete metric inducing the strong topology.

For every strong-continuous dual functional $\ell$, there is a
unique $u=\Pi(\ell)$ and normal residual $n$ such that
\[
 \ell=J(u)+n,\qquad n(v)\leq0\ (v\in\mathscr C),\qquad n(u)=0.
 \tag{M10}
\]
This is the unique maximizer of $\ell(v)-e(v)/d$, not automatically
a source solution.

Let $\mu$ be a literal positive Radon datum with a strong-continuous
energy functional and at least one completed subsolution.
Let $H_\mu$ be the greatest relative subsolution furnished by
the frozen Perron theorem, and let $u_\mu=\Pi(\ell_\mu)$. Then
\[
 \Delta_S(\mu)=e(H_\mu)-e(u_\mu)\geq0,\qquad
 I(H_\mu,u_\mu)\leq q\Delta_S(\mu).
 \tag{M11}
\]
Moreover $\Delta_S(\mu)=0$ if and only if $\mu$ is the original
top measure of a member of $\mathscr C$. When it exists, this
potential is unique and equals both $H_\mu$ and $u_\mu$.
\end{theorem}

\noindent\textbf{Status: \texttt{NEW\_PROVED} (restatement).}
The proofs are in
Sections~\ref{sec:r12-energy-difference},
\ref{sec:r12-nonlinear-dual-source}, and
\ref{sec:r12-perron-variational-gap}.
A capacity modulus with
$\int_0^1 F(s)s^{-1-1/d}\,ds<\infty$ supplies an explicit
class of literal signed data in this dual, without requiring
domination by a single source.

There is also a distinct order-energy Banach lattice $\mathbb L_S$,
with
\[
 \|h\|_{\mathbb L_S}
       =r(P(-|h|)),\qquad
 P(f)=\sup_{\rm rel}\{v\in\mathscr C:v\preceq f\}.
 \tag{M12}
\]
It provides new literal scalar multipliers for every original
completed mixed current. The potential cone stays closed and
unchanged. The old weighted total-variation source distance is
the dual norm on this lattice; it is not the signed distance
in (M9).

\begin{theorem}[One geometry-only probe detects literal source solvability]
\label{thm:r12-main-single-normal-probe}
Under (M2), there is one bounded continuous original classical
finite-mass potential $\psi_0$, depending only on the analytic box,
with $-1\leq\psi_0<0$. Let $\mu$ be a literal positive Radon datum
whose energy functional is strong-continuous, and put
$u=\Pi(\ell_\mu)$ and $n=\ell_\mu-J(u)$. Then
\[
 \Xi_{S,\psi_0}(\mu)=-n(\psi_0)\geq0,
 \qquad
 \Xi_{S,\psi_0}(\mu)=0
       \ \Longleftrightarrow\ \widehat\mu_u=\mu
       \ \Longleftrightarrow\ \mu\in\widehat\mu(\mathscr C).
 \tag{M17}
\]
For every original real compact smooth scalar $\phi$, there is
$\varepsilon_\phi>0$, depending on the original finite test stage
and $\phi$ but not on $S$ or $\mu$, such that
\[
 \left|\int\phi\,d(\widehat\mu_u-\mu)\right|
              \leq \varepsilon_\phi^{-1}\Xi_{S,\psi_0}(\mu).
 \tag{M18}
\]
No subsolution family is assumed. For a nonliteral energy dual,
zero in (M17) implies only zero compact-test shadow, not zero
functional on the entire energy space.
\end{theorem}

\noindent\textbf{Status: \texttt{NEW\_PROVED} (restatement).}
Section~\ref{sec:r12-single-probe-normal-residual} constructs
$\psi_0$ so that $\psi_0\pm\varepsilon_\phi\phi$ are actual
admissible potentials. Applying the normal inequalities to these
two perturbations proves (M18). Literalness then identifies the
zero shadow with the entire original source equation. This proves
a testable exact criterion, not automatic vanishing of $\Xi$.

\subsection{A complete specified parameter module, not an ambient replacement}

Assume (M2) with $R$ as background, with the unit parameters $s_i$
defined on an ordinary neighborhood of $K_0$ and with (M6) holding
there. Let $R$ have degree $(P,P)$ and set
$k=n-P\geq1$, $d_R=k+1$. The geometry-only probe may be chosen
for this degree as well. Write
\[
 \omega=(D\psi_0)^k\wedge R,\qquad
 \theta=s_*((-\psi_0)\omega),\qquad
 h_c=\psi_0 c(s),\quad c\in L^{d_R}_+(\theta).
 \tag{M19}
\]
These are explicitly defined relative Borel parameter potentials;
they need not be ambient classical psh.

\begin{theorem}[Parameter energy, compatible sources, and attained capacity]
\label{thm:r12-main-parameter-module}
The mixed products of the potentials (M19) are intrinsic positive
separately closed original currents on every original test, for
$0\leq j\leq k$:
\[
 T_{R,j}(h_{c_1},\ldots,h_{c_j})
    =\left(\prod_{i=1}^j c_i(s)\right)(D\psi_0)^j\wedge R,
 \qquad
 B_{s,\psi_0}(h_{c_0},\ldots,h_{c_k})
    =\int\prod_{i=0}^k c_i\,d\theta.
 \tag{M20}
\]
The module is strongly complete with energy $\|c\|_{d_R}^{d_R}$,
has the printed relative comparison, domination, quasicontinuity
and obstacle properties, and solves each compatible datum
$b(s)\omega$, $b\geq0$, $b\in L^{d_R/k}(\theta)$ uniquely within
the module by $h=\psi_0b(s)^{1/k}$.

For a Borel set $A\subset W\Subset V$, where $W$ is ordinary open,
put $g_A=d s_*(\mathbf1_A\omega)/d\theta$. Its energy-normalized
capacity has the exact formula
\[
 \mathfrak c_{\psi_0}(A)
    :=\sup_{\substack{c\geq0\\\|c\|_{d_R}\leq1}}
               \int_A c(s)^k\,d\omega
     =\|g_A\|_{d_R}.
 \tag{M21}
\]
If $\omega(A)>0$, the unique maximizing parameter coefficient is
$c_A=g_A/\|g_A\|_{d_R}$ and is bounded. The corresponding
potential has ordinary zero boundary value. This capacity is
raw/ordinary-outer regular on every Borel subset of $W$ and has
exactly the $\omega$-null ideal there.
\end{theorem}

\noindent\textbf{Status: \texttt{NEW\_PROVED} (restatement).}
The complete proofs are in
Sections~\ref{sec:r12-parameter-energy-Dirichlet} and
\ref{sec:r12-parameter-energy-capacity}.
For unbounded coefficients the boundary is the stated
bounded-parameter completion boundary, not an asserted ordinary
boundary limit. Uniqueness among all ambient potentials and
equality with unrestricted original BT capacity are not claimed.
Every arbitrary $S$ enters through its proved mixed-slot current
$R=\Omega_s\wedge S$, without imposing conormality on $S$ itself.

\subsection{A general compact family of full-tester energy data}

\begin{theorem}[Weakly compact data and uniform strong test continuity]
\label{thm:r12-main-full-tester-data}
Retain the arbitrary background and strict-box hypotheses (M2).
For a compact $H\Subset V$ and finite $m_0,M\geq0$, let
\[
 \begin{split}
 \mathcal G(H,m_0,M)=\{\lambda\geq0:\ &\operatorname{supp}\lambda\subset H,
 \quad\lambda(V)\leq m_0,\\
 &\lambda(E)\leq M C_S^*(E;V)\quad\text{for every Borel }E\}.
 \end{split}
 \tag{M23}
\]
This is a weakly compact family of original Radon measures, and
its literal integration functionals form a weak-$*$ compact subset
of $\mathbb D_S^*$ contained in the strong-continuous dual.
For each energy-root bound $R$, there is a constant $C$ depending
only on the printed data such that, if $r(u),r(v)\leq R$ and
$0<\varepsilon=I(u,v)^{2^{-q}}\leq1$, then
\[
 \sup_{\lambda\in\mathcal G(H,m_0,M)}
              \int|u-v|\,d\lambda
               \leq C\varepsilon\bigl(1+\log(1/\varepsilon)\bigr).
 \tag{M24}
\]
The integral is zero when $I(u,v)=0$. If $\lambda_j\to\lambda$
weakly in this family and $u_j\to u$ strongly in $\mathscr C$,
then $\int u_j\,d\lambda_j\to\int u\,d\lambda$.

Every weak limit of $\chi(Da_j)^q\wedge S$, where
$0\leq\chi\leq M$ is one fixed original compact smooth scalar
and every $a_j$ is an arbitrary original normalized classical
tester on the whole $V$, belongs to such a family. No special
moment geometry or source-compatibility condition is imposed.
\end{theorem}

\noindent\textbf{Status: \texttt{NEW\_PROVED} (restatement),}
with the cluster membership a \texttt{DIRECT\_COROLLARY} of
the full-tester capacity bound. The full proof and explicit
constant are in Section~\ref{sec:r12-capacity-cluster-strong-dual}.
The theorem is not operator-norm compactness and does not give
actual normalized tester realization of every weak limit.
Indeed a constant tester row with a compact cutoff can have
strictly positive global source normal $\Xi$ while having no
raw-capacity gain at all. Source compatibility and compact
no-gain are therefore separate conditions.

A complementary stability theorem proves that every fixed measure
order interval $0\leq\mu\leq M\widehat\mu_v$ is a weak-$*$ compact,
norm-closed set of \emph{literal} strong-dual data. Its proof
uses $L^1(\widehat\mu_v)$ approximation on the full difference
space, not merely agreement on compact tests. It neither places
all finite-energy data in such intervals nor forces zero normal.

\subsection{A positive current records nonlinear weak defects}

\begin{theorem}[Positive Jensen primitive and a single no-defect probe]
\label{thm:r12-main-Jensen}
For every arbitrary allowed $S$, every finite bounded classical
tuple $h_i$ and convex weights $\lambda_i$, put
$w=\sum_i\lambda_i h_i$. For $2\leq k\leq q$ the original
positive coefficient current
\[
 V_k=\sum_i\lambda_i\sum_{a=0}^{k-2}(k-1-a)
       \Gamma_{(Dh_i)^a(Dw)^{k-2-a}S}(h_i-w)
 \tag{M13}
\]
satisfies
\[
                \sum_i\lambda_i T_{S,k}(h_i^{[k]})
                        =T_{S,k}(w^{[k]})-DV_k.
 \tag{M14}
\]
It need not be closed, and $-DV_k$ has no fixed sign.
The formulas extend to the full first-energy completion under (M2).

For $q\geq2$, there is one bounded continuous actual finite-mass potential
$\psi_0$, depending only on the analytic box and not on $S$,
with the following property. For any energy-bounded weak top-source
cluster $\nu$, let $u$ be its minimum-energy convex-tail center.
Then some positive original current $V_q$ satisfies
\[
 \nu=\widehat\mu_u-DV_q,\qquad
 \|\mu_{V_q,\eta}\|_{\rm TV}
       \leq C_\eta
       \left(\int(-\psi_0)\,d\nu-B(\psi_0,u^{[q]})\right)
 \tag{M15}
\]
for every original compact coefficient test $\eta$.
The number in parentheses is nonnegative and is zero exactly
when $\nu=\widehat\mu_u$ on the full original test sheaf.
\end{theorem}

\noindent\textbf{Status: \texttt{NEW\_PROVED} (restatement).}
See Sections~\ref{sec:r12-positive-jensen-defects},
\ref{sec:r12-weak-cluster-defects}, and
\ref{sec:r12-universal-jensen-probe}.
The frozen R3 moving-atom example shows that this scalar loss
need not vanish, even for uniformly bounded families with
common boundary controls. Its old counterexample is not counted anew.

\subsection{What an unrestricted final theorem would get wrong}

For a proper full-rank moment $t:V\to Q\subset\mathbf R^r$ and the
explicit factorization $S=(t^*\beta_0)^{r-q}\wedge R$, with $q<r$,
every weak bounded-source limit has a closed positive real-current
lift:
\[
                     t_*\mu=T\wedge\beta_0^{r-q},
                 \qquad d'T=d''T=0,\quad T\geq0.
 \tag{M16}
\]
This is a necessary condition, not a sufficient reconstruction.
It implies a collar inequality, rules out nonzero compact
tropical source carriers, and makes every proper moment fiber
outer-capacity null. On the original split torus, these statements
give nonsolvable dominated data in every $2\leq q<r$.
Their variational maximizers exist, but their source, contact,
and Perron--variational defects persist.

There is also a direct original-smooth perturbation obstruction
in every $n>q\geq1$. For the explicit allowed radial background
of Section~\ref{sec:r12-radial-pole-order-dual}, a bounded actual
source $\mu_0$ and a bounded nonnegative original smooth $G$ satisfy:
\[
 \begin{gathered}
 \mu_\epsilon=(1+\epsilon G)\mu_0,\\
 \begin{cases}
  \mu_\epsilon\text{ has a unique full-cone source},&-c<\epsilon\leq0,\\
  \mu_\epsilon\text{ has no full-cone source},&\epsilon>0.
 \end{cases}
 \end{gathered}
 \tag{M22}
\]
Here $c>0$ ensures strictly positive densities. All these measures
have the same Borel null ideal, are dominated by a bounded source
up to a constant, and converge in signed energy-dual norm to
$\mu_0$. Thus compatibility cannot be removed merely by assuming
strict positive density or arbitrarily small smooth perturbation.
Section~\ref{sec:r12-smooth-density-source-instability} proves the
threshold by an exact covariance normal over the full cone, not
only radial competitors.

Two further distinctions are proved by exact examples. A continuous
Banach-dual functional need not be strongly continuous, even if
nonnegative on the potential cone. Conversely, a positive
strong-continuous order functional can restrict to a Dirac measure
on all original smooth tests while failing to be literal integration
on the relative Borel quotient. Arbitrary relative-current positivity
also does not supply a hyperbolic Hessian cone for an enlarged
potential class. These are failures of specific proposed routes,
not failures of the bounded arbitrary-$S$ calculus.

\subsection{Status convention and dependency order}

Each formal result retains one of
\texttt{EXTERNAL\_KNOWN}, \texttt{INTERNAL\_INPUT},
\texttt{DIRECT\_COROLLARY}, \texttt{NEW\_PROVED},
\texttt{CONDITIONAL}, \texttt{OPEN}, \texttt{FALSE},
or \texttt{FAILED\_ATTEMPT}.
Here \texttt{NEW\_PROVED} means proved in this record relative
to the frozen research chain; it is not a claim of a completed
worldwide priority search.

The main theorems above are restatements of the detailed results,
not additional counted discoveries. The conormal calculus uses
only frozen bounded products, coefficient measures and real
measure theory. The energy constructions use the frozen strong
cone before defining either Banach space. Jensen currents use the
bounded recursion and then the already proved completion.
Source obstructions use positive lifts and compact Stokes,
independently of any analytic quotient claim. The GQ section
retains the missing actual-zero-to-cap-zero implication.

The final unrestricted analytic PS/PL relation quotient, low
real-degree signed constant-order transfer, full raw-capacity
regularity, and appropriately formulated general Dirichlet
existence remain separate questions. Their exact surviving
forms and the disproved shortcuts are listed at the end.


% END REVIEWED MODULE research_r12/scope_and_main_results.tex


\part{Closed parameter bands and original conditional currents}

% BEGIN REVIEWED MODULE research_r12/gq/closed_parameter_bands.tex
\section{Closed parameter bands and an original-current Radon calculus}
\label{sec:r12-closed-parameter-bands}

Fix the original finite extension $K/\mathbf Q_p$ and its original
absolute value. Let $U$ be an ordinary open in a good, Hausdorff,
boundaryless, quasi-smooth, pure $n$-dimensional $K$-analytic space.
All currents and tests in this section are original CLD objects.
In particular, an equality is required on every original compact
smooth complementary test, allowing every actual moment refinement.
No equality on an enlarged GK delta-test sheaf is asserted.

Let $f_1,\ldots,f_c$ be original analytic units, set
\[
 s_i=\log|f_i|,\qquad s=(s_1,\ldots,s_c):U\longrightarrow J\subset\mathbf R^c,
 \qquad
 \Omega_s=\bigwedge_{i=1}^c(d's_i\wedge d''s_i),
 \tag{PB1}
\]
where $J$ is a real open containing the image under consideration.
No rank, independence, submersion, saturation, or product hypothesis
on $s$ is imposed in this section. The original smooth functions
$s_i$ satisfy $Ds_i=0$ and their prime and double-prime differentials
are separately closed.

\begin{definition}[Parameter-annihilated currents]
\label{def:r12-parameter-annihilated}
A coefficient-Radon current $R$ on $U$ is \emph{$s$-annihilated} if
\[
                 d's_i\wedge R=d''s_i\wedge R=0
                         \qquad(1\le i\le c).
 \tag{PB2}
\]
An allowed $s$-background is an $s$-annihilated, original symmetric
middle-positive, separately $d'$- and $d''$-closed diagonal current.
The adjective describes a proved extra structure of a derived
current; it does not restrict the background $S$ in OPEN-GQ.
\end{definition}

For every arbitrary original allowed $(p,p)$-background $S$ with
$p+c\le n$, the derived current $\Omega_s\wedge S$ is an allowed
$s$-background. Indeed $\Omega_s$ is smooth, strongly positive and
separately closed, so its wedge preserves the original middle cone;
every conormal in (PB2) is repeated in this product. This uses no
regularization or replacement of $S$. It also covers all cases where
$\Omega_s$ happens to vanish because of moment-rank degeneracy.

\subsection{The Borel parameter algebra preserves closedness}

\begin{lemma}[Bounded smooth approximation for finitely many parameter measures]
\label{lem:r12-parameter-L1-approximation}
Let $\nu_1,\ldots,\nu_N$ be finite positive Radon measures on a
compact subset of $J$, and let $b:J\to\mathbf R$ be Borel and
bounded on a neighborhood of that subset. There are smooth real
functions $b_m$ on $J$, with one common bound, such that
\[
                   \sum_{j=1}^N\int|b_m-b|\,d\nu_j\longrightarrow0.
 \tag{PB3}
\]
If $b\ge0$, the $b_m$ can be chosen nonnegative. No measure in this
statement is assumed to be nonatomic.
\end{lemma}

\noindent\textbf{Status: \texttt{EXTERNAL\_KNOWN}} as elementary
Radon regularity; the short proof fixes precisely the version used.

\begin{proof}
Replace the finite family by its sum. Truncate $b$ to its fixed
bound on a compact neighborhood and approximate it in $L^1$ by
bounded simple functions. Inner and outer regularity approximate
each indicator by a continuous function between zero and one.
Continuous compactly supported functions on $J$ can be uniformly
approximated by smooth functions by convolution after a compact
cutoff. Truncation with a smooth approximation to the scalar
clipping map keeps a common bound and, when required, nonnegativity.
A diagonal choice proves (PB3).
\end{proof}

\begin{theorem}[Closed Borel bands and differential commutation]
\label{thm:r12-closed-parameter-bands}
Let $R$ be an allowed $s$-background. For every locally bounded
Borel $b:J\to\mathbf R$, the intrinsic coefficient multiplication
$b(s)R$ is separately closed and $s$-annihilated. It is middle-positive
if $b\ge0$. In particular, for every Borel $E\subset J$,
\[
                            R_E:=\mathbf1_{s^{-1}(E)}R
 \tag{PB4}
\]
is a genuine positive separately closed original current.

More generally, let $u$ be an original locally bounded Borel scalar
such that $T=D(uR)$ is coefficient-Radon. Then $T$ is $s$-annihilated
and, on all original tests,
\[
                         D\bigl(u\,b(s)R\bigr)=b(s)D(uR).
 \tag{PB5}
\]
Thus (PB5) applies in particular to every frozen locally bounded
classical BT insertion into $R$.
\end{theorem}

\noindent\textbf{Status: \texttt{NEW\_PROVED}} for the original
closed-band calculus. The coefficient-Radon nature and positivity of
the frozen bounded BT insertion are \texttt{INTERNAL\_INPUT}.

\begin{proof}
First take $b$ smooth. Its differentials are combinations of the
$d's_i$ and $d''s_i$. Consequently
\[
 d'b(s)\wedge R=d''b(s)\wedge R=Db(s)\wedge R=0.
 \tag{PB6}
\]
The distributional product rule and separate closedness of $R$
give separate closedness of $b(s)R$. Conormality follows immediately
from (PB2). For a compact original test $\eta$ complementary to $T$,
the smooth identity
\[
 D(b(s)\eta)=b(s)D\eta+d'b(s)\wedge d''\eta
                    -d''b(s)\wedge d'\eta+Db(s)\wedge\eta
\]
and (PB6), still valid after scalar multiplication by $u$, imply
\[
 \langle b(s)D(uR),\eta\rangle
   =\langle uR,D(b(s)\eta)\rangle
   =\langle uR,b(s)D\eta\rangle.
 \tag{PB7}
\]
This is (PB5) for smooth $b$.

For Borel $b$, fix one arbitrary compact original test stage. Choose
finite coefficient-Radon frames controlling $R$ on the supports of
the test and its required derivatives. Push the total variations of
these finitely many coefficient measures forward by $s$; they are
finite Radon measures on the real parameter space. Apply
Lemma~\ref{lem:r12-parameter-L1-approximation}. The resulting
$b_m(s)R$ converge to $b(s)R$ against every pairing needed for the
closedness equation of that test. Distributional differentiation is
continuous, so closedness passes to the limit. Multiplication of the
already zero coefficient measures in (PB2) proves conormality.
Positive scalar multiplication proves middle positivity when $b\ge0$.

For (PB5), include in the same finite controlling family the
coefficient measures of $T$ on $\operatorname{supp}\eta$ and those
of $R$ on $\operatorname{supp}D\eta$. The common compact bound on
$u$ then gives convergence of both sides of (PB7). This proves
(PB5) without inferring convergence against a discontinuous weight
from weak current convergence alone. The approximating functions
may depend on the test: since that test was arbitrary, the resulting
intrinsic equality holds on the whole original test sheaf. No
countability assumption on the collection of all refinements is
needed.

Finally, separate closedness of $d's_i,d''s_i$ gives directly, as
distributions,
\[
 d's_i\wedge D(uR)=D(u\,d's_i\wedge R)=0,
 \qquad
 d''s_i\wedge D(uR)=D(u\,d''s_i\wedge R)=0.
 \tag{PB8}
\]
This proves the asserted conormality of $T$. All arguments involve
intrinsic original forms and coefficient measures, and hence are
unchanged after any actual refinement containing the chosen data.
\end{proof}

\begin{corollary}[A closed current-valued Radon resolution]
\label{cor:r12-current-valued-parameter-resolution}
The assignment $E\mapsto R_E$ is countably additive on Borel sets
in every original compact-test pairing. Its local coefficient
variation is controlled by that of $R$. Every $R_E$ is separately
closed and positive. For all permitted bounded classical tuples,
\[
 T_{R_E,k}(u_1,\ldots,u_k)
       =\mathbf1_{s^{-1}(E)}T_{R,k}(u_1,\ldots,u_k).
 \tag{PB9}
\]
If $E_1,\ldots,E_N$ form a finite Borel partition on the parameter
range under consideration, then $R=\sum_iR_{E_i}$ and every bounded
mixed BT product has the corresponding exact finite decomposition.
The same conclusions hold for countable partitions whenever a
compact test is fixed and the sums are interpreted in coefficient
variation.
\end{corollary}

\noindent\textbf{Status: \texttt{DIRECT\_COROLLARY}}.

\begin{proof}
Coefficient-Radon multiplication gives countable additivity and the
variation bound. The theorem gives positivity and closedness of each
band. Apply (PB5) repeatedly: every positive mixed tail is separately
closed, coefficient-Radon, and $s$-annihilated by (PB8). This proves
(PB9). Sum the resulting coefficient measures over a partition.
\end{proof}

\paragraph{What is and is not a fiber construction.}
The resolution (PB4) retains every Borel set, including individual
parameters and atoms. It does not discard an exceptional parameter
set. It is a current-valued measure on the parameter space, not a
construction of BT products over a completed residue field. Its band
currents live on the original $U$, are carried measure-theoretically
by $s^{-1}(E)$, and use every original test there. Their topological
supports are only asserted to lie in the closure of that carrier.
No conditional fiber current on a new analytic field is asserted.

The conormal hypothesis is essential. A general Borel restriction
of a closed current need not be closed. Moreover, the frozen R10
annular example shows that arbitrary Borel scalar coefficients
cannot commute with curvature insertion: for the unit Gauss point
$\xi$, $\mathbf1_{\{\xi\}}[U]=0$ but
$\mathbf1_{\{\xi\}}D((\log|T|)_+[U])=\delta_\xi$.
That example does not contradict (PB5), because $[U]$ is not
annihilated by $d'\log|T|$ and $d''\log|T|$.

\paragraph{Frozen inputs and exact scope.}
Only the frozen original coefficient-Radon calculus, smooth form
Leibniz rules, positivity under strongly positive smooth wedges,
and bounded arbitrary-background BT theorem are used. The R10
source \path{signed_descent/pws_differential_curvature_module.tex}
is used for its explicit warning against unrestricted Borel
coefficient linearity; its finite PS module is not enlarged
without the separate proof above. No R/B/A2 proof is re-audited.

% END REVIEWED MODULE research_r12/gq/closed_parameter_bands.tex


% BEGIN REVIEWED MODULE research_r12/gq/parameter_family_bt_calculus.tex
\section{A continuous parameter-family Bedford--Taylor calculus}
\label{sec:r12-parameter-family-bt}

We retain the original analytic category, the original field and norm,
and the full original test sheaf of
Section~\ref{sec:r12-closed-parameter-bands}. Let $s:U\to J$ be the
same unit-logarithm parameter map. In addition let $Y$ be an allowed
original $K$-analytic space and let $\pi:U\to Y$ be an original
analytic map. No product hypothesis is required.

\begin{definition}[Continuous slice-psh families]
\label{def:r12-continuous-slice-psh}
Let $H:J\times Y\to\mathbf R$ be a finite continuous function,
where this product denotes the ordinary topological parameter space.
Assume that, for every $a\in J$, the actual function
$H_a(y)=H(a,y)$ is a locally bounded original classical psh function
on $Y$. Associate the actual continuous function on $U$
\[
                      h(x)=H\bigl(s(x),\pi(x)\bigr).
 \tag{PF1}
\]
Write $\mathcal P_{s,\pi}(U)$ for the cone of functions admitting
such a presentation locally in the ordinary topology. A chosen
presentation is not part of the function. Different functions in
a tuple may initially use different reference spaces and maps;
a finite product of their reference spaces and maps puts them in
one common notation. Only compact local stages are used.
\end{definition}

Such an $h$ need not be psh on $U$. Even on a product with a
one-dimensional parameter, the example $H(a,y)=-a^2$, constant in
$y$, is slice-psh while $h=-s^2$ has negative horizontal curvature.
The new products will be defined on parameter-annihilated backgrounds,
not on every positive background for these enlarged potential slots.

\begin{lemma}[Finite parameter step functions]
\label{lem:r12-finite-parameter-step-products}
Let $R$ be an allowed $s$-background of degree $(P,P)$, let
$P+k\le n$, and let $E_1,\ldots,E_N$ be a finite Borel partition
of the relevant parameter region. For every $i,j$ let $u_{i,j}$ be
a locally bounded original classical psh function on $U$, and put
\[
               h_i=\sum_{j=1}^N\mathbf1_{s^{-1}(E_j)}u_{i,j}
                          \qquad(1\le i\le k).
 \tag{PF2}
\]
The recursive distributions $D(h_1D(h_2\cdots D(h_kR)\cdots))$
are well-defined coefficient-Radon currents and satisfy
\[
 D(h_1D(h_2\cdots D(h_kR)\cdots))
              =\sum_{j=1}^N T_{R_{E_j},k}(u_{1,j},\ldots,u_{k,j}).
 \tag{PF3}
\]
They are symmetric in the slots, positive, separately closed, and
$s$-annihilated. They agree with the frozen products when the slots
are independent of the parameter partition.
\end{lemma}

\noindent\textbf{Status: \texttt{DIRECT\_COROLLARY}} of the closed
band theorem and the frozen arbitrary-background bounded BT theorem.

\begin{proof}
Every $R_{E_j}$ is a genuine allowed positive closed background.
Equation (PB9) says that each of its BT tails remains carried by
$s^{-1}(E_j)$. Multiplication of such a tail by $h_i$ in (PF2)
therefore equals multiplication by $u_{i,j}$. Starting with
$R=\sum_jR_{E_j}$ and applying $D$ successively proves (PF3).
Frozen mixed symmetry and positivity apply on each band. Separate
closedness and (PB8) give the remaining assertions. No derivative
of a Borel indicator has been ignored: the fact that these particular
band restrictions are closed was proved before using them.
\end{proof}

\begin{theorem}[Mixed products for continuous slice-psh potentials]
\label{thm:r12-parameter-family-bt-calculus}
Let $R$ be an allowed $s$-background of degree $(P,P)$, and let
$h_1,\ldots,h_k\in\mathcal P_{s,\pi}(U)$ with $P+k\le n$.
Then the intrinsic recursion
\[
 T^s_{R,0}=R,\qquad
 T^s_{R,k}(h_1,\ldots,h_k)
                  =D\bigl(h_1T^s_{R,k-1}(h_2,\ldots,h_k)\bigr)
 \tag{PF4}
\]
has all the following properties.
\begin{enumerate}
\item Its values are symmetric, original middle-positive,
separately closed, coefficient-Radon, and $s$-annihilated.
\item The differentiated slots are mixed symmetric and multilinear
on the cone, and the construction depends only on their actual
functions. It commutes with ordinary restriction and gluing.
\item It has the frozen relative CLN bounds on each fixed compact
test stage, with the same type of background controls and constants
depending on uniform bounds of the chosen parameter slices on the
fixed enlarged compact stage, not on any parameter partition or
its derivatives. For a fixed continuous $H_i$ these bounds follow
from compactness of the parameter and reference-space images.
\item If the slots converge locally uniformly to continuous
members of this class and their chosen families have uniform slice
bounds on each fixed enlarged compact stage, their mixed products
converge on every original compact smooth test. The same holds for
mixed finite-step approximants (PF2) converging locally uniformly to
the given continuous slots, provided every chosen classical
constituent $u_{i,j}$ has a common bound on that enlarged compact.
This is a bound on the constituents, not just on the assembled
step function along its own parameter bands. The approximants (PF6)
have it automatically.
\item For every locally bounded Borel $b$ on the parameter space,
\[
            T^s_{b(s)R,k}(\mathbf h)=b(s)T^s_{R,k}(\mathbf h),
 \tag{PF5}
\]
where signed $b$ is interpreted by the difference of its positive
and negative parts. In particular the full closed-band Radon
resolution persists under every mixed product.
\end{enumerate}
All assertions use every original complementary test and every
actual moment refinement. No smooth approximation of $R$ is made.
\end{theorem}

\noindent\textbf{Status: \texttt{NEW\_PROVED}} as the enlarged
relative family calculus. The bounded classical BT products on each
closed band, their relative CLN estimate, and continuous coefficient
Radon convergence are the exact frozen \texttt{INTERNAL\_INPUT}s.

\begin{proof}[Positivity of one insertion]
Work on a fixed compact test neighborhood, enlarged slightly inside
an ordinary chart. The parameter image and the $\pi$-image of this
larger compact are compact. Continuity of $H$ makes
$a\mapsto H_a$ uniformly continuous in the sup norm on the latter
compact when $a$ ranges over a compact parameter neighborhood.
Choose finer finite Borel parameter partitions, with samples $a_j$
in their cells, and set
\[
                 h_m=\sum_j\mathbf1_{s^{-1}(E_{m,j})}
                                      (H_{a_j}\circ\pi).
 \tag{PF6}
\]
The slices pull back to original classical psh functions, by their
all-curve definition. These are functions on the original $U$;
no current is pulled back. The approximants have a common local
sup bound and converge uniformly to $h$ on the chosen compact.
By (PF3), $D(h_mR)$ is positive and closed. Since $R$ is
coefficient-Radon, $h_mR\to hR$ against all pairings needed for a
distribution derivative. Thus $D(h_mR)\to D(hR)$ as original
currents. The positive cone is closed under this convergence, so
$D(hR)$ is positive and hence coefficient-Radon. Separate closedness
is automatic for $D(hR)$, and (PB8) proves conormality.

Apply this argument recursively. Each already constructed tail is
again an allowed $s$-background, so (PF4) is well-defined and has
the assertions in (1). It is already intrinsic in the actual scalar
functions, before any symmetry or choice-independence argument.
\end{proof}

\begin{proof}[Uniform controls and convergence]
For a simultaneous finite-step tuple, take one common refinement of
the finite parameter partitions. Suppose its sliced potentials have
absolute values at most $M_i$ on the fixed larger compact stage.
The frozen relative CLN estimate applies on every positive closed
band background $R_{E_j}$. In the notation of its finite positive
control forms it yields
\[
 \sum_j\bigl\langle
       T_{R_{E_j},k}(u_{1,j},\ldots,u_{k,j}),\Theta\bigr\rangle
   \le C\prod_{i=1}^k(1+M_i)
                        \sum_j\sum_\ell\langle R_{E_j},\Xi_\ell\rangle
   = C\prod_{i=1}^k(1+M_i)\sum_\ell\langle R,\Xi_\ell\rangle.
 \tag{PF7}
\]
Use the usual harmless scalar normalization if the frozen estimate
is written with $M_i$ instead of $1+M_i$. The constant and the
finite controls are those of the one fixed compact stage. In
particular no number of cells and no derivative of a partition
function occurs. This step uses genuinely closed backgrounds,
not an unproved closedness claim about arbitrary restrictions of $R$.

We prove convergence of simultaneous step approximants by induction
on $k$. The case $k=1$ was just proved. By induction the tails
$Q_m$ converge weakly to the intrinsic tail $Q$ in (PF4), and
(PF7) gives uniform finite coefficient-variation controls on each
fixed original test stage. If $h_{1,m}\to h_1$ uniformly there,
\[
             (h_{1,m}-h_1)Q_m\longrightarrow0,
                   \qquad h_1Q_m\longrightarrow h_1Q.
 \tag{PF8}
\]
The first convergence follows from uniform variation bounds. The
second is the frozen continuous multiplier theorem (R, Lemma 3.1,
transported by the accepted ordinary algebraic charts). Equivalently,
on every coefficient measure it is weak Radon convergence against
a fixed continuous multiplier; original smooth scalar density and
the same uniform controls prove it directly. Only a continuous
limiting multiplier is used here. Taking $D$ proves convergence
to (PF4), and passing (PF7) to the limit proves the relative CLN
bound for the continuous family slots.

For a locally uniformly convergent tuple of continuous family
potentials with the stated uniform slice bounds, the just proved
CLN bounds again give uniform controls for its tails. The same
induction using (PF8) proves (4).
Thus this part does not depend on the particular finite-step
approximations chosen above.
\end{proof}

\begin{proof}[Symmetry, locality, and the parameter algebra]
The finite-step products are symmetric by (PF3). Approximate all
slots simultaneously by (PF6), with a common finite partition,
and pass to the proved limits. This gives mixed symmetry in (2).
The recursion is linear in its first slot; symmetry gives linearity
in every slot whenever sums remain in the cone. Signed differences
then extend by multilinearity, with independence from their
presentations because a zero actual scalar function gives a zero
recursive current. Ordinary locality follows directly from (PF4).
It also permits the finite compact-stage constructions to glue.

Apply (PB5) with $u=h_1$ to every recursively constructed tail.
That lemma requires only local boundedness and a coefficient-Radon
derivative, both already proved here; it does not require $h_1$
itself to be an ordinary psh function on $U$. Induction proves
(PF5). The argument first uses nonnegative $b$ so that $b(s)R$
is allowed, and then takes the difference of two such identities.
Every expression used in these proofs is an original function,
current, or test. Reexpressing them on any common actual refinement
therefore preserves all equalities.
\end{proof}

\begin{corollary}[A new all-dimensional arbitrary-background construction]
\label{cor:r12-parametric-arbitrary-S}
Let $S$ be any original symmetric middle-positive separately closed
$(p,p)$-current on $U$. If $p+c+k\le n$ and the $h_i$ satisfy
Definition~\ref{def:r12-continuous-slice-psh}, then
\[
               T^s_{\Omega_s\wedge S,k}(h_1,\ldots,h_k)
 \tag{PF9}
\]
is a canonical positive separately closed original current, with
all the properties in the theorem. For one slot,
\[
                 T^s_{\Omega_s\wedge S,1}(h)
                                =\Omega_s\wedge D(hS).
 \tag{PF10}
\]
The original arbitrary $S$ is fixed and unchanged. There is no
matroidal, simplicial-polytopal, model-current, or $W_{\rm loc}$
hypothesis on it or on the reference space $Y$.
\end{corollary}

\noindent\textbf{Status: \texttt{DIRECT\_COROLLARY}} of the two
new theorems; it is not the unrestricted signed analytic quotient.

\begin{proof}
The paragraph following (PB2) proves that $\Omega_s\wedge S$ is
an allowed $s$-background. Apply the theorem. Since $\Omega_s$ is
smooth and separately closed, distributional Leibniz gives (PF10).
\end{proof}

\paragraph{Limits of this enlargement.}
No arbitrary Borel function of the analytic point is admitted as a
central coefficient: only Borel functions of the annihilated
parameters are. A bounded decreasing family whose limit is not
continuous is not covered by the continuity assertion above.
Nor does the theorem infer slice-psh data from a general analytic
PS/PL zero relation. That latter inference, without the additional
product/trace structure proved separately, remains an OPEN-GQ/CTO
gate. This calculus is a new derived construction available for
every original $S$, not a replacement of $S$ by a special class in
the statement of the general goal.

% END REVIEWED MODULE research_r12/gq/parameter_family_bt_calculus.tex


% BEGIN REVIEWED MODULE research_r12/gq/product_parameter_order_transfer.tex
\section{Product-parameter order transfer with every original test retained}
\label{sec:r12-product-parameter-order}

The new closed-band calculus has an unconditional original-current
order consequence when the parameter family has a fixed original
reference fiber. The product hypothesis below concerns the analytic
data, never the allowed arbitrary background $S$.

Let $Y$ be good, Hausdorff, boundaryless, quasi-smooth, pure of
dimension $d\ge1$ over the original $K$. Let $J\subset\mathbf R^c$
be open, $c\ge1$, and put
\[
 A_J=\{z\in(\mathbf G_m^c)^{\mathrm{an}}:
                         (\log|z_1|,\ldots,\log|z_c|)\in J\},
 \qquad U=A_J\times_KY,
 \qquad n=c+d.
 \tag{PT1}
\]
Write $s_i=\log|z_i|$, $\pi:U\to Y$, and use $\Omega_s$ from
(PB1). Let $H:J\times Y\to\mathbf R$ be finite continuous and
let $h(x)=H(s(x),\pi(x))$. No psh hypothesis on $h$ or on the
slices $H_a$ is imposed initially.

\subsection{The precise fixed-field regularization input}

We use only the following two authorized \texttt{INTERNAL\_INPUT}
statements, without re-auditing their proofs.
\begin{itemize}
\item R, \emph{Continuous CLD-positive potentials and adelic mixed
measures}, 10 November 2026, Theorem 2.9, printed p.~4: on an ordinary
open of a smooth quasi-projective algebraic $K$-variety, a finite
continuous function with positive original curvature current is
locally uniformly approximable by smooth functions with positive
original curvature forms on a fixed ordinary neighborhood.
\item B, \emph{Locally bounded classical plurisubharmonic functions
and mixed currents}, 10 November 2026, Lemma D.4, printed p.~29:
every point of an allowed original analytic space has, inside any
prescribed ordinary neighborhood, a smaller ordinary neighborhood
$K$-analytically isomorphic to an ordinary open of a smooth affine
$K$-variety. No change of $K$ occurs.
\end{itemize}
The archived PDFs have respective SHA256 values
{\small\nolinkurl{123269c9727efb8f530ba8ff79b30a44d49b5dc8068e15e305843099af84e5db}}
and
{\small\nolinkurl{6e9f599abf514df0f53d87be4d705d2830f8a8ff9e7f6c9c2c6f31595c33a5fe}}.
These are the exact R4 input copies, reached through the verified
R11--R3 chain. A finite extension of $\mathbf Q_p$ can be realized
as a number-field completion without altering the prescribed norm,
as already fixed in the frozen scope register.

\begin{lemma}[Continuous first-order transfer along an original map]
\label{lem:r12-continuous-original-map-transfer}
Let $v$ be a finite continuous function on an allowed original
analytic space $Y$ with $D[v]\ge0$. For any original analytic map
$\pi:V\to Y$ and any original symmetric middle-positive separately
closed current $T$ on $V$, in every permitted degree,
\[
                           D\bigl((v\circ\pi)T\bigr)\ge0.
 \tag{PT2}
\]
If $D[v]=0$, the current in (PT2) is zero.
\end{lemma}

\noindent\textbf{Status: \texttt{DIRECT\_COROLLARY}} of the exact
two internal inputs above, the original smooth positive wedge rule,
and coefficient-Radon multiplication. This is the already frozen
continuous first-order mechanism, restated to pin its scope.

\begin{proof}
Use B, Lemma D.4, to choose the fixed ordinary algebraic neighborhood
of a point of $Y$, and apply R, Theorem 2.9, there. Pull back its
smooth positive-curvature approximants along $\pi$. Smooth positive
$(1,1)$-forms retain positivity under original analytic pullback, so
their wedges with the arbitrary $T$ are positive. The scalar
approximants converge locally uniformly, and $T$ is coefficient-Radon.
Their distribution derivatives therefore converge to the current
in (PT2), which is positive. All tests on $V$ are used in this
passage, including those represented by moments not pulled back from
$Y$. If $D[v]=0$, apply the argument to $v$ and to $-v$; a current
and its negative cannot both be positive unless they are zero.
No background current, and no BT quotient, is pulled to a new field.
\end{proof}

\subsection{A necessary detector using product tests}

\begin{lemma}[Every-parameter product detection]
\label{lem:r12-every-parameter-product-detector}
For every $\rho\in C_c^\infty(J)$ and every original compact
smooth $(d-1,d-1)$-test $\eta$ on $Y$,
\[
 \bigl\langle\Omega_s\wedge D[h],\rho(s)\pi^*\eta\bigr\rangle
             =\int_J\rho(a)\,\langle D[H_a],\eta\rangle\,da.
 \tag{PT3}
\]
The function $a\mapsto\langle D[H_a],\eta\rangle$ is continuous.
Consequently
\[
 \Omega_s\wedge D[h]\ge0\quad\Longrightarrow\quad
                            D[H_a]\ge0\ \text{for every }a\in J,
 \tag{PT4}
\]
and the same implication holds for zero in place of positivity.
\end{lemma}

\noindent\textbf{Status: \texttt{NEW\_PROVED}} as the explicit
necessary detector and every-parameter upgrade. Its integration
input is the original calibrated product rule, not a restriction
of arbitrary currents to a Gauss fiber.

\begin{proof}
The support of $\rho(s)\pi^*\eta$ is compact in $U$: a compact
parameter box gives compact polyannuli, and the remaining support
lies over a compact subset of $Y$, covered by finitely many compact
analytic neighborhoods. The original product calibration gives,
for every compact top-degree original test $\beta$ on $Y$,
\[
 \langle[h],\rho(s)\Omega_s\wedge\pi^*\beta\rangle
       =\int_J\rho(a)\langle[H_a],\beta\rangle\,da.
 \tag{PT5}
\]
For clarity, this is checked on each finite moment presentation of
$\beta$: adjoining the $c$ independent torus coordinates takes its
calibrated tropical integration cells to their products with
the base parameter cells. The base coordinate lattice has index one,
and its original logarithmic volume form is $\Omega_s$. Ordinary
Fubini on each calibrated product cell gives (PT5). Finite smooth
localization on $Y$ and continuous coefficient multiplication prove
it for $H$. Both sides retain the original $K$ calibration, including
any finite constant-field multiplicities already in $[Y]$.

Since $\Omega_s$ is separately closed and every differential of
$\rho(s)$ is killed by its wedge, distributional integration by
parts in (PT5), with $\beta=D\eta$, gives exactly (PT3).
Continuity of $H$ is uniform in $a$ on the compact support of
$D\eta$. The finite coefficient measure $[Y]\wedge D\eta$
therefore makes the scalar pairing in (PT3) continuous.

If $\eta$ is original middle-positive and $\rho\ge0$, its displayed
pullback test is positive. Positivity of the left side implies
nonnegativity of the scalar distribution in $a$ on the right. This
first yields nonnegativity almost everywhere. Continuity upgrades
it to every parameter: a negative value would persist on an open
parameter ball and would be detected by a nonnegative smooth
$\rho$ supported there. Apply this argument separately to each
original positive $\eta$. It proves (PT4) on the complete original
test sheaf of $Y$, with no common exceptional set or countable
test assumption. For zero, use arbitrary signed $\rho$ and arbitrary
$\eta$, or apply the positive argument to both signs.
\end{proof}

The product tests in this lemma are only necessary detectors. They
are not asserted to span all original tests on $U$. In particular,
an original unit on $U$ can depend genuinely on both $z$ and $y$;
such a unit need not come from either factor. Sufficiency in the
next theorem is supplied by the intrinsic all-test band construction,
not by a product-test density assertion.

\subsection{Full arbitrary-background transfer in this parameter class}

\begin{theorem}[Product-parameter middle-order transfer]
\label{thm:r12-product-parameter-order-transfer}
Under (PT1), suppose
\[
                           \Omega_s\wedge D[h]\ge0
 \tag{PT6}
\]
as an original current on $U$. Then for every allowed $s$-background
$R$ of every degree $(P,P)$ with $P+1\le n$,
\[
                                    D(hR)\ge0.
 \tag{PT7}
\]
If the current in (PT6) is zero, the current in (PT7) is zero.
In particular, for every arbitrary original middle-positive
separately closed $(p,p)$-background $S$, with $p+c+1\le n$,
\[
 \Omega_s\wedge D[h]\ge0
      \quad\Longrightarrow\quad \Omega_s\wedge D(hS)\ge0,
 \qquad
 \Omega_s\wedge D[h]=0
      \quad\Longrightarrow\quad \Omega_s\wedge D(hS)=0.
 \tag{PT8}
\]
These are ordinary-current statements on all original compact
smooth complementary tests and all actual moment refinements.
There is no matroidal, simplicial-polytopal, or special-current
hypothesis on $Y$ or $S$.
\end{theorem}

\noindent\textbf{Status: \texttt{NEW\_PROVED}} relative to the
exact original R/B inputs. The geometric product and the
parameter-family form of $h$ are essential printed hypotheses.

\begin{proof}
The detector gives $D[H_a]\ge0$ on the original $Y$ for every
$a\in J$. Lemma~\ref{lem:r12-continuous-original-map-transfer}
then gives
\[
             D\bigl((H_a\circ\pi)T\bigr)\ge0
\]
for every original allowed positive closed $T$ on $U$.
For the particular $s$-annihilated $R$, approximate $h$ uniformly
on a fixed compact test neighborhood by the finite parameter
step functions (PF6). Its band currents $R_{E_j}$ are positive
and genuinely separately closed. Thus
\[
                    D(h_mR)=\sum_j
                            D\bigl((H_{a_j}\circ\pi)R_{E_j}\bigr)
                                      \ge0.
 \tag{PT9}
\]
Uniform convergence and the coefficient-Radon property of $R$
give $D(h_mR)\to D(hR)$ on every original test. This proves (PT7).
It requires no simultaneous smooth regularization of infinitely
many parameter slices, and no assertion about BT products over
the completed residue field of a Gauss point.

If (PT6) is zero, apply the same proof to $H$ and to $-H$.
Both signed currents in (PT7) are then positive, hence zero.
For (PT8), take the derived background $R=\Omega_s\wedge S$.
It is allowed by (PB2), and smooth closed Leibniz identifies
$D(hR)$ with $\Omega_s\wedge D(hS)$. The ambient test in this
argument was arbitrary; the assertion is therefore not restricted
to product tests or to any selected skeleton.
\end{proof}

\begin{corollary}[A common-parameter actual relation sector]
\label{cor:r12-product-parameter-actual-relation}
Suppose in addition that $h$ is an actual continuous PS function,
locally a finite difference of permitted bounded classical psh
potentials, so that $D_\delta h$ and its prescribed action are
defined. If
\[
                             \Omega_s\wedge D_\delta h=0
 \tag{PT10}
\]
as an actual analytic delta-form, then for every allowed original
$S$ in the permitted degree its prescribed response
$\Omega_s\wedge D(hS)$ is zero on every original test.
The same conclusion compares two such $h$ when the corresponding
actual forms agree.
\end{corollary}

\noindent\textbf{Status: \texttt{DIRECT\_COROLLARY}}. This is an
actual relation sector with the displayed common smooth parameter
factor; it is not the full degree-two constant-cycle quotient.

\begin{proof}
Compatibility of the actual delta derivative with the original
current map sends (PT10) to the zero premise in (PT8). The
prescribed action in this sector is exactly the printed intrinsic
distribution. Apply the theorem. Subtract the two continuous
family functions for the equality statement.
\end{proof}

\paragraph{What the proof does not permit.}
For a general moment map, an actual zero form does not provide a
fixed original reference space $Y$ with the full parameter-family
description (PT1) and (PF1). In particular, replacing $Y$ here by
the complete residue-field fiber of a general map would change the
scope of R, Theorem 2.9; that substitution has not been made.
Likewise, closed Borel bands do not identify an arbitrary current
on $A_J\times_KY$ with a current on $Y$: their original tests may
contain units depending on both factors. The theorem supplies no
such disintegration or inverse pushforward assertion.

The common-factor result consequently leaves OPEN-GQ for arbitrary
finite actual PS/PL presentations and OPEN-CTO for general signed
constant tensors, especially lower real degree without the stated
product parameter structure. An ordinary equality proved here is
not equality on all GK delta-test forms. These limitations are
part of the theorem's scope, not omitted hypotheses.

% END REVIEWED MODULE research_r12/gq/product_parameter_order_transfer.tex


% BEGIN REVIEWED MODULE research_r12/gq/original_current_disintegration.tex
\section{Disintegration into original ambient currents}
\label{sec:r12-original-current-disintegration}

The closed parameter bands admit a genuine local disintegration,
but its conditional currents remain on the original ambient
$K$-analytic space. They are not currents on a newly introduced
residue-field analytic space. This distinction is essential for
using the frozen arbitrary-background BT theorem without changing
its field or its test sheaf.

Throughout this section $U$ is second countable in addition to the
standing original hypotheses. The conclusion is ordinary-local on
every space in the project: every point has an ordinary neighborhood
of this kind, as proved below. Let $s:U\to J\subset\mathbf R^c$ be
the original parameter map, and let $R$ be an allowed $s$-background.

\subsection{A countably cofinal system of original smooth tests}

\begin{lemma}[Local countability does not discard actual refinements]
\label{lem:r12-countably-cofinal-original-tests}
Every point of an allowed original analytic space has a second
countable ordinary neighborhood. On any such $U$, the original
compact smooth CLD test spaces in all bidegrees admit a countable
cofinal family of finite moment-and-cutoff presentation stages.
Each stage is a quotient of a separable smooth coefficient space,
has a countable dense subset, and has a countable dense subset in
its original middle-positive cone. All equality kernels between
the stages can likewise be checked on countably many dense
coefficient relations.

More precisely, every original compact test, including a test
initially represented after an arbitrary actual refinement, is a
finite sum of tests in these stages. Approximation within each
stage is in its usual smooth coefficient topology on one fixed
compact support. No assertion about all GK delta-tests is made.
\end{lemma}

\noindent\textbf{Status: \texttt{NEW\_PROVED}} for the exact
original-test cofinality used in the conditional-current argument.
The original smooth cutoff/partition and common-moment properties
are frozen \texttt{INTERNAL\_INPUT}s.

\begin{proof}
The finite extension $K/\mathbf Q_p$ is separable as a metric field.
Every $K$-affinoid algebra is a quotient of a generalized Tate
algebra in finitely many variables. Polynomials with coefficients
in a fixed countable dense subfield, followed by passage to the
quotient, give a countable dense subset of that Banach algebra.
This argument also applies to non-strict radii: finite polynomials
are dense in the defining weighted completion.

The Berkovich spectrum of a separable affinoid algebra is compact
metrizable. Indeed a countable dense set of functions determines
each bounded multiplicative seminorm, and all remaining evaluations
are uniform limits of these evaluations, since
$\bigl||f(x)|-|g(x)|\bigr|\le\|f-g\|$. The spectrum embeds
continuously and injectively in a countable product of compact
real intervals. Compactness and the Hausdorff property identify
its topology with that induced metric topology. By goodness, an
ordinary point has an affinoid neighborhood; its ordinary interior
is the required second countable neighborhood.

On a second countable good $U$, choose a countable collection of
affinoid neighborhoods whose ordinary interiors form a basis.
To construct it, for each member of a countable topological basis
cover its points by affinoid interiors contained in that member,
and choose a countable subcover. Choose a countable dense subset
$D_B$ of each resulting affinoid algebra $\mathcal O(B)$.
Whenever $g$ is an analytic unit on a neighborhood of $B$, choose
$g_0\in D_B$ with
\[
 \|g-g_0\|_B<\min\left\{\inf_{x\in B}|g(x)|,\|g^{-1}\|_B^{-1}\right\}.
 \tag{DI1}
\]
The right side is positive by compactness and invertibility. The
ultrametric inequality gives $|g_0|=|g|$ on $B$, and the Neumann
series applied to $1+g^{-1}(g_0-g)$ makes $g_0$ a unit there. Their actual
logarithmic moment functions and their original differentials
therefore agree exactly on the ordinary interior. Thus every finite
actual moment tuple on a compact local stage is replaceable, after
finite ordinary localization, by a tuple from a countable family.
This applies equally to arbitrary extra units in a refinement.

Fix countably many original smooth cutoffs for pairs of basis
neighborhoods with nested compact closures, equal to one on the
smaller compact and supported in the larger. The frozen cutoff
theorem supplies each such choice. Include the finitely many
moments used by each cutoff in the countable family; finite unions
and finite common refinements remain countable. A finite collection
of these cutoffs covers the support of any test. Normalizing their
sum on a neighborhood of that support yields a finite smooth
partition and reduces the test to compact stages in the stated
countable family.

For a fixed tuple and cutoff, write its coefficient arrays on
real coordinate neighborhoods. Use a countable exhaustion by
rational real balls and compact coordinate cores. The corresponding
$C^\infty$ coefficient spaces are separable metrizable spaces;
finite products and their quotients have the same property. The
actual zero-form conditions on a fixed moment image are closed
linear conditions, obtained by restrictions to its relevant tangent
spaces. The smooth coefficient relations comparing two stages can
be imposed on their countable common refinements; their kernels
are separable subspaces. Choose a countable dense subset in each
such kernel. The preimage of the middle-positive cone is likewise
a subspace subset of a second countable metric coefficient space,
and hence is separable; choose dense positive representatives.
These choices prove all countability assertions. No rationality of
the real coefficients of a test is required: only density within
its fixed smooth presentation is used.
\end{proof}

\begin{lemma}[A recalled Radon disintegration construction]
\label{lem:r12-Radon-conditional-measures}
Let $M$ be a locally compact second countable Hausdorff space,
$s:M\to J$ continuous, and $\mu$ a finite positive Radon measure.
Put $\nu=s_*\mu$. There is a measurable family of positive Radon
probability measures $\mu_a$, for $\nu$-almost every $a$, carried
by $s^{-1}(a)$, such that for every nonnegative Borel $F$,
\[
                       \int_M F\,d\mu
                          =\int_J\left(\int_M F\,d\mu_a\right)d\nu(a).
 \tag{DI2}
\]
The identity also holds after insertion of any bounded Borel
parameter weight. Integrable signed functions are allowed.
\end{lemma}

\noindent\textbf{Status: \texttt{EXTERNAL\_KNOWN}} as the usual
Radon conditional-measure theorem. The construction below records
the countable version needed here, using Radon--Nikodym and Riesz
representation only.

\begin{proof}
Take a countable dense rational vector lattice in $C_0(M)$, enlarged
to include a compact approximate unit $\chi_m\uparrow1$, and products
with pullbacks of a countable determining algebra in $C_c(J)$.
For each chosen $f$, the signed measure
$E\mapsto\int_{s^{-1}(E)}f\,d\mu$ is absolutely continuous with
respect to $\nu$. Choose its Radon--Nikodym derivative $g_f$.
Outside one null set the maps $f\mapsto g_f(a)$ are rational-linear,
positive, bounded by $\|f\|_\infty$, and satisfy
\[
                  g_{\phi(s)f}(a)=\phi(a)g_f(a).
\]
There are only countably many identities and inequalities to impose.
Extend by continuity to $C_0(M)$ and apply Riesz representation to
obtain $\mu_a$ of mass at most one. The approximate unit and
monotone convergence show that its mass equals one almost everywhere:
the increasing limit of $g_{\chi_m}$ has integral $\nu(J)$ and is
bounded above by one. The displayed multiplication identity then
gives $s_*\mu_a=\delta_a$, first on the determining continuous
algebra and then on Borel sets. Thus $\mu_a$ is carried by
$s^{-1}(a)$.

The derivatives $g_f$ supply measurability and (DI2) on the chosen
algebra. Uniform approximation extends these assertions to
$C_0(M)$; the monotone-class theorem extends them to nonnegative
Borel functions. Multiplication by a bounded Borel function of
$s$ gives the weighted version. The signed assertion follows from
positive and negative parts. No regularity of a conditional density
in $a$ is asserted.
\end{proof}

\subsection{Conditional currents retain the entire original test sheaf}

\begin{theorem}[Disintegration of a parameter-annihilated original current]
\label{thm:r12-original-current-disintegration}
Let $R$ be an allowed $s$-background on second countable $U$.
There is a finite positive Radon control measure $\nu$ on $J$ and
a measurable family of original currents $R_a$ on the same $U$
such that, for $\nu$-almost every $a$,
\begin{enumerate}
\item $R_a$ is symmetric, middle-positive, separately $d'$- and
$d''$-closed, coefficient-Radon, and $s$-annihilated;
\item its coefficient measures are carried by $s^{-1}(a)$;
\item for every Borel $E\subset J$ and every original compact
smooth test $\eta$,
\[
               \langle R_E,\eta\rangle
                    =\int_E\langle R_a,\eta\rangle\,d\nu(a).
 \tag{DI3}
\]
\end{enumerate}
Local coefficient variations of the family are $\nu$-integrable.
The exceptional set can be chosen once for all original tests and
all actual moment refinements. For the chosen control measure $\nu$,
the family is unique up to one $\nu$-null set as original currents.
At every parameter atom of positive $\nu$-mass,
\[
                       R_a=\nu(\{a\})^{-1}R_{\{a\}}.
 \tag{DI4}
\]
The normalization measure is not claimed to be canonical; the
closed-band current-valued measure $E\mapsto R_E$ is canonical.
\end{theorem}

\noindent\textbf{Status: \texttt{NEW\_PROVED}} for original
all-test currents. This is not an auxiliary-field current theorem.

\begin{proof}
Choose the countably cofinal stages of
Lemma~\ref{lem:r12-countably-cofinal-original-tests}. On each stage,
use a finite constant exterior coefficient frame with a fixed
compact cutoff. The coefficient-Radon property gives finitely many
finite signed Radon measures representing all its test pairings.
List the resulting countable collection as $(\mu_j)_{j\ge1}$,
including finitely many positive domination controls on every stage.
The measure
\[
            \mu=\sum_{j\ge1}
                    \frac{2^{-j}}{1+\|\mu_j\|_{\mathrm{TV}}}|\mu_j|
 \tag{DI5}
\]
is finite positive Radon and dominates every listed coefficient
measure by a fixed constant depending on its index. Thus every
fixed compact coefficient frame is bounded by $C\mu$, not merely
absolutely continuous without quantitative control.
If $R=0$ the theorem is trivial; otherwise discard any
zero summands and put $\nu=s_*\mu$.

Write $\mu_j=g_j\mu$ using Borel versions of their Radon--Nikodym
densities, and disintegrate $\mu$ by the preceding lemma. Outside
one $\nu$-null set all $\int|g_j|\,d\mu_a$ are finite. Put
$\mu_{j,a}=g_j\mu_a$. On a given stage these finitely many signed
Radon measures define an order-zero functional on its smooth
coefficient arrays.

We check that these functionals give an original current, rather
than independent unrelated functions on the chosen stages. For
each of the countably many dense zero-form relations from the
cofinality lemma, its original coefficient measure is zero, hence
its conditional pairing is zero almost everywhere. Impose all
these equalities outside one null set. Each conditional stage
functional is continuous for the smooth coefficient topology,
by its finite coefficient-variation bound, so it kills the entire
closed relation kernel. The same reasoning on common stages
identifies any two representations. Consequently the functionals
extend to one coefficient-Radon original current $R_a$ on all
compact original smooth tests. The measure construction proves
(DI3), first on the countable stages and then on every test.
Its finite stage variations are integrable because
\[
                 \int_J\int|g_j|\,d\mu_a\,d\nu
                                           =\int|g_j|\,d\mu<\infty.
 \tag{DI6}
\]

For every positive test in the chosen countable dense family,
positivity of every $R_E$ in (DI3) forces
$\langle R_a,\eta\rangle\ge0$ almost everywhere. Impose these
countably many inequalities simultaneously; continuity on each
stage gives positivity on its full middle-positive cone. Thus
$R_a$ is middle-positive, not merely positive on decomposable or
strongly positive tests. The same countable argument applied to
the symmetry identities gives symmetry.

Every $R_E$ is separately closed by the new closed-band theorem.
For each test $\alpha$ in a countable dense source stage,
(DI3) applied to $d'\alpha$ gives
\[
                      \int_E\langle R_a,d'\alpha\rangle\,d\nu=0
                           \qquad\text{for all Borel }E.
\]
Its integrand is therefore zero almost everywhere. Impose the
countably many prime and double-prime equations and extend by
continuity of the differentials on each source stage. This proves
separate closedness on all original tests. The identical argument
for $d's_i\wedge R_E$ and $d''s_i\wedge R_E$ proves (PB2) for
$R_a$. The support assertion follows from that of $\mu_a$.
Set $R_a=0$ on the one exceptional null set if a family indexed by
every parameter is desired; no potential function is changed.

Uniqueness follows from (DI3) on a countable determining family:
two families have equal pairings there almost everywhere, and
continuity then identifies their original currents. Applying (DI3)
to $E=\{a\}$ proves (DI4). A null exceptional set cannot contain
an atom of positive $\nu$-mass. All refinements were included in
the exact cofinality lemma, so the conclusion is not a statement
only on the initial countable moment names.
\end{proof}

\paragraph{The scope of the conditional object.}
Each $R_a$ is an original $K$-current on $U$, tested by the entire
ambient original sheaf, including moments with genuine mixed
dependence on any parameter and fiber coordinates. The theorem
does not identify it with a current on $Y$, on a norm-level
polyhedron, or on a completed-residue-field analytic fiber.
It asserts no inverse to a geometric pushforward. Its use in the
next section is legitimate precisely because the frozen BT theorem
already permits every original positive separately closed current
on $U$, including these conditional ones.

% END REVIEWED MODULE research_r12/gq/original_current_disintegration.tex


% BEGIN REVIEWED MODULE research_r12/gq/bounded_borel_family_bt.tex
\section{Bounded Borel parameter families and joint decreasing continuity}
\label{sec:r12-bounded-borel-family-bt}

The conditional currents of the preceding section allow actual
bounded Borel parameter dependence, including jumps and atoms.
The potentials in each fixed parameter slice remain in the frozen
classical psh class over the original $K$. No positivity or quotient
theorem over a completed residue field is invoked.

Let $U,s,J,\pi:U\to Y$ be as in the continuous-family construction,
and let $R$ be any allowed $s$-background of degree $(P,P)$.
All statements are ordinary-local; first work on a second countable
ordinary neighborhood and choose the disintegration (DI3).

\begin{definition}[Locally bounded Borel slice families]
\label{def:r12-bounded-Borel-slice-families}
A slice family is a jointly Borel function
$H:J\times Y\to\mathbf R$ such that:
\begin{enumerate}
\item each $H_a$ is an actual locally bounded original classical
psh function on $Y$;
\item for every compact $Q\subset J$ and compact $L\subset Y$,
the actual bound $\sup_{Q\times L}|H|$ is finite.
\end{enumerate}
Its associated actual Borel function is
$h(x)=H(s(x),\pi(x))$. Functions are not modified on null sets,
and the second condition is a pointwise bound, not an essential
supremum. Let $\mathcal P^{\mathrm{Bor}}_{s,\pi}$ denote the cone
of such associated functions, locally in the ordinary topology.
\end{definition}

The ambient function $h$ need not be upper semicontinuous or psh.
Arbitrary locally bounded Borel functions of $s$ alone belong to
this cone, because their fixed slices are constants. Their
curvature on an $s$-annihilated background is zero by (PB5).
This enlargement is therefore relative to the conormal current;
it is not a redefinition of the classical psh class on $U$.

\begin{lemma}[Measurable bounded BT families on original conditional currents]
\label{lem:r12-measurable-original-BT-families}
For $H_1,\ldots,H_k$ as above, put
$u_{i,a}=H_i(a,\pi(\cdot))$ on the original $U$. If $P+k\le n$,
the original currents
\[
                  Q_{a,k}=T_{R_a,k}(u_{1,a},\ldots,u_{k,a})
 \tag{BF1}
\]
form a measurable coefficient-Radon family, and their local
coefficient variations are integrable in $a$. They are positive,
separately closed, $s$-annihilated, and carried by $s^{-1}(a)$.
The same measurability and local integrability hold after a bounded
jointly Borel scalar multiplier $G(a,x)$ on compact parameter and
ambient stages.
\end{lemma}

\noindent\textbf{Status: \texttt{NEW\_PROVED}} relative to the
original conditional-current theorem and the frozen bounded BT
theorem with relative CLN controls.

\begin{proof}
All $R_a$ outside the one fixed exceptional set are allowed original
backgrounds on $U$. Set them to zero on that exceptional set. The
functions $u_{i,a}$ are original classical psh by analytic pullback,
so every current in (BF1) is a frozen product on the original
ambient space. Its conormality follows by (PB8); its measure carrier
remains in $s^{-1}(a)$ by ordinary support locality of the recursion.

The coefficient measures of $R_a$ form measurable kernels by their
construction from $g_j\mu_a$. Assume this has been proved for the
tail of length $k-1$. For every fixed compact original test $\eta$,
the recursion reads
\[
                \langle Q_{a,k},\eta\rangle
                  =\langle u_{1,a}Q_{a,k-1}^{\mathrm{tail}},D\eta\rangle.
 \tag{BF2}
\]
The multiplier $(a,x)\mapsto H_1(a,\pi(x))$ is jointly Borel.
Integration against the finitely many measurable coefficient kernels
of the tail therefore makes this pairing measurable. This assertion
also holds after multiplying a test by any original smooth scalar.
On a fixed compact stage, use a countable dense collection of
continuous scalar tests and the finite coefficient domination of
the positive current $Q_{a,k}$. The usual monotone-class argument
then gives measurable Radon kernels for every coefficient measure,
including their integrals against bounded jointly Borel functions.
This establishes the induction step; it is not an appeal to weak
convergence against an arbitrary discontinuous scalar.

For integrability, choose the fixed enlarged compact stage and its
finite positive controls in the frozen relative CLN theorem. Only
parameters in a compact neighborhood of the $s$-image of that stage
can contribute, because each $R_a$ is carried by its level. The
uniform slice bounds in the definition give numbers $M_i$ independent
of $a$ on the corresponding compact product. The CLN estimate gives
\[
 \langle Q_{a,k},\Theta\rangle
       \le C\prod_{i=1}^k(1+M_i)
                                      \sum_\ell\langle R_a,\Xi_\ell\rangle.
 \tag{BF3}
\]
The right side is integrable by (DI3). Original finite middle-matrix
domination bounds all other coefficient variations by finitely many
such controls. A bounded multiplier adds only its compact bound.
This proves all local integrability assertions. No uniform bound
on the conditional currents themselves is imposed as a hypothesis.
\end{proof}

\begin{theorem}[The bounded Borel parameter-family BT calculus]
\label{thm:r12-bounded-Borel-family-BT}
Let $h_1,\ldots,h_k\in\mathcal P^{\mathrm{Bor}}_{s,\pi}$ and
$P+k\le n$. Then the intrinsic recursion
\[
 \mathcal T^s_{R,0}=R,\qquad
 \mathcal T^s_{R,k}(h_1,\ldots,h_k)
          =D\bigl(h_1\mathcal T^s_{R,k-1}(h_2,\ldots,h_k)\bigr)
 \tag{BF4}
\]
is well-defined using the actual Borel values of all potentials.
It is represented on every original test by
\[
               \mathcal T^s_{R,k}(\mathbf h)
                      =\int_J T_{R_a,k}(\mathbf u_a)\,d\nu(a).
 \tag{BF5}
\]
Its currents are symmetric, middle-positive, separately closed,
coefficient-Radon, and $s$-annihilated. The differentiated slots
are mixed symmetric and multilinear on the cone. The products are
independent of the chosen parameter-family descriptions and of the
control measure/disintegration, commute with ordinary restriction
and gluing, satisfy the same local CLN controls (BF3) after
integration, and commute with every locally bounded Borel parameter
coefficient as in (PF5).

Suppose for $0\le i\le k$ that jointly Borel slice families
$H_{i,m}$ decrease pointwise on $J\times Y$ to slice families
$H_i$ satisfying the same local boundedness conditions. Let
$h_{i,m},h_i$ be the associated actual functions. Then the full
weighted joint decreasing continuity holds:
\[
 h_{0,m}\,\mathcal T^s_{R,k}(h_{1,m},\ldots,h_{k,m})
          \longrightarrow
 h_0\,\mathcal T^s_{R,k}(h_1,\ldots,h_k)
 \tag{BF6}
\]
against every original compact smooth complementary test. In top
degree this is local weak convergence of the corresponding Radon
measures against every continuous compact scalar test. The
undifferentiated weighted current need not be positive or closed.
\end{theorem}

\noindent\textbf{Status: \texttt{NEW\_PROVED}} for these relative
families. This is a full bounded-family theorem over general
original conormal backgrounds, not the unrestricted analytic
PS/PL relation quotient.

\begin{proof}[The intrinsic formula and its properties]
The measurable-family lemma makes the right side of (BF5) an
original coefficient-Radon current, with all the stated positivity,
closedness, conormality and mixed symmetry inherited from the
conditional frozen products. We identify it with (BF4) by induction.
On the carrier of $R_a$ and of every conditional tail,
\[
                       h_i(x)=H_i(a,\pi(x))=u_{i,a}(x)
 \tag{BF7}
\]
as an equality of actual Borel functions. The coefficient-kernel
Fubini theorem and (BF2) therefore give, for any original test $\eta$,
\[
 \begin{split}
 \left\langle D\left(h_1\int_JQ^{\mathrm{tail}}_{a,k-1}d\nu(a)\right),
                                                           \eta\right\rangle
   &=\int_J\langle h_1Q^{\mathrm{tail}}_{a,k-1},D\eta\rangle\,d\nu(a)\\
   &=\int_J\langle u_{1,a}Q^{\mathrm{tail}}_{a,k-1},D\eta\rangle\,d\nu(a)\\
   &=\int_J\langle Q_{a,k},\eta\rangle\,d\nu(a).
 \end{split}
 \tag{BF8}
\]
All integrals are absolutely controlled on this compact stage by
the lemma. This proves the recursive identity and in particular
shows that the result depends only on the actual functions and $R$.
Hence it is independent of the auxiliary scalar control measure,
the conditional normalization, and the family descriptions.

Symmetry gives multilinearity in every differentiated slot from
the intrinsic first-slot linearity. It also makes signed-difference
extensions independent of their descriptions whenever such an
extension is used. Ordinary locality follows from the recursion;
therefore constructions on second countable neighborhoods agree
on overlaps and give the stated result in the full original
analytic category. Integrating (BF3) gives relative CLN with no
parameter derivatives or partition constants. The Borel parameter
coefficient rule follows directly from (PB5) and the recursion,
or from (DI3) and (BF5).
\end{proof}

\begin{proof}[Weighted joint decreasing continuity]
For every good parameter $a$, the actual original functions
$u_{i,m,a}=H_{i,m}(a,\pi(\cdot))$ decrease to the locally bounded
classical $u_{i,a}$ on the same original $U$. Apply the frozen
weighted joint bounded BT theorem to the fixed original background
$R_a$. It gives the required convergence of each conditional
pairing against any fixed original test $\eta$.

On the compact stage supporting that test, the endpoint inequalities
$H_i\le H_{i,m}\le H_{i,1}$ give one uniform bound on the entire
compact parameter/reference-space product. Use (BF3), finite
coefficient domination, and the bound for $H_{0,m}$ to dominate
every conditional pairing by
\[
                         C_\eta\sum_\ell\langle R_a,\Xi_\ell\rangle,
 \tag{BF9}
\]
an integrable function independent of $m$. Dominated convergence
in (BF5), with its weighted coefficient-kernel version, proves
(BF6). No exceptional parameter is ignored without its measure
being zero: the disintegration exceptional set is fixed once,
and every positive-mass atom has its original band current (DI4).
All actual potential values remain unchanged.

For top-degree measures, their compact total variations have the
same uniform bounds. Original smooth scalar tests are uniformly
dense in continuous compact scalars on a slightly enlarged compact
stage, by the frozen cutoff/density theorem. The usual uniform
approximation and variation estimate extends (BF6) to those tests.
This is not a convergence assertion against arbitrary Borel scalar
tests. Every step has used original ambient tests, so it remains
valid after every actual moment refinement.
\end{proof}

\begin{corollary}[Every arbitrary background supplies the relative construction]
\label{cor:r12-bounded-Borel-family-arbitrary-S}
For every original arbitrary allowed $(p,p)$-background $S$ and
$p+c+k\le n$, set $R=\Omega_s\wedge S$. Then (BF4)--(BF6)
construct the mixed parameter-family currents over this derived
background, with all original tests, all actual refinements, and
the unchanged original $K$ and norm. The initial case $k=1$ is
$\Omega_s\wedge D(hS)$.
\end{corollary}

\noindent\textbf{Status: \texttt{DIRECT\_COROLLARY}}.

\begin{proof}
Apply (PB2) and the theorem. Smooth closed Leibniz proves the last
identity. This retains, rather than replaces, the original arbitrary
$S$ in the displayed formula.
\end{proof}

\paragraph{The unresolved actual-relation step.}
For each parameter the functions in (BF1) are already classical
psh functions on the original ambient $U$, obtained from specified
original slice data. A general actual relation on a complete Gauss
fiber does not yet supply such original ambient representatives,
nor an equality of their prescribed products. Disintegration alone
cannot manufacture that missing implication. In particular this
theorem must not be used to change an actual fiber relation into a
relation over every conditional current. Full OPEN-GQ and the
general lower-degree signed OPEN-CTO remain open. The new result is
the closed original conditional-current mechanism and its bounded
family calculus, not the old coarea formula for the special initial
current $\Omega_s\wedge S$.

% END REVIEWED MODULE research_r12/gq/bounded_borel_family_bt.tex


% BEGIN REVIEWED MODULE research_r12/gq/level_trace_calculus_and_relation_gate.tex
\section{Ordinary level traces and a precise relation-lifting gate}
\label{sec:r12-level-trace-calculus}

The conditional currents in Section~\ref{sec:r12-original-current-disintegration}
remain original currents on the ambient space. This permits a useful
extension of the family construction without constructing any analytic
space over a residue field. It also makes the remaining relation gate
more precise: a relation of analytic delta-forms must still supply
ordinary ambient witnesses, or an independent argument replacing them.

Keep the original $U,s,J$ of Section~\ref{sec:r12-closed-parameter-bands},
and write $F_a=s^{-1}(a)$. The set $F_a$ is a closed subset of $U$;
we do not give it a new analytic structure. In particular, every
``neighborhood'' in the following definition is an ordinary open
of the original ambient $K$-analytic space.

\begin{definition}[Ordinary psh and harmonic level traces]
\label{def:r12-ordinary-level-traces}
A locally bounded actual Borel function $h:U\to\mathbf R$ is an
\emph{ordinary psh level trace} for $s$ if, for every $a\in J$ and
every $x\in F_a$, there are an ordinary neighborhood $O$ of $x$ and
an original locally bounded classical psh function $u$ on $O$ such
that
\[
                            h|_{F_a\cap O}=u|_{F_a\cap O}
                            \quad\hbox{pointwise}.                 \tag{LT1}
\]
Let $\mathcal P_s^{\mathrm{tr}}(U)$ be this cone. The witnesses
$O,u$ may depend on $a,x$; neither a measurable choice of them nor
a common neighborhood for all parameters is required.

A locally bounded actual Borel function $h$ is an \emph{ordinary
harmonic level trace} if the witnesses in (LT1) may be chosen
finite continuous with $D[u]=0$ as an original current on $O$.
Denote this real vector space by $\mathcal H_s^{\mathrm{tr}}(U)$.
The frozen continuous potential quotient implies that $D(uQ)=0$
for every original allowed positive closed background $Q$ on $O$.
This last implication is an \texttt{INTERNAL\_INPUT}, not a new
regularization claim for arbitrary currents.
\end{definition}

The Borel slice-family class of
Definition~\ref{def:r12-bounded-Borel-slice-families} is contained in
$\mathcal P_s^{\mathrm{tr}}$: take $u=H_a\circ\pi$. Every locally
bounded Borel function of $s$ alone belongs to
$\mathcal H_s^{\mathrm{tr}}$, using constant witnesses. Neither
class is asserted to consist of ambient upper semicontinuous
functions. The definition is background-independent.

\begin{lemma}[A one-step estimate for positive curvature recursions]
\label{lem:r12-positive-recursion-one-step-control}
Let $T$ be an original positive separately closed coefficient-Radon
current, let $h$ be an actual locally bounded Borel scalar, and
suppose $Q=D(hT)$ is positive. On each fixed compact original test
stage there are finitely many positive coefficient controls
$\Xi_\ell$ on a fixed larger compact stage and a constant $C$,
independent of $T,h$, such that
\[
             \|Q\|_{K,\mathrm{coeff}}
               \le C\|h\|_{L}\sum_\ell\langle T,\Xi_\ell\rangle.
                                                                    \tag{LT2}
\]
Here $K\Subset L$ are the fixed compact stages, and the left side
is any specified finite coefficient-variation seminorm on $K$.
If the degree of $Q$ is top degree, this includes its mass on $K$.
\end{lemma}

\noindent\textbf{Status: \texttt{DIRECT\_COROLLARY}} of the
frozen finite middle-cone coefficient domination and the recursion.

\begin{proof}
Choose finitely many positive smooth compact controls $\Theta_j$
which dominate the desired coefficient variations of a positive
current on $K$. For each $j$, the recursion gives
\[
                 \langle Q,\Theta_j\rangle
                         =\langle hT,D\Theta_j\rangle.
\]
On the fixed enlarged moment stage, the finitely many coefficients
of $D\Theta_j$ are bounded. The frozen finite-dimensional
middle-matrix domination bounds the total variation of their
pairings with $T$ by $C\sum_\ell\langle T,\Xi_\ell\rangle$.
Multiplication by $h$ adds its actual bound on $L$. Summing the
finitely many controls proves (LT2). Positivity of $Q$ is essential
for reducing its coefficient variations to the controls.
\end{proof}

\begin{theorem}[BT calculus for ordinary level traces]
\label{thm:r12-level-trace-BT}
Let $R$ be any allowed $s$-background of degree $(P,P)$, and let
$h_1,\ldots,h_k\in\mathcal P_s^{\mathrm{tr}}(U)$, where $P+k\le n$.
The intrinsic actual-value recursion
\[
       T^{\mathrm{tr}}_{R,0}=R,\qquad
       T^{\mathrm{tr}}_{R,k}(h_1,\ldots,h_k)
           =D\bigl(h_1T^{\mathrm{tr}}_{R,k-1}(h_2,\ldots,h_k)\bigr)
                                                                    \tag{LT3}
\]
is a positive separately closed symmetric coefficient-Radon
original current. It is $s$-annihilated and symmetric in the slots.
It satisfies ordinary locality, addition and nonnegative scaling
in each slot, and relative CLN bounds using only the actual local
bounds of the $h_i$ on fixed enlarged compact stages.

For every $h\in\mathcal H_s^{\mathrm{tr}}(U)$ and every allowed
$s$-background $T$,
\[
                                  D(hT)=0.                         \tag{LT4}
\]
Consequently two psh level traces differing by a harmonic level
trace have the same insertion on every allowed $s$-background,
including every positive mixed tail in (LT3). All identities hold
on every original compact smooth test and all actual moment
refinements, over the original $K$ and its original norm.
\end{theorem}

\noindent\textbf{Status: \texttt{NEW\_PROVED}} relative to the
original conditional-current theorem and the frozen bounded BT
calculus. No unrestricted decreasing-continuity assertion for
$\mathcal P_s^{\mathrm{tr}}$ is included in this theorem.

\begin{proof}
First work on a second countable ordinary neighborhood and use
$R=\int R_a\,d\nu(a)$ from
Theorem~\ref{thm:r12-original-current-disintegration}. For an actual
locally bounded Borel $h$, the conditional formula
\[
       \langle D(hR),\eta\rangle
           =\int_J\langle hR_a,D\eta\rangle\,d\nu(a)                \tag{LT5}
\]
follows directly from the coefficient kernels and their local
integrability. No chosen witness in (LT1) enters this measurability
or Fubini step.

If $h\in\mathcal P_s^{\mathrm{tr}}$, then $D(hR_a)$ is positive.
Indeed near each point of $F_a$ it equals $D(uR_a)$ for one witness
$u$ in (LT1), because $R_a$ is carried by $F_a$ and the functions
agree there in their actual values. Frozen BT gives positivity.
Outside $F_a$ the current is zero. Ordinary localization proves
positivity on all of $U$. Thus (LT5) makes $D(hR)$ positive. A
positive original current is coefficient-Radon by the frozen
middle-cone coefficient theorem. Separate closedness follows from
the form $D(hR)$, and conormality follows from (PB8).

For precision, the conditional currents $D(hR_a)$ also have
integrable coefficient variations. Their pairings with each
fixed test are measurable by the right side of (LT5). Positivity
and a countable dense family on each of the cofinal original
stages promote these pairings to measurable coefficient kernels,
as in Lemma~\ref{lem:r12-measurable-original-BT-families}.
Estimate (LT2), with the actual bound of $h$ on the enlarged
compact, bounds their variations by integrable coefficient
controls of $R_a$. Thus they may be used as a conditional family
in the next insertion, without choosing a new null set separately
for each original test.

Inductively construct
\[
      Q_{a,0}=R_a,\qquad
      Q_{a,j}=D(h_{k-j+1}Q_{a,j-1}).
                                                                    \tag{LT6}
\]
They are positive closed conormal currents carried by $F_a$.
The preceding argument gives measurable coefficient kernels and
integrable local variations at every step. Their integrals are
exactly (LT3), by the recursive test formula. Iterating (LT2)
gives the asserted CLN estimates with actual local bounds; no
uniform bound for chosen ordinary witnesses is used.

To prove symmetry, fix $a$ outside the common exceptional set and
$x\in F_a$. Choose the finitely many witnesses for the $h_i$ and
intersect their ordinary neighborhoods. On that intersection,
every conditional tail in (LT6) is carried by $F_a$, so every
$h_i$ can be replaced by its classical witness $u_i$. The resulting
current is the frozen original product
$T_{R_a,k}(u_1,\ldots,u_k)$ and hence is symmetric. This is local
near every point of $F_a$; off $F_a$ both sides vanish. Integrating
the resulting original-current identity proves symmetry in (LT3).
Addition and scaling follow from the recursion and symmetry.

For a harmonic trace $h$, the same local replacement gives
$D(hR_a)=D(uR_a)=0$, using the frozen continuous potential quotient.
Formula (LT5) yields (LT4). Subtracting two insertions proves the
last assertion. The arguments are local and intrinsic in actual
functions and currents, so they glue on arbitrary allowed $U$.
At no point is the test sheaf replaced by tests on $F_a$ or by
tests on an auxiliary analytic fiber.
\end{proof}

\paragraph{A precise limitation on continuity.}
For one fixed allowed background $R$, bounded pointwise convergence
$h_m\to h$ gives $D(h_mR)\to D(hR)$ distributionally by dominated
convergence of its original coefficient measures. This observation
does not prove joint decreasing continuity for several moving
level-trace slots. The ordinary neighborhoods and psh witnesses in
(LT1) may depend on $m$ with no common ordinary neighborhood and
no monotone sequence of witnesses. The stronger hypotheses in the
Borel slice-family theorem supply exactly such original psh
sequences and remain in force for its decreasing-continuity result.
The general trace-class continuity problem is \texttt{OPEN} here.

\subsection{Carrier relations are stronger than ambient pointwise relations}

\begin{lemma}[A carrier version of the zero-product relation]
\label{lem:r12-carrier-zero-product}
Let $Q$ be an allowed original background whose coefficient measures
are carried by a closed set $F\subset U$. Let $f_1,\ldots,f_m$ be
continuous local bounded classical differences, and suppose
$\prod_i f_i=0$ pointwise on $F$. Then
\[
                        L_{f_1}\cdots L_{f_m}Q=0.                  \tag{LT7}
\]
The conclusion holds on each signed mixed tail over $Q$.
\end{lemma}

\noindent\textbf{Status: \texttt{DIRECT\_COROLLARY}} extending
the proof of frozen R6 Lemma~\path{r6-zero-function-product}
from $F=U$ to the actual coefficient carrier.

\begin{proof}
Every signed mixed tail is coefficient-Radon, closed, and carried
by $F$, by ordinary support locality. For $m=1$, $f_1Q=0$ by
actual coefficient multiplication. For $m>1$, put
$T=L_{f_2}\cdots L_{f_m}Q$. On the ordinary open
$G=\{f_1\ne0\}$, the remaining product vanishes on $F\cap G$.
Induction, applied to $Q|_G$, gives $T|_G=0$. Outside $G$ the
multiplier $f_1$ is zero, hence $f_1T=0$. Taking $D$ proves (LT7).
Finite positive-tail expansions give the signed assertion.
\end{proof}

For example, on a finite complete simplicial fan $\Delta$ with
Courant functions, let $C$ be a subfan containing $t(F)$. Frozen
R6 identifies the ideal
$\ker(\mathrm{PP}(\Delta)\to\mathrm{PP}(C))$ with its nonface
monomials. Lemma~\ref{lem:r12-carrier-zero-product} shows that this
whole ideal acts by zero on $Q$, even if $t(U)$ is not contained
in $|C|$. Linear functions whose pullbacks agree with a constant
on $F$ likewise act by zero. This is a carrier-sensitive sufficient
criterion and does not identify every actual analytic relation
with such a pointwise ideal.

\paragraph{The remaining non-tautological lifting question.}
The preceding theorems prove the following sufficient reduction.
After a fixed conormal disintegration, a finite actual PS/PL
presentation can be simplified, at almost every parameter, by:
\begin{enumerate}
\item carrier zero-product relations of continuous classical
differences as in (LT7);
\item replacement of an allowed potential by another whose
difference is an ordinary harmonic level trace as in (LT4);
\item the already frozen ambient real-cycle, PS coefficient,
symmetry, and finite signed multilinearity relations.
\end{enumerate}
If such a finite local simplification reduces it to zero on each
original conditional current, with locally integrable coefficient
bounds, then its action is zero after integration. The background
is never regularized. Parameter atoms are included by (DI4).

The missing implication is not Radon disintegration: it is whether
an actual analytic delta-form relation supplies these ordinary
carrier relations, harmonic trace witnesses, or a strictly stronger
proved original-test certificate. In particular, actual equality
on a norm fiber does not by itself produce an ordinary harmonic
extension of a function, or a piecewise-polynomial representative
in the displayed restriction ideal. These implications remain
\texttt{OPEN}. No assertion is made that the three listed types
generate the whole analytic relation kernel. The statement is
uniform in every original allowed background when its hypotheses
hold at every parameter; an almost-everywhere certificate relative
to one disintegration proves only the corresponding fixed-current
conclusion.

\paragraph{Exact dependencies.}
The original bounded BT recursion, signed symmetry, coefficient
domination, and continuous harmonic potential quotient are frozen
R4 inputs. The carrier zero-product proof is the stated extension
of R6 \path{real_cycle_stanley_reisner_action.tex}, SHA256
{\small\nolinkurl{48e9930cba841d394b33634bbc99970186f8755a9cdef67472155b2cd5dc7d20}}.
The required R6 archive and its successive R10--R9--R8--R7 parents
were verified against their frozen manifests; no old result is
relabelled as new.

% END REVIEWED MODULE research_r12/gq/level_trace_calculus_and_relation_gate.tex


% BEGIN REVIEWED MODULE research_r12/capacity/closed_band_capacity_regularization.tex
\section{Full-tester capacity regularity on closed parameter bands}
\label{sec:r12-closed-band-capacity}

Fix the original finite extension $K/\mathbf Q_p$ and its original
absolute value. Let $V$ be an ordinary open in a good, Hausdorff,
boundaryless, quasi-smooth, pure $n$-dimensional $K$-analytic space.
Let $W\Subset V$ be an \emph{ordinary open} set. No assumption is
made that $\partial W$ is null for any curvature measure. Choose
original analytic units $f_1,\ldots,f_c$ on $V$, and put
\[
 s=(\log|f_1|,\ldots,\log|f_c|):V\longrightarrow J\subset\mathbf R^c,
 \qquad \Omega_s=\bigwedge_{i=1}^c(d's_i\wedge d''s_i),
 \tag{CB1}
\]
where $J$ is a real open containing the image in question.

Let $R$ be an original symmetric middle-positive, separately
$d'$- and $d''$-closed $(P,P)$-current, where
$k=n-P\geq1$. Assume the precise conormal identities
\[
                    d's_i\wedge R=d''s_i\wedge R=0
                                  \qquad(1\leq i\leq c).
 \tag{CB2}
\]
All currents in this section act on the full original compact
smooth complementary test sheaf and allow every actual moment
refinement. The derived-current interpretation for an unrestricted
background $S$ is stated separately below; (CB2) is not imposed
on that arbitrary $S$ without proof.

Write $\mathcal A(V)$ for the \emph{full} family of actual
classical psh functions $a$ on $V$ with $0\leq a\leq1$.
For Borel $A\subset V$, set
\[
 \begin{aligned}
 C_R(A;V)&=\sup_{a\in\mathcal A(V)}\mu^R_a(A),
       &\mu^R_a&=(Da)^k\wedge R,\\
 C_R^*(A;V)&=\inf_{\substack{G\supset A\\G\text{ ordinary open in }V}}
                                                C_R(G;V).
 \end{aligned}
 \tag{CB3}
\]
These are the original raw and ordinary-open outer capacities
for the specified background and full tester domain. No countable,
fixed-moment, smooth, or prescribed-product tester subfamily has
replaced $\mathcal A(V)$.

\begin{lemma}[A fixed Radon control for every full-tester parameter image]
\label{lem:r12-band-uniform-Radon-control}
There is a finite positive Radon measure $\nu$ on $J$, depending
on $R,W,V,s$ but not on $a\in\mathcal A(V)$, such that
\[
 \mu^R_a\bigl(W\cap s^{-1}(E)\bigr)\leq\nu(E)
                      \qquad(E\subset J\text{ Borel}).
 \tag{CB4}
\]
The measure $\nu$ may be chosen as the $s$-pushforward of a finite
sum of compactly supported original positive coefficient traces
of $R$. In particular, no properness of $s$ and no auxiliary-field
fiber construction is required. The same measure, after enlarging
its geometric constant, controls all mixed products
$(Da_1)\wedge\cdots\wedge(Da_k)\wedge R$ with
$a_i\in\mathcal A(V)$.
\end{lemma}

\noindent\textbf{Status: \texttt{NEW\_PROVED}} for this uniform
full-tester domination on the parameter sigma algebra. Its inputs
are the proved closed-band identity (PB9) and the frozen
\texttt{INTERNAL\_INPUT} arbitrary-background relative CLN theorem.

\begin{proof}
The compact set $\overline W$ has a finite ordinary neighborhood
cover whose enlarged compact stages are contained in $V$. The
homogeneous relative CLN estimate on those stages gives a finite
family of compactly supported original positive trace forms
$\Xi_1,\ldots,\Xi_N$ of complementary degree $(k,k)$ and a
finite geometric constant $C$ such that, for every original
middle-positive separately closed background $Q$ of degree
$(P,P)$ and every normalized classical tuple,
\[
 \int_W (Da_1)\wedge\cdots\wedge(Da_k)\wedge Q
                  \leq C\sum_{l=1}^N\langle Q,\Xi_l\rangle.
 \tag{CB5}
\]
The forms and the constant are independent of the tuple and of
$Q$. To recall the exact scope of this frozen estimate, use an
original compact cutoff equal to one near $\overline W$, apply
the recursion and integration by parts, and bound each derivative
of a cutoff by the positive trace on the next compact stage.
Repeat through the $k$ insertions. The scalar bounds
$\|a_i\|_\infty\leq1$ leave only finitely many traces of the
background on the final enlarged compact. Original finite
middle-coefficient domination supplies these positive traces;
they can be taken as finite sums of nonnegative scalar multiples
of strongly positive trace powers. This is the same finite-control
version of relative CLN used in (BF3). It does not bound only a
preselected moment family of potentials.

For Borel $E\subset J$, (PB4) gives the positive separately
closed original background
$R_E=\mathbf1_{s^{-1}(E)}R$. Apply (CB5) to this $Q$ and
use the exact original identity (PB9). The left side is
\[
 \int_{W\cap s^{-1}(E)}(Da_1)\wedge\cdots\wedge(Da_k)\wedge R,
\]
while its right side is
\[
 C\sum_l\langle R_E,\Xi_l\rangle
       =\int_{s^{-1}(E)}d\tau,
 \qquad \tau=C\sum_l\Xi_l\wedge R.
 \tag{CB6}
\]
The original trace measure $\tau$ is positive and has compact
support in $V$, hence finite. Its continuous pushforward
$\nu=s_*\tau$ is a finite Radon measure on $J$. This proves
(CB4) and its mixed version. Closedness of $R_E$ is essential
for this application of CLN; an arbitrary Borel restriction of
$R$ would not justify it. Every step uses the full original
bounded calculus, so its identity and its estimate survive all
actual moment refinements.
\end{proof}

\begin{theorem}[Raw equals outer on every Borel parameter band]
\label{thm:r12-band-raw-outer-capacity}
Under (CB1)--(CB2), for every Borel $E\subset J$ one has
\[
 \boxed{\quad
 C_R\bigl(W\cap s^{-1}(E);V\bigr)
       =C_R^*\bigl(W\cap s^{-1}(E);V\bigr).
 \quad}
 \tag{CB7}
\]
The same assertion holds if ordinary outer regularization is
taken relative to $W$, with the same fixed full tester family
on $V$. It requires no finite-energy representative of a weak
curvature cluster, no vanishing Jensen defect, and no regularity
of the boundary of $E$ or $W$.

Put
\[
       c_{R,s,W}(E)=C_R\bigl(W\cap s^{-1}(E);V\bigr)
                                  \qquad(E\subset J\text{ Borel}).
 \tag{CB8}
\]
Then, for all Borel $E,F\subset J$,
\[
 |c_{R,s,W}(E)-c_{R,s,W}(F)|\leq\nu(E\mathbin\triangle F).
 \tag{CB9}
\]
In particular this finite parameter capacity is inner regular
by compact real sets, outer regular by real open sets, and
continuous both on increasing and on decreasing Borel sequences.
Its real-open outer extension is a finite Choquet capacity on
all subsets of $J$.
\end{theorem}

\noindent\textbf{Status: \texttt{NEW\_PROVED}} for full BT
capacity on these analytic saturated sets. The measure-theoretic
regularization following (CB4) is elementary. The distinct R11
result (EC12)--(EC14) already treated scalar images of \emph{energy}
currents by a coarea-density argument; it is not claimed anew.

\begin{proof}
Set $A=W\cap s^{-1}(E)$. Finite Radon outer regularity supplies
a real open $O\supset E$ with $\nu(O\setminus E)<\epsilon$.
The set $G=W\cap s^{-1}(O)$ is ordinary open in $V$ and
contains $A$. For every actual normalized tester,
\[
 \mu^R_a(G)
    \leq\mu^R_a(A)+\mu^R_a\bigl(W\cap s^{-1}(O\setminus E)\bigr)
    \leq\mu^R_a(A)+\epsilon.
 \tag{CB10}
\]
Take the supremum over the full tester family and then let
$\epsilon\downarrow0$. This gives $C_R^*(A;V)\leq C_R(A;V)$;
the reverse inequality is always monotonicity. This proves (CB7).

For a set $A\subset W$, the infimum over open neighborhoods
inside $W$ equals the ambient infimum in (CB3): intersect any
ambient neighborhood with the ordinary open $W$, and use
monotonicity of raw capacity. Thus the relative-outer formulation
has not silently replaced the requested ordinary-open topology.
In particular no mass on $\partial W$ has been discarded in a
limit. The open window $W$ itself is used in every neighborhood.

Subadditivity of the supremum capacity and (CB4) give (CB9).
Finite Radon regularity of $\nu$ then yields the two regularity
statements and continuity on decreasing Borel sequences. For
increasing Borel sequences continuity also follows immediately
by commuting the two suprema over the sequence index and tester.

For completeness, let $c^{\mathrm o}$ be the real-open outer
extension to arbitrary parameter sets. Every such set has a Borel
hull of the same outer value: intersect a sequence of open
neighborhoods whose capacities tend to the infimum and use the
already proved Borel outer regularity. If $A_j$ increase, take
such hulls $B_j$ and put $H_j=\bigcap_{l\geq j}B_l$. These
increase, contain $A_j$, and satisfy
$c(H_j)=c^{\mathrm o}(A_j)$. Borel continuity from below gives
continuity from below of $c^{\mathrm o}$. Continuity from above
on compact sets was proved by (CB9). These are the Choquet
axioms claimed here. This last assertion concerns parameter
subsets; it does not identify analytic and parameter outer
regularizations for every non-Borel inverse image.
\end{proof}

\begin{proposition}[A Borel-stable weak-star relaxation]
\label{prop:r12-band-weak-star-relaxation}
For $a\in\mathcal A(V)$ define the finite parameter measure
\[
                    \lambda_a=s_*(\mathbf1_W\mu^R_a).
 \tag{CB11}
\]
Then $0\leq\lambda_a\leq\nu$. If $g_a=d\lambda_a/d\nu$ and
$\mathcal H$ is the weak-star closure of the full density family
in $L^\infty(\nu)$, then $\mathcal H$ is compact,
$0\leq g\leq1$ for $g\in\mathcal H$, and
\[
                c_{R,s,W}(E)=\max_{g\in\mathcal H}\int_Eg\,d\nu
                                    \quad(E\subset J\text{ Borel}).
 \tag{CB12}
\]
Moreover, for Borel $b_j,b$ with $b_j\to b$ in $L^1(\nu)$,
\[
 \sup_{a\in\mathcal A(V)}
 \int_W |b_j(s)-b(s)|\,d\mu^R_a
                              \leq\|b_j-b\|_{L^1(\nu)}\longrightarrow0.
 \tag{CB13}
\]
Thus this parameter relaxation creates no gain on \emph{any}
Borel parameter set. It need not be the pushforward of an actual
tester, and it does not give a representation of a full analytic
weak current cluster.
\end{proposition}

\noindent\textbf{Status: \texttt{DIRECT\_COROLLARY}} of (CB4),
Radon--Nikodym and Banach--Alaoglu; the last two functional-analytic
facts are \texttt{EXTERNAL\_KNOWN}.

\begin{proof}
Domination gives the densities in the weak-star compact order
interval $[0,1]\subset L^\infty(\nu)$. For each Borel $E$,
its indicator is in $L^1(\nu)$; hence integration over $E$
is weak-star continuous. Passing to the closure therefore changes
neither its supremum nor any of its finite Borel linear pairings,
and compactness attains the supremum. This proves (CB12).
The same measure domination gives (CB13). Unlike ordinary weak
convergence of Radon measures, this specified topology tests the
possibly discontinuous indicator $\mathbf1_E$ itself.
\end{proof}

\begin{corollary}[An unrestricted-background mixed-slot theorem]
\label{cor:r12-arbitrary-S-band-capacity}
Let $S$ be \emph{any} original symmetric middle-positive,
separately closed $(p,p)$-current on $V$. Put $q=n-p$, and
choose any $1\leq c\leq q-1$ original unit parameters as in
(CB1). The derived current
\[
                             R=\Omega_s\wedge S
 \tag{CB14}
\]
satisfies (CB2), and all conclusions (CB4)--(CB13) hold for
its remaining degree $k=q-c\geq1$. Equivalently they hold
for the full original mixed-slot family
\[
                    (Da)^{q-c}\wedge\Omega_s\wedge S,
                                  \qquad a\in\mathcal A(V).
 \tag{CB15}
\]
No smoothness, model origin, algebraic-cycle representation, or
membership in an already completed-background class is imposed
on $S$. Degenerate moment ranks are allowed; if $R=0$, the
statements have their exact zero interpretation.
\end{corollary}

\noindent\textbf{Status: \texttt{DIRECT\_COROLLARY}} of the
new closed-band theorem and the original smooth-factor calculus.

\begin{proof}
The original smooth form $\Omega_s$ is strongly positive and
separately closed. Its wedge preserves middle positivity and
separate closedness of the arbitrary original $S$. Each
$d's_i$ and $d''s_i$ repeats in that wedge, proving (CB2).
Apply the preceding results to this derived current. The proof
does not replace $S$ by $R$ in any claim about unrestricted
original-$S$ capacity; it constructs and controls precisely the
mixed-slot family (CB15).
\end{proof}

\paragraph{Remaining gates and failed extensions.}
The new regularity class consists of arbitrary Borel parameter
bands in ordinary relatively compact \emph{open} windows. It is
not limited to intervals, finitely many levels, or completed
potential sublevels. Nevertheless the following passages are
not consequences of the theorem.
\begin{enumerate}
\item Replacing the open window by an arbitrary compact set $K$
would make $K\cap s^{-1}(O)$ non-open. The argument then lacks
control of every ordinary neighborhood of $K$. This is precisely
an additional compact-boundary gate, not an innocuous notational
change.
\item Dropping (CB2) destroys the separately closed background
$R_E$ used in (CB5). The frozen annular Borel-scalar failure
quoted after (PB9) shows that unrestricted Borel restriction
cannot be inserted into that proof. The proposed unrestricted
extension is a \texttt{FAILED\_ATTEMPT}, not a proof of raw
capacity regularity for arbitrary $S$.
\item The measure $\nu$ controls the parameter image of the
curvature family, not its complete analytic measures. Neither
one finite Radon measure dominating the entire full tester
family on $V$ nor equality of the energy and BT null ideals
has been established.
\item No measurable choice of independently optimized fiber
potentials is asserted. In particular (CB12) is not an equality
between capacity and an integral of fiber capacities. All
suprema in (CB3) and (CB8) retain actual globally admissible
testers on the fixed domain $V$.
\end{enumerate}
The remaining unrestricted raw/outer capacity and compact
representative questions, as well as OPEN-GQ and OPEN-CTO, remain
\texttt{OPEN}. This theorem supplies a new full-dimensional,
arbitrary-background mixed-slot mechanism without using any of
those conclusions as an input.

\paragraph{Exact source dependencies.}
The only new current-theoretic dependency is
\path{gq/closed_parameter_bands.tex}, specifically (PB4)--(PB9).
The finite-control form of the frozen relative CLN estimate is
also printed in \path{gq/bounded_borel_family_bt.tex}, (BF3).
The earlier scalar energy-capacity comparison is the frozen R11
\path{capacity/uniform_coarea_capacity.tex}, (EC12)--(EC14).
No old proof or old result status is rewritten here.

% END REVIEWED MODULE research_r12/capacity/closed_band_capacity_regularization.tex


% BEGIN REVIEWED MODULE research_r12/capacity/band_stable_tester_capacity.tex
\section{A band-stable tester class and an exact Radon selection formula}
\label{sec:r12-band-stable-capacity}

Retain all hypotheses (CB1)--(CB2), including the original
conormal background $R$, its remaining degree $k\geq1$, and
the ordinary open window $W\Subset V$. Let $\nu$ be the finite
parameter control measure in (CB4). The enlarged admissible class
defined here is explicit. It is \emph{not} identified with the
original full family $\mathcal A(V)$ of classical psh testers.
Nor is it identified with every jointly Borel slice-classical
family from the broader parameter-family calculus.

\begin{definition}[Countable band pastings of actual testers]
\label{def:r12-band-pasted-testers}
Let $\mathcal A_s^{\rm band}(V)$ be the family of actual Borel
functions $h:V\to[0,1]$ for which there exist a countable Borel
partition $(E_i)_{i\geq1}$ of $J$ and functions
$a_i\in\mathcal A(V)$ such that
\[
                         h(x)=a_i(x)\quad\text{if }s(x)\in E_i.
 \tag{BA1}
\]
Empty pieces are permitted. A finite pasting is a member with a
finite such representation. The domain of every $a_i$ remains
the fixed original $V$; independent functions merely defined on
an auxiliary fiber are not part of this definition.
\end{definition}

\begin{proposition}[Intrinsic positive products for band-pasted testers]
\label{prop:r12-band-pasted-products}
For $h_1,\ldots,h_j\in\mathcal A_s^{\rm band}(V)$ and
$0\leq j\leq k$, the original recursion
\[
 \widetilde T_{R,0}=R,\qquad
 \widetilde T_{R,j}(h_1,\ldots,h_j)
      =D\bigl(h_1\widetilde T_{R,j-1}(h_2,\ldots,h_j)\bigr)
 \tag{BA2}
\]
is well-defined and gives a symmetric middle-positive,
separately closed, $s$-annihilated coefficient-Radon current.
If $(E_l)$ is a common countable refinement of the partitions
in (BA1), and $h_i=a_{i,l}$ on $s^{-1}(E_l)$, then
\[
 \widetilde T_{R,j}(h_1,\ldots,h_j)
       =\sum_l T_{R_{E_l},j}(a_{1,l},\ldots,a_{j,l}).
 \tag{BA3}
\]
The series converges in local coefficient total variation.
The product is independent of every presentation of the
\emph{same actual Borel functions}. It agrees with the frozen
bounded BT product when the entries are actual classical psh.
Every identity holds on every original compact complementary
test and every actual moment refinement.
\end{proposition}

\noindent\textbf{Status: \texttt{DIRECT\_COROLLARY}} of the
new closed-band calculus (PB4)--(PB9) and the frozen relative
CLN estimates. This countably pasted subclass is also covered
by the broader Borel-family theorem where that theorem's
explicit parameter presentation is used.

\begin{proof}
Each $R_{E_l}$ is an allowed original positive closed background.
Thus every summand in (BA3) exists by the frozen bounded theorem
and is positive, separately closed and conormal. On any fixed
compact test stage, the normalized scalar bounds and relative
CLN control its coefficient variations by finitely many fixed
positive traces of $R_{E_l}$ on an enlarged compact stage.
Their sums are the corresponding traces of $R$, by the countable
closed-band decomposition. Hence (BA3) converges in local
coefficient variation and is an original positive order-zero
current. Separate closedness and conormality pass on original
tests.

Each tail summand is carried measure-theoretically by
$s^{-1}(E_l)$. Multiplication by the bounded actual Borel $h_1$
therefore agrees on that summand with multiplication by
$a_{1,l}$. Local absolute variation convergence allows this
bounded multiplication to commute with the sum, and distributional
$D$ then gives (BA2). Induction from $j=0$ proves the stated
intrinsic recursion. It proves independence of presentations,
because (BA2) uses the actual Borel values themselves rather than
an unidentified quotient by background coefficient-null sets.
The usual symmetry on each original bounded tuple gives symmetry
after summation. For one common classical tuple (PB9) sums back
to its frozen original product.
\end{proof}

Write $\widetilde\mu_h=\widetilde T_{R,k}(h^{[k]})$. For a
Borel $A\subset W$ define the band-stable raw capacity
\[
       C_R^{\rm band}(A;V,s)
           =\sup_{h\in\mathcal A_s^{\rm band}(V)}\widetilde\mu_h(A).
 \tag{BA4}
\]
The original capacity satisfies $C_R(A;V)\leq C_R^{\rm band}(A;V,s)$.
This inequality does not assert equality of the admissible classes
or equality of the capacities.

\begin{lemma}[A countable essential supremum for a dominated density family]
\label{lem:r12-countable-essential-supremum}
Let $(J,\nu)$ be a finite measure space, and let
$\mathcal G\subset L^\infty(\nu)$ consist of functions between
zero and one. There are $g_i\in\mathcal G$ such that
\[
                   g^*=\sup_i g_i
                       =\mathop{\rm ess\,sup}_{g\in\mathcal G}g
                                   \quad\nu\text{-almost everywhere}.
 \tag{BA5}
\]
For real Borel $J$ and a finite Radon $\nu$, the selected functions
and their finite-maximum argmax partitions can be chosen Borel.
\end{lemma}

\noindent\textbf{Status: \texttt{EXTERNAL\_KNOWN}} as finite
measure-lattice theory; the following proof records the exact
selection used and avoids an uncountable pointwise supremum.

\begin{proof}
Put $L=\sup_{F\subset\mathcal G\text{ finite}}\int\max_{g\in F}g\,d\nu$.
Choose finite families whose integrals tend to $L$, and enumerate
their union as $(g_i)$. The increasing finite maxima tend to a
function $g^*$ with integral $L$. For any fixed $g\in\mathcal G$,
the finite-maximum definition of $L$ and monotone convergence give
$\int\max(g,g^*)\,d\nu\leq L=\int g^*\,d\nu$.
Consequently $g\leq g^*$ almost everywhere. Conversely every
essential upper bound of $\mathcal G$ bounds this countable
supremum, so it is the essential supremum. On a real Radon space
each selected measurable density has a Borel version. The sets
on which the first index attaining a finite maximum is $i$ are
then Borel and form the required partition.
\end{proof}

\begin{theorem}[Exact band selection and unchanged raw null sets]
\label{thm:r12-band-selection-capacity}
For each Borel $A\subset W$ and each actual tester
$a\in\mathcal A(V)$, form the parameter measure
\[
          \lambda_a^A=s_*(\mathbf1_A\mu_a^R),
          \qquad g_a^A=\frac{d\lambda_a^A}{d\nu},
          \qquad
          g_A^*=\mathop{\rm ess\,sup}_{a\in\mathcal A(V)}g_a^A.
 \tag{BA6}
\]
The densities exist and lie between zero and one. The enlarged
capacity has the exact formula
\[
 \boxed{\qquad
             C_R^{\rm band}(A;V,s)=\int_J g_A^*\,d\nu.
 \qquad}
 \tag{BA7}
\]
Already finite band pastings have the supremum on the left.
The value is independent of the choice of dominating parameter
measure $\nu$.

Moreover, on every such Borel $A$,
\[
 C_R(A;V)=0
       \quad\Longleftrightarrow\quad
 C_R^{\rm band}(A;V,s)=0.
 \tag{BA8}
\]
This is an equality of the two \emph{raw Borel null ideals}.
It is not an equality of their outer capacities, and it does
not identify either null ideal with a previously defined energy
null ideal.
\end{theorem}

\noindent\textbf{Status: \texttt{NEW\_PROVED}} for the exact
selection formula on the explicit band-stable original-current
tester class and for (BA8). No fiberwise admissibility theorem
is assumed in this statement.

\begin{proof}
Since $A\subset W$, (CB4) gives $0\leq\lambda_a^A\leq\nu$.
Thus (BA6) is well-defined as an equivalence class in
$L^\infty(\nu)$; no simultaneous pointwise choice for the full
uncountable tester family is necessary.

For a countable pasting $h$ with pieces $E_i$ and actual entries
$a_i$, the top identity (BA3) gives
\[
 \widetilde\mu_h(A)
    =\sum_i\mu^R_{a_i}(A\cap s^{-1}(E_i))
    =\sum_i\int_{E_i}g_{a_i}^A\,d\nu
    \leq\int_Jg_A^*\,d\nu.
 \tag{BA9}
\]
The final inequality needs only the countably many entries in
this one pasting; the union of their exceptional null sets is
still null. This proves the upper bound in (BA7).

Apply Lemma~\ref{lem:r12-countable-essential-supremum} to choose
actual testers $a_i$ with $g_A^*=\sup_i g_{a_i}^A$ almost
everywhere. For each $N$, take the finite Borel partition
selecting the first argmax among $g_{a_1}^A,\ldots,g_{a_N}^A$
and paste the corresponding original $a_i$. Formula (BA9) for
this $h_N$ is the exact equality
\[
                   \widetilde\mu_{h_N}(A)
                        =\int_J\max_{1\leq i\leq N}g_{a_i}^A\,d\nu.
 \tag{BA10}
\]
Monotone convergence gives the lower bound and proves that finite
pastings suffice. The intrinsic definition on the left of (BA7)
also proves independence of the auxiliary dominating measure.

If the original raw capacity of $A$ is zero, every nonnegative
density $g_a^A$ has integral zero. Hence zero is an essential
upper bound for the family, and $g_A^*=0$ as a measure-algebra
class. Equivalently, the countable family selected above consists
of zero classes. Formula (BA7) gives zero band capacity. The
reverse implication follows from inclusion of original testers.
This proves (BA8) without asserting a common pointwise exceptional
set for every actual tester on every conditional fiber.
\end{proof}

\begin{corollary}[The band envelope is a measure on saturated sets]
\label{cor:r12-band-capacity-measure-envelope}
Put $g_W^*=\mathop{\rm ess\,sup}_a g_a^W$ and
\[
                              \Lambda_W=g_W^*\nu.
 \tag{BA11}
\]
This finite positive Radon measure is intrinsic, and for every
Borel $E\subset J$,
\[
         C_R^{\rm band}\bigl(W\cap s^{-1}(E);V,s\bigr)
                                      =\Lambda_W(E).
 \tag{BA12}
\]
It is the least finite positive measure in ordinary measure order dominating
all $s_*(\mathbf1_W\mu^R_a)$ with $a\in\mathcal A(V)$.
The ordinary-open outer band capacity has the same value on
these saturated sets. The measure $\Lambda_W$ is not required
to be the curvature image of one admissible tester.
\end{corollary}

\noindent\textbf{Status: \texttt{DIRECT\_COROLLARY}}.

\begin{proof}
For $A=W\cap s^{-1}(E)$, its density family is
$g_a^A=\mathbf1_E g_a^W$, hence
$g_A^*=\mathbf1_Eg_W^*$. This proves (BA12) and its
intrinsic character. To check the least-measure assertion, let
$\zeta$ be any finite positive measure dominating all these
parameter measures. Relative to $\theta=\nu+\zeta$, domination
says each density is at most $d\zeta/d\theta$. Their countable
essential supremum, which represents $\Lambda_W$, has the same
upper bound. Thus $\Lambda_W\leq\zeta$.

For real open $O\supset E$, the corresponding open analytic
band $W\cap s^{-1}(O)$ has band capacity $\Lambda_W(O)$.
Radon outer regularity gives the upper outer bound; monotonicity
gives the reverse. As in (CB7), ordinary relative and ambient
outer neighborhoods agree because $W$ is open. No boundary-null
condition or actual maximizing tester is used.
\end{proof}

\paragraph{Exact interpretation and remaining minimal selection gate.}
Equations (BA7) and (BA12) deliberately use a Radon--Nikodym
\emph{essential} supremum. If an original-current disintegration
$R=\int R_t\,d\nu_0(t)$ is also chosen, then for each fixed
actual tester the relevant density can be expressed by its
conditional curvature, after the appropriate change of control
measure. That does not turn (BA7) into an integral of a
simultaneous raw pointwise fiber capacity: the exceptional set
for a conditional identity may depend on an uncountable tester.
No common full-test fiber supremum is asserted.

The finite argmax proof supplies exactly the admissibility and
measurable selection needed for the \emph{defined} band-pasting
class. It does not show that the pasted Borel potential has an
ambient classical psh representative. Upgrading (BA7) to an
equality for $C_R(A;V)$ would require such a representative or
another theorem showing no gain under these parameter-dependent
switches. That extra gate remains \texttt{OPEN}. The theorem
also does not address an arbitrary Borel set's ordinary outer
regularity; outside saturated sets, neither raw capacity has
been equated to its outer regularization.

For every unrestricted original $S$, the construction applies
to the mixed-slot background $R=\Omega_s\wedge S$ of (CB14).
This is a new calculus and capacity mechanism built from the
same $S$, not a replacement of the arbitrary background in
OPEN-GQ or in the unrestricted classical BT target. All field,
absolute-value and original-test requirements are retained.

% END REVIEWED MODULE research_r12/capacity/band_stable_tester_capacity.tex


% BEGIN REVIEWED MODULE research_r12/capacity/band_quasicontinuity_and_continuity.tex
\section{Quasicontinuity and intrinsic decreasing continuity for band testers}
\label{sec:r12-band-quasicontinuity}

Retain (CB1)--(CB2) and the precise band-pasting class
$\mathcal A_s^{\rm band}(V)$ of (BA1). In this section the same
formula (BA4) is used on all Borel subsets of $V$, with value
possibly $+\infty$ outside compact subsets. Write
\[
 c(A)=C_R^{\rm band}(A;V,s),\qquad
 c^*(A)=\inf_{\substack{G\supset A\\G\text{ ordinary open in }V}}c(G).
 \tag{BQ1}
\]
Both retain the complete band-pasting tester family on the fixed
domain $V$. Local relative CLN makes $c$ finite on compact sets.
It is monotone and countably subadditive, being a supremum of
positive Radon measures. Its outer regularization has the same
properties. No equality $c=c^*$ on all Borel sets is assumed.

Let $\mathcal A_s^{\rm band,ct}(V)$ be the subclass of (BA1)
which has a presentation whose actual classical branches $a_i$
are \emph{continuous on $V$}. The resulting Borel function $h$
need not itself be ordinarily continuous, and need not be ambient
classical psh. Its positive original products have already been
constructed by (BA2)--(BA3).

\begin{theorem}[Two precise quasicontinuity assertions]
\label{thm:r12-band-quasicontinuity}
Every $h\in\mathcal A_s^{\rm band}(V)$ is locally
ordinary-open outer $C_R$-quasicontinuous, where $C_R$ uses
the full original classical testers on $V$.

Every $h\in\mathcal A_s^{\rm band,ct}(V)$ is locally
ordinary-open outer $c$-quasicontinuous for its own enlarged
band tester family. More explicitly, for any compact $H\Subset V$
and $\epsilon>0$ there is an ordinary open $G\subset V$ such
that
\[
                     c(G)<\epsilon,
                  \qquad h|_{H\setminus G}\text{ is continuous}.
 \tag{BQ2}
\]
The open exceptional set can be chosen simultaneously for any
countable collection of continuous-branch band testers.
\end{theorem}

\noindent\textbf{Status: \texttt{NEW\_PROVED}} for the two
printed classes and their distinct capacities. Frozen bounded
relative quasicontinuity, specifically frozen R11
\path{dirichlet/analytic_laurent_collars.tex},
\path{cor:r11-analytic-bounded-relative-QC}, is
\texttt{INTERNAL\_INPUT}; finite
Radon Lusin approximation is \texttt{EXTERNAL\_KNOWN}.

\begin{proof}
Choose an ordinary open $W$ with $H\subset W\Subset V$ and
the finite parameter control $\nu$ from (CB4). It controls not
only original testers but \emph{every band-pasted tester} on
parameter bands. Indeed if $b$ has pieces $E_i$ and branches
$b_i$, then (BA3) and (CB4) give, for every Borel $E$,
\[
 \begin{split}
 \widetilde\mu_b\bigl(W\cap s^{-1}(E)\bigr)
 &=\sum_i\mu^R_{b_i}\bigl(W\cap s^{-1}(E\cap E_i)\bigr)\\
 &\leq\sum_i\nu(E\cap E_i)=\nu(E).
 \end{split}
 \tag{BQ3}
\]
Consequently $c(W\cap s^{-1}(E))\leq\nu(E)$.

For the given $h$, let $i:J\to\mathbf N$ be its Borel branch
index. Lusin's theorem for this map into the discrete countable
space gives a compact $F\subset J$ with
$\nu(J\setminus F)<\epsilon/2$ on which $i$ is continuous.
Its image on $F$ is finite. The set
$G_0=W\cap s^{-1}(J\setminus F)$ is ordinary open, and
(BQ3) gives $c(G_0)<\epsilon/2$. On
$H\setminus G_0$, the branch index is a continuous
finite-valued function.

If the branches are continuous, this already makes $h$ continuous
there, proving (BQ2). For general actual bounded classical
branches, apply their frozen outer $C_R$-quasicontinuity to the
finitely many branches occurring on $F$. Remove finitely many
ordinary open sets whose total original capacity is less than
$\epsilon/2$. Together with $G_0$ they give the first assertion.
This latter step claims a bound only in the original capacity,
not in the larger band capacity.

For a countable collection in the continuous-branch class, apply
(BQ2) with exceptional bounds $\epsilon2^{-j-1}$ and take the
union. Countable subadditivity bounds its capacity, and restricting
a continuous function to the smaller complement preserves
continuity. No common finite partition for that countable
collection is needed.
\end{proof}

\subsection{All original mixed coefficients are controlled by band capacity}

For the rest of this section impose the precise frozen geometric
compensator hypothesis (DG1): there is one compact original
$K$-affinoid domain $B$ with
\[
                         \overline V\Subset\operatorname{Int}(B).
 \tag{BQ4}
\]
The background $R$ need only be defined on $V$; no extension of
it across the boundary is required here. All original actual
moment refinements are still permitted.

\begin{lemma}[Uniform original coefficient domination]
\label{lem:r12-band-capacity-coefficients}
For $0\leq j\leq k$ and every original compact smooth
complementary test form $\eta$ of degree $(k-j,k-j)$, there is
a finite constant $C_\eta$, independent of the normalized band
tuple, such that for every Borel $E\subset V$,
\[
 \left|\mu_{\widetilde T_{R,j}(h_1,\ldots,h_j),\eta}\right|(E)
                              \leq C_\eta c(E).
 \tag{BQ5}
\]
In particular the variations are uniformly bounded on compact
sets and ignore all band-outer-capacity-null sets.
\end{lemma}

\noindent\textbf{Status: \texttt{NEW\_PROVED}} for the enlarged
tester family; the full original compensator (DG2) is frozen
\texttt{INTERNAL\_INPUT}.

\begin{proof}
Apply (DG2) to the positive original current
$T=\widetilde T_{R,j}(h_1,\ldots,h_j)$ and the chosen test.
It gives a normalized original smooth classical psh $a$ on $V$
and $A_\eta<\infty$ with
\[
                  |\mu_{T,\eta}|
                        \leq A_\eta\mathbf1_{\operatorname{supp}\eta}
                                                (Da)^{k-j}\wedge T.
\]
On each common parameter band, the original smooth-factor
compatibility identifies the right side with the corresponding
mixed top product. The countable sum (BA3) gives the same
identification globally. The average of the $k$ normalized
entries $h_1,\ldots,h_j,a,\ldots,a$ belongs to
$\mathcal A_s^{\rm band}(V)$: only finitely many countable
partitions are refined, and the averages of actual classical
branches are actual normalized testers. Its positive top
expansion contains $k!$ times the product of those labeled
entries, divided by $k^k$. Therefore
\[
               (Da)^{k-j}\wedge T(E)
                         \leq\frac{k^k}{k!}\,c(E).
\]
This proves (BQ5) with $C_\eta=A_\eta k^k/k!$.
For $j=0$ use only the $k$ smooth slots; for $j=k$ the
ordinary scalar estimate suffices. The compensator includes any
additional actual refinement of $\eta$, so this is a full
original-test domination, not merely a moment-trace statement.
\end{proof}

\begin{lemma}[Weighted weak passage under the specified capacity]
\label{lem:r12-band-weighted-weak-passage}
Let $T_m\to T$ on all original smooth compact tests, with
every current positive and of the same diagonal degree. Suppose
the coefficient bounds (BQ5) hold uniformly for $T_m$ and $T$.
Let $g$ be a bounded Borel function which is locally ordinary
outer $c$-quasicontinuous. Then
\[
                              gT_m\longrightarrow gT
                         \quad\text{on every original compact test}.
 \tag{BQ6}
\]
If the Borel functions $g_m$ have one common local bound and
converge locally in outer $c$-capacity to $g$, then also
$g_mT_m\to gT$. In this last statement
$(g_m-g)T_m$ tends to zero in local coefficient total variation.
No coefficient-total-variation convergence of $T_m$ itself is
asserted.
\end{lemma}

\noindent\textbf{Status: \texttt{DIRECT\_COROLLARY}} of
quasicontinuity, (BQ5), ordinary continuous extension and original
scalar density. The proof fixes the capacity used in the passage.

\begin{proof}
Fix an original test $\eta$ and a compact neighborhood of its
support. Remove an ordinary open $G$ of $c$-capacity less than
$\epsilon$ so that $g$ is continuous on the compact complement.
Extend that restriction to a bounded continuous function on the
compact neighborhood, without increasing its bound, and multiply
by a compact original cutoff equal to one near
$\operatorname{supp}\eta$. Such continuous extension is valid
on the local compact metrizable affinoid stage over the original
separable field. Call the resulting continuous compact scalar
$g_\epsilon$. On the support of $\eta$ it differs from $g$
only in $G$, and its absolute difference is at most twice the
common bound. Formula (BQ5) therefore bounds the error in each
pairing by $2\|g\|_\infty C_\eta\epsilon$, uniformly in $m$
and for $T$.

Original smooth scalar density approximates $g_\epsilon$
uniformly with a fixed compact support. Uniform compact
coefficient variation and the given original weak convergence
then imply $g_\epsilon T_m\to g_\epsilon T$.
Letting $\epsilon\downarrow0$ proves (BQ6).

For a varying multiplier and a compact coefficient support $H$,
split into the sets where $|g_m-g|\leq\delta$ and where it is
larger. If $M$ bounds the multipliers locally, (BQ5) gives
\[
 \int_H|g_m-g|\,d|\mu_{T_m,\eta}|
 \leq\delta\sup_m|\mu_{T_m,\eta}|(H)
       +2M C_\eta c^*(H\cap\{|g_m-g|>\delta\}).
 \tag{BQ7}
\]
First let $m\to\infty$ and then $\delta\downarrow0$.
This proves local coefficient variation convergence of the
error, and the fixed-multiplier result finishes the proof.
\end{proof}

\begin{theorem}[Capacity continuity of the band products]
\label{thm:r12-band-capacity-continuity}
Under (BQ4), let $h_{i,m},h_i\in\mathcal A_s^{\rm band}(V)$
for $1\leq i\leq j\leq k$. Suppose each $h_i$ is locally
ordinary outer $c$-quasicontinuous and
\[
 c^*\bigl(H\cap\{|h_{i,m}-h_i|>\delta\}\bigr)\longrightarrow0
 \quad\text{for every }H\Subset V\text{ compact and }\delta>0.
 \tag{BQ8}
\]
Then
\[
 \widetilde T_{R,j}(h_{1,m},\ldots,h_{j,m})
       \longrightarrow\widetilde T_{R,j}(h_1,\ldots,h_j)
 \tag{BQ9}
\]
on every original compact smooth complementary test, including
every actual moment refinement. The same assertion remains
valid after multiplication by bounded locally outer-$c$-
quasicontinuous weights, and by bounded varying weights that
converge locally in outer $c$-capacity to such a weight.
\end{theorem}

\noindent\textbf{Status: \texttt{NEW\_PROVED}} for the new
band tester class and its own capacity. The hypotheses do not
substitute original raw capacity for the larger band capacity.

\begin{proof}
Induct on $j$. At $j=0$ the current is the fixed $R$.
Assume convergence for the tail of length $j-1$ and denote
its currents by $Q_m\to Q$. The uniform bounds (BQ5)
apply to both. The preceding lemma with multiplier
$h_{1,m}\to h_1$ gives
$h_{1,m}Q_m\to h_1Q$ on all original tests. Apply
distributional $D$ and the intrinsic recursion (BA2) to obtain
(BQ9). This is legitimate on a test $\eta$ because the
preceding convergence applies to the original compact test
$D\eta$ as well. This completes the induction. Applying the
same weighted weak-passage lemma to the resulting mixed
currents proves the last assertion.
\end{proof}

\begin{corollary}[Intrinsic joint decreasing continuity without monotone witnesses]
\label{cor:r12-intrinsic-band-decreasing-continuity}
Under (BQ4), suppose
$h_{i,m},h_i\in\mathcal A_s^{\rm band,ct}(V)$ and
$h_{i,m}\downarrow h_i$ pointwise on $V$, for the finitely
many slots $1\leq i\leq j\leq k$. Then (BQ9) and its
stated weighted versions hold.

Only the actual Borel functions on $V$ must decrease. Their
countable band presentations may vary with $m$; no common
countable partition for the whole row and no monotone ambient
classical representatives on every parameter slice are assumed.
\end{corollary}

\noindent\textbf{Status: \texttt{NEW\_PROVED}} for this
continuous-branch class. It extends the direct monotone-slice
argument precisely by replacing that hypothesis with the
proved self-capacity mechanism.

\begin{proof}
On a fixed compact $H$, the simultaneous version of (BQ2)
removes one ordinary open $G$ of arbitrarily small $c$-capacity
on whose compact complement every $h_{i,m}$ and $h_i$ is
continuous. Dini's theorem on $H\setminus G$ makes each
decreasing row converge uniformly there. Thus, for every
$\delta>0$ and all sufficiently large $m$, the set in
(BQ8) is contained in $G$. Its outer capacity is at most
$c(G)$, which can be arbitrarily small. This proves (BQ8)
and permits application of the preceding theorem. Dini is
applied to actual Borel values restricted to this compact set;
it does not impose a new psh representative or new fiber field.
\end{proof}

\paragraph{The remaining full band-class gate.}
The first theorem proves original-$C_R$ quasicontinuity for
every countable band pasting, but self-$c$ quasicontinuity only
for the continuous-branch subclass. The raw-null-ideal equality
(BA8) gives no quantitative bound of $c$ by $C_R$ on small
open sets. Consequently it cannot convert the former
quasicontinuity assertion into the latter. This is a recorded
\texttt{FAILED\_ATTEMPT}; the missing uniform passage remains
\texttt{OPEN} for arbitrary discontinuous actual branches.
The intrinsic decreasing theorem above must not be relabeled
as a theorem for that larger class or for the full Borel
slice-classical class.

Every arbitrary original $S$ gives the allowed mixed-slot
background $R=\Omega_s\wedge S$ as in (CB14), so these
quasicontinuity and continuity statements have full-dimensional
arbitrary-$S$ instances without changing the original field or
test category. They do not resolve unrestricted original-$S$
raw capacity, OPEN-GQ, or OPEN-CTO.

% END REVIEWED MODULE research_r12/capacity/band_quasicontinuity_and_continuity.tex


% BEGIN REVIEWED MODULE research_r12/capacity/band_capacitary_completion.tex
\section{A capacitary completion with an intrinsic bounded current calculus}
\label{sec:r12-band-capacitary-completion}

Assume (CB1)--(CB2) and the containing-affinoid hypothesis
(BQ4). Thus $R$ is still an original arbitrary middle-positive
separately closed conormal current of remaining degree $k\geq1$.
Put
\[
 \mathcal A_0=\mathcal A_s^{\rm band,ct}(V),\qquad
 \gamma_0(E)=\sup_{h\in\mathcal A_0}\widetilde\mu_h(E),\qquad
 \gamma_1(E)=C_R^{\rm band}(E;V,s).
 \tag{BC1}
\]
The distinction is essential: $\mathcal A_0$ uses continuous
actual branches, whereas $\gamma_1$ allows all bounded actual
classical branches. Both capacities use their entire displayed
tester families, and both are finite on compact sets. Write
$\gamma_i^*$ for ordinary-open outer regularization, $i=0,1$.
Then $\gamma_0\leq\gamma_1$, but equality is not assumed.

We give one construction for either fixed choice $i\in\{0,1\}$.
The self-generated case $i=0$ has the exact unchanged-outer-capacity
conclusion. The stronger full-band control $i=1$ has only the
one-sided comparison printed below.

The bounded recursive completion method itself is an explicitly
frozen input: R11 \path{capacity/barrier_relative_completion.tex},
\path{lem:r11-analytic-bounded-completion}, transfers R2
\path{capacitary_closure.tex}, \path{prop:r2-completion} and
\path{thm:r2-completed-calculus}, to the original analytic setting.
The proofs below verify that method for the new band seed and
its two distinct capacities. They do not claim the first bounded
arbitrary-background capacitary completion.

\begin{definition}[The two explicitly normalized capacitary completions]
\label{def:r12-band-capacitary-completion}
Let $\mathfrak B_i(V)$ be the closure of $\mathcal A_0$ among
Borel functions $V\to[0,1]$, modulo equality outside a locally
$\gamma_i^*$-null set, for local convergence in outer
$\gamma_i$-capacity. Thus a defining row satisfies
\[
 \gamma_i^*\bigl(H\cap\{|h_m-h|>\delta\}\bigr)\longrightarrow0
       \quad(H\Subset V\text{ compact},\ \delta>0),
       \qquad h_m\in\mathcal A_0.
 \tag{BC2}
\]
No ambient classical psh representative of $h$ is required or
asserted. Affine normalization extends the construction to any
fixed bounded real range when that is needed.
\end{definition}

\begin{theorem}[Concrete completeness and quasicontinuity]
\label{thm:r12-band-completion-metric}
Each $\mathfrak B_i(V)$ is a complete metrizable space of actual
Borel classes. Its members are locally ordinary-open outer
$\gamma_i$-quasicontinuous. It is closed under finite convex
combinations, maxima and the bounded operation
$\max(f,g-\epsilon)$ for $\epsilon\geq0$.

For an increasing compact exhaustion
$H_l\subset\operatorname{Int}(H_{l+1})$ of $V$, one compatible
metric is
\[
 \begin{split}
 d_i(f,g)&=\sum_{l\geq1}2^{-l}\,\kappa_{i,l}(f,g),\\
 \kappa_{i,l}(f,g)&=\inf\{\epsilon>0:
  \gamma_i^*(H_l\cap\{|f-g|>\epsilon\})\leq\epsilon\}.
 \end{split}
 \tag{BC3}
\]
The terms lie in $[0,1]$ because all functions are normalized.
Only a countable compact exhaustion of the ordinary geometry is
used here; there is no countable-cofinal replacement of the
capacity tester family.
\end{theorem}

\noindent\textbf{Status: \texttt{NEW\_PROVED}} for the concrete
completions of this seed class.

\begin{proof}
The affinoid space is metrizable over the original separable
field, as in the exact test-countability argument (DI1), so its
ordinary open $V$ has the stated compact exhaustion. The union
bound for the two error sets and countable subadditivity of
$\gamma_i^*$ give the triangle inequality for each
$\kappa_{i,l}$, and hence for $d_i$. Vanishing distance is
equivalent to equality outside a locally outer-null set: test
the difference at thresholds $1/m$ and take the countable union.
The same threshold argument shows that metric convergence is
exactly (BC2).

The entire space of normalized Borel classes is complete for
this metric. Indeed take a Cauchy row and select a subsequence,
still denoted $f_j$, such that
\[
 \gamma_i^*(H_j\cap\{|f_{j+1}-f_j|>2^{-j}\})<2^{-j}.
\]
Choose ordinary open covers $G_j$ of these error sets with
$\gamma_i(G_j)<2^{1-j}$. The late-tail unions have capacity
tending to zero. Off their limsup, the successive differences
are eventually bounded by the summable sequence $2^{-j}$ on
each compact core. Thus $f_j$ has a finite pointwise Borel
limit there, lying in $[0,1]$. Set the limit equal to zero on
the outer-null limsup. Uniform tail convergence outside the
late open union proves local capacity convergence, and the
Cauchy property returns the whole row to the same limit.
The closed seed closure $\mathfrak B_i(V)$ is therefore complete.

Each seed member is outer-$\gamma_1$-quasicontinuous by (BQ2),
and hence is outer-$\gamma_0$-quasicontinuous. For a capacity
limit of seed members, use the same fast-subsequence construction.
On a fixed compact remove a late union of the error covers and
also open quasicontinuity exceptions for the countably many
approximants, assigning summable arbitrarily small capacities.
Outside their union the approximants are continuous and converge
uniformly. Their limit is continuous there. This proves the
asserted quasicontinuity for either completion, and also proves
that changing values on an outer-null set has no effect on it.

The seed is closed under the printed finite operations: refine
finitely many countable parameter partitions and apply the
corresponding maximum or convex operation to their continuous
actual classical branches. The outputs remain normalized. These
operations are continuous for local capacity convergence, by
their pointwise Lipschitz inequalities and a finite union of
error sets. Taking seed limits proves the closure statements.
\end{proof}

\subsection{Construction on all original tests, before any compact restriction}

\begin{theorem}[Intrinsic bounded products in the completion]
\label{thm:r12-completed-band-products}
For $h_1,\ldots,h_j\in\mathfrak B_i(V)$ and $0\leq j\leq k$,
the recursion
\[
 \widehat T^{,i}_{R,0}=R,\qquad
 \widehat T^{,i}_{R,j}(h_1,\ldots,h_j)
    =D\bigl(h_1\widehat T^{,i}_{R,j-1}(h_2,\ldots,h_j)\bigr)
 \tag{BC4}
\]
is well-defined on the actual Borel quotient and gives original
symmetric middle-positive, separately closed, conormal
coefficient-Radon currents. It is independent of all defining
rows. For every original compact complementary test $\eta$,
\[
 \left|\mu_{\widehat T^{,i}_{R,j}(\mathbf h),\eta}\right|(E)
                \leq C_\eta\gamma_i^*(E)
                         \qquad(E\subset V\text{ Borel}),
 \tag{BC5}
\]
where the constant is independent of the normalized completed
tuple. For a diagonal top product the exact constant is one:
\[
       \widehat\mu^{,i}_h(E)
           :=\widehat T^{,i}_{R,k}(h^{[k]})(E)
                          \leq\gamma_i^*(E).
 \tag{BC6}
\]
The products are symmetric and multilinear in the usual bounded
positive-cone sense after affine normalization. They agree with
the seed products, annihilate a constant differentiated slot,
and have original smooth-factor compatibility. They are local
on ordinary opens for equality of the Borel entries modulo the
specified outer-null sets.

Every equality is on every original compact smooth complementary
test and every actual moment refinement. Ordinary-current
equality is not promoted to equality on extra GK delta tests.
\end{theorem}

\noindent\textbf{Status: \texttt{NEW\_PROVED}} for this completed
bounded current calculus. Frozen bounded products enter only at
the actual seed level.

\begin{proof}
First observe that (BQ5) for seed tuples holds with the smaller
$\gamma_0$, not merely with $\gamma_1$. Its proof averages a
finite seed tuple and one normalized original smooth compensator;
that average is still in $\mathcal A_0$. Thus every seed
coefficient has the bound $C_\eta\gamma_0(E)$, and hence the
bound $C_\eta\gamma_i(E)$ for either choice of $i$.
The same argument at degree zero controls the original background
coefficients, so multiplication of $R$ by a bounded Borel
$\gamma_i$-class is intrinsic.

Induct on $j$. Assume the tail current $Q$ is constructed and
all its seed approximating tails $Q_m$ converge to it on original
tests, with (BC5). Each $h_1$ has the proved quasicontinuity.
The proof of the weighted weak-passage lemma (BQ6)--(BQ7) applies
unchanged when the limit coefficient bound is $C_\eta\gamma_i^*$:
its fixed exceptional sets are ordinary open, on which
$\gamma_i^*=\gamma_i$, and its varying-multiplier estimate uses
outer capacity already. Therefore
$h_{1,m}Q_m\to h_1Q$ on all original tests. Applying
distributional $D$ proves convergence of the seed products to
the intrinsic right side of (BC4).

The seed coefficient variations are uniformly bounded on compact
sets. Their distributional limit consequently has order-zero
Radon coefficients and is the same original current. Positivity,
symmetry, separate closedness and the conormal identities pass
on arbitrary original tests. For each fixed coefficient measure,
variation lower semicontinuity on an ordinary open $G$ gives
\[
        |\mu_{\widehat T,\eta}|(G)
                 \leq\liminf_m|\mu_{T_m,\eta}|(G)
                 \leq C_\eta\gamma_i(G).
 \tag{BC7}
\]
One first tests continuous compact scalar factors; original
smooth scalar density and the uniform coefficient bounds justify
this passage from current convergence. Individual Radon outer
regularity then proves (BC5) on every Borel $E$. In particular
the new current ignores the specified outer-null sets, validating
all Borel multiplications at the next induction step. For a
diagonal top tuple the seed measure bound is exactly
$\widetilde\mu_{h_m}\leq\gamma_i$; the identical open passage
and Radon outer step prove (BC6).

The recursion proves independence of defining rows. Symmetry and
the bounded affine-cone multilinear identities follow by passing
the corresponding seed identities through simultaneous defining
rows. For a smooth slot the distributional Leibniz identity and
separate closedness of the tail give the usual smooth wedge.
A constant slot gives zero for the same reason. Ordinary locality
follows by induction from the local distributional recursion and
coefficient nullity of the exceptional set. No varying weak limit
has been restricted to a compact set or to a discontinuous
equality indicator in this construction.
\end{proof}

\begin{theorem}[Capacity and intrinsic joint decreasing continuity]
\label{thm:r12-completed-band-continuity}
For tuples in $\mathfrak B_i(V)$, local outer-$\gamma_i$-capacity
convergence in every slot implies convergence of every mixed
current (BC4) on all original tests. Bounded Borel weights may
also vary in local outer $\gamma_i$-capacity, provided their
limit is locally outer-$\gamma_i$-quasicontinuous; the weighted
current convergence then holds as well.

In particular, if finitely many normalized completed rows
$h_{a,m}\downarrow h_a$ pointwise on their Borel representatives
and each $h_a\in\mathfrak B_i(V)$, all mixed products and
the stated weighted versions converge. No monotonicity of any
chosen actual or band-pasted defining witnesses is required.
\end{theorem}

\noindent\textbf{Status: \texttt{NEW\_PROVED}} for the completed
class with its printed topology and limit-membership condition.

\begin{proof}
Every completed entry is quasicontinuous and every completed
mixed tail has the outer bound (BC5). The weighted weak-passage
proof therefore applies exactly as in the construction of (BC4).
Induction on the tuple length proves the first assertion, now
for completed rather than merely seed rows. The weighted
conclusion is the same lemma applied to the resulting mixed
currents.

For decreasing rows, remove one ordinary open set of arbitrarily
small $\gamma_i$-capacity so that the countably many entries
and their finitely many limits are all continuous on the fixed
compact complement. This uses summable error budgets and the
proved quasicontinuity. Dini's theorem gives uniform convergence
there. Hence the rows converge in local outer $\gamma_i$-capacity,
and the first assertion applies. Membership of the limit is an
explicit hypothesis; an arbitrary pointwise infimum is not
declared a member without a separate capacity-limit argument.
\end{proof}

\begin{theorem}[Strict maximum carriers and the equality-side inequality]
\label{thm:r12-completed-band-maximum-carriers}
Let $u,v,z_1,\ldots,z_b\in\mathfrak B_i(V)$, let
$w=\max(u,v)$, and suppose $a\geq1$, $a+b\leq k$. Then
\[
 \mathbf1_{\{u>v\}}\widehat T^{,i}_{R,a+b}(w^{[a]},\mathbf z)
    =\mathbf1_{\{u>v\}}\widehat T^{,i}_{R,a+b}(u^{[a]},\mathbf z).
 \tag{BC8}
\]
Moreover the full middle-positive current inequality holds:
\[
 \begin{split}
 \widehat T^{,i}_{R,a+b}(w^{[a]},\mathbf z)
 \geq{}&\mathbf1_{\{u\geq v\}}
                    \widehat T^{,i}_{R,a+b}(u^{[a]},\mathbf z)\\
       &+\mathbf1_{\{u<v\}}
                    \widehat T^{,i}_{R,a+b}(v^{[a]},\mathbf z).
 \end{split}
 \tag{BC9}
\]
No restricted current in these formulas is asserted separately
closed. In top degree, finite maxima therefore preserve any
fixed original positive-measure subsolution lower bound.
\end{theorem}

\noindent\textbf{Status: \texttt{NEW\_PROVED}} for this completion.
The proof parallels the frozen carrier method, but uses the
newly established band-capacity continuity rather than assuming
it for the larger class.

\begin{proof}
At the seed level refine finitely many countable band partitions.
On each band the continuous actual branches have an ordinary
open strict comparison locus. Frozen ordinary locality gives
the maximum identity there; (PB9) and the positive band sum give
the seed version of (BC8).

Choose seed defining rows for the completed tuple and their
maxima. Let $\chi_\epsilon$ be a continuous Lipschitz scalar
cutoff which is zero on $(-\infty,\epsilon]$, one on
$[2\epsilon,\infty)$, and nondecreasing between. Multiply the
seed strict-carrier identity by
$\chi_\epsilon(u_m-v_m)$. These bounded weights converge in
local capacity to $\chi_\epsilon(u-v)$, whose limit is
quasicontinuous. The completed weighted continuity theorem
passes the identity. Let $\epsilon\downarrow0$ using a fixed
rescaled cutoff. Bounded convergence for every original
coefficient measure gives (BC8).

Put $w_\epsilon=\max(u,v-\epsilon)$. This belongs to the
completion and converges uniformly to $w$, hence in local
capacity. Apply (BC8) on its two strict carriers. The shift
can be normalized by replacing $u,v$ by
$(u+\epsilon)/(1+\epsilon),v/(1+\epsilon)$, so the identity
uses only already constructed normalized entries and their
proved affine compatibility. Positivity on the remaining
restriction gives, before any limiting operation,
\[
 \widehat T^{,i}(w_\epsilon^{[a]},\mathbf z)
 \geq\mathbf1_{\{u\geq v\}}\widehat T^{,i}(u^{[a]},\mathbf z)
       +\mathbf1_{\{u<v-\epsilon\}}\widehat T^{,i}(v^{[a]},\mathbf z).
\]
The first right-hand current is fixed, retaining any equality
mass. The second increases on every positive original test as
$\epsilon\downarrow0$. The left side converges by the proved
capacity theorem. This yields (BC9) on the full middle-positive
dual cone. In top degree the two indicators partition the
finite Borel locus, and all relevant measures ignore its
outer-null complement; the subsolution assertion follows.
\end{proof}

\subsection{The exact capacity statement and the remaining cofinality gate}

For Borel $E$ define
\[
                \widehat\gamma_i(E)
                  =\sup_{h\in\mathfrak B_i(V)}\widehat\mu_h^{,i}(E).
 \tag{BC10}
\]
\begin{theorem}[Outer capacity is unchanged for the exact seed]
\label{thm:r12-band-completion-outer-capacity}
For every Borel $E\subset V$,
\[
 \gamma_0(E)\leq\widehat\gamma_0(E)\leq\gamma_0^*(E),
 \qquad
 \widehat\gamma_0(G)=\gamma_0(G)\quad(G\text{ ordinary open}),
 \qquad
 \widehat\gamma_0^*=\gamma_0^*.
 \tag{BC11}
\]
For the full-band-control completion only the following is
asserted without an extra cofinality theorem:
\[
 \gamma_0(E)\leq\widehat\gamma_1(E)\leq\gamma_1^*(E),
 \qquad
 \gamma_0(G)\leq\widehat\gamma_1(G)\leq\gamma_1(G)
                       \quad(G\text{ ordinary open}).
 \tag{BC12}
\]
In particular, (BC11) does not prove that completing the seed
creates no mass gain on compact sets.
\end{theorem}

\noindent\textbf{Status: \texttt{NEW\_PROVED}} for the exact
capacity comparisons. The equality with the outer capacity of
the \emph{entire original or full band} tester class remains
\texttt{OPEN} unless its required cofinality is separately proved.

\begin{proof}
The seed is included in both completions and its products agree
with the completed products. This gives both lower bounds.
The upper bounds are (BC6). On an ordinary open set raw and
outer capacity agree by definition, proving the open equalities
or inequalities as printed. Taking infima over ordinary open
neighborhoods proves the last equality in (BC11). No limit has
been restricted to a compact set in this argument.
\end{proof}

\begin{proposition}[All globally continuous Borel slice families enter the completion]
\label{prop:r12-continuous-slice-completion}
Let $H:J\times V\to[0,1]$ be jointly Borel, and suppose that
for every $t\in J$ the function $H_t=H(t,\cdot)$ is a
continuous actual classical psh function on the entire fixed
$V$. Then the actual diagonal Borel function
\[
                                h(x)=H(s(x),x)
 \tag{BC13}
\]
belongs to both $\mathfrak B_0(V)$ and $\mathfrak B_1(V)$.
It admits locally uniformly convergent approximants from
$\mathcal A_0$. Therefore the preceding bounded products,
maximum carriers and intrinsic decreasing theorem apply to
such slice families, even when the actual diagonal rows
decrease without decreasing ambient slice witnesses.
\end{proposition}

\noindent\textbf{Status: \texttt{NEW\_PROVED}} for this explicit
inclusion and its intrinsic continuity consequence. It is not an
uncountable pointwise fiber-capacity selection theorem.

\begin{proof}
On the locally compact second countable ordinary space $V$,
$C(V)$ with its compact-open metric is separable. This follows,
for example, from separability of the continuous function spaces
on a countable compact exhaustion and their countable product.
Its subset consisting of continuous normalized actual classical
psh functions is therefore separable as well. Choose a countable
dense family $(a_l)$ in that subset; these are actual functions
on $V$, not conditional fiber representatives.

Fix a compact exhaustion $H_m$ as in (BC3). For each fixed
$a_l$, the function
$t\mapsto\sup_{x\in H_m}|H(t,x)-a_l(x)|$ is Borel:
the supremum can be taken over a countable dense subset of
$H_m$, because every slice is continuous. For each $m$, choose
the first index $l=l_m(t)$ for which this supremum is less than
$1/m$. Density ensures existence for every $t$, and the first
index defines a countable Borel parameter partition. Put
$h_m(x)=a_{l_m(s(x))}(x)$. Then $h_m\in\mathcal A_0$ and
\[
                     \sup_{x\in H_m}|h_m(x)-h(x)|<1/m.
\]
Thus $h_m\to h$ locally uniformly, which implies (BC2) for
both capacities. The completed conclusions now follow. This
selection uses the compact-open topology of \emph{continuous
global functions}; it makes no simultaneous exceptional-set
assertion for an uncountable family of conditional current tests.
\end{proof}

\begin{corollary}[The empty parameter tuple retains the original arbitrary background]
\label{cor:r12-empty-parameter-completion}
Permit the empty parameter tuple, with $J=\mathbf R^0$ a single
point and $\Omega_s=1$. Then (CB2) is vacuous and $R=S$ is
the \emph{unchanged arbitrary original} middle-positive separately
closed background. The seed $\mathcal A_0$ is exactly the
normalized continuous actual classical psh family on $V$,
$\gamma_0$ is its raw tester capacity, and $\gamma_1=C_S$ is
the full original normalized bounded classical tester capacity.
All preceding constructions and printed bounds apply under
(BQ4), without any additional structural restriction on $S$.

Where the frozen R11 bounded completion hypotheses overlap,
the $i=1$ construction is a subclass of that frozen completion,
and its products agree with the frozen products. Neither equality
of the two tester capacities nor equality of the two completed
potential classes is inferred.
\end{corollary}

\noindent\textbf{Status: \texttt{DIRECT\_COROLLARY}} and an
explicit recovery/transfer of the frozen arbitrary-background
bounded mechanism, not a new claim of its first construction.

\begin{proof}
There is only one nonempty band, so the definitions reduce to
the displayed actual seed and its two capacities. The proofs
above require no conormal equation when the tuple is empty.
For the frozen comparison, a continuous actual defining row
is also an allowed row of bounded actual classical functions.
The intrinsic recursions agree at degree zero and hence at
every successive degree. This identifies their products on
all original tests wherever both constructions are defined.
\end{proof}

\paragraph{The precise failed inclusion and the remaining theory.}
The original first attempt at the open-capacity assertion used
the words ``the original testers are included.'' That is valid
for the exact seed $\mathcal A_0$ and invalid for the larger
class of possibly discontinuous classical branches without an
additional approximation theorem. It is retained here as a
\texttt{FAILED\_ATTEMPT}, repaired by the two separate capacities
(BC1) and the exact inequalities (BC11)--(BC12). No counterexample
to the missing cofinality is claimed.

The completion supplies positive closed products, the original
recursion, symmetry and bounded multilinearity, original smooth
compatibility, locality, maximum carriers, and weighted joint
decreasing continuity with its printed limit-membership condition.
It does not assert every classical capacity or Dirichlet property:
unrestricted comparison with boundary data, arbitrary compact
raw/outer equality, identification with all ambient classical psh
functions, and analytic PS/PL presentation descent require their
own proofs. The two completions use different possible null
quotients; they are not silently identified with one another or
with the energy lattice. For arbitrary original $S$, applying
the construction to the derived $R=\Omega_s\wedge S$ retains
all original-field and all-original-test requirements without
turning it into a proof of OPEN-GQ or OPEN-CTO.

Domain restriction is a separate issue from locality of the
constructed currents. Capacity on a smaller tester domain can
be larger, so convergence in the capacity of $V$ alone does not
prove membership in the corresponding completion on an arbitrary
smaller open. No such domain-capacity equivalence or gluing theorem
for the potential classes is assumed. Whenever the class
constructions coexist on an overlap and their Borel values agree
in the relevant null quotients, their original currents restrict
and agree by the local recursion. That conditional compatibility
does not assert the missing class membership.

% END REVIEWED MODULE research_r12/capacity/band_capacitary_completion.tex


\part{Energy duality and the unchanged potential cone}

% BEGIN REVIEWED MODULE research_r12/energy/energy_difference_banach.tex
\section{A Banach space of differences of strong energy potentials}
\label{sec:r12-energy-difference}

\paragraph{Exact standing hypotheses and frozen inputs.}
Let $K/\mathbf Q_p$ be the original finite extension with its original
absolute value. Let $X$ be good, Hausdorff, boundaryless, quasi-smooth,
and pure of dimension $n$. Choose a compact original affinoid domain
$B\subset X$, functions $f_a,g_b\in\mathcal O(B)$ and positive real
thresholds $R_a,r_b$, such that
\[
 K_0=\{|f_a|\leq R_a,\ |g_b|\geq r_b\}
                         \Subset\operatorname{Int}_X(B),\qquad
 V=\{|f_a|<R_a,\ |g_b|>r_b\}\ne\varnothing .
 \tag{ED1}
\]
Let $S$ be an arbitrary original symmetric CLD middle-positive,
separately $d'$- and $d''$-closed $(p,p)$-current on an ordinary
neighborhood of $K_0$. Put $q=n-p\geq1$ and $d=q+1$.
No strong positivity, geometric representation or approximation of
$S$ is assumed. Every current below is tested on the full original
CLD test sheaf, with every actual moment refinement.

Write $\mathscr C=\widehat{\mathcal E}^{1}_{\rm str}(V)$ for the
precise R11 strong completion, with its encoded approximation
boundary condition. Its elements are nonpositive Borel potentials
modulo local ordinary-open outer $S$-capacity zero. We use the R11
inputs (GE10)--(GE18), (IC10)--(IC22), (CC1)--(CC11), (CI1)--(CI8)
and the completed extension of (EV1). They provide, respectively,
the symmetric global mixed energy, its positive energy matrices,
the complete strong cone, its intrinsic original currents, and its
monotone closure. All these statements are \texttt{INTERNAL\_INPUT}
in their printed R11 scopes.

For $u_0,\ldots,u_q\in\mathscr C$, abbreviate
\[
 B(u_0,\ldots,u_q)=\int_V(-u_0)\,\widehat N_q(u_1,\ldots,u_q),
 \quad e(u)=B(u^{[d]}),\quad r(u)=e(u)^{1/d}.
 \tag{ED2}
\]
The frozen inequalities are $B(\mathbf u)\leq\prod_i r(u_i)$ and
$r(u+v)\leq r(u)+r(v)$. The zero-energy class is zero. Completed
potentials are finite outside outer-null sets; all differences below
are formed there and assigned any finite Borel value on the joint
exceptional set. In particular no expression $(-\infty)-(-\infty)$
is used as a pointwise number.

\begin{definition}[Energy-difference space and its norm]
\label{def:r12-energy-difference}
Let $\mathbb D_S=\mathscr C-\mathscr C$ be the real vector space of
such Borel differences modulo local outer-capacity equality. Define
\[
 \|h\|_{\rm ed}
   =\inf\{r(a)+r(b):a,b\in\mathscr C,\ h=a-b\}.
 \tag{ED3}
\]
This is a norm defined by the energy cone, not the discrepancy $I$.
It is not declared to be the ordinary supremum norm or a Hilbert norm.
\end{definition}

\begin{lemma}[Continuous extension of the mixed energy]
\label{lem:r12-difference-mixed-extension}
There is a unique symmetric real $d$-linear form on $\mathbb D_S$
extending $B$, and it satisfies
\[
       |B(h_0,\ldots,h_q)|\leq\prod_{i=0}^q\|h_i\|_{\rm ed}.
 \tag{ED4}
\]
Moreover $\|u\|_{\rm ed}=r(u)$ for every $u\in\mathscr C$.
\end{lemma}

\noindent\textbf{Status: \texttt{NEW\_PROVED}} for the full
unbounded cone-difference construction; the mixed energy bound on
the positive cone is the explicit frozen input.

\begin{proof}
Expand each difference into its two cone terms and extend by the
finite multilinear expansion. If $a-b=a'-b'$, then
$a+b'=a'+b$ in the original relative cone. Separate additivity of
the frozen $B$ therefore makes the expansion independent of that
presentation, one slot at a time. Positive real homogeneity and
changing the signs of differences give real multilinearity.
For specified decompositions $h_i=a_i-b_i$, the triangle inequality
and cone Holder give the bound $\prod_i(r(a_i)+r(b_i))$.
Taking independent infima proves (ED4).

The definition gives $\|u\|_{\rm ed}\leq r(u)$. Applying (ED4)
to the diagonal tuple gives $e(u)\leq\|u\|_{\rm ed}^{d}$, so the
reverse inequality follows. This argument also covers zero energy.
\end{proof}

\begin{theorem}[Actual Borel realization and Banach completeness]
\label{thm:r12-energy-difference-Banach}
Formula (ED3) is a norm, $\mathbb D_S$ is complete for that norm,
and $\mathscr C$ is a closed convex cone in this Banach space.
Norm convergence implies convergence in the full local outer
$S$-capacity topology. More precisely, with
$C_0=d2^d$, one has
\[
 C_S^*(\{|h|>t\};V)\leq C_0 t^{-d}\|h\|_{\rm ed}^{d}
                    \qquad(h\in\mathbb D_S,\ t>0).
 \tag{ED5}
\]
For $u,v\in\mathscr C$,
\[
 I(u,v)\leq\bigl(r(u)^q+r(v)^q\bigr)\|u-v\|_{\rm ed}.
 \tag{ED6}
\]
The result does not identify this norm topology with the strong
discrepancy topology on $\mathscr C$.
\end{theorem}

\noindent\textbf{Status: \texttt{NEW\_PROVED}}.

\begin{proof}
The infimum in (ED3) is finite. Minkowski on $\mathscr C$ proves
the triangle inequality; switching the two cone terms proves
symmetry, and homogeneity follows by rescaling decompositions.
The R11 capacity integral bound $J_1^*(u)\leq2^d e(u)$ and
monotonicity of the sublevels give
\[
 \frac{t^d}{d}C_S^*(\{u<-t\};V)
       \leq\int_0^t s^{d-1}C_S^*(\{u<-s\};V)\,ds
       \leq2^d e(u).
\]
For $h=a-b$, the set $\{|h|>t\}$ is contained, modulo an
outer-null set, in $\{a<-t\}\cup\{b<-t\}$. Hence its outer
capacity is at most $C_0t^{-d}(r(a)^d+r(b)^d)$, bounded by
$C_0t^{-d}(r(a)+r(b))^d$. Infimizing proves (ED5), and shows
that a zero seminorm has the zero relative Borel class. Local
outer-null sets here also have zero outer capacity in the fixed
domain: use a countable compact exhaustion and the R11 domain
comparison on its cores. The same comparison transfers norm
convergence to every smaller-domain original capacity.

We prove completeness inside the space of actual Borel differences,
rather than introducing an abstract completion with unidentified
elements. For a norm-Cauchy sequence choose a subsequence $(h_j)$
whose successive increments have decompositions
\[
 h_{j+1}-h_j=a_j-b_j,\qquad
 \sum_{j\geq1}(r(a_j)+r(b_j))<\infty.
 \tag{ED7}
\]
Write $h_1=a_0-b_0$. The partial sums
$A_N=a_0+\sum_{j=1}^N a_j$ and
$B_N=b_0+\sum_{j=1}^N b_j$ decrease as Borel potentials,
outside one countable union of outer-null sets. Their energy
roots are bounded by the corresponding sums of roots, so the
R11 decreasing-closure theorem (CI8) gives strong limits
$A,B\in\mathscr C$. Apply the same theorem to each tail series.
It yields $A-A_N,B-B_N\in\mathscr C$ and
\[
 r(A-A_N)+r(B-B_N)
          \leq\sum_{j>N}(r(a_j)+r(b_j))\longrightarrow0.
 \tag{ED8}
\]
The identities with the tails hold on the common finite locus,
hence in the relative Borel quotient. Thus $h=A-B\in\mathbb D_S$
is the norm limit of the subsequence and of the original Cauchy
sequence. No abstract vector invisible to Borel capacity has been
adjoined.

For $u,v\in\mathscr C$, multilinearity gives
\[
 I(u,v)=B(u-v,u^{[q]})-B(u-v,v^{[q]}).
\]
Apply (ED4) to obtain (ED6). If a sequence in $\mathscr C$
converges in norm to $h\in\mathbb D_S$, its energy roots are
bounded by the norm triangle inequality and the preceding lemma.
Equation (ED6) makes it strongly Cauchy. R11 completeness supplies
a limit $u\in\mathscr C$. The norm limit and the strong limit
have the same local capacity class by (ED5) and the R11 strong
topology theorem. Thus $h=u$, proving norm closedness.
Convexity is the frozen cone addition property.
\end{proof}

\begin{definition}[Curvature as an energy-dual functional]
\label{def:r12-dual-curvature}
For $u\in\mathscr C$ put
\[
       J(u)(h)=B(h,u^{[q]})=-\int_Vh\,d\widehat\mu_u
                      \qquad(h\in\mathbb D_S).
 \tag{ED9}
\]
The displayed integral is literal and absolutely convergent: for
$h=a-b$, $|h|\leq(-a)+(-b)$ off a common null set, and the
two mixed first energies are finite. The current measure ignores
that exceptional set. Thus $J(u)$ is a genuine functional in
$\mathbb D_S^*$, with its original current interpretation.
\end{definition}

\begin{theorem}[Exact dual norm and strong continuity of curvature]
\label{thm:r12-dual-curvature-continuity}
For every continuous linear functional $\ell$ on $\mathbb D_S$,
\[
 \|\ell\|_*=
     \sup_{\substack{a\in\mathscr C\\r(a)\leq1}}|\ell(a)|.
 \tag{ED10}
\]
For $u\in\mathscr C$ one has $\|J(u)\|_*=r(u)^q$.
For $u,v\in\mathscr C$ and $R=\max(r(u),r(v))$,
\[
 \begin{split}
 \|J(u)-J(v)\|_*&\leq qR^{q-1}\|u-v\|_{\rm ed},\\
 \|J(u)-J(v)\|_*&\leq\sqrt q\,R^{(q-1)/2}I(u,v)^{1/2}.
 \end{split}
 \tag{ED11}
\]
At $R=0$ both sides mean zero; for $q=1$ the power $R^0$ is
one. In particular strong convergence of potentials gives norm
convergence of their source functionals, in every remaining degree.
\end{theorem}

\noindent\textbf{Status: \texttt{NEW\_PROVED}} for the precise
Banach dual and quantitative all-degree source map. The integrated
positive energy matrices used in the proof are frozen inputs.

\begin{proof}
If the supremum in (ED10) is $L$, a decomposition $h=a-b$
gives $|\ell(h)|\leq L(r(a)+r(b))$. Taking the infimum yields
$\|\ell\|_*\leq L$; the reverse follows from
$\|a\|_{\rm ed}=r(a)$. Equation (ED4) bounds the norm of
$J(u)$ by $r(u)^q$. For $r(u)>0$, evaluation on $u/r(u)$
gives equality; if $r(u)=0$, the class and functional are zero.

The finite telescoping expansion gives
\[
 (J(u)-J(v))(h)
    =\sum_{l=0}^{q-1}B(h,u-v,u^{[l]},v^{[q-1-l]}).
 \tag{ED12}
\]
Equation (ED4) proves the first bound in (ED11).
For the sharper strong-topology bound define
\[
 H(f,g)=\sum_{l=0}^{q-1}
                 B(f,g,u^{[l]},v^{[q-1-l]}).
\]
This is a positive semidefinite bilinear form on $\mathbb D_S$.
Indeed each summand is the integrated energy form with the
fixed completed positive closed mixed background. Positivity on
finite differences follows from R11 (GE7), passed through the
global mixed-energy convergence (IC13). It is not a pointwise
Alexandrov--Fenchel assertion. We have
\[
 H(u-v,u-v)=I(u,v),\qquad
 H(a,a)\leq q r(a)^2R^{q-1}\quad(a\in\mathscr C).
\]
Cauchy--Schwarz for this positive semidefinite form and (ED12)
bound the functional on every $a$ with $r(a)\leq1$.
Equation (ED10) proves the second bound. Strongly convergent
rows have bounded energy roots by R11, so the final assertion
follows. Every construction used the intrinsic original products;
no special tropical trace was substituted for their test sheaf.
\end{proof}

\paragraph{Two different source distances.}
When signed measure integration defines the corresponding continuous
functionals, their difference norm is
\[
 \delta_*(\mu,\nu)
     =\sup_{\substack{a\in\mathscr C\\r(a)\leq1}}
                    \left|\int_V(-a)\,d(\mu-\nu)\right|.
 \tag{ED13}
\]
This is not the R11 absolute-variation distance (VF11), which has
$\int(-a)\,d|\mu-\nu|$ inside the supremum. The latter dominates
(ED13) whenever it is defined, and can be strictly stronger.
The signed cancellation in (ED13) is essential to (ED11).
The Banach space and its curvature map do not prove an analytic
PS/PL presentation quotient, or equality on additional GK delta tests.

% END REVIEWED MODULE research_r12/energy/energy_difference_banach.tex


% BEGIN REVIEWED MODULE research_r12/energy/nonlinear_dual_source_geometry.tex
\section{Closed nonlinear source images and a dual normal decomposition}
\label{sec:r12-nonlinear-dual-source}

Keep exactly (ED1), the unchanged arbitrary original $S$, and the
Banach space $\mathbb D_S$ of the preceding section. Its positive
potential cone is denoted $\mathscr C$ even though its Borel
representatives are nonpositive functions. Write $J$ for (ED9),
$r(u)=e(u)^{1/d}$, and $\|\cdot\|_*$ for the genuine operator
dual norm (ED10). In particular this is not the absolute-variation
distance of R11 (VF11).

\begin{theorem}[A complete curvature-dual metric and closed source image]
\label{thm:r12-closed-full-source-image}
The image $J(\mathscr C)$ is norm closed in $\mathbb D_S^*$.
The formula
\[
                 d_J(u,v)=\|J(u)-J(v)\|_*
 \tag{NI1}
\]
defines a complete metric on $\mathscr C$, inducing precisely the
R11 strong discrepancy topology. For $u,v\in\mathscr C$,
\[
 \begin{split}
 I(u,v)&\leq(r(u)+r(v))d_J(u,v),\\
 d_J(u,v)&\leq\sqrt q\,
     \max(r(u),r(v))^{(q-1)/2}I(u,v)^{1/2}.
 \end{split}
 \tag{NI2}
\]
No common dominating measure and no fixed actual lower barrier
are required. The assertion concerns all $q\geq1$ and the full
original test category on the stated domain.
\end{theorem}

\noindent\textbf{Status: \texttt{NEW\_PROVED}}. R11 (OI1)--(OI3)
proved closedness only inside a fixed dominated interval for its
stronger absolute-variation distance; neither restriction occurs here.

\begin{proof}
Multilinearity gives the exact identity
$I(u,v)=(J(u)-J(v))(u-v)$. The dual inequality and
$\|u-v\|_{\rm ed}\leq r(u)+r(v)$ prove the first estimate.
The second is (ED11). Thus $d_J=0$ gives $I=0$ and equal
relative classes. Symmetry and the triangle inequality come
from the dual norm, not from an asserted triangle inequality
for $I$.

For completeness, mixed Holder also gives the explicit root bound
\[
 I(u,v)\geq\bigl(r(u)-r(v)\bigr)
                   \bigl(r(u)^q-r(v)^q\bigr)
                \geq|r(u)-r(v)|^{q+1}.
\]
The first inequality follows by bounding the two mixed terms in
$e(u)+e(v)-B(u,v^{[q]})-B(v,u^{[q]})$ separately. The second
is the elementary inequality $x^q-y^q\geq(x-y)^q$ for
$x\geq y\geq0$. Hence convergence to a fixed potential in
$I$ automatically bounds the energy roots.

Suppose $J(u_j)$ is norm Cauchy. Its norms are bounded, so
$r(u_j)^q=\|J(u_j)\|_*$ bounds the energy roots. The first
estimate in (NI2) makes $(u_j)$ strongly Cauchy. R11 strong
completeness gives $u\in\mathscr C$ with $I(u_j,u)\to0$.
Its energy roots remain bounded, and the second estimate gives
$J(u_j)\to J(u)$ in norm. Every norm limit of image elements
is consequently in the image, proving closedness and metric
completeness. The same inequalities prove equivalence of the
two convergences: strong convergence bounds the roots by the
frozen energy-root estimate, whereas $d_J$ convergence bounds
them by the exact dual norm identity. These arguments also
identify their Cauchy sequences.
\end{proof}

\begin{corollary}[Source approximation on every original test]
\label{cor:r12-original-source-approximation}
Let $u_j\in\mathscr C$ and suppose their source functionals
$J(u_j)$ converge in $\mathbb D_S^*$ to $\ell$. Then there is
a unique $u\in\mathscr C$ with $J(u)=\ell$, and $u_j\to u$
strongly. If a specified original Radon measure $\mu$ satisfies
\[
                 \ell(\phi)=-\int_V\phi\,d\mu
            \quad\text{for every original compact smooth scalar }\phi,
 \tag{NI3}
\]
then $\widehat\mu_u=\mu$ on all such tests, after every actual
moment refinement. All other strong current and weighted limits
of R11 (IC11)--(IC13) hold as well.
\end{corollary}

\noindent\textbf{Status: \texttt{DIRECT\_COROLLARY}}.

\begin{proof}
Closedness supplies $u$ and the quantitative estimates supply
strong convergence. Every original compact smooth scalar is an
element of $\mathbb D_S$: R11 (HC14)--(HC15) writes it as the
exact difference of two actual bounded finite-mass common-collar
potentials. That construction is valid for all $q$, not only
for its separate Hilbert-space consequence. Thus (ED9) and
(NI3) give equality on the complete original scalar test sheaf.
Original smooth scalar density and Radon regularity identify
the same measures. The strong calculus supplies the other
limits. This is an approximation criterion, not a claim that
every source admits such image approximants.
\end{proof}

\begin{definition}[The strong-continuous energy dual]
\label{def:r12-strong-continuous-dual}
Let $\mathbb D_{S,\mathrm{s}}^*$ consist of those
$\ell\in\mathbb D_S^*$ for which
\[
  u_j,u\in\mathscr C,\ I(u_j,u)\to0
                 \quad\Longrightarrow\quad\ell(u_j)\to\ell(u).
 \tag{NI4}
\]
This is a continuity requirement on the actual strong cone.
It is not automatically imposed on every Banach-dual functional.
\end{definition}

\begin{lemma}[A closed dual class containing every curvature]
\label{lem:r12-strong-dual-closed}
The class $\mathbb D_{S,\mathrm{s}}^*$ is a norm-closed linear
subspace of $\mathbb D_S^*$, and $J(\mathscr C)$ is contained
in it. Consequently the norm-closed linear span of all curvature
functionals is contained in this class.
\end{lemma}

\noindent\textbf{Status: \texttt{DIRECT\_COROLLARY}} of the
proved Banach realization and frozen strong tester estimate.

\begin{proof}
Linearity is immediate. A strongly convergent row has bounded
$r(u_j)$, hence is bounded in $\mathbb D_S$. If $\ell_m\to\ell$
in norm, the error $|(\ell-\ell_m)(u_j-u)|$ is uniformly small
for large $m$. Fixing $m$ then invokes (NI4), proving closure.
For fixed $v\in\mathscr C$, the completed R11 tester estimate
(EV1) bounds $|J(v)(u_j-u)|$ by
$\int|u_j-u|\,d\widehat\mu_v\to0$. Its constants are finite
on the bounded-energy family. This proves the curvature assertion.
\end{proof}

Let
\[
 \mathscr C^\circ
   =\{n\in\mathbb D_S^*:n(v)\leq0\text{ for every }v\in\mathscr C\}.
 \tag{NI5}
\]
This is the polar of the potential cone in the specified dual,
not an identification of any two current-positivity cones.

\begin{theorem}[All-degree variational normal decomposition]
\label{thm:r12-nonlinear-normal-decomposition}
For every $\ell\in\mathbb D_{S,\mathrm{s}}^*$, the functional
\[
                F_\ell(v)=\ell(v)-e(v)/d\qquad(v\in\mathscr C)
 \tag{NI6}
\]
has a unique maximizer $u=\Pi(\ell)$. Every maximizing sequence
converges strongly to $u$. Equivalently, there is a unique pair
$(u,n)$ with
\[
       \ell=J(u)+n,\qquad
       u\in\mathscr C,\quad n\in\mathscr C^\circ,
                           \quad n(u)=0.
 \tag{NI7}
\]
The residual also lies in $\mathbb D_{S,\mathrm{s}}^*$.
Define the nonnegative support gauge
\[
       p(\ell)=\sup_{\substack{v\in\mathscr C\\r(v)\leq1}}\ell(v).
 \tag{NI8}
\]
Then
\[
 \begin{split}
 e(u)&=\ell(u),\qquad r(u)^q=p(\ell)\leq\|\ell\|_*,\\
 \sup F_\ell&=\frac qd\,p(\ell)^{d/q}.
 \end{split}
 \tag{NI9}
\]
If $\ell$ is nonnegative on $\mathscr C$, then
$p(\ell)=\|\ell\|_*$. No assertion that the residual $n$ vanishes
is part of this theorem.
\end{theorem}

\noindent\textbf{Status: \texttt{NEW\_PROVED}} for arbitrary
strong-continuous signed dual data and every remaining degree.
The midpoint existence mechanism is explicitly inherited from
R11; the Banach source class and exact normal decomposition
are the new extension. This is not unrestricted source solvability.

\begin{proof}
Since $|\ell(v)|\leq\|\ell\|_*r(v)$,
$F_\ell(v)\leq\|\ell\|_*r(v)-r(v)^d/d$. Thus its supremum
$L\geq F_\ell(0)=0$ is finite, and every maximizing sequence
has bounded energies. In the midpoint gap the linear term
cancels. The frozen sharp inequality gives
\[
 c_q I(v_j,v_k)
     \leq L-\tfrac12(F_\ell(v_j)+F_\ell(v_k))\longrightarrow0,
 \quad c_q=\frac{1-2^{-q}}{2d}.
\]
Strong completeness provides $u\in\mathscr C$. Its energy
converges, and (NI4) passes the datum term, so $F_\ell(u)=L$.
The same midpoint estimate proves uniqueness and convergence
of every maximizing sequence. Actual classical strong density
also gives the same supremum over the frozen actual class.

The continuous $d$-linear extension $B$ makes $e/d$ a polynomial
on $\mathbb D_S$ whose differential at $u$ is $J(u)$.
Only segments inside $\mathscr C$ are used for maximization.
Their first variation gives
\[
                   (\ell-J(u))(v-u)\leq0\qquad(v\in\mathscr C).
 \tag{NI10}
\]
Conversely convexity of $e$ on the cone makes this sufficient.
Taking the competitors $0,2u$ gives $(\ell-J(u))(u)=0$.
Then (NI10) is precisely the polar condition in (NI7).
Conversely (NI7) implies (NI10). The pair is unique because
the maximizer is unique and $n=\ell-J(u)$. Its strong continuity
follows from the preceding lemma.

Scaling along $tu$ gives $e(u)=\ell(u)$, also when $u=0$.
For each $v\in\mathscr C$, the polar inequality gives
$\ell(v)\leq J(u)(v)\leq r(v)r(u)^q$. Thus
$p(\ell)\leq r(u)^q$. If $u\ne0$, evaluation on $u/r(u)$
gives the reverse inequality; if $u=0$, the polar condition
and the allowed zero tester give $p(\ell)=0$. This proves
(NI9). Finally (ED10) says
$\|\ell\|_* =\max(p(\ell),p(-\ell))$; if $\ell\geq0$ on
the cone, the latter second term is zero.
\end{proof}

\begin{theorem}[A quantitative nonlinear retraction of source functionals]
\label{thm:r12-source-retraction}
Put $\mathcal P(\ell)=J(\Pi(\ell))$ on
$\mathbb D_{S,\mathrm{s}}^*$. This is a retraction onto the norm-closed
image $J(\mathscr C)$, and is positively homogeneous of degree one.
For $\|\ell\|_*,\|m\|_*\leq R$, with $R>0$, one has
\[
 \begin{split}
 I(\Pi(\ell),\Pi(m))&\leq2R^{1/q}\|\ell-m\|_*,\\
 \|\mathcal P(\ell)-\mathcal P(m)\|_*
                         &\leq\sqrt{2qR}\,\|\ell-m\|_*^{1/2}.
 \end{split}
 \tag{NI11}
\]
The first inequality has the more precise right side
$(\ell-m)(\Pi(\ell)-\Pi(m))$. If
$n=\ell-\mathcal P(\ell)$ and $R=\|\ell\|_*>0$, then
\[
 \left(\sqrt{qR+\|n\|_*}-\sqrt{qR}\right)^2
 \leq\operatorname{dist}_*(\ell,J(\mathscr C))\leq\|n\|_*.
 \tag{NI12}
\]
In particular a nonzero residual is a quantitative obstruction to
source-image approximation, not merely a failure of one chosen
maximization algorithm.
\end{theorem}

\noindent\textbf{Status: \texttt{NEW\_PROVED}} for this nonlinear
all-degree dual retraction and the exact image-approximation gap.
It is not asserted to be a Hilbert metric projection.

\begin{proof}
If $\ell=J(a)$, the pair $(a,0)$ satisfies (NI7), so uniqueness
gives the retraction. Homogeneity of $J$ gives
$\Pi(t\ell)=t^{1/q}\Pi(\ell)$ and
$\mathcal P(t\ell)=t\mathcal P(\ell)$ for $t\geq0$;
the case $t=0$ is the zero class.

Write $u=\Pi(\ell)$ and $v=\Pi(m)$. Use $v,u$ as the
competitors in their two inequalities (NI10), and add. It yields
$I(u,v)\leq(\ell-m)(u-v)$. Duality bounds this by
$\|\ell-m\|_*(r(u)+r(v))$; (NI9) bounds the roots by
$R^{1/q}$. The second line of (NI11) now follows from (ED11).
If $R=0$ both projections are zero.

Let $\delta$ be the distance in (NI12). Since $J(0)=0$,
$\delta\leq R$. Choose $\eta_j\in J(\mathscr C)$ with
$\delta_j=\|\ell-\eta_j\|_*\to\delta$. Their norms are
at most $R+\delta_j$. Retraction and (NI11) imply
\[
 \|n\|_*\leq\delta_j+
               \sqrt{2q(R+\delta_j)\delta_j}.
\]
Passing to the limit and using $\delta\leq R$ gives
$\|n\|_*\leq\delta+2\sqrt{qR\delta}$. Solving this
quadratic inequality in $\sqrt\delta$ gives the lower bound.
The upper bound uses the particular image element $\mathcal P(\ell)$.
\end{proof}

\paragraph{Source equations and scope.}
For an actual measure datum whose literal integration defines $\ell$,
the equation in the completed class holds exactly when the residual
is zero. The decomposition does not assert that an arbitrary positive
Radon datum has that property. The R11 compact-source obstruction
still gives nonzero residuals, and the new higher-degree cap-lift
obstruction supplies further examples. No ambient classical
representative, ordinary pointwise boundary trace, analytic GQ,
or low-real-degree CTO is inferred.

% END REVIEWED MODULE research_r12/energy/nonlinear_dual_source_geometry.tex


% BEGIN REVIEWED MODULE research_r12/energy/capacity_controlled_dual_data.tex
\section{Capacity-controlled signed data in the strong energy dual}
\label{sec:r12-capacity-dual-data}

We retain exactly the strict Laurent-box and extension hypotheses
(ED1), the arbitrary original middle-positive separately closed
background $S$, and all original tests. Put $d=q+1\geq2$,
$\mathscr C=\widehat{\mathcal E}^{1}_{\rm str}(V)$, and
$\mathbb D_S=\mathscr C-\mathscr C$. No moment presentation or
special form of $S$ is imposed in this section. The outer capacity
$C_S^*(E;V)$ always uses the full original bounded classical tester
class on the specified whole domain $V$.

The datum criterion below is a sufficient condition for membership
in the strong-continuous dual, not a nonlinear solvability theorem.
It does not require domination by one actual curvature measure.

\begin{definition}[An integrable capacity modulus]
\label{def:r12-integrable-capacity-modulus}
An integrable capacity modulus in degree $d$ is a nondecreasing
function $F:[0,\infty]\to[0,\infty]$, finite on every finite
positive argument, with $F(0)=0$, such that
\[
                 \int_0^1 F(s)s^{-1-1/d}\,ds<\infty.
 \tag{CD1}
\]
The value at infinity is its increasing limit. Let $\mu$ be a
finite signed original Radon measure on $V$. We say that its
variation has modulus $F$ if
\[
                    |\mu|(E)\leq F\bigl(C_S^*(E;V)\bigr)
                    \quad\text{for every Borel }E\subset V.
 \tag{CD2}
\]
The total variation $m=|\mu|(V)$ is finite independently of
whether $C_S^*(V;V)$ is finite.
\end{definition}

Put $C_0=d2^d$, as in (ED5), and define
\[
       A_{F,m}:=\int_0^\infty
              \min\{m,F(C_0s^{-d})\}\,ds.
 \tag{CD3}
\]
If $m=0$, set $A_{F,0}=0$. Substitution $x=C_0s^{-d}$ on a
tail of the integral shows that (CD1) makes (CD3) finite.

\begin{theorem}[Literal signed data and strong continuity]
\label{thm:r12-capacity-modulus-strong-dual}
Suppose $\mu$ satisfies (CD2) for an integrable capacity modulus.
Then every $h\in\mathbb D_S$ is absolutely integrable for
$|\mu|$, and
\[
       \ell_\mu(h)=-\int_V h\,d\mu,
       \qquad
       \|\ell_\mu\|_*\leq A_{F,m}
 \tag{CD4}
\]
defines an element of $\mathbb D_{S,\mathrm s}^*$.
More explicitly, if $u_j,u\in\mathscr C$ and $I(u_j,u)\to0$,
then
\[
                    \int_V|u_j-u|\,d|\mu|\longrightarrow0.
 \tag{CD5}
\]
All differences and integrals are unchanged by locally outer-null
modifications. The assertion holds for every $q\geq1$ and every
allowed original $S$, including degenerate backgrounds after the
specified relative quotient.
\end{theorem}

\noindent\textbf{Status: \texttt{NEW\_PROVED}} as a
full-degree signed datum class in the newly constructed Banach
dual. The sublevel and all-local strong capacity estimates are
the explicitly inherited R11 inputs (CI3), (EV7)--(EV8).

\begin{proof}
For $a\in\mathscr C$, (ED5) and (CD2) give
\[
       |\mu|\{a<-t\}
       \leq \min\{m,F(C_0r(a)^d t^{-d})\},\qquad t>0.
 \tag{CD6}
\]
The actual Borel representative of $a$ is nonpositive and finite
outside an outer-null set. Condition (CD2) annihilates this set.
The layer-cake formula, followed by $t=r(a)s$, therefore gives
\[
                        \int_V(-a)\,d|\mu|
                                      \leq A_{F,m}r(a).
 \tag{CD7}
\]
When $r(a)=0$, the class is zero outside a local outer-null set,
which is globally outer-null in the fixed domain by the R11
countable compact exhaustion and domain comparison. Thus (CD7)
holds in that case too. If $h=a-b$, then
$|h|\leq(-a)+(-b)$ off the joint null set. Consequently
$\int|h|\,d|\mu|\leq A_{F,m}(r(a)+r(b))$. Taking the infimum
over its presentations proves absolute integrability and (CD4).
This is a literal integral on the actual Borel difference space;
no unproved extension of a measure to abstract completion points
is used.

Now let $u_j\to u$ strongly. The energy-root estimate gives a
finite $R$ with $r(u_j),r(u)\leq R$. If $R=0$ the assertion
is immediate. Otherwise, for every $L>0$, (CD6) gives
\[
 \sup_j\int_V\bigl((-u_j)-L\bigr)_+\,d|\mu|
       \leq\int_L^\infty
                  \min\{m,F(C_0R^d t^{-d})\}\,dt
                   \longrightarrow0\quad(L\to\infty).
 \tag{CD8}
\]
The same bound holds for $u$. This is a uniform integrability
statement for the full finite datum, including its boundary tail.

Fix $L$ and write $u_j^{(L)}=\max(u_j,-L)$ and
$u^{(L)}=\max(u,-L)$. On each compact $K\Subset V$, the
R11 strong estimate makes
$C_S^*(K\cap\{|u_j^{(L)}-u^{(L)}|>\varepsilon\};V)$ tend
to zero. An integrable nondecreasing modulus satisfies
$F(s)\to0$ as $s\downarrow0$: otherwise (CD1) would diverge.
Hence (CD2) gives convergence in $|\mu|$-measure on $K$.
The clipped differences are bounded by $L$, so their integrals
over $K$ tend to zero. Radon inner regularity and finiteness of
$|\mu|$ allow a compact $K$ with arbitrarily small
$|\mu|(V\setminus K)$; the same bound $L$ then passes this
convergence to the whole domain. Finally (CD8) removes the
clipping uniformly. This proves (CD5) and the defining
strong-dual continuity (NI4).

Every capacity in the argument is the frozen full original
capacity. Actual moment refinements alter neither the measure
integrals nor the classes of tests defining these estimates.
No assertion of raw outer regularity on arbitrary Borel sets
is needed.
\end{proof}

\begin{corollary}[Power and logarithmic endpoint moduli]
\label{cor:r12-power-capacity-dual-data}
Let $m=|\mu|(V)<\infty$. If, for constants $A>0$ and
$\gamma>1/d$,
\[
                 |\mu|(E)\leq A\,C_S^*(E;V)^\gamma
                    \quad(E\text{ Borel}),
 \tag{CD9}
\]
then $\ell_\mu\in\mathbb D_{S,\mathrm s}^*$ and
\[
 \|\ell_\mu\|_*
 \leq \frac{d\gamma}{d\gamma-1}
       m^{1-1/(d\gamma)}
                  (A C_0^\gamma)^{1/(d\gamma)}.
 \tag{CD10}
\]
There are also endpoint admissible moduli: any nondecreasing
extension, for $s$ bounded away from zero, of
\[
             F(s)=A s^{1/d}\bigl(\log(e/s)\bigr)^{-\alpha}
                   \quad(0<s\leq1),\qquad \alpha>1,
 \tag{CD11}
\]
satisfies (CD1). Neither the unmodified critical power
$s^{1/d}$ nor the range $\alpha\leq1$ satisfies this
integral criterion. This last statement is about the criterion,
not a universal necessity claim for all possible measures.
\end{corollary}

\noindent\textbf{Status: \texttt{DIRECT\_COROLLARY}}.

\begin{proof}
For (CD9), (CD3) becomes the integral of
$\min(m,A C_0^\gamma s^{-d\gamma})$. Splitting at
$s_0=(A C_0^\gamma/m)^{1/(d\gamma)}$ yields (CD10) when
$m>0$, and the zero case is immediate. For (CD11), substitution
$y=\log(e/s)$ turns (CD1), up to the harmless constant $A$,
into $\int_1^\infty y^{-\alpha}\,dy$. This proves all printed
assertions without an endpoint counterexample being inferred
from the failure of a sufficient estimate.
\end{proof}

\begin{proposition}[Uniform capacity control converts variation to dual convergence]
\label{prop:r12-capacity-dini-data-stability}
Let $\mu_j,\mu$ be finite signed Radon measures such that
$m_j=|\mu_j-\mu|(V)\to0$. Suppose the variations
$|\mu_j-\mu|$ all satisfy (CD2) with one common integrable
modulus $F$. Then
\[
         \|\ell_{\mu_j-\mu}\|_*
                \leq A_{F,m_j}\longrightarrow0.
 \tag{CD12}
\]
This unconditional statement defines only the functional of the
controlled difference. If the individual datum functionals are
defined, it is their difference by literal integration.
If additionally each $\ell_{\mu_j}$ and $\ell_\mu$ lies in
$\mathbb D_{S,\mathrm s}^*$, their canonical variational
potentials from (NI7) converge strongly, and their normal
residuals converge in dual norm. For power domination the
right side in (CD12) has the explicit bound (CD10) with
$m=m_j$. If $\ell_{\mu_j}=J(u_j)$ are actual completed source
functionals, then the limit is itself $J(u)$; in that special
case $\widehat\mu_u=\mu$ on every original scalar test.
\end{proposition}

\noindent\textbf{Status: \texttt{DIRECT\_COROLLARY}} of
the new datum criterion and nonlinear closed-image theorem.

\begin{proof}
Apply (CD4) to $\mu_j-\mu$. Since $m_j$ is bounded, the
integrands in (CD3) are dominated by one integrable function
$\min(M,F(C_0s^{-d}))$ and tend pointwise to zero. Dominated
convergence proves (CD12). The quantitative retraction estimate
(NI11) supplies strong convergence of the variational potentials
and dual convergence of their curvatures; subtracting from the
data gives convergence of the residuals. The final assertion is
the closed-image criterion (NI1)--(NI3). It assumes source
approximants and does not create them from capacity domination.
\end{proof}

\paragraph{Relation to the frozen first-order theory.}
R11 already solves a different linear transposition equation for
arbitrary finite diffuse data in a first-order Dirichlet-form
space. That result does not say that the solution is a classical
psh strong-completion member. The present theorem places signed
data in an all-degree nonlinear energy dual and supplies its
variational decomposition. It neither supersedes that linear
result nor removes the nonlinear normal residual. In particular,
bounded-source domination implies (CD9) with $\gamma=1$, but
the compact-source examples in Section~\ref{sec:r12-positive-moment-lifts}
show that even this datum class does not imply the unrestricted
source equation.

% END REVIEWED MODULE research_r12/energy/capacity_controlled_dual_data.tex


% BEGIN REVIEWED MODULE research_r12/energy/strong_dual_gate_counterexample.tex
\section{Why strong-continuous dual data cannot be replaced by arbitrary dual data}
\label{sec:r12-strong-dual-counterexample}

This section gives one exact allowed example. It is used to disprove
universal identifications of topologies and a missing continuity
hypothesis, not as a substitute for the full-degree theorems of the
preceding sections. The field and its norm remain the original ones.

Let $X=(\mathbf G_{m,K})^{\mathrm{an}}$, let $t=\log|T|$, and let
\[
       V=t^{-1}((-1,1)),\qquad S=[V],\qquad q=1,
       \sigma:(-1,1)\longrightarrow V
 \tag{CG1}
\]
be the original annulus and its Gauss section. The closed annulus
embeds in a larger compact original annulus, so (ED1) holds. Write
$\mathscr C$ and $\mathbb D_S$ for precisely the same strong cone
and Banach difference space as before. The current $S$ is used here
only as an explicit permitted background for a counterexample.

For $-1<s<1$ put
\[
 \begin{split}
 G(x,s)&=\frac{(1-\max(x,s))(1+\min(x,s))}{2},\\
 u_s&=-G(t,s).
 \end{split}
 \tag{CG2}
\]
The function $u_s$ is a bounded continuous convex PL moment
potential, negative on $V$ and zero at its two boundary ends.
Its one-variable slope has jump one at $s$. The frozen original
torus calibration therefore gives, on every original test,
\[
        \mu_{u_s}=\delta_{\sigma(s)},\qquad
        B(u_s,u_a)=G(s,a),\qquad
        e(u_s)=G(s,s)=\frac{1-s^2}{2}.
 \tag{CG3}
\]
All these potentials are actual members of the finite-mass
zero-boundary class. Arbitrary real $s$ gives an actual real PL
potential over the same field; no value-group enlargement is made.

\begin{proposition}[A Banach-dual atom functional outside the strong dual]
\label{prop:r12-atom-functional-not-strong}
There is a continuous linear functional $L\in\mathbb D_S^*$,
nonnegative on $\mathscr C$, such that
\[
             L(u_0)=1,\qquad L(u_s)=0\quad(s\ne0),
             \qquad \|L\|_*\leq\sqrt2.
 \tag{CG4}
\]
Nevertheless $I(u_s,u_0)\to0$ as $s\to0$. In particular,
\[
                     \mathbb D_{S,\mathrm s}^*
                                      \ne\mathbb D_S^*.
 \tag{CG5}
\]
The energy-difference norm topology on the cone is strictly
stronger than its discrepancy topology in this example.
\end{proposition}

\noindent\textbf{Status: \texttt{NEW\_PROVED}} for the
new dual and norm. The universal assertion obtained by omitting
the strong-dual condition (NI4) is \texttt{FALSE}.

\begin{proof}
In degree one the completed first current is additive and
positively homogeneous on $\mathscr C$. For $h=a-b\in\mathbb D_S$
its linear extension
\[
                    M_h=\widehat\mu_a-\widehat\mu_b
 \tag{CG6}
\]
is well defined as a locally finite signed original measure.
Indeed two presentations give $a+b'=a'+b$ in the completed
relative cone, hence the same sum of completed first currents.
Define
\[
                          L(h)=M_h(\{\sigma(0)\}).
 \tag{CG7}
\]
This definition uses the proved linear completed current, not
an assumed distributional multiplication of every Borel
difference by $S$.

Since $-u_0\geq0$ and $-u_0(\sigma(0))=1/2$, for $a\in\mathscr C$
one has
\[
 \widehat\mu_a(\{\sigma(0)\})
       \leq2\int_V(-u_0)\,d\widehat\mu_a
       =2B(u_0,a)
       \leq\sqrt2\,r(a).
 \tag{CG8}
\]
For $h=a-b$, the variation of $M_h$ is dominated by
$\widehat\mu_a+\widehat\mu_b$ locally. Thus
$|L(h)|\leq\sqrt2(r(a)+r(b))$, and taking the infimum proves
continuity with the stated norm bound. Positivity on the cone
and (CG4) follow from (CG3).

For $0<s<1$, direct substitution in (CG3) gives
\[
                 I(u_s,u_0)
       =G(s,s)+G(0,0)-2G(s,0)=s-\frac{s^2}{2}\longrightarrow0.
 \tag{CG9}
\]
By contrast $L(u_s-u_0)=-1$. Therefore $L$ fails (NI4) and
\[
                    \|u_s-u_0\|_{\rm ed}\geq1/\sqrt2
                    \qquad(0<s<1).
 \tag{CG10}
\]
This proves both strictness assertions. The current formula
(CG3), the definition (CG6), and all limits are equalities
of original currents, so using refined non-toric test moments
does not remove the example.
\end{proof}

\begin{theorem}[Positive dual data with an unattained variational supremum]
\label{thm:r12-dual-nonattainment}
Fix $0<\varepsilon<1/2$ and let
\[
                         \ell=J(u_0)-\varepsilon L.
 \tag{CG11}
\]
This is a norm-continuous linear functional on $\mathbb D_S$,
nonnegative on the whole cone $\mathscr C$. The finite supremum
\[
           \sup_{v\in\mathscr C}
                     \bigl(\ell(v)-e(v)/2\bigr)=\frac14
 \tag{CG12}
\]
is not attained. An explicit maximizing sequence is $u_s$ as
$s\downarrow0$, and it converges strongly to $u_0$, where the
functional has value $1/4-\varepsilon$.

Consequently boundedness in the energy-difference dual, even
with positivity on the potential cone, cannot replace the
strong-continuity hypothesis in (NI7). This does not contradict
any theorem for literal positive Radon data proved to satisfy
that hypothesis.
\end{theorem}

\noindent\textbf{Status: \texttt{NEW\_PROVED}}; unrestricted
dual-data variational attainment is \texttt{FALSE}.

\begin{proof}
Norm continuity is immediate from (CG8). For $v\in\mathscr C$,
symmetry and literal mixed energy integration give
\[
        J(u_0)(v)=B(v,u_0)=\int_V(-u_0)\,d\widehat\mu_v
                          \geq\tfrac12 L(v).
\]
Thus $\ell(v)\geq(1/2-\varepsilon)L(v)\geq0$.
In degree one the exact quadratic identity is
\[
 J(u_0)(v)-e(v)/2
                =e(u_0)/2-I(v,u_0)/2
                =\frac14-I(v,u_0)/2.
 \tag{CG13}
\]
Since $I\geq0$ and $L(v)\geq0$, the variational expression
in (CG12) is at most $1/4$. Equality would require
$I(v,u_0)=0$ and $L(v)=0$. The first condition means
$v=u_0$ as a relative class, whereas $L(u_0)=1$.
Thus equality is impossible.

For $s>0$, (CG3)--(CG4) instead give
\[
      \ell(u_s)-e(u_s)/2
            =\frac{1-s}{2}-\frac{1-s^2}{4}
            =\frac14-\frac s2+\frac{s^2}{4}
                      \longrightarrow\frac14.
 \tag{CG14}
\]
At $u_0$ the value is $1/4-\varepsilon$. This proves the
exact supremum, nonattainment, and failure of passage to the
strong limit. There is no failure of completeness: it is
precisely the datum term which is not strongly continuous.
\end{proof}

\begin{corollary}[The signed dual distance is not weighted total variation]
\label{cor:r12-dual-versus-variation}
Let the R11 absolute-variation distance be
\[
 \delta_1(\mu,\nu)
     =\sup_{\substack{a\in\mathscr C\\e(a)\leq1}}
                     \int_V(-a)\,d|\mu-\nu|.
 \tag{CG15}
\]
For the exact moving sources above,
\[
 \begin{split}
  \|J(u_s)-J(u_0)\|_*&\longrightarrow0,\\
  \liminf_{s\downarrow0}\delta_1(\mu_{u_s},\mu_{u_0})
                                                  &\geq\sqrt2.
 \end{split}
 \tag{CG16}
\]
Thus the new closed-source-image theorem genuinely uses a
weaker signed-cancellation dual distance, not a renamed
version of the old variation distance.
\end{corollary}

\noindent\textbf{Status: \texttt{DIRECT\_COROLLARY}} with
an explicit counterexample to equality of the two topologies.

\begin{proof}
The first statement follows from (ED11) and (CG9). For $s>0$,
the two distinct Gauss atoms have variation
$\delta_{\sigma(s)}+\delta_{\sigma(0)}$. The cone tester
$a=u_0/r(u_0)=\sqrt2u_0$ has energy one. Its variation pairing
is
\[
            \sqrt2\bigl(G(s,0)+G(0,0)\bigr)
                         =\sqrt2(1-s/2)\longrightarrow\sqrt2.
\]
Taking the supremum in (CG15) proves the assertion. Both
distances are evaluated on the same allowed analytic annulus
and original measures.
\end{proof}

% END REVIEWED MODULE research_r12/energy/strong_dual_gate_counterexample.tex


% BEGIN REVIEWED MODULE research_r12/nonlinear/energy_obstacles_and_contact_defects.tex
\section{The energy obstacle map and its contact defect}
\label{sec:r12-energy-obstacles}

Keep exactly the strict original analytic Laurent box and background
extension \textup{(ED1)}.  Thus $S$ remains an arbitrary original
symmetric CLD middle-positive separately $d'$- and $d''$-closed
$(p,p)$-current, $q=n-p\ge1$, and $d=q+1$.  No factorization of
$S$ is imposed in this section.  Write
\[
 \mathscr C=\widehat{\mathcal E}^{1}_{\rm str}(V),\qquad
 \mathbb D_S=\mathscr C-\mathscr C,\qquad
 r(u)=e(u)^{1/d},\qquad E(u)=e(u)/d.
 \tag{OC1}
\]
The Banach norm $\|\cdot\|_{\rm ed}$, its literal Borel
realization, and the curvature functional
$J(u)(h)=-\int h\,d\widehat\mu_u$ are the proved R12
\textup{(ED2)--(ED11)}.  Their proof does not use the obstacle
map constructed below.  The R11 strong relative Perron theorem
\textup{(SE1)--(SE6)} is an \texttt{INTERNAL\_INPUT}.
All orders are relative Borel orders modulo local ordinary-open
outer $S$-capacity zero, denoted by $\preceq$.  All measures are
the original measures on the full original test sheaf.

\begin{definition}[An obstacle in the energy-difference space]
\label{def:r12-energy-obstacle}
For $f\in\mathbb D_S$ put
\[
 P(f)=\sup_{\rm rel}\{v\in\mathscr C:v\preceq f\},
 \qquad \mathfrak E(f)=E(P(f)).
 \tag{OC2}
\]
The supremum is the R11 relative least upper bound among Borel
classes.  It is not the ordinary upper semicontinuous
regularization of a pointwise supremum.
\end{definition}

\begin{theorem}[A global obstacle map and quantitative continuity]
\label{thm:r12-global-energy-obstacle}
The definition is valid for every $f\in\mathbb D_S$.  Its value
$P(f)$ is the greatest member of $\mathscr C$ below $f$, and
\[
 P(f)\preceq f,\qquad P(u)=u\quad(u\in\mathscr C),\qquad
 r(P(f))\le\|f\|_{\rm ed}.
 \tag{OC3}
\]
The map $P$ is order preserving, positively homogeneous and
superadditive in the relative order, hence order concave.
For $f,g\in\mathbb D_S$, set
$R=\max(r(P(f)),r(P(g)))$.  Then
\[
 \begin{split}
 |r(P(f))-r(P(g))|&\le\|f-g\|_{\rm ed},\\
 I(P(f),P(g))&\le R^q\|f-g\|_{\rm ed},\\
 \|J(P(f))-J(P(g))\|_*
       &\le\sqrt q\,R^{(2q-1)/2}\|f-g\|_{\rm ed}^{1/2}.
 \end{split}
 \tag{OC4}
\]
At $R=0$ the second and third left sides are zero.  In
particular, norm convergence of obstacles gives strong
convergence of their envelopes and dual-norm convergence of
their original curvature functionals.  Norm continuity of
$P$ with values in $\mathbb D_S$ is not claimed.
\end{theorem}

\noindent\textbf{Status: \texttt{NEW\_PROVED}} for the
global energy-difference obstacle domain and its quantitative
estimates, in every degree and for arbitrary original $S$.

\begin{proof}
Write $f=a-b$ with $a,b\in\mathscr C$.  Since $b\preceq0$,
one has $a\preceq f$.  The family in \textup{(OC2)} is therefore
nonempty.  R11 \textup{(SE2)} gives its relative envelope in
$\mathscr C$ and keeps it below the Borel upper constraint $f$.
Thus the envelope itself belongs to the family and is its
greatest member.  In particular $P(u)=u$ on the cone.
Energy monotonicity gives $r(P(f))\le r(a)$; infimizing over
decompositions proves the last assertion of \textup{(OC3)}.

Increasing the obstacle increases its competitor family.
Positive rescaling maps competitor families bijectively onto
each other, and $P(0)=0$ handles scale zero.  Also
$P(f)+P(g)\in\mathscr C$ lies below $f+g$, proving
superadditivity and hence concavity.

For the estimates write $f-g=a-b$, $a,b\in\mathscr C$, and
put $u=P(f)$, $v=P(g)$.  Since
\[
 u+b\preceq f+b=g+a\preceq g,\qquad
 v+a\preceq g+a=f+b\preceq f,
\]
both are legitimate cone competitors.  Consequently
\[
                        a\preceq u-v\preceq-b.
 \tag{OC5}
\]
In particular $u\succeq v+a$ and $v\succeq u+b$.
Energy monotonicity and Minkowski give
$r(u)\le r(v)+r(a)$ and $r(v)\le r(u)+r(b)$.
Infimizing proves the first estimate in \textup{(OC4)}.

The literal mixed integrals and \textup{(OC5)} give the sharper
signed estimate, without inserting an absolute-variation norm,
\[
 \begin{split}
 I(u,v)&=J(u)(u-v)-J(v)(u-v)\\
       &\le J(u)(a)+J(v)(b)\\
       &\le r(a)r(u)^q+r(b)r(v)^q
        \le R^q(r(a)+r(b)).
 \end{split}
 \tag{OC6}
\]
Taking the infimum gives the second estimate.  Combine it with
the proved integrated Cauchy--Schwarz bound
$\|J(u)-J(v)\|_*\le\sqrt q R^{(q-1)/2}I(u,v)^{1/2}$
to obtain the third.  The bound in \textup{(OC3)} controls
$R$ on a norm-convergent obstacle row.  All subsequent
strong current and weighted limits are therefore the full
original limits of R11, not only selected moment pairings.
\end{proof}

\begin{proposition}[A finite convex obstacle energy]
\label{prop:r12-convex-obstacle-energy}
The function $\mathfrak E:\mathbb D_S\to[0,\infty)$ is
finite, convex, order decreasing and positively homogeneous
of degree $d$.  It is locally Lipschitz, with the precise bound
\[
       |\mathfrak E(f)-\mathfrak E(g)|
           \le\max(r(P(f)),r(P(g)))^q\|f-g\|_{\rm ed}.
 \tag{OC7}
\]
Its restriction to $\mathscr C$ is the original first energy
$e/d$.  Its convexity is a statement on all of $\mathbb D_S$;
it does not declare the polynomial extension $B(f^{[d]})/d$
convex outside the potential cone.
\end{proposition}

\noindent\textbf{Status: \texttt{NEW\_PROVED}}.

\begin{proof}
Finiteness, order decrease and homogeneity follow from
\textup{(OC3)} and the corresponding cone energy properties.
For $0\le s\le1$, the order concavity of $P$ gives
\[
 P((1-s)f+sg)\succeq(1-s)P(f)+sP(g).
\]
Apply energy monotonicity, followed by convexity of $E$ on
the cone, to obtain convexity of $\mathfrak E$.  Finally,
$|x^d-y^d|/d\le\max(x,y)^{d-1}|x-y|$ for $x,y\ge0$.
Use the first line of \textup{(OC4)} and $d-1=q$ to obtain
\textup{(OC7)}.  The norm bound \textup{(OC3)} makes this a
local Lipschitz estimate on the entire Banach space.
\end{proof}

\subsection{Contact defects and exact subgradients}

Call a continuous linear functional $\ell\in\mathbb D_S^*$
\emph{order reversing} if $h\succeq0$ implies $\ell(h)\le0$.
This is order on Borel scalar differences, not an identification
of positivity cones for differential forms.  Every literal
positive measure datum $\ell_\mu(h)=-\int h\,d\mu$ with finite
integrals and a continuous extension to $\mathbb D_S$ has this
property.  Write
\[
 \partial\mathfrak E(f)=\{\ell\in\mathbb D_S^*:
 \mathfrak E(g)\ge\mathfrak E(f)+\ell(g-f)
                      \text{ for every }g\in\mathbb D_S\}
 \tag{OC8}
\]
for the Banach subdifferential, and retain the cone polar
$\mathscr C^\circ$ of \textup{(NI5)}.

\begin{definition}[Relative contact defect]
\label{def:r12-relative-contact-defect}
For $f\in\mathbb D_S$, with $u=P(f)$, put
\[
 \mathfrak Q(f)=\int_V(f-u)\,d\widehat\mu_u
                         =J(u)(u-f)\ge0.
 \tag{OC9}
\]
The integral is literal and finite: $f-u$ is a finite sum of
cone differences, and all its absolute mixed first-energy
integrals are finite.  Its nonnegativity uses $u\preceq f$
and the fact that the original completed measure ignores the
relative null set.  Vanishing means the relative Borel
contact identity
$(f-u)\widehat\mu_u=0$; no ordinary support regularity is
part of this definition.
\end{definition}

\begin{theorem}[The obstacle subdifferential and contact criterion]
\label{thm:r12-obstacle-subdifferential}
Let $f\in\mathbb D_S$ and $u=P(f)$.  A functional
$\ell\in\mathbb D_S^*$ belongs to $\partial\mathfrak E(f)$
if and only if all three conditions hold:
\[
 \begin{gathered}
 \ell\text{ is order reversing},\\
 \ell-J(u)\in\mathscr C^\circ,
                   \qquad (\ell-J(u))(u)=0,\\
                         \ell(f-u)=0.
 \end{gathered}
 \tag{OC10}
\]
Every such functional has $\|\ell\|_*=r(u)^q$.
In particular,
\[
       \boxed{\quad J(P(f))\in\partial\mathfrak E(f)
                        \quad\Longleftrightarrow\quad
                         \mathfrak Q(f)=0.\quad}
 \tag{OC11}
\]
For $u\in\mathscr C$, one always has
$J(u)\in\partial\mathfrak E(u)$.

If $\ell$ is additionally in the strong-continuous dual
$\mathbb D_{S,\mathrm s}^*$, the middle line of
\textup{(OC10)} is equivalent to $u=\Pi(\ell)$ by the
proved R12 normal-decomposition theorem.  Strong continuity
is not being asserted for every Banach subgradient.
\end{theorem}

\noindent\textbf{Status: \texttt{NEW\_PROVED}} for the
precise original arbitrary-background obstacle energy.  This is
a contact/variational characterization; it does not assert
that the contact defect vanishes.

\begin{proof}
Suppose first that $\ell\in\partial\mathfrak E(f)$.
For $h\succeq0$ and $t>0$, order decrease gives
$\mathfrak E(f+th)\le\mathfrak E(f)$.  The subgradient
inequality then gives $t\ell(h)\le0$, proving order reversal.
Since $P(u)=u$ and $E(u)=\mathfrak E(f)$, the competitor
$g=u$ in \textup{(OC8)} gives $\ell(u-f)\le0$.
But $u-f\preceq0$ and order reversal give the opposite
inequality.  Hence $\ell(f-u)=0$.

Restricting \textup{(OC8)} to $v\in\mathscr C$ now gives
\[
                        E(v)\ge E(u)+\ell(v-u).
 \tag{OC12}
\]
Along the cone segment $u+t(v-u)$ the polynomial energy
has derivative $J(u)(v-u)$ at zero.  Applying
\textup{(OC12)}, dividing by $t>0$ and letting $t\downarrow0$
gives $(\ell-J(u))(v-u)\le0$.  The competitors $0$ and
$2u$ yield $(\ell-J(u))(u)=0$, and then the same inequality
is the polar condition.  This proves necessity.

Conversely assume \textup{(OC10)}.  Cone convexity and its
first derivative give
\[
 E(v)\ge E(u)+J(u)(v-u)\ge E(u)+\ell(v-u)
                                     \quad(v\in\mathscr C).
\]
For arbitrary $g\in\mathbb D_S$ use $v=P(g)\preceq g$.
Order reversal gives $\ell(P(g)-u)\ge\ell(g-u)$, so
\[
 \mathfrak E(g)\ge E(u)+\ell(g-u)
                 =\mathfrak E(f)+\ell(g-f).
\]
The equality uses $\ell(f-u)=0$.  This proves sufficiency.

Order reversal makes $\ell(v)\ge0$ on the nonpositive cone.
The polar condition gives
$0\le\ell(v)\le J(u)(v)\le r(v)r(u)^q$.
The exact dual norm formula \textup{(ED10)} proves
$\|\ell\|_*\le r(u)^q$.  If $u\ne0$, evaluation on
$u/r(u)$ gives equality because $\ell(u)=e(u)$.
For $u=0$, the upper bound gives the same conclusion.
Finally $J(u)$ is order reversing, its normal residual is
zero, and its last condition in \textup{(OC10)} is exactly
$-\mathfrak Q(f)=0$.  This proves \textup{(OC11)} and
the assertion at cone points.  The final equivalence is
precisely the already proved normal-decomposition theorem
with its stated continuity gate.
\end{proof}

\begin{corollary}[Zero-mass slack does not move a cone envelope]
\label{cor:r12-zero-mass-obstacle-slack}
If $u\in\mathscr C$, $g\in\mathbb D_S$, $g\succeq0$, and
$\int g\,d\widehat\mu_u=0$, then
$P(u+g)=u$ and $\mathfrak Q(u+g)=0$.
\end{corollary}

\noindent\textbf{Status: \texttt{DIRECT\_COROLLARY}} of
the frozen ordered energy separation and the new obstacle definition.
It is recorded as a valid nontrivial contact mechanism, not as
unrestricted orthogonality.

\begin{proof}
Put $v=P(u+g)$.  Then $u\preceq v\preceq u+g$.
Thus $\int(v-u)\,d\widehat\mu_u=0$ and convexity gives
$E(v)\ge E(u)+J(u)(v-u)=E(u)$.  Energy monotonicity gives
the reverse inequality.  The frozen ordered discrepancy bound
$I(u,v)\le e(u)-e(v)$ makes $u=v$ as relative classes.
The contact defect is then the assumed zero integral.
\end{proof}

\subsection{Directional loss, source equations and failure of orthogonality}

For $u\in\mathscr C$ and $h\in\mathbb D_S$, define the
one-sided defect
\[
 \mathfrak d_u(h)=
 \lim_{t\downarrow0}
       \frac{\mathfrak E(u+th)-E(u)}{t}-J(u)(h).
 \tag{OC13}
\]
Convexity and local Lipschitz continuity make the limit finite.
Since $J(u)$ is a subgradient, the defect is nonnegative.
It is positively homogeneous and subadditive, being a convex
directional derivative minus a linear functional.

\begin{theorem}[A first-order contact gate and its exact loss]
\label{thm:r12-first-order-contact-gate}
For $u,h$ as above, put $v_t=P(u+th)$ and
$G_t=u+th-v_t\succeq0$ for $t>0$.  Then
\[
 \begin{split}
 &\mathfrak E(u+th)-E(u)-tJ(u)(h)\\
 &\quad=
 \underbrace{E(v_t)-E(u)-J(u)(v_t-u)}_{\ge0}
                       +\int_VG_t\,d\widehat\mu_u\ge0,
 \end{split}
 \tag{OC14}
\]
and
\[
 \mathfrak E(u+th)-E(u)-tJ(u)(h)
       \le t(J(v_t)-J(u))(h)+\mathfrak Q(u+th).
 \tag{OC15}
\]
Consequently
\[
 \mathfrak d_u(h)
       \le\liminf_{t\downarrow0}\frac{\mathfrak Q(u+th)}{t}.
 \tag{OC16}
\]

Let $\mu$ be a positive Radon datum whose literal integration
functional $\ell_\mu\in\mathbb D_{S,\mathrm s}^*$ is defined
on every energy difference, and let $u=\Pi(\ell_\mu)$.
If, for every original compact smooth scalar $\phi$ and both
signs, one has
\[
                  \mathfrak Q(u\pm t\phi)=o(t)
                                     \quad(t\downarrow0),
 \tag{OC17}
\]
then $\widehat\mu_u=\mu$ on every original compact smooth
scalar test, after every actual moment refinement.
Condition \textup{(OC17)} is a sufficient first-order contact
criterion, not a condition proved for general obstacles or data.
\end{theorem}

\noindent\textbf{Status: \texttt{NEW\_PROVED}} for the
defect identities and \texttt{CONDITIONAL} for the source
conclusion under the explicit geometric contact criterion.
This condition is stated using the independently constructed
envelope and its own curvature, not by assuming the desired
source equality.

\begin{proof}
All integrals are finite by the mixed first-energy inequality.
Subtract $tJ(u)(h)$ from the convex tangent remainder at $u$.
Since $v_t-u-th=-G_t$ and $J(u)=-\int(\cdot)\,d\widehat\mu_u$,
this gives exactly \textup{(OC14)}.  The tangent inequality
at $v_t$ instead gives
$E(v_t)-E(u)\le J(v_t)(v_t-u)$.
Substitute $v_t-u=th-G_t$ to obtain \textup{(OC15)}.
The third estimate of \textup{(OC4)} gives
$J(v_t)\to J(u)$ in the actual dual norm.  Divide
\textup{(OC15)} by $t$ and let $t\downarrow0$; the error
$(J(v_t)-J(u))(h)$ tends to zero.  This proves
\textup{(OC16)}, including the possibility that its right
side is infinite.

For the datum assertion, order reversal and the normal
decomposition give $\ell_\mu\in\partial\mathfrak E(u)$
by \textup{(OC10)}.  Thus
\[
 \mathfrak d_u(h)\ge(\ell_\mu-J(u))(h)
                                    \quad(h\in\mathbb D_S).
 \tag{OC18}
\]
Every original compact smooth $\phi$ belongs to $\mathbb D_S$
by the exact all-degree finite-mass difference construction
R11 \textup{(HC14)--(HC15)}.  Assumption \textup{(OC17)},
\textup{(OC16)} and nonnegativity make both
$\mathfrak d_u(\phi)$ and $\mathfrak d_u(-\phi)$ zero.
Apply \textup{(OC18)} in both signs to get
$\ell_\mu(\phi)=J(u)(\phi)$.  The literal integral
definitions give the original source identity on every
such test.  The choice of $\phi$ was unrestricted in the
original sheaf, so actual refinements are included.
\end{proof}

\begin{corollary}[A nonzero variational residual forces an obstacle kink]
\label{cor:r12-normal-residual-obstacle-kink}
For the positive datum and maximizer of the preceding theorem,
put $n=\ell_\mu-J(u)$.  For every $h\in\mathbb D_S$ the
right and left directional slopes of $\mathfrak E$ at $u$
have a gap at least $|n(h)|$.  If $n(h)>0$, then
\[
 \mathfrak d_u(h)\ge n(h)>0,
 \qquad
 \liminf_{t\downarrow0}\frac{\mathfrak Q(u+th)}{t}\ge n(h).
 \tag{OC19}
\]
If the source equation fails on an original compact smooth
test, the direction $h$ can be chosen to be such a test,
with its sign changed if necessary.  In that case universal
relative obstacle orthogonality and two-sided differentiability
of this obstacle energy at the maximizer are both false.
\end{corollary}

\noindent\textbf{Status: \texttt{NEW\_PROVED}} as a
quantitative contact interpretation of the normal residual.

\begin{proof}
Both $J(u)$ and $\ell_\mu$ are subgradients at $u$.
The right slope in direction $h$ is at least the larger of
$J(u)(h)$ and $\ell_\mu(h)$, whereas the left slope is at
most their smaller value.  Their difference is at least
$|n(h)|$.  Equation \textup{(OC18)} gives the first bound
in \textup{(OC19)}, and \textup{(OC16)} gives the second.
Distinct original Radon measures are distinguished by an
original compact smooth scalar test; the exact test-difference
construction places that direction in $\mathbb D_S$.
No ordinary upper envelope or stronger delta-test equality
is used in this separation.
\end{proof}

\paragraph{Which proposed classical step fails.}
The map $P$ exists on every energy difference, its energy is
finite and convex, and the original curvature is strongly
continuous.  These facts alone do not make the curvature a
subgradient at the \emph{obstacle} $f$: by \textup{(OC11)},
the missing statement is precisely its nonnegative contact
defect.  The new all-degree compact-source examples force
that defect to persist at first order along an original
smooth direction through a genuine variational maximizer.
Thus the unrestricted contact step is \texttt{FALSE} in
this classical-approximation strong class.  The qualified
criterion \textup{(OC17)} remains useful on a separately
proved obstacle/data subclass; no such universal subclass
identification is claimed.

% END REVIEWED MODULE research_r12/nonlinear/energy_obstacles_and_contact_defects.tex


% BEGIN REVIEWED MODULE research_r12/nonlinear/fenchel_order_gauge.tex
\section{The exact Fenchel dual and the order-energy gauge}
\label{sec:r12-Fenchel-order-gauge}

Retain all of \textup{(ED1)} and the arbitrary original $S$.
Let $P$, $\mathfrak E$, $\mathscr C$, $\mathbb D_S$,
$r=e^{1/d}$ and $d=q+1$ be exactly the objects of
Section~\ref{sec:r12-energy-obstacles}.  In particular no
ordinary upper regularization or larger relative potential
cone has been substituted for the R11 strong completion.

Define the closed order-dual cone
\[
 \mathscr O_S=\{\ell\in\mathbb D_S^*:
                  h\succeq0\Longrightarrow\ell(h)\le0\},
 \qquad m(f)=r(P(f)).
 \tag{FO1}
\]
The order dual is not the polar $\mathscr C^\circ$ of the
potential cone: its sign is imposed on all nonnegative Borel
differences, not only on nonpositive potentials.

\begin{theorem}[An attained dual formula for the obstacle gauge]
\label{thm:r12-order-gauge-duality}
The function $m$ is nonnegative, sublinear, order decreasing
and $1$-Lipschitz for $\|\cdot\|_{\rm ed}$.  It has the
exact primal and dual descriptions
\[
 \begin{split}
 m(f)&=\min\{r(a):a\in\mathscr C,\ a\preceq f\},\\
 m(f)&=\max\{\ell(f):\ell\in\mathscr O_S,
                                      \|\ell\|_*\le1\}.
 \end{split}
 \tag{FO2}
\]
The first minimum is attained by $P(f)$, and the second
maximum is attained for every $f$.  Thus every obstacle has
an order-reversing supporting functional, although its own
curvature need not be one at that obstacle.
\end{theorem}

\noindent\textbf{Status: \texttt{NEW\_PROVED}} for this
exact obstacle duality.  The real Hahn--Banach theorem for
a continuous sublinear majorant is \texttt{EXTERNAL\_KNOWN};
the application and all energy identifications are proved below.

\begin{proof}
The greatest-minorant property of $P$ and cone energy
monotonicity prove the first formula.  Superadditivity of
$P$, followed by energy monotonicity and Minkowski, gives
\[
 m(f+g)=r(P(f+g))\le r(P(f)+P(g))\le m(f)+m(g).
\]
Positive homogeneity, nonnegativity, order decrease and the
Lipschitz bound have already been proved in \textup{(OC3)--(OC4)}.
Also $m(h)=0$ whenever $h\succeq0$, since the zero
potential is then the greatest cone minorant.

For $\ell\in\mathscr O_S$ with $\|\ell\|_*\le1$,
order reversal and $P(f)\preceq f$ give
$\ell(f)\le\ell(P(f))\le r(P(f))=m(f)$.
For the reverse direction, if $f\ne0$, define a linear
functional on $\mathbf R f$ by $L(tf)=t m(f)$.
It is dominated by $m$: for $t\ge0$ this is equality,
and for $t<0$ it follows from $-m(f)\le m(-f)$.
Hahn--Banach extends it to a real linear $\ell$ on
$\mathbb D_S$ with $\ell(h)\le m(h)$ for every $h$.
Since $m(\pm h)\le\|h\|_{\rm ed}$, this functional
is continuous with norm at most one.  For $h\succeq0$,
$m(h)=0$ makes it order reversing.  Finally
$\ell(f)=m(f)$.  When $f=0$, the zero functional suffices.
This proves the attained dual formula without invoking
an unidentified measure representation of $\ell$.
\end{proof}

\begin{theorem}[Exact Fenchel conjugate]
\label{thm:r12-obstacle-Fenchel-conjugate}
For $\ell\in\mathbb D_S^*$, define
\[
 \mathfrak E^*(\ell)=
       \sup_{f\in\mathbb D_S}\{\ell(f)-\mathfrak E(f)\}.
\]
Then
\[
 \boxed{\quad
 \mathfrak E^*(\ell)=
 \begin{cases}
 \displaystyle\frac qd\|\ell\|_*^{d/q},&\ell\in\mathscr O_S,\\
 +\infty,&\ell\notin\mathscr O_S.
 \end{cases}\quad}
 \tag{FO3}
\]
If $\ell\in\mathscr O_S\cap\mathbb D_{S,\mathrm s}^*$,
the supremum is attained.  Put $u=\Pi(\ell)$.  Its complete
set of maximizers is
\[
 \{f\in\mathbb D_S:P(f)=u,\ \ell(f-u)=0\}.
 \tag{FO4}
\]
In particular the potential maximizer $u$ is unique, even
when the full obstacle maximizer need not be unique.
For arbitrary $\ell\in\mathscr O_S$, formula
\textup{(FO3)} is a value formula and does not silently
assert attainment or strong continuity.
\end{theorem}

\noindent\textbf{Status: \texttt{NEW\_PROVED}} for the
full Banach-space conjugate and its exact order domain.

\begin{proof}
If $\ell$ is not order reversing, choose $h\succeq0$
with $\ell(h)>0$.  Then $P(th)=0$ for every $t\ge0$,
so the conjugate supremum is infinite.

Suppose now that $\ell\in\mathscr O_S$.  For every
$f\in\mathbb D_S$,
\[
 \ell(f)-\mathfrak E(f)
       \le\ell(P(f))-E(P(f)).
 \tag{FO5}
\]
Conversely cone members are obstacles fixed by $P$.  Hence
the conjugate supremum is exactly the supremum of
$\ell(v)-r(v)^d/d$ over $v\in\mathscr C$.
Order reversal makes $\ell$ nonnegative on this nonpositive
cone, and the exact dual norm formula \textup{(ED10)} says
\[
 \|\ell\|_*=
          \sup_{v\in\mathscr C,\ r(v)\le1}\ell(v).
\]
Optimize the scalar expression $at-t^d/d$ over $t\ge0$:
its supremum is $q a^{d/q}/d$.  Applying this to cone
directions approaching the displayed support supremum,
and using the matching upper bound, proves \textup{(FO3)}.
It remains valid when the norm is zero.

Under the additional strong-dual hypothesis, the proved
normal-decomposition theorem gives the unique cone
maximizer $u$.  Equality in \textup{(FO5)} requires
$\ell(f-P(f))=0$, while its right side attains the
supremum precisely when $P(f)=u$.  These conditions are
also sufficient.  This proves \textup{(FO4)} and the
attainment assertion with exactly its stated hypothesis.
\end{proof}

\begin{corollary}[Supporting order functionals at every obstacle]
\label{cor:r12-obstacle-support-existence}
For every $f\in\mathbb D_S$, the subdifferential
$\partial\mathfrak E(f)$ is nonempty.  If $m(f)>0$,
choose a maximizer $\lambda$ in the second line of
\textup{(FO2)}.  Then
\[
          \ell=m(f)^q\lambda\in\partial\mathfrak E(f),
          \qquad\|\ell\|_*=m(f)^q.
 \tag{FO6}
\]
If $m(f)=0$, the zero functional is a subgradient.
Every such functional satisfies the exact normal and contact
conditions \textup{(OC10)}.  Neither literal positive
Radon representation on all completed Borel potentials
nor membership in the strong-continuous dual is inferred
solely from this supporting-functional construction.
\end{corollary}

\noindent\textbf{Status: \texttt{DIRECT\_COROLLARY}}.

\begin{proof}
For $m(f)>0$, the convex scalar function $s^d/d$ gives
\[
 \mathfrak E(g)\ge\mathfrak E(f)
                 +m(f)^q(m(g)-m(f)).
\]
Since $\lambda(g)\le m(g)$ and $\lambda(f)=m(f)$,
this proves the subgradient inequality for \textup{(FO6)}.
The norm follows from \textup{(OC10)} or from evaluation
on $P(f)/m(f)$.  At zero gauge the nonnegativity of
$\mathfrak E$ makes zero a subgradient.  The qualifications
follow because no continuity stronger than the Banach norm
was used in Hahn--Banach.
\end{proof}

\subsection{A lattice norm, kept distinct from the original difference norm}

The Borel difference space is a vector lattice.  Indeed if
$f=a-b$ and $g=c-e$, with all four terms in $\mathscr C$,
then
\[
 \max(f,g)=\max(a+e,c+b)-(b+e)\in\mathbb D_S.
 \tag{FO7}
\]
In particular $|f|\in\mathbb D_S$.  This algebraic
lattice property does not make $\|\cdot\|_{\rm ed}$
a lattice norm.

\begin{definition}[Order-energy norm]
\label{def:r12-order-energy-norm}
For $f\in\mathbb D_S$ put
\[
               \|f\|_{\rm o}=m(-|f|)=r(P(-|f|)).
 \tag{FO8}
\]
This norm measures the least energy root of a cone
minorant of the negative absolute value, not the least
sum of roots in an exact difference decomposition.
\end{definition}

\begin{proposition}[The precise order-norm bounds]
\label{prop:r12-order-energy-norm}
Formula \textup{(FO8)} is a lattice norm, and
\[
 \|f\|_{\rm o}\le\|f\|_{\rm ed},\qquad
 \|u\|_{\rm o}=r(u)\quad(u\in\mathscr C).
 \tag{FO9}
\]
For every $\ell\in\mathscr O_S$,
\[
                |\ell(f)|\le\|\ell\|_*\|f\|_{\rm o}.
 \tag{FO10}
\]
Its operator norm for the order norm equals its operator
norm for the difference norm.  For $u,v\in\mathscr C$,
\[
 \begin{split}
 |r(u)-r(v)|&\le\|u-v\|_{\rm o},\\
 I(u,v)&\le(r(u)^q+r(v)^q)\|u-v\|_{\rm o}.
 \end{split}
 \tag{FO11}
\]
Finally, for $f\in\mathbb D_S$ and $t>0$,
\[
 C_S^*(\{|f|>t\};V)
                  \le C_0t^{-d}\|f\|_{\rm o}^d,
                         \qquad C_0=d2^d.
 \tag{FO12}
\]
No converse strong-to-order-norm estimate is asserted.
\end{proposition}

\noindent\textbf{Status: \texttt{NEW\_PROVED}}.

\begin{proof}
The gauge is homogeneous and order decreasing.  The
inequality $-|f+g|\succeq-|f|-|g|$ and subadditivity
of $m$ give the triangle inequality.  Symmetry is clear.
If the norm is zero, $P(-|f|)$ has zero energy and is
the zero relative class.  Since it is below $-|f|$,
this forces $f=0$.  If $|f|\preceq|g|$, the order
property gives $\|f\|_{\rm o}\le\|g\|_{\rm o}$;
thus the norm is a lattice norm.

If $f=a-b$, then $|f|\preceq-a-b$, so
$P(-|f|)\succeq a+b$.  Consequently
$\|f\|_{\rm o}\le r(a+b)\le r(a)+r(b)$.
Infimizing proves the first assertion of \textup{(FO9)}.
The second follows because $u\preceq0$ and $P(u)=u$.

For an order-reversing $\ell$,
$|\ell(f)|\le\ell(-|f|)$ by the two pointwise
inequalities $-|f|\preceq\pm f$.  Formula
\textup{(FO2)}, applied to $-|f|$, proves
\textup{(FO10)}.  Testing on unit cone members and
using \textup{(ED10)} gives equality of the two
operator norms for these order-reversing functionals.

Put $w=P(-|u-v|)$.  Then $u\succeq v+w$ and
$v\succeq u+w$, yielding the root estimate by
energy monotonicity and Minkowski.  Also
\[
 I(u,v)\le\int|u-v|(d\widehat\mu_u+d\widehat\mu_v)
       \le B(w,u^{[q]})+B(w,v^{[q]}),
\]
which gives the second line of \textup{(FO11)}.
For arbitrary $f$, the inequality $|f|\preceq-w$
with $w=P(-|f|)$ gives
$\{|f|>t\}\subset\{w<-t\}$ modulo an outer-null
set.  The frozen first-energy sublevel bound proves
\textup{(FO12)}.
\end{proof}

\paragraph{The positive-dual gate is false in general.}
Theorem~\ref{thm:r12-pole-order-dual-not-strong} proves
that an order-reversing functional in $\mathbb D_S^*$
need not belong to $\mathbb D_{S,\mathrm s}^*$ under
the present general hypotheses.  The stronger implication
\[
 I(u_j,u)\longrightarrow0
       \quad\Longrightarrow\quad
 r\bigl(P(-|u_j-u|)\bigr)\longrightarrow0.
 \tag{FO13}
\]
is also \texttt{FALSE}, by
Theorem~\ref{thm:r12-strong-not-order-topology}; only the
reverse implication follows from \textup{(FO11)}.
Both failures occur in every remaining degree, with
uniform bounds in every fixed finite power class.
The new counterexample has positivity on the full
Borel order cone, not merely on $\mathscr C$.
Its restriction to all original compact smooth scalar
tests is zero but its value on an unbounded actual
finite-energy potential is one.  Thus positivity on
compact original tests does not give literal integration
on all relative Borel representatives.  None of these
false passages is used in \textup{(FO1)--(FO12)}.
The earlier \texttt{OPEN} formulation is retained in
\path{fenchel_order_gauge_before_pole_counterexample.tex}
as a dated failed research gate, not as current status.

% END REVIEWED MODULE research_r12/nonlinear/fenchel_order_gauge.tex


% BEGIN REVIEWED MODULE research_r12/nonlinear/order_energy_lattice_completion.tex
\section{An actual Borel order-energy lattice and original current multipliers}
\label{sec:r12-order-energy-lattice}

Retain exactly \textup{(ED1)} and the arbitrary original
middle-positive separately closed background $S$.  The
potential cone remains the frozen
$\mathscr C=\widehat{\mathcal E}^1_{\rm str}(V)$.
We now complete its difference space for the distinct lattice
norm \textup{(FO8)}.  This introduces a larger \emph{scalar
multiplier space}, not a larger potential cone.  Its Borel
realization and the closedness of the old cone are proved
rather than assumed.

\begin{definition}[The order-energy lattice]
\label{def:r12-order-energy-lattice}
Let $\mathbb L_S$ be the norm completion of
$(\mathbb D_S,\|\cdot\|_{\rm o})$.
\end{definition}

The notation $P(f)$ will also be used for a Borel obstacle
$f$ not initially known to lie in $\mathbb D_S$, whenever
at least one member of $\mathscr C$ lies below $f$.
The R11 relative Perron theorem then gives exactly the same
greatest-minorant construction.  This is an explicitly
nonempty-family use of the frozen theorem.

\begin{theorem}[Borel realization and an exact majorant norm]
\label{thm:r12-Borel-order-lattice}
The completion $\mathbb L_S$ is a Banach lattice whose
elements have canonical Borel classes, finite outside a
locally outer $S$-capacity-null set.  These classes are
injectively determined by the completion elements.  Lattice
operations agree with the actual Borel operations, and
\[
 \boxed{\qquad
 \|f\|_{\mathbb L_S}
    =r(P(-|f|))
    =\min\{r(a):a\in\mathscr C,\ a\preceq-|f|\}.
 \qquad}
 \tag{OL1}
\]
In particular every member admits an actual cone majorant
of its absolute value.  The capacity estimate
\textup{(FO12)} holds with this norm.  The original potential
cone $\mathscr C$ is closed in $\mathbb L_S$, and its norm
is still exactly $r$.

More precisely, each $f\in\mathbb L_S$ has a sequence
$f_j\in\mathbb D_S$ and cone functions $W_j\in\mathscr C$
such that
\[
 |f-f_j|\preceq-W_j,\qquad r(W_j)\longrightarrow0.
 \tag{OL2}
\]
The construction does not assert that every Borel function
with some finite-energy cone majorant belongs to $\mathbb L_S$.
\end{theorem}

\noindent\textbf{Status: \texttt{NEW\_PROVED}} as an
identified Borel Banach lattice over the unchanged full
strong potential cone.  The decreasing-series theorem R11
\textup{(CI8)} is an \texttt{INTERNAL\_INPUT}.

\begin{proof}
Take a Cauchy row in the order norm and pass to a subsequence,
still denoted $(f_j)$, with
\[
                 \sum_j\|f_{j+1}-f_j\|_{\rm o}<\infty.
\]
Put $w_j=P(-|f_{j+1}-f_j|)\in\mathscr C$.  Then
$|f_{j+1}-f_j|\preceq-w_j$ and
$\sum_jr(w_j)<\infty$.  The partial sums of the nonpositive
$w_j$ decrease, and Minkowski bounds their energy roots.
R11 \textup{(CI8)} gives their Borel cone sum and every
tail sum
\[
 W_j=\sum_{l\ge j}w_l\in\mathscr C,\qquad
 r(W_j)\le\sum_{l\ge j}r(w_l)\longrightarrow0.
 \tag{OL3}
\]
Outside one countable union of relative null sets, this is
an ordinary series of nonpositive numbers with finite sum.
Thus the series of successive absolute differences is
pointwise convergent there.  It defines a finite Borel
limit $f$ and gives \textup{(OL2)}.  The same argument
identifies the limit in local capacity for the entire
original Cauchy row, by \textup{(FO12)} and the Cauchy
property.

If this limit is the zero Borel class, then
$|f_j|\preceq-W_j$.  The greatest-minorant property gives
$\|f_j\|_{\rm o}\le r(W_j)\to0$.  Thus the completion
element itself is zero.  Applying this observation to
the difference of two Cauchy rows proves both independence
of choices and injectivity of the Borel realization.  This
step excludes an abstract norm vector invisible to Borel
capacity.

Put $a_j=P(-|f_j|)$.  Since
$|f|\preceq-a_j-W_j$, the Borel obstacle $-|f|$ has
a cone minorant $a_j+W_j$.  Its envelope $a=P(-|f|)$
therefore exists, and
\[
 r(a)\le r(a_j+W_j)\le\|f_j\|_{\rm o}+r(W_j).
 \tag{OL4}
\]
Conversely $|f_j|\preceq|f|+|f_j-f|\preceq-a-W_j$,
so
\[
                  \|f_j\|_{\rm o}\le r(a)+r(W_j).
 \tag{OL5}
\]
Passing to the limit proves \textup{(OL1)}.  The minimum
statement follows from energy monotonicity and the greatest
minorant property.  The same majorant gives
\textup{(FO12)} for the completed element.

The pointwise inequalities
$||f|-|g||\le|f-g|$ and the corresponding inequalities
for maxima show that lattice operations are continuous
in the order norm on $\mathbb D_S$.  They extend to the
completion and agree with the Borel limits just constructed.
Formula \textup{(OL1)} also proves directly that the extended
norm is monotone in absolute Borel order.  Thus the completion
is a Banach lattice with the asserted concrete realization.

Finally let a cone row $u_j$ converge in $\mathbb L_S$.
Its roots $r(u_j)=\|u_j\|_{\mathbb L_S}$ are bounded.
The second line of \textup{(FO11)} makes it strongly
Cauchy, and the frozen strong completeness supplies a
limit $u\in\mathscr C$.  Both its strong limit and its
order-norm limit have the same local capacity class.
Injectivity of the Borel realization identifies them.
This proves closedness of $\mathscr C$, without declaring
it equal to every nonpositive member of the larger lattice.
\end{proof}

\begin{proposition}[Relative quasicontinuity in the lattice]
\label{prop:r12-lattice-quasicontinuity}
Every member of $\mathbb L_S$ has a finite representative
which is locally ordinary-open outer $S$-quasicontinuous,
with the same fixed-domain outer capacity as the original
completion.  This does not identify raw and outer capacity
on arbitrary Borel sets.
\end{proposition}

\noindent\textbf{Status: \texttt{DIRECT\_COROLLARY}} of
\textup{(OL2)--(OL3)}, the capacity estimate and the frozen
quasicontinuity of strong cone members.

\begin{proof}
In constructing \textup{(OL2)}, pass to a faster subsequence
so that
$C_0 2^{jd}r(W_j)^d\le2^{-2j}$.
The sets where $|f-f_j|>2^{-j}$ have outer capacity at
most $2^{-2j}$.  Cover a sufficiently late tail of these
sets by ordinary open sets with arbitrarily small total
capacity.  Outside their union the approximants converge
uniformly.  Each $f_j$ is a difference of two strong cone
members and is therefore finite and continuous off a
further ordinary open set of arbitrarily small capacity
on the fixed compact core.  Choose and sum these exceptional
capacities over the countably many approximants.  The uniform
limit is finite and continuous off their union.  A countable
compact exhaustion gives the stated local result.  Only
outer capacity subadditivity and its defining open covers
are used.
\end{proof}

\subsection{The variation norm is the dual norm of this lattice}

\begin{theorem}[Literal curvature integration and weighted total variation]
\label{thm:r12-lattice-dual-variation}
For each $u\in\mathscr C$, the functional $J(u)$ extends
uniquely to $\mathbb L_S$ and satisfies
\[
 J(u)(f)=-\int_V f\,d\widehat\mu_u,\qquad
 \|J(u)\|_{\mathbb L_S^*}=r(u)^q.
 \tag{OL6}
\]
The integral is literal and absolutely convergent for every
completed lattice element.  If $u,v\in\mathscr C$ and
$\sigma=\widehat\mu_u-\widehat\mu_v$, then
\[
 \begin{split}
 \|J(u)-J(v)\|_{\mathbb L_S^*}
   &=\sup_{\substack{a\in\mathscr C\\r(a)\le1}}
                               \int_V(-a)\,d|\sigma|\\
   &=\delta_1(\widehat\mu_u,\widehat\mu_v).
 \end{split}
 \tag{OL7}
\]
Here $\delta_1$ uses the R11 weighted absolute-variation
formula, not the new signed operator distance $\delta_*$
on $\mathbb D_S^*$.  In \textup{(OL7)} the supremum is
unchanged if it is restricted to actual first-energy
potentials, as in the frozen formula.  Applying that
formula to these general completed source measures is
not an import of R11's more restricted dominated-data
stability theorem.
\end{theorem}

\noindent\textbf{Status: \texttt{NEW\_PROVED}} as a
norm-dual interpretation of the two distinct source distances.

\begin{proof}
Since $J(u)$ is order reversing, \textup{(FO10)} gives its
unique continuous extension, with the same norm.  Let
$f\in\mathbb L_S$ and $a=P(-|f|)$.  Then
\[
 \int|f|\,d\widehat\mu_u
       \le\int(-a)\,d\widehat\mu_u
       =B(a,u^{[q]})\le\|f\|_{\mathbb L_S}r(u)^q.
 \tag{OL8}
\]
All relative null sets are ignored by the original curvature
measure.  The approximants in \textup{(OL2)} have integral
error at most $r(W_j)r(u)^q\to0$, proving literal agreement
with the functional extension.  Testing on $u/r(u)$ gives
the norm equality, with the zero case interpreted separately.

The difference $\sigma$ is a locally finite signed
Radon measure; its global total variation may be infinite.
All weighted variation integrals below are finite.
Put $\delta$ equal to the supremum in \textup{(OL7)}.
It is finite, at most $r(u)^q+r(v)^q$.
For any $f\in\mathbb L_S$ of norm at most one, its
minorant $a=P(-|f|)$ has $r(a)\le1$.  Therefore
\[
 \left|\int f\,d\sigma\right|
          \le\int|f|\,d|\sigma|
          \le\int(-a)\,d|\sigma|\le\delta.
 \tag{OL9}
\]
This is one inequality in \textup{(OL7)}.

For the reverse fix $a\in\mathscr C$, $r(a)\le1$.
The finite positive measure $(-a)|\sigma|$ ignores the
joint pole-null set.  By Radon regularity choose disjoint
compact subsets of a Hahn positive and negative carrier
that omit arbitrarily little of this finite weighted measure.
Choose an original compact smooth scalar $\chi$ equal to
$1$ on the first compact and $-1$ on the second, with
$|\chi|\le1$.  The original smooth cutoff and finite
ordinary-open gluing calculus supplies such a function.
For $M>0$ define the lattice clipping
\[
                     f_M=\max\bigl(a,\min(-a,M\chi)\bigr).
 \tag{OL10}
\]
Every entry is in $\mathbb D_S$, including $M\chi$ by
R11 \textup{(HC14)--(HC15)}.  Hence $f_M\in\mathbb D_S$.
It satisfies $|f_M|\preceq-a$, so its order norm is at
most $r(a)\le1$.  On the two chosen compact sets,
$f_M$ tends respectively to $-a$ and $a$ as $M\to\infty$,
outside the harmless pole set.  The error everywhere is
bounded by $2(-a)$, which is integrable for $|\sigma|$.
First let $M\to\infty$ and then let the omitted weighted
mass tend to zero.  The pairings $\int f_M\,d\sigma$
approach $\int(-a)\,d|\sigma|$ from within arbitrarily
small error.  Taking the supremum over $a$ proves the
reverse inequality.  No product $a\chi$ or density of
background currents was required.

Finally, to compare the tester class with the frozen
absolute-variation formula, choose for any $a\in\mathscr C$
an actual strong defining row $a_j$.  The completed tester
estimate \textup{(EV1)} gives
\[
 \int|a_j-a|\,d|\sigma|
  \le\int|a_j-a|\,(d\widehat\mu_u+d\widehat\mu_v)
                                      \longrightarrow0.
\]
Their energy roots converge.  If $r(a)\le1$, replacing
$a_j$ by $a_j/\max(1,r(a_j))$ gives actual unit-root
testers with the same limiting weighted variation
pairing.  When $r(a)=0$, its pairing is zero.  The
actual and completed suprema therefore agree.  This
argument needs neither a common actual lower barrier
nor a common actual dominating source measure.
\end{proof}

\paragraph{Topology qualification.}
The positive functionals $J(u)$ have the same individual
norm in the two dual spaces, but differences of them need
not.  In the original annulus example \textup{(CG1)--(CG16)},
$J(u_s)\to J(u_0)$ in the signed $\mathbb D_S^*$ norm,
while their $\mathbb L_S^*$ distance has lower limit at
least $\sqrt2$.  Thus strong curvature continuity must not
be upgraded to lattice-dual norm continuity.  Moreover,
Theorems~\ref{thm:r12-pole-order-dual-not-strong}
and~\ref{thm:r12-strong-not-order-topology} show that
the positive-dual continuity gate and \textup{(FO13)}
are both \texttt{FALSE}: the order lattice is strictly
stronger than the strong potential topology, and it has
nonzero positive functionals invisible to every original
compact smooth scalar test.  The multiplier conclusions
below need only order-norm convergence and are unaffected.

\subsection{Multipliers of every completed mixed original current}

\begin{lemma}[A finite-energy trace-completing potential]
\label{lem:r12-trace-completing-potential}
Let $0\le k\le q$ and let $\eta$ be an original compact
smooth complementary $(q-k,q-k)$-test form for an original
middle-positive $(p+k,p+k)$-current.  There are a bounded
actual classical $\psi\in\mathscr C$ of finite total top
mass and a finite constant $C_\eta$ such that, for every
tuple $u_1,\ldots,u_k\in\mathscr C$, the coefficient
measure of $T=\widehat N_k(u_1,\ldots,u_k)$ satisfies
\[
 |\mu_{T,\eta}|
        \le C_\eta\,
             \widehat N_q(u_1,\ldots,u_k,\psi^{[q-k]}).
 \tag{OL11}
\]
For $k=q$, one may omit $\psi$ and take
$C_\eta=\|\eta\|_\infty$ for the scalar test.
The domination is local on the coefficient support and
is an inequality of positive Radon measures, not a
closedness assertion for their Borel restrictions.
\end{lemma}

\noindent\textbf{Status: \texttt{DIRECT\_COROLLARY}} of
the frozen original middle-coefficient trace bound, analytic
compensators R11 \textup{(AE2)--(AE5)}, and finite-mass
collar construction \textup{(TC7)--(TC8)}.

\begin{proof}
Use finitely many original moment stages and a compact
smooth partition near the support of $\eta$.  On each
stage a fixed multiple of the full moment trace power
dominates every coefficient of that test in the middle
matrix convention.  The original analytic compensator
lemma gives a global smooth classical psh function whose
curvature dominates this stage's positive $(1,1)$ trace
near the relevant compact piece.  Sum the finitely many
compensators, obtaining one smooth classical $g$ with
all these local domination properties.  Choose $C$ with
$g-C<-1$ on $V$, and choose $A$ large enough that
\[
                         \psi=\max(g-C,A\rho)
\]
equals $g-C$ near the whole test support.  Here $\rho$
is the strict-box barrier.  The exact collar construction
proves that $\psi$ is bounded, actual, zero-boundary in
the frozen energy sense, and has finite total top mass.
Thus it belongs to every positive power class.

On each compact test stage, $D\psi=Dg$ dominates its
full moment trace.  Taking powers preserves the necessary
strong-form domination; the middle-positive coefficient
matrix bound then controls $|\mu_{T,\eta}|$ by a
constant times $(D\psi)^{q-k}\wedge T$ there.
Original smooth-factor compatibility identifies the latter
with the restriction of the mixed top measure in
\textup{(OL11)}.  Sum the finitely many stage bounds.
All moment stages may include any specified actual test
refinement, so this is a complete original-test assertion.
\end{proof}

\begin{theorem}[Order-energy multipliers and closed distributional cap currents]
\label{thm:r12-lattice-current-multipliers}
Let $0\le k\le q$, $u_1,\ldots,u_k\in\mathscr C$, and
$T=\widehat N_k(u_1,\ldots,u_k)$, with $T=S$ when $k=0$.
Every $f\in\mathbb L_S$ defines a literal order-zero
original coefficient product $fT$.  For every compact
original complementary test $\eta$ and the potential
$\psi$ from Lemma~\ref{lem:r12-trace-completing-potential},
\[
 \int_V|f|\,d|\mu_{T,\eta}|
       \le C_\eta\|f\|_{\mathbb L_S}
                   \prod_{i=1}^k r(u_i)\,r(\psi)^{q-k}.
 \tag{OL12}
\]
For $k=q$, the last factor is one.  Thus order-norm
convergence of multipliers gives local coefficient-total-
variation convergence against every original test.

If $k<q$, the formula
\[
                 \mathcal D_T(f)=D(fT)
 \tag{OL13}
\]
defines a linear continuous original current-valued map
whose values are separately $d'$- and $d''$-closed.
For $f\in\mathscr C$ it is exactly the frozen positive
BT current
$\widehat N_{k+1}(f,u_1,\ldots,u_k)$.
For general $f\in\mathbb L_S$, neither positivity nor
order zero of \textup{(OL13)} is asserted.
Every equality and convergence is on the full original
compact smooth complementary test sheaf and every actual
moment refinement.
\end{theorem}

\noindent\textbf{Status: \texttt{NEW\_PROVED}} for the
new completed scalar multiplier domain and its continuous
closed-current action, in every dimension and for arbitrary
original $S$.  The already constructed BT products are
\texttt{INTERNAL\_INPUT}, not newly defined by this theorem.

\begin{proof}
For $f\in\mathbb L_S$ put $a=P(-|f|)$, so
$|f|\preceq-a$ and $r(a)=\|f\|_{\mathbb L_S}$.
The coefficient domination \textup{(OL11)} and the
literal mixed-energy formula give
\[
 \begin{split}
 \int|f|\,d|\mu_{T,\eta}|
 &\le C_\eta\int(-a)\,
            d\widehat N_q(u_1,\ldots,u_k,\psi^{[q-k]})\\
 &=C_\eta B(a,u_1,\ldots,u_k,\psi^{[q-k]})\\
 &\le C_\eta r(a)\prod_i r(u_i)\,r(\psi)^{q-k}.
 \end{split}
\]
The top mixed measure ignores every relative null set;
the coefficient domination therefore shows that the
same is true of the coefficient measure.  Hence the
product is independent of the Borel representative of
$f$.  These coefficient products agree on overlaps by
their ordinary measure definition, and their local
bounds give an original order-zero current $fT$.
Applying the estimate to a difference of multipliers
proves the stated coefficient variation convergence.

Distributional $D$ is continuous on original currents.
It is therefore legitimate to define \textup{(OL13)};
the identities $d'D=0$ and $d''D=0$ give separate
closedness.  At a cone multiplier the frozen original
completed recurrence identifies it with the stated BT
product.  At an exact difference of cone members it is
the corresponding signed difference.  At an arbitrary
lattice limit it is its original current limit, but
no closed positive cone argument has been supplied to
make that distribution a positive order-zero current.
The proof retains this distinction explicitly.
\end{proof}

\paragraph{No relation-quotient shortcut.}
The scalar Borel quotient used in this construction is
already the proved relative potential quotient.  It is
not the unresolved actual analytic PS/PL presentation
quotient.  Extending $f\mapsto D(fT)$ to the lattice
does not show that two actual delta-form presentations
with the same analytic value have the same arbitrary-$S$
action.  Nor does ordinary current equality become an
equality on all additional GK delta tests.  These are
separate gates from the original-test multiplier theorem.

% END REVIEWED MODULE research_r12/nonlinear/order_energy_lattice_completion.tex


% BEGIN REVIEWED MODULE research_r12/nonlinear/radial_pole_order_dual_obstruction.tex
\section{An order-reversing dual which is not strongly continuous}
\label{sec:r12-radial-pole-order-dual}

This section resolves negatively the two universal topology
questions posed after \textup{(FO12)}.  The counterexample
works in every remaining degree $q\ge1$, with an explicitly
constructed allowed background on the original field.  All
potentials involved even belong to every finite power-energy
class for that fixed background.  The example does not replace
the arbitrary background in any preceding positive theorem.

Fix integers $n>q\ge1$, put $m=n-q$, and choose a real
$\delta>0$.  On the original affine $n$-space set
\[
 \begin{gathered}
 W=\{\max_i|T_i|<e^\delta\},\qquad
 V=\{\max_i|T_i|<1\},\\
 t(x)=-\log\max_i|T_i(x)|\quad(x\in V),\qquad
 u=-t=\log\max_i|T_i|,\qquad u(o)=-\infty,
 \end{gathered}
 \tag{PO1}
\]
where $o$ is the rigid origin.  Let $\sigma(s)$, $s>0$,
be the original diagonal radius-$e^{-s}$ Gauss point.  For
every $s>0$, define on the whole larger open $W$
\[
             \psi_s=\max\{\log\max_i|T_i|,-s\}.
 \tag{PO2}
\]
These are actual bounded continuous classical psh functions
on compact subsets of $W$, and on $V$ they are
$-\min(t,s)$ with ordinary boundary value zero.

The exact calibration input is the original radial-shell
formula in frozen R6
\path{dirichlet_radial_measure_carriers.tex}, recalled
as R11 \textup{(RE2)} in
\path{comparison/radial_unbounded_energy_and_loss.tex}.
It is an equality on every original scalar test: a bounded
radial slope $A$ has top measure $\sigma_*\tau$ with
$\tau([s,\infty))=A(s)^n$ at continuity points.  Its unit
radial-shell normalization is one.  No field extension or
value-group restriction is made.

\begin{lemma}[An allowed radial-mixture background and its exact bounded sources]
\label{lem:r12-radial-mixture-background}
The weak coefficient integral
\[
                    S=\int_0^\infty e^{-s}(D\psi_s)^m\,ds
 \tag{PO3}
\]
defines on $W$ an original symmetric middle-positive
separately closed $(m,m)$-current.  Its restriction to $V$
satisfies all strict-box and extension hypotheses
\textup{(ED1)}.  For every $R>0$, its original bounded
top source is
\[
 \mu_R:=(D\psi_R)^q\wedge S
     =\sigma_*\bigl(\mathbf1_{(0,R)}e^{-s}\,ds
                                         +e^{-R}\delta_R\bigr).
 \tag{PO4}
\]
Every $\psi_R$ is an actual bounded finite-mass
zero-boundary member of every finite positive power class
for this fixed $S$.
\end{lemma}

\noindent\textbf{Status: \texttt{NEW\_PROVED}} for the
fixed background and its all-original-test source formula.
The bounded arbitrary-background calculus, relative CLN,
and radial-shell calibration are \texttt{INTERNAL\_INPUT}.

\begin{proof}
The functions in \textup{(PO2)} are the original radial
PL truncations of a logarithmic maximum.  Their bounded
classicality, continuity across the origin and all coordinate
zero loci, and their extension to $W$ are the frozen
maximum/curvewise calculus.  On a fixed compact subset
of $W$, their absolute values are bounded by $s+C$ for
one fixed $C$.  Relative CLN for the integration background
therefore bounds every compact coefficient trace of
$(D\psi_s)^m$ by a constant times $(1+s)^m$.
This bound is integrable against $e^{-s}ds$.
The functions $\psi_s$ vary locally uniformly with $s$,
indeed $|\psi_s-\psi_{s'}|\le|s-s'|$.  Frozen bounded
continuous-potential stability makes every current pairing
measurable in $s$.  Consequently \textup{(PO3)} has
locally finite Radon coefficients, is middle-positive and
symmetric, and is separately closed by integration of the
two distributional closedness equations.

The same CLN domination applies to
$T_{[W],m+k}(\psi_s^{[m]},\psi_R^{[k]})$ for each
fixed $R$ and $0\le k\le q$.  Induction in the intrinsic
recursion, multiplying by the fixed continuous bounded
$\psi_R$ on the compact test stage and applying coefficient
Fubini, gives
\[
 T_{S,q}(\psi_R^{[q]})
    =\int_0^\infty e^{-s}
       T_{[W],n}(\psi_R^{[q]},\psi_s^{[m]})\,ds.
 \tag{PO5}
\]
This passage does not use an unproved product of two
arbitrary currents: all intermediate products are the
frozen products of the displayed bounded potentials,
and the action over their integrated background is
identified by its defining recursion.

On $V$, polarization of the frozen radial-shell formula
gives exactly
\[
          T_{[V],n}(\psi_R^{[q]},\psi_s^{[m]})
                                    =\delta_{\sigma(\min(R,s))}.
 \tag{PO6}
\]
To check both the normalization and the original-test
scope, assume first $R\le s$ and take $a,b\ge0$.
The bounded radial potential $a\psi_R+b\psi_s$ has
slope $a+b$ below depth $R$, slope $b$ between depths
$R$ and $s$, and zero afterwards.  Its original top
measure is therefore
\[
       ((a+b)^n-b^n)\delta_{\sigma(R)}
                                      +b^n\delta_{\sigma(s)}.
\]
Comparing the coefficient of $a^qb^m$ with bounded
mixed multilinearity proves \textup{(PO6)}.  The case
$s\le R$ is symmetric, and equality of the two depths
gives the same formula.  The binomial coefficient
cancels on the two sides, leaving unit mass rather
than an extra factorial.  Since the radial calibration
and polarization are original current equalities, this
comparison holds on every actual refinement as well.

Substituting \textup{(PO6)} into \textup{(PO5)} and
splitting the integral at $s=R$ gives \textup{(PO4)}.
Its total mass is one.  On $V$, $\psi_R$ is bounded,
has compact strict negative sublevels and boundary value
zero.  Thus its finite mass gives every finite positive
power energy.  Finally choose a compact polyannulus-free
closed polydisc of radius greater than one inside $W$.
The closed unit polydisc is compactly contained in its
ordinary interior, and $V$ is the strict upper Laurent
box $|T_i|<1$.  The current \textup{(PO3)} is already
defined on its ordinary neighborhood.  This verifies
\textup{(ED1)} without changing the original norm.
\end{proof}

\begin{proposition}[A logarithmic pole in every relative power class]
\label{prop:r12-logarithmic-all-power-relative-pole}
For the background \textup{(PO3)}, the actual logarithmic
potential $u$ in \textup{(PO1)} has no truncation loss and
\[
 \begin{gathered}
 \mu_u=\sigma_*(e^{-s}\,ds),\qquad
 u\in\bigcap_{\alpha>0}\mathcal E^\alpha_{S,\partial}(V),\\
 e_1(u)=1,\qquad e_1(\psi_R)=1-e^{-R},\qquad
                     I(u,\psi_R)=e^{-R}.
 \end{gathered}
 \tag{PO7}
\]
In particular $\psi_R\downarrow u$ strongly as $R\to\infty$,
on the same actual representatives and on every original test.
For every fixed $\alpha>0$ the entire row has the uniform bound
\[
 e_\alpha(\psi_R)
       =\alpha\int_0^R s^{\alpha-1}e^{-s}ds
       \le\Gamma(\alpha+1)=e_\alpha(u).
 \tag{PO7a}
\]
\end{proposition}

\noindent\textbf{Status: \texttt{NEW\_PROVED}} for this
fixed-background example.  It does not assert that a logarithmic
pole has finite energy for every allowed background.

\begin{proof}
The potential is negative extended classical psh, and
every strict negative sublevel is a compactly contained
smaller polydisc.  Its canonical truncations are precisely
$\psi_j$, $j\in\mathbf N$.  Formula \textup{(PO4)}
identifies their restriction to $\{u>-j\}$ as
$\sigma_*(\mathbf1_{(0,j)}e^{-s}ds)$ and their discarded
part as $e^{-j}\delta_{\sigma(j)}$.  The latter has
total mass tending to zero.  Thus $u$ belongs to the
frozen no-loss boundary class and has the displayed
top measure.  For every $\alpha>0$,
\[
                   e_\alpha(u)=\int_0^\infty s^\alpha e^{-s}ds<\infty.
\]
This is exactly the frozen definition of its power class.
In particular it is an actual member of the first-energy
class, hence of the strong cone $\mathscr C$.

Direct integration in \textup{(PO4)} gives
\[
 e_1(\psi_R)=\int_0^R s e^{-s}ds+Re^{-R}=1-e^{-R}.
\]
The same integration by parts with $s^\alpha$ proves
\textup{(PO7a)}.  Thus the later topology failures persist
even under uniform bounds in every fixed finite power class.
Moreover $\psi_R-u=(t-R)_+$ on the finite locus.
Its integral against $\mu_R$ is zero and its integral
against $\mu_u$ is $e^{-R}$.  The frozen discrepancy
identity therefore gives the last formula in
\textup{(PO7)}.  Pole values are ignored only by the
already proved no-loss measure; no representative is
changed in deriving these identities.
\end{proof}

\subsection{A pole-slope functional on the full completed cone}

\begin{lemma}[Every radial Gauss point is detected and every completed trace is convex]
\label{lem:r12-radial-positive-capacity-traces}
For every $t_0>0$ the singleton $\{\sigma(t_0)\}$ has
positive full outer $S$-capacity.  Every completed
$a\in\mathscr C$ has a finite nonpositive convex trace
\[
                           g_a(t)=a(\sigma(t)),\qquad t>0,
 \tag{PO8}
\]
with $\lim_{t\downarrow0}g_a(t)=0$.
Here these are actual values of the relative class:
no outer-null set can contain any one of the displayed
Gauss points.  The trace is the pointwise limit of any
bounded actual strong defining row, and its sum and
positive scaling agree with cone addition and scaling.
\end{lemma}

\noindent\textbf{Status: \texttt{DIRECT\_COROLLARY}} of
\textup{(PO4)}, the frozen mixed energy and strong tester
estimate, and original classical Gauss convexity R6
\path{lem:r6-joint-gauss-convexity}.

\begin{proof}
Formula \textup{(PO4)} gives an atom of mass $e^{-t_0}$
at $\sigma(t_0)$ in the source of the actual bounded
$\psi_{t_0}$.  This source ignores every relative
outer-null set by the frozen completed capacity calculus.
Hence the singleton is not outer-null.  Equivalently,
normalizing $\psi_{t_0}$ by its finite oscillation gives
a full raw capacity tester with positive mass there,
and its domination by outer capacity gives the same
conclusion.

The literal mixed first energy yields, for every
$a\in\mathscr C$,
\[
 0\le -a(\sigma(t))
     \le e^t B(a,\psi_t^{[q]})
     \le e^t r(a)(1-e^{-t})^{q/d}.
 \tag{PO9}
\]
Thus these point values are finite and their upper bound
tends to zero as $t\downarrow0$.  They cannot change
under a relative-null modification.  For a bounded
actual defining row $a_j\to a$, the strong tester
estimate against the fixed $\mu_{\psi_t}$ gives
\[
 |a_j(\sigma(t))-a(\sigma(t))|
       \le e^t\int|a_j-a|\,d\mu_{\psi_t}\longrightarrow0.
\]
The actual bounded traces are convex: restrict to the
coordinate-unit torus inside the polydisc and apply the
frozen joint Gauss convexity lemma along the diagonal
logarithmic ray.  Finiteness and pointwise convergence
preserve convexity on every finite interval.  This
proves \textup{(PO8)} and its boundary limit.  The
linearity assertions follow from the point values,
which are unchanged on the common relative quotient.
\end{proof}

\begin{theorem}[A norm-one order-reversing functional outside the strong dual]
\label{thm:r12-pole-order-dual-not-strong}
For $a\in\mathscr C$ the limit
\[
                    L_\infty(a)=
                         \lim_{t\to\infty}\frac{-g_a(t)}t
 \tag{PO10}
\]
exists, is finite and nonnegative.  It extends linearly
to an order-reversing functional
\[
 L_\infty\in\mathscr O_S\subset\mathbb D_S^*,\qquad
 \|L_\infty\|_*=1,\qquad
 L_\infty(u)=1,\qquad L_\infty(\psi_R)=0.
 \tag{PO11}
\]
It is not strongly continuous.  Consequently, in every
remaining degree $q\ge1$, the universal implication
\[
          \ell\in\mathscr O_S\quad\Longrightarrow\quad
                   \ell\in\mathbb D_{S,\mathrm s}^*
 \tag{PO12}
\]
is false.

The restriction of $L_\infty$ to every original compact
smooth scalar test is zero.  Thus it is not literal
integration against any Radon measure on the relative
Borel difference space, even though it has the stated
sign on the full scalar order cone.  In the standard
Banach-lattice convention, $-L_\infty$ is positive.
It is an asymptotic pole
functional, not a positive current on the original
smooth test sheaf.
\end{theorem}

\noindent\textbf{Status: \texttt{NEW\_PROVED}}.  The
universal implication \textup{(PO12)} is \texttt{FALSE},
not merely unproved.  This strengthens the earlier
counterexample which had positivity only on the
potential cone.

\begin{proof}
The function $f_a=-g_a$ is nonnegative and concave on
$[0,\infty)$, with $f_a(0)=0$.  Its quotients
$f_a(t)/t$ are nonincreasing and nonnegative, so they
have a finite limit $L_\infty(a)$.  Moreover
\[
                    f_a(t)\ge tL_\infty(a)\qquad(t>0).
\]
Using the exact source of $u$ and the mixed energy
inequality gives
\[
 L_\infty(a)
  \le\int_0^\infty f_a(t)e^{-t}dt
  =B(a,u^{[q]})\le r(a)r(u)^q=r(a).
 \tag{PO13}
\]
Addition and positive homogeneity follow by taking the
finite limits of the corresponding trace identities.
If two cone difference presentations agree, their
values agree at every radial Gauss point by the
preceding lemma, so their slope differences agree.
Thus $L_\infty$ extends to a linear functional on
$\mathbb D_S$.  For $h=a-b$, the bound
$|L_\infty(h)|\le r(a)+r(b)$ proves continuity after
infimizing.  The potential $u$ has energy root one and
slope one, so the norm is exactly one.

If $h\succeq0$, its actual value at every $\sigma(t)$
is nonnegative.  Therefore
\[ L_\infty(h)=\lim_{t\to\infty}-h(\sigma(t))/t\le0. \]
This proves order reversal on the full scalar order
cone, not only nonnegativity on $\mathscr C$.
The values on $u$ and $\psi_R$ are immediate.  Since
$\psi_R\to u$ strongly by \textup{(PO7)}, those
values disprove strong continuity.

An original compact smooth scalar $\phi$ is bounded
on $V$.  It belongs to $\mathbb D_S$ by the frozen
exact common-collar difference theorem.  Thus
$L_\infty(\phi)=\lim_{t\to\infty}-\phi(\sigma(t))/t=0$.
Any Radon measure with these original smooth test
integrals would be zero.  It could not also represent
$L_\infty(u)=1$.  This proves the final assertion
without assigning a fictitious point mass to the pole.
\end{proof}

\begin{theorem}[Strong convergence need not converge in the order norm]
\label{thm:r12-strong-not-order-topology}
For every $R>0$ one has the exact identities
\[
 \boxed{\quad
 P(-|u-\psi_R|)=u,
       \qquad\|u-\psi_R\|_{\rm o}
              =\|u-\psi_R\|_{\mathbb L_S}=1,
       \qquad I(u,\psi_R)=e^{-R}.
 \quad}
 \tag{PO14}
\]
Hence the order norm topology on $\mathscr C$ is
strictly stronger than its strong discrepancy topology
in every remaining degree.  The order-energy Banach
lattice norm is not order continuous: the positive
functions $\psi_R-u$ decrease to zero outside the pole
null set while their norms remain one.
\end{theorem}

\noindent\textbf{Status: \texttt{NEW\_PROVED}}.  The
universal strong-to-order-norm implication considered
in \textup{(FO13)} is \texttt{FALSE}.

\begin{proof}
Put $f_R=-|u-\psi_R|=u-\psi_R$ and $w=P(f_R)$.
Since $\psi_R\preceq0$, the original potential $u$
is below $f_R$, so $u\preceq w$ and $r(w)\le1$.
The order-dual norm inequality \textup{(FO10)}, applied
to the norm-one $L_\infty$, gives
\[
       1=|L_\infty(f_R)|\le\|f_R\|_{\rm o}=r(w).
\]
Thus $e(w)=e(u)=1$.  Ordered energy separation gives
$I(u,w)\le e(u)-e(w)=0$, hence $w=u$ as relative
classes.  This proves the first two formulas in
\textup{(PO14)}; the discrepancy formula was computed
in \textup{(PO7)}.  On the finite locus of $u$,
$(\psi_R-u)=(t-R)_+$ decreases pointwise to zero.
The pole is a relative outer-null set by the frozen
power-class capacity theorem.  This proves the
failure of order continuity with precisely the
actual relative order of the lattice.
\end{proof}

\begin{corollary}[A purely unbounded component of the order dual]
\label{cor:r12-pole-dual-on-lattice}
The norm-one order-reversing functional $L_\infty$
extends to $\mathbb L_S$.  On that lattice it has the
same concrete formula
\[
                     L_\infty(f)=
                         \lim_{t\to\infty}\frac{-f(\sigma(t))}{t}.
 \tag{PO15}
\]
In particular it vanishes on every bounded lattice
element.  The norm closure of the bounded elements
of $\mathbb L_S$ is therefore a proper subspace,
and does not contain $u$.
\end{corollary}

\noindent\textbf{Status: \texttt{DIRECT\_COROLLARY}} of
the actual Borel lattice realization and the pole-slope
estimate; no identification with a Radon measure is made.

\begin{proof}
The order-dual norm theorem gives the extension with
the same norm.  For $f\in\mathbb L_S$ let
$a=P(-|f|)$.  Its values at each $\sigma(t)$ are
well defined and finite, because the fixed source
$\mu_{\psi_t}$ has an atom there and every lattice
element is literally integrable for that source.
For $t\ge1$, concavity of $-g_a$ gives
\[
 \frac{|f(\sigma(t))|}{t}
    \le\frac{-g_a(t)}t\le-g_a(1)
    \le e(1-e^{-1})^{q/d}\|f\|_{\mathbb L_S}.
\]
Thus the normalized evaluation functionals are
uniformly bounded on $\mathbb L_S$ for $t\ge1$.
They converge on the dense subspace $\mathbb D_S$
to $L_\infty$.  Approximation in norm and this uniform
bound extend the convergence to every lattice element,
which proves \textup{(PO15)}.  A bounded element has
zero limit in that formula.  Since $L_\infty(u)=1$,
the closure of those elements cannot contain $u$.
\end{proof}

\paragraph{Revised gate status and nonclaims.}
The full order-dual continuity and strong-to-order-norm
implications are now \texttt{FALSE} under the project's
allowed arbitrary-background hypotheses, even with the
uniform all-finite-power bounds in \textup{(PO7a)}.  They cannot
be used to remove the strong-dual hypothesis from the
general signed-data variational theorem.  This example
does not prove that every order-reversing datum fails
variational attainment: $L_\infty$ itself attains its
cone variational maximum at a positive scaling of $u$.
Whether a useful larger positive datum class has a
separate upper-semicontinuity or attainment theorem
requires an additional proof.  Literal capacity-controlled
data and all previously proved source/contact theorems
are unchanged.  No GQ, CTO, or extra delta-test conclusion
follows from the asymptotic pole functional.

% END REVIEWED MODULE research_r12/nonlinear/radial_pole_order_dual_obstruction.tex


% BEGIN REVIEWED MODULE research_r12/energy/compact_test_shadow_and_core.tex
\section{The compact-test core and the Radon shadow of an energy dual}
\label{sec:r12-compact-test-core}

Assume exactly the strict-box hypotheses (M2), with the arbitrary
original background unchanged. Let $\mathbb L_S$ be the actual Borel
Banach lattice of Section~\ref{sec:r12-order-energy-lattice}. A
\emph{positive functional} on this lattice means a continuous linear
map $F$ such that $F(f)\geq0$ for $f\geq0$ in its relative Borel
order. Thus the corresponding order-reversing energy functional
is $\ell=-F$. This sign convention will be used throughout this
section.

\begin{definition}[The compact-test core]
Let
\[
 \mathbb L_{S,0}
   =\overline{C_c^\infty(V,\mathbf R)}^{\,\|\cdot\|_{\mathbb L_S}} .
 \tag{CS1}
\]
The original smooth test functions in (CS1) include every actual
moment refinement. The closure is in the proved scalar multiplier
lattice, not an ambient Sobolev space and not a new potential cone.
\end{definition}

\begin{lemma}[Continuous compact tests and the zero-boundary finite-mass core]
\label{lem:r12-continuous-compact-core}
Every actual continuous compactly supported scalar function belongs
canonically to $\mathbb L_{S,0}$. For a compact set $H\Subset V$ there
is a finite $A_H$ such that
\[
       \|f\|_{\mathbb L_S}\leq A_H\|f\|_\infty
                   \qquad(f\in C_c(V),\ \operatorname{supp}f\subset H).
 \tag{CS2}
\]
The space $\mathbb L_{S,0}$ is a closed vector sublattice.

If $v<0$ is bounded, continuous, actual classical psh, has compact
negative sublevels and finite total top mass, then
$v\in\mathbb L_{S,0}\cap\mathscr C$. More precisely, for
$v^\epsilon=\max(v,-\epsilon)$,
\[
 \begin{split}
 v-v^\epsilon&\in C_c(V),\\
 \|v^\epsilon\|_{\mathbb L_S}
   &=r(v^\epsilon)\leq
            \bigl(\epsilon\,\mu_v(V)\bigr)^{1/d}.
 \end{split}
 \tag{CS3}
\]
No assertion about all bounded discontinuous psh functions is
contained in this statement.
\end{lemma}

\noindent\textbf{Status: \texttt{NEW\_PROVED}} for this identified
core and norm estimate, using the frozen smooth scalar density,
strict-box barrier, comparison, and the new order norm.

\begin{proof}
Let $\rho$ be the bounded strictly negative continuous barrier
of (TC3)--(TC4). On $H$ choose $c_H>0$ with $\rho\leq-c_H$.
For an original smooth $f$ supported in $H$,
\[
 |f|\leq \|f\|_\infty(-\rho)/c_H
 \quad\Longrightarrow\quad
 \|f\|_{\mathbb L_S}\leq r(\rho)\|f\|_\infty/c_H.
\]
The inequality holds outside $H$ as well because $f=0$ there.
Original smooth scalar functions approximate continuous compact
functions uniformly, with support in one fixed slightly larger
compact set. This is the frozen original cutoff and scalar-density
calculus, not density of background currents. The displayed estimate
makes such approximations Cauchy in $\mathbb L_S$; its capacity
realization identifies their limit with the given continuous
function. Applying the estimate with any fixed larger compact
gives (CS2). In particular the identification does not depend on
the chosen smooth approximations.

Maxima and absolute values of continuous compact functions are
again continuous compact functions. They are thus in the same
closure. The lattice operations are norm continuous in
$\mathbb L_S$, so its closed subspace (CS1) is a sublattice.

The finite-mass potential $v$ belongs to the frozen actual
first-energy class. So does $v^\epsilon$. For every $c>1$,
strict negativity gives $cv<v^\epsilon$ everywhere, and each
positive-margin comparison gap is compact by the boundary
condition. The frozen comparison theorem gives
$\mu_{v^\epsilon}(V)\leq c^q\mu_v(V)$. Letting $c\downarrow1$
yields the same inequality without $c$. Since
$0\leq-v^\epsilon\leq\epsilon$,
\[
                e(v^\epsilon)
                     \leq\epsilon\,\mu_{v^\epsilon}(V)
                     \leq\epsilon\,\mu_v(V).
\]
This proves the second line of (CS3). The first line follows from
continuity and the compactness of the closure of
$\{v<-\epsilon\}$. Therefore $v-v^\epsilon\in\mathbb L_{S,0}$
and these functions converge to $v$ in the order norm. Closedness
proves $v\in\mathbb L_{S,0}$. The zero potential may be added
separately. This proof does not replace bounded psh functions
by continuous ones without an approximation theorem.
\end{proof}

\begin{definition}[Admissible compact-test Radon data]
Let $\mathcal M_{\rm ct}^+(S,V)$ be the positive Radon measures
$\mu$ on $V$ for which
\[
 \|\mu\|_{\rm ct}
     =\sup_{\substack{f\in C_c(V),\ f\geq0\\
                         \|f\|_{\mathbb L_S}\leq1}}
                                   \int_V f\,d\mu <\infty.
 \tag{CS4}
\]
A nonnegative continuous compact function of zero order norm must
therefore have zero $\mu$-integral; otherwise (CS4) is infinite
by scaling. This is a condition on compact continuous functions,
not on all representatives of relative Borel classes. It must not
be confused with literal finite-energy integration on
$\mathbb L_S$.
\end{definition}

\begin{theorem}[An exact compact-test shadow map]
\label{thm:r12-positive-dual-shadow}
Every positive $F\in\mathbb L_S^*$ has a unique positive Radon
measure $\mu_F$ characterized by
\[
             F(\varphi)=\int_V\varphi\,d\mu_F
                 \qquad(\varphi\in C_c^\infty(V,\mathbf R)).
 \tag{CS5}
\]
Equality also holds for every actual continuous compact function.
The measure belongs to $\mathcal M_{\rm ct}^+(S,V)$, with
\[
        \|\mu_F\|_{\rm ct}
             =\|F|_{\mathbb L_{S,0}}\|
             \leq\|F\|.
 \tag{CS6}
\]

Conversely every $\mu\in\mathcal M_{\rm ct}^+(S,V)$ defines a
unique positive continuous functional $F_0$ on $\mathbb L_{S,0}$.
It has at least one positive extension $F$ to $\mathbb L_S$
satisfying
\[
                       \|F\|=\|F_0\|=\|\mu\|_{\rm ct}.
 \tag{CS7}
\]
For two positive continuous functionals $F_1,F_2$,
\[
 \mu_{F_1}=\mu_{F_2}
   \quad\Longleftrightarrow\quad
      F_1-F_2\in\mathbb L_{S,0}^{\perp}.
 \tag{CS8}
\]
Thus compact original tests determine the restriction to the
core, not automatically the functional on the whole energy
lattice. None of (CS5)--(CS8) asserts literal integration on
all relative Borel representatives.
\end{theorem}

\noindent\textbf{Status: \texttt{NEW\_PROVED}} as the exact
shadow/core statement for the new energy lattice.
The positive Riesz representation theorem and Hahn--Banach are
\texttt{EXTERNAL\_KNOWN}. The extension argument is included to
specify its positivity and norm rather than invoking an unspecified
positive extension principle.

\begin{proof}
For every compact $H$, (CS2) bounds the restriction of $F$ to
continuous functions supported in $H$ in the uniform norm.
Its positivity is the pointwise positivity of such functions
passed to relative classes. The positive Riesz representation
theorem on the locally compact Hausdorff ordinary space gives
a unique positive Radon measure $\mu_F$. Local compactness
follows from goodness; a finite cover of each compact stage
suffices, so no countability at infinity is being inferred for
an arbitrary $X$. Smooth scalar density proves uniqueness using
only the original test functions, and proves (CS5) for all
continuous compact functions.

A positive functional on a normed lattice has its norm computed
on nonnegative elements: $|F(f)|\leq F(|f|)$ and
$\||f|\|=\|f\|$. The nonnegative continuous compact functions
are dense in the nonnegative cone of $\mathbb L_{S,0}$, by
taking positive parts of continuous compact approximations.
This proves (CS6).

Conversely the integration functional on $C_c(V)$ is bounded
in the order norm by (CS4), since
\[
                 \left|\int f\,d\mu\right|
                      \leq\int |f|\,d\mu
                      \leq\|\mu\|_{\rm ct}\|f\|_{\mathbb L_S}.
\]
It annihilates zero-norm functions, so it extends continuously
and uniquely to the core. Positivity passes by positive-part
approximation, and its norm is exactly (CS4).

Set $M=\|F_0\|$ and define on the real vector space
$\mathbb L_S$ the finite sublinear functional
\[
                            p(f)=M\|f_+\|_{\mathbb L_S}.
 \tag{CS9}
\]
Indeed $(f+g)_+\leq f_++g_+$ and the norm is a lattice norm.
For $f\in\mathbb L_{S,0}$ positivity implies
$F_0(f)\leq F_0(f_+)\leq p(f)$. Hahn--Banach supplies a linear
extension $F\leq p$. If $f\geq0$, then $p(-f)=0$, hence
$F(-f)\leq0$ and $F(f)\geq0$. Applying the domination to both
$f$ and $-f$ gives $|F(f)|\leq M\|f\|$. Therefore
$\|F\|\leq M$; the restriction gives the reverse inequality.
This proves (CS7), including the zero case. Finally (CS5) and
the density defining the core show that equal shadows are
equivalent to vanishing of the difference on that core. This
is (CS8).
\end{proof}

\begin{proposition}[Literal source functionals remain distinguishable]
\label{prop:r12-source-shadow-injectivity}
For $u\in\mathscr C$, let $F_u=-J(u)$ on $\mathbb L_S$.
Its compact-test shadow is the original measure
$\widehat\mu_u$. The restriction of the shadow map to the
source family $\{F_u:u\in\mathscr C\}$ is injective:
\[
                    \mu_{F_u}=\mu_{F_v}\quad\Longrightarrow\quad u=v.
 \tag{CS10}
\]
This does not require the shadow map to be injective on all
positive functionals.
\end{proposition}

\noindent\textbf{Status: \texttt{DIRECT\_COROLLARY}} of literal
curvature integration (OL6) and frozen strong-cone uniqueness.

\begin{proof}
Formula (OL6) identifies the shadow on every original smooth
compact test and hence as a Radon measure. If the two shadows
are equal, then $\widehat\mu_u=\widehat\mu_v$ on the full
original test sheaf. The frozen completed comparison/uniqueness
theorem, with the same relative boundary and null quotient,
gives $u=v$. Equivalently the first-energy discrepancy vanishes.
\end{proof}

\paragraph{Two distinct failures of a reconstruction shortcut.}
The proper-moment example in
Section~\ref{sec:r12-null-fibers-nonliteral-duals} gives a positive
strong-continuous dual whose compact-test shadow is a Dirac
measure on an outer-capacity-null fiber. Its functional cannot
be literal integration on arbitrary representatives of the
relative Borel class. This is failure of literalizability even
when the shadow is known.

Separately, Section~\ref{sec:r12-radial-pole-order-dual} constructs
a norm-one order-reversing $L_\infty$. The functional
$-L_\infty$ is a nonzero positive member of
$\mathbb L_{S,0}^{\perp}$, because it vanishes on every
bounded lattice element. In that example the bounded continuous
finite-mass potentials $\psi_R$ belong to this core by (CS3),
but their strong limit $u$ does not. Thus the core is not closed
for the strong topology on its intersection with the potential
cone. This is a \texttt{DIRECT\_COROLLARY} of the exact
radial proof, not an assumption in (CS5)--(CS10).
The functional is invisible to every original compact scalar
test without being zero as an energy dual. It is not a nonzero
original current with zero tests. An energy dual and an original
current are different objects, and the distinction is essential.


% END REVIEWED MODULE research_r12/energy/compact_test_shadow_and_core.tex


% BEGIN REVIEWED MODULE research_r12/nonlinear/literal_dual_continuity_rooftop_gate.tex
\section{The remaining literal-data continuity question and its finite-rooftop gate}
\label{sec:r12-literal-dual-rooftop-gate}

Retain the full strict-box arbitrary-background hypotheses and
the unchanged strong cone $\mathscr C$.  The counterexample
$L_\infty$ is not literal integration against a Radon measure.
It therefore does not answer the following narrower question:
\[
 \begin{gathered}
 \mu\ge0\text{ is Radon and literal integration defines a}\\
 \text{continuous functional }\ell_\mu\in\mathbb D_S^*
 \quad\overset{?}{\Longrightarrow}\quad
                   \ell_\mu\in\mathbb D_{S,\mathrm s}^*
 \end{gathered}
 \tag{LG1}
\]
Here ``literal'' includes independence of all relative Borel
representatives; in particular $\mu$ ignores every locally
ordinary-open outer-capacity-null Borel set.  Finiteness of the
unweighted total mass of $\mu$ is not assumed.  The question
\textup{(LG1)} remains \texttt{OPEN} in this round.  The
capacity-Dini and actual-source-dominated data theorems already
prove useful sufficient cases and are not being reopened.

For finitely many cone members define their lower rooftop
\[
       P(a_0,\ldots,a_j)=P\bigl(\min(a_0,\ldots,a_j)\bigr).
 \tag{LG2}
\]
The right side uses the proved Borel obstacle construction.
It exists: the finite sum $a_0+\cdots+a_j$ is a cone member
below every entry, since all entries are nonpositive.
Thus this definition introduces no existence hypothesis.

\begin{proposition}[A common lower bound is exactly a bounded finite-rooftop row]
\label{prop:r12-finite-rooftop-common-bound}
Let $u_j\to u$ strongly in $\mathscr C$.  The following
properties of a chosen subsequence $(u_{j_l})$ are equivalent:
\begin{enumerate}
\item there is $w\in\mathscr C$ with
$w\preceq u$ and $w\preceq u_{j_l}$ for every $l$;
\item the decreasing finite rooftops
\[
              w_l=P(u,u_{j_1},\ldots,u_{j_l})
 \tag{LG3}
\]
have uniformly bounded first energy.
\end{enumerate}
If they hold, the actual decreasing relative limit of
$w_l$ is itself a common lower bound in $\mathscr C$,
and $w_l$ converges strongly to it.
\end{proposition}

\noindent\textbf{Status: \texttt{DIRECT\_COROLLARY}}
of the proved obstacle existence and frozen decreasing
energy-bounded closure, with the equivalence recorded
to isolate the exact finite-dimensional obstruction.

\begin{proof}
If a common lower bound $w$ exists, it is admissible
for every rooftop, so $w\preceq w_l$.  Energy
monotonicity gives $e(w_l)\le e(w)$.
Conversely the rooftops decrease, because their
constraints increase.  The frozen decreasing closure
for an energy-bounded row gives a strong limit
$w\in\mathscr C$ represented by their actual
decreasing Borel limit outside the fixed countable
union of relative-null sets.  For each fixed $l$,
all subsequent rooftops are below $u_{j_l}$ and
below $u$.  Their limit has the same inequalities.
No uncountable-union or ordinary usc envelope is used.
\end{proof}

\begin{proposition}[The common-bound mechanism is valid on the full radial energy family]
\label{prop:r12-radial-common-lower-subsequence}
For the explicit background \textup{(PO3)}, every strongly
convergent row of actual radial potentials from
Theorem~\ref{thm:r12-exact-radial-flux-inverse} has a
subsequence with a common \emph{actual radial} lower
bound in the same first-energy class.  Consequently
every literal positive datum with bounded energy-difference
functional is continuous along all such strong radial rows,
even if it is not known to be strongly continuous on the
whole ambient cone.
\end{proposition}

\noindent\textbf{Status: \texttt{NEW\_PROVED}} on the
all-degree radial class, not on arbitrary ambient potentials.
This supplies a positive test of the proposed mechanism
without promoting it to the general gate.

\begin{proof}
Write the row as $u_{A_j}\to u_A$.  The exact discrepancy
theorem gives $A_j\to A$ in $L^d(e^{-t}dt)$.
Choose a subsequence with
$\sum_l\|A_{j_l}-A\|_d^d<\infty$ and put
\[
                  A_*(t)=\sup\{A(t),A_{j_1}(t),A_{j_2}(t),\ldots\}.
 \tag{LG3a}
\]
This is nonnegative, nonincreasing and right continuous.
Pointwise,
\[
  A_*\le A+\left(\sum_l|A_{j_l}-A|^d\right)^{1/d},
  \qquad
  \|A_*\|_d
        \le\|A\|_d+\left(\sum_l\|A_{j_l}-A\|_d^d\right)^{1/d}.
\]
Hence it is finite at every positive depth: an infinite
value at one depth would, by monotonicity, force an
infinite weighted integral on an interval.  Thus
$A_*\in\mathscr A_q$.  The actual potential $u_{A_*}$
is below every selected potential and below $u_A$,
because its integrated slope is larger at every depth.
For a literal positive datum, $-u_{A_*}$ is integrable.
A further quasi-everywhere convergent subsequence is
dominated by that function, so dominated convergence
applies.  Repeating the construction from every
subsequence proves the last assertion.  The pointwise supremum
is taken in the proved radial slope space, not in an
unconstructed lattice of arbitrary current Hessians.
\end{proof}

\begin{proposition}[The common-lower-bound subsequence property would settle literal continuity]
\label{prop:r12-common-bound-literal-continuity}
Suppose that every strongly convergent sequence in
$\mathscr C$ has a subsequence satisfying the equivalent
conditions of Proposition~\ref{prop:r12-finite-rooftop-common-bound}.  Then
\textup{(LG1)} holds.  More precisely, for every literal
positive datum as in \textup{(LG1)} and every strong
row $u_j\to u$,
\[
                         \int_V|u_j-u|\,d\mu\longrightarrow0.
 \tag{LG4}
\]
\end{proposition}

\noindent\textbf{Status: \texttt{CONDITIONAL}}.  Its
printed subsequence premise is \texttt{OPEN}, not a
consequence silently imported from the ordinary
Perron theorem or from the order norm.

\begin{proof}
Start from any subsequence of the given row and use
the assumed property to pass to a further subsequence
with a common lower bound $w\in\mathscr C$.
Strong convergence implies local outer-capacity
convergence.  A fast further subsequence, chosen on
a countable compact exhaustion of the fixed box,
converges pointwise off one locally outer-null Borel
set.  Countable subadditivity proves this familiar
capacity Borel--Cantelli passage, without raw/outer
identification.  Literalness makes the exceptional
set $\mu$-null.

All selected potentials and $u$ lie between $w$ and
zero, outside a further countable union of null sets.
Thus $|u_{j_l}-u|\le -w$ there.  The dominating
function is integrable because
$\int(-w)d\mu=\ell_\mu(w)\le\|\ell_\mu\|_*r(w)$.
It also shows that its pole set is $\mu$-null.
Dominated convergence gives \textup{(LG4)} along
that further subsequence.  Since the argument starts
from every subsequence, the full row converges in
$L^1(\mu)$.  In particular its scalar pairings
converge, proving strong continuity.
\end{proof}

\begin{proposition}[An ordered two-entry continuity gate is sufficient]
\label{prop:r12-ordered-rooftop-reduction}
Consider the following assertion for the same cone:
\[
 \begin{gathered}
 w\preceq u,\quad u_j\longrightarrow u\text{ strongly}
 \quad\Longrightarrow\quad
        e\bigl(P(w,u_j)\bigr)\longrightarrow e(w).
 \end{gathered}
 \tag{LG5}
\]
If \textup{(LG5)} holds, every strong row admits a
subsequence with a common lower bound.  Hence it
implies \textup{(LG1)}.  Only this ordered-target
two-entry rooftop statement is needed; no full
continuity theorem for arbitrary moving obstacles
is required.
\end{proposition}

\noindent\textbf{Status: \texttt{CONDITIONAL}} with
the new minimal rooftop gate \textup{(LG5)} explicitly
\texttt{OPEN}.

\begin{proof}
Fix a positive summable sequence $(\epsilon_l)$ and
put $w_0=u$.  Having selected $w_{l-1}\preceq u$,
apply \textup{(LG5)} with this fixed lower member.
Choose $j_l>j_{l-1}$ so large that
\[
          e(P(w_{l-1},u_{j_l}))
                                 \le e(w_{l-1})+\epsilon_l.
\]
Set $w_l=P(w_{l-1},u_{j_l})$.  Greatestness of
the rooftops shows that it equals the cumulative
rooftop in \textup{(LG3)}.  Indeed being below
$w_{l-1}$ is equivalent, for a cone member, to
being below all its preceding constraints.
The row decreases and
$e(w_l)\le e(u)+\sum_l\epsilon_l$.
The previous two propositions apply.
\end{proof}

\paragraph{What the existing estimates do prove.}
The same ordered gate holds if strong convergence
is strengthened to energy-difference norm convergence.
If $u_j-u=a_j-b_j$ with $a_j,b_j\in\mathscr C$,
then $w+a_j\preceq w$ and
$w+a_j\preceq u+a_j\preceq u_j$.  Hence
\[
  w+a_j\preceq P(w,u_j)\preceq w,
 \qquad
 0\le r(P(w,u_j))-r(w)\le r(a_j).
 \tag{LG6}
\]
Taking decompositions with
$r(a_j)+r(b_j)\to0$ proves the claim.  A fast
energy-difference convergent subsequence also has
the direct common lower bound
$u+\sum_j a_j$, by the summable-root decreasing
series construction.  These are
\texttt{DIRECT\_COROLLARY} statements in the
stronger norm, not proofs of \textup{(LG5)}.

\paragraph{Why the attempted general proof stops.}
The frozen full Perron theorem supplies least upper
bounds and energy-minimizing finite maxima.  The
rooftops in \textup{(LG3)} move in the opposite
direction.  Their individual existence supplies
only the rough finite bound
\[
                   r(P(w,u_j))\le r(w)+r(u_j),
 \tag{LG7}
\]
not a vanishing increment above $r(w)$.  Iterating
that bound over an infinite row need not keep energy
bounded.  The obstacle theorem's proved continuity
uses energy-difference control; the current strong
topology does not provide that control.  Replacing
the latter by an order-norm assertion would also
be invalid, by the pole counterexample.

The pole counterexample does \emph{not} disprove
\textup{(LG5)} or the common-lower-bound property:
its row $\psi_R\downarrow u$ already has the
common lower bound $u$.  It does show why strong
convergence must not be replaced by order-norm
convergence in a purported proof.  Nor do the
smooth-density source obstructions disprove
\textup{(LG1)}: their literal data already lie
in the strong dual and fail at the separate
normal-defect/source-compatibility step.

Accordingly the route through \textup{(LG5)} is
recorded as a \texttt{FAILED\_ATTEMPT} at its
unproved rooftop-continuity step, with its exact
conditional consequences proved above.  No
unrestricted comparison, orthogonality, or source
existence theorem has been used to bypass that step.

\paragraph{A further counterstress, not a counterexample.}
A possible test in degree one and dimension two is to
take rigid centers $p_j=(a_j,0)\to0$, a summable
positive mixture of translated radial curvature
backgrounds, and the isolated logarithmic poles
$\log\max(|T_1-a_j|,|T_2|)$.  One would have to
prove their precise strong convergence for that fixed
mixture and then prove that every completed common
minorant is forced to have a divisorial pole along
$T_2=0$ of forbidden relative energy.  Finite classical
products suggest such a forcing phenomenon, but
ordinary curvewise intuition is not a proof for the
completed Borel cone and its original ambient tests.
In particular no Siu-type or completed-current
divisorial-forcing theorem has been established in
this argument.  This stress test is retained as a
\texttt{FAILED\_ATTEMPT}, with both the exact strong
convergence calculation and the completed minorant
forcing step unproved; it does not change the
\texttt{OPEN} status of \textup{(LG1)} or \textup{(LG5)}.

% END REVIEWED MODULE research_r12/nonlinear/literal_dual_continuity_rooftop_gate.tex


% BEGIN REVIEWED MODULE research_r12/nonlinear/fixed_domination_literal_closure.tex
\section{Fixed domination closes the literal energy-dual interval}
\label{sec:r12-fixed-domination-literal-closure}

Retain exactly the strict-box hypotheses \textup{(ED1)}, the
arbitrary original symmetric middle-positive separately closed
background $S$, and the unchanged full strong cone $\mathscr C$.
The degree $q=n-p\geq1$ is unrestricted. Write
$\widehat\mu_v=\widehat N_q(v^{[q]})$ for the original completed
source of $v\in\mathscr C$, and retain the energy-difference
Banach space $\mathbb D_S$ and its dual norm $\|\cdot\|_*$.
Every literal integral below uses actual Borel representatives
modulo the same local ordinary-open outer-capacity null sets.

\begin{lemma}[Fixed-source integrability and original scalar density]
\label{lem:r12-fixed-source-L1-density}
For $v\in\mathscr C$, every $h\in\mathbb D_S$ is literally
integrable against $\widehat\mu_v$, and
\[
 \int_V |h|\,d\widehat\mu_v
       \leq r(v)^q\|h\|_{\rm ed}.
 \tag{FD1}
\]
For each such $h$ there are original smooth compactly supported
scalar functions $\phi_k$ such that
\[
       \int_V|h-\phi_k|\,d\widehat\mu_v\longrightarrow0.
 \tag{FD2}
\]
No convergence in the energy-difference norm is asserted in
\textup{(FD2)}.
\end{lemma}

\noindent\textbf{Status: \texttt{DIRECT\_COROLLARY}} of literal
completed energy integration, mixed H\"older, and original
scalar density. The usual $C_c$ density in $L^1$ of a Radon
measure is \texttt{EXTERNAL\_KNOWN}; the passage to the original
CLD scalar tests is given explicitly below.

\begin{proof}
Write $h=a-b$ with $a,b\in\mathscr C$. Outside a source-null
set containing their pole and relative exceptional sets,
$|h|\leq-a-b$. Consequently
\[
 \int_V|h|\,d\widehat\mu_v
 \leq B(a,v^{[q]})+B(b,v^{[q]})
 \leq(r(a)+r(b))r(v)^q.
\]
Taking the infimum over the exact Borel decompositions of $h$
proves \textup{(FD1)}, including independence of representatives.
The measure $\widehat\mu_v$ is Radon and locally finite, although
its total mass need not be finite. Standard Radon $L^1$ density
first approximates $h$ by continuous compactly supported functions.
For each one, the frozen original scalar-density and cutoff
calculus used in \textup{(CS2)} gives uniform original smooth
approximations with support in one slightly larger compact set.
That compact has finite $\widehat\mu_v$ mass, so the uniform
approximation is an $L^1$ approximation. A diagonal choice proves
\textup{(FD2)}. Each $\phi_k$ belongs to $\mathbb D_S$ by the
frozen exact common-collar construction \textup{(HC14)--(HC15)}.
This argument uses neither density of $S$ nor an unproved density
of compact tests in the energy-difference norm.
\end{proof}

Fix $v\in\mathscr C$ and a finite constant $M\geq0$. Set
\[
 \mathfrak I_{v,M}
  =\left\{\ell_\mu:\ 0\leq\mu\leq M\widehat\mu_v\right\},
 \qquad
 \ell_\mu(h)=-\int_Vh\,d\mu\quad(h\in\mathbb D_S).
 \tag{FD3}
\]
The domination in this definition is an inequality of original
Radon measures on every Borel subset of $V$. It is not merely
an inequality tested on the cone of potentials.

\begin{theorem}[A norm-closed interval of literal strong-dual data]
\label{thm:r12-fixed-domination-literal-closure}
The set $\mathfrak I_{v,M}$ is convex, weak-star compact in
$\sigma(\mathbb D_S^*,\mathbb D_S)$, and hence norm closed in
$\mathbb D_S^*$. Its elements are literal positive Radon data
in the strong-continuous dual, with
\[
 \mathfrak I_{v,M}\subset\mathbb D_{S,\mathrm s}^*\cap
       \{\ell:\|\ell\|_*\leq M r(v)^q\}.
 \tag{FD4}
\]
In particular, if
\[
 0\leq\mu_j\leq M\widehat\mu_v,\qquad
       \|\ell_{\mu_j}-\ell\|_*\longrightarrow0,
 \tag{FD5}
\]
then there is a unique original positive Radon measure
$\mu\leq M\widehat\mu_v$ for which
\[
                    \ell(h)=-\int_Vh\,d\mu
                       \qquad(h\in\mathbb D_S).
 \tag{FD6}
\]
The same conclusion holds if norm convergence in
\textup{(FD5)} is replaced by weak-star convergence.
\end{theorem}

\noindent\textbf{Status: \texttt{NEW\_PROVED}} as a literal
closure and compactness statement for the new Banach dual.
Strong integration under fixed completed-source domination is a
\texttt{DIRECT\_COROLLARY} of the frozen completed tester
estimate \textup{(EV1)}. The proof of literalness on the full
energy space, rather than only on compact tests, is essential.

\begin{proof}
For a dominated measure, \textup{(FD1)} gives
\[
       |\ell_\mu(h)|\leq M r(v)^q\|h\|_{\rm ed}.
 \tag{FD7}
\]
It also shows that all the displayed integrals exist literally.
The measure ignores all relative outer-null sets because
$\widehat\mu_v$ does. If $u_j\to u$ strongly in $\mathscr C$,
the completed frozen \textup{(EV1)} estimate gives
\[
 |\ell_\mu(u_j-u)|
 \leq M\int_V|u_j-u|\,d\widehat\mu_v\longrightarrow0.
 \tag{FD8}
\]
Thus \textup{(FD4)} follows, and convexity is immediate.

We first prove the literal conclusion directly. Suppose that
$\ell_{\mu_j}\to\ell$ weak-star; norm convergence is a special
case. On original smooth compact tests define
$T(\phi)=-\ell(\phi)$. Passing to the limit in the positive
domination inequalities gives
\[
 T(\phi)\geq0\quad(\phi\geq0),\qquad
            |T(\phi)|\leq M\int_V|\phi|\,d\widehat\mu_v.
 \tag{FD9}
\]
Uniform original scalar density with compact support extends $T$
uniquely to a positive functional on $C_c(V)$. On each compact
support it is bounded in the uniform norm. The positive Riesz
theorem supplies a unique original Radon measure $\mu$ with
$T(\phi)=\int\phi\,d\mu$. Testing nonnegative continuous compact
functions in \textup{(FD9)}, and using Radon regularity, gives
$0\leq\mu\leq M\widehat\mu_v$ on all Borel sets.

This identification on original compact tests does not yet prove
\textup{(FD6)}. Fix an arbitrary $h\in\mathbb D_S$ and take
the approximations \textup{(FD2)}. For each fixed $k$, weak-star
convergence may be evaluated on the actual fixed vector
$h-\phi_k\in\mathbb D_S$. Therefore
\[
 \begin{split}
 |\ell(h-\phi_k)|
   &=\lim_j|\ell_{\mu_j}(h-\phi_k)|
       \leq M\int_V|h-\phi_k|\,d\widehat\mu_v,\\
 \int_V|h-\phi_k|\,d\mu
   &\leq M\int_V|h-\phi_k|\,d\widehat\mu_v.
 \end{split}
 \tag{FD10}
\]
Both errors tend to zero. Equality on $\phi_k$ consequently
passes to equality on $h$, proving \textup{(FD6)} on the entire
space. This passage does not assume that the approximations
converge in its norm. Uniqueness follows already on original
compact scalar tests. Estimate \textup{(FD8)} applies to the
limit measure, so its functional is strongly continuous.

For compactness, the fixed strict box has a countable compact
exhaustion and $\widehat\mu_v$ is $\sigma$-finite. Every dominated
measure has a unique density $b\in L^\infty(\widehat\mu_v)$
with $0\leq b\leq M$, modulo measure-zero equality. The map
\[
 b\longmapsto \ell_{b\widehat\mu_v}
 \tag{FD11}
\]
is weak-star continuous from the weak-star compact interval
$\{0\leq b\leq M\}\subset L^\infty(\widehat\mu_v)$ into
$\mathbb D_S^*$: each $h\in\mathbb D_S$ is an $L^1$ tester
by \textup{(FD1)}. Banach--Alaoglu, applied to this closed order
interval, shows its image is compact. It is injective because
equality of its images on original compact tests identifies
the Radon measures, hence their densities. Its image is exactly
$\mathfrak I_{v,M}$. A compact subset of the Hausdorff weak-star
dual is closed; it is consequently also norm closed. The case
$\widehat\mu_v=0$ is the singleton zero interval and satisfies
all the same conclusions.
\end{proof}

\begin{proposition}[Vague convergence becomes full energy-test convergence inside the interval]
\label{prop:r12-fixed-domination-moving-testers}
Let $\mu_j,\mu$ all satisfy
$0\leq\mu_j,\mu\leq M\widehat\mu_v$. The following are
equivalent:
\begin{enumerate}
\item $\mu_j\to\mu$ on all original smooth compact scalar tests;
\item $\mu_j\to\mu$ vaguely on $C_c(V)$;
\item $\ell_{\mu_j}\to\ell_\mu$ in
$\sigma(\mathbb D_S^*,\mathbb D_S)$.
\end{enumerate}
In this case, for every strongly convergent cone row $u_j\to u$,
\[
                  \int_Vu_j\,d\mu_j
                         \longrightarrow\int_Vu\,d\mu.
 \tag{FD12}
\]
No dual-norm convergence is inferred from these three equivalent
weak convergence statements.
\end{proposition}

\noindent\textbf{Status: \texttt{NEW\_PROVED}} for the
fixed dominated literal interval.

\begin{proof}
The common domination bounds all masses on each fixed compact,
so uniform original scalar density makes the first two statements
equivalent. The third implies the first because every original
compact scalar belongs to $\mathbb D_S$.
Conversely, fix $h\in\mathbb D_S$ and use \textup{(FD2)}.
The two approximation errors in its integrals against $\mu_j$
and $\mu$ are bounded uniformly in $j$ by
$M\int|h-\phi_k|\,d\widehat\mu_v$. Pass first in $j$ on the
fixed compact test $\phi_k$, then let $k\to\infty$.
This proves the third statement. Finally
\[
 \left|\int_Vu_j\,d\mu_j-\int_Vu\,d\mu\right|
 \leq M\int_V|u_j-u|\,d\widehat\mu_v
         +\left|\int_Vu\,d(\mu_j-\mu)\right|.
\]
The first term tends to zero by completed \textup{(EV1)}, and
the second by the just-proved convergence on the fixed energy
tester $u$. This proves \textup{(FD12)}.
\end{proof}

\begin{proposition}[Pointwise finite energy integrals already have a uniform bound]
\label{prop:r12-pointwise-finite-literal-dual}
Let $\lambda$ be a positive original Radon measure which
annihilates every locally outer-capacity-null Borel set. The
following are equivalent:
\begin{enumerate}
\item $\int_V(-a)\,d\lambda<\infty$ for every $a\in\mathscr C$;
\item every $h\in\mathbb D_S$ belongs literally to $L^1(\lambda)$;
\item there is a finite $A$ such that
$\int_V|h|\,d\lambda\leq A\|h\|_{\rm ed}$ for every
$h\in\mathbb D_S$.
\end{enumerate}
When these conditions hold, the optimal constant is
\[
 A_\lambda
   =\sup_{\substack{a\in\mathscr C\\r(a)\leq1}}
                  \int_V(-a)\,d\lambda
   =\|\ell_\lambda\|_*<\infty,
 \qquad \ell_\lambda(h)=-\int_Vh\,d\lambda.
 \tag{FD13}
\]
This proves bounded literal dual membership, not membership
in $\mathbb D_{S,\mathrm s}^*$ and not the source equation.
\end{proposition}

\noindent\textbf{Status: \texttt{DIRECT\_COROLLARY}} of the
proved summable-cone-series construction. The positive Radon
measure is not assumed to have finite unweighted total mass.

\begin{proof}
Only the implication from the first condition to the third
requires proof. Suppose that the displayed supremum defining
$A_\lambda$ were infinite. A zero-energy potential is zero
outside a locally outer-null set, so it has zero integral.
Rescaling nonzero cone members therefore produces $a_k\in
\mathscr C$ with
\[
                r(a_k)=1,\qquad
                \int_V(-a_k)\,d\lambda\geq4^k.
\]
The partial sums $w_N=\sum_{k=1}^N2^{-k}a_k$ decrease
and have energy roots at most $\sum_k2^{-k}=1$. The
summable-root cone construction in \textup{(ED7)--(ED8)},
equivalently the frozen decreasing energy-bounded closure,
gives $w\in\mathscr C$ whose relative Borel representative
is their pointwise decreasing limit off one locally outer-null
set. The measure annihilates that set. Monotone convergence
of the nonnegative functions $-w_N$ yields
\[
          \int_V(-w)\,d\lambda
              =\sum_{k\geq1}2^{-k}
                         \int_V(-a_k)\,d\lambda
              \geq\sum_{k\geq1}2^k=\infty,
\]
contradicting the first condition. Thus $A_\lambda<\infty$.
For $h=a-b$, the pointwise bound $|h|\leq-a-b$ gives
$\int|h|\,d\lambda\leq A_\lambda(r(a)+r(b))$.
Taking the infimum over decompositions proves the third
condition with that constant. The remaining implications
are immediate. Its optimality and the dual-norm identity
follow by testing $h=a$ on the normalized cone, while
the just-proved estimate supplies the reverse bound.
Representative independence uses the printed null-set
assumption throughout; it is not deduced from finite
integrals on one arbitrary choice of representatives.
\end{proof}

\paragraph{Exact consequence and limits.}
The intersection
$\mathfrak I_{v,M}\cap J(\mathscr C)$ is dual-norm closed by
Theorem~\ref{thm:r12-closed-full-source-image}. A dual-norm
convergent row in the interval therefore has a literal strong-dual
limit, and its variational maximizers converge strongly by
\textup{(NI11)}. If every datum in the row is a source, its limit
is a source as well. These are \texttt{DIRECT\_COROLLARY}
statements; no equality between the signed dual norm and the
absolute-variation distance is used.

It would be \texttt{FALSE} to conclude that every member of a
fixed interval is a source. The original smooth-density examples
of Section~\ref{sec:r12-smooth-density-source-instability} are
contained in one such interval, have positive densities bounded
above and below, and have a nonzero normal defect. Likewise this
theorem does not assert that all literal finite-energy data belong
to a union of fixed-source dominated intervals. It repairs
literalness under fixed domination, not the still-open general
literal-continuity implication \textup{(LG1)}.

The nonliteral norm limits of the proper-moment mollifier example
cannot occur in \textup{(FD5)}: its approximating domination
constants are not bounded for one fixed completed source. Compact
original-test agreement alone still does not determine an arbitrary
energy functional; the uniform domination and the full
$L^1(\widehat\mu_v)$ approximation in \textup{(FD10)} are the
additional ingredients that remove that obstruction here.

\paragraph{Sources, refinements and dependency check.}
The only internal inputs are the original strong-cone literal
calculus and tester continuity \textup{(IC), (EV1)}, the mixed
energy inequalities, and the exact original compact scalar
realization \textup{(HC14)--(HC15)}. Original scalar density is
the same frozen local calculus used in \textup{(CS2)}. The
external functional-analytic inputs are the Radon $L^1$ density,
positive Riesz theorem, Radon--Nikodym theorem and Banach--Alaoglu.
All arguments concern the original ordinary space over the fixed
field and the complete original scalar test sheaf. A new actual
moment refinement changes neither the Borel measure nor any
integral; its smooth compact tests are already included. No
Gauss-residue current, auxiliary field, unrestricted source
existence, or raw/outer-capacity equality is invoked.

% END REVIEWED MODULE research_r12/nonlinear/fixed_domination_literal_closure.tex


\part{Positive Jensen currents and full-test weak defects}

% BEGIN REVIEWED MODULE research_r12/capacity/positive_jensen_defect_currents.tex
\section{Positive Jensen defect currents in every original degree}
\label{sec:r12-positive-jensen-defects}

The field is the fixed finite extension $K/\mathbf Q_p$, with its
original absolute value.  Let $X$ be good, Hausdorff, boundaryless,
quasi-smooth and pure of dimension $n$.  On an ordinary open set
$U\subset X$ fix an arbitrary original symmetric middle-positive
$(p,p)$-current $S$ which is separately $d'$- and $d''$-closed, and
put $q=n-p\geq1$ and $D=d'd''$.  Every equality below is an equality
of original currents on every compact smooth complementary CLD test,
including tests on every actual moment refinement.  No equality on
additional GK delta tests is inferred.

The bounded mixed calculus and positive energy matrices are
\texttt{INTERNAL\_INPUT}.  Specifically, if $Q$ is a permitted bounded
mixed product over $S$ and $f=a-b$ is a real difference of bounded
classical psh functions, then
\[
 \Gamma_Q(f)=\tfrac12D(f^2Q)-fD(fQ)
 \tag{JD1}
\]
is a positive coefficient-Radon current.  Signed expressions involving
$Df$ below always mean the finite multilinear expansion of $Da-Db$;
they never mean the product of arbitrary currents.

\begin{lemma}[The differential of an energy current]
\label{lem:r12-differential-energy-current}
If $Q$ contains $k-2$ differentiated bounded slots over $S$, where
$2\leq k\leq q$, then
\[
                 D\Gamma_Q(f)=-(Df)^2\wedge Q.
 \tag{JD2}
\]
The right side is an original signed coefficient-Radon current.
In particular, an energy current is not in general separately closed.
\end{lemma}

\noindent\textbf{Status: \texttt{DIRECT\_COROLLARY}} of the frozen
bounded recursion and energy identity, with the sign recorded explicitly.

\begin{proof}
Apply $D$ distributionally to (JD1).  Since $D^2=0$, its first
term vanishes.  Expanding the remaining term in $f=a-b$, bounded
recursion and mixed symmetry give
\[
 \begin{split}
 D\bigl(fD(fQ)\bigr)
 &=D\bigl(aD(aQ)\bigr)-D\bigl(aD(bQ)\bigr)
       -D\bigl(bD(aQ)\bigr)+D\bigl(bD(bQ)\bigr)\\
 &=(Da)^2Q-2Da\wedge Db\wedge Q+(Db)^2Q.
 \end{split}
\]
All four multipliers are bounded and all four recursions are within
the frozen degree range.  This proves (JD2) without differentiating
a Borel restriction.  As an independent smooth sign check,
$\Gamma_Q(f)=d'f\wedge d''f\wedge Q$ and separate closedness gives
$d''\Gamma_Q(f)=-Df\wedge d''f\wedge Q$, hence
$D\Gamma_Q(f)=-Df\wedge Df\wedge Q$.
\end{proof}

\begin{theorem}[An exact positive primitive for the Jensen defect]
\label{thm:r12-positive-jensen-current}
Let $h_1,\ldots,h_N$ be locally bounded classical psh functions on
$U$, let $\lambda_i\geq0$ with $\sum_i\lambda_i=1$, and put
\[
 w=\sum_i\lambda_i h_i,\qquad f_i=h_i-w.
\]
For $2\leq k\leq q$ define the finite current
\[
 V_k(\boldsymbol\lambda,\mathbf h)
  =\sum_{i=1}^N\lambda_i\sum_{a=0}^{k-2}(k-1-a)
       \Gamma_{(Dh_i)^a(Dw)^{k-2-a}\wedge S}(f_i).
 \tag{JD3}
\]
It is a symmetric middle-positive coefficient-Radon current of degree
$(p+k-1,p+k-1)$, and
\[
 \boxed{\quad
 \sum_i\lambda_i T_{S,k}(h_i^{[k]})
       =T_{S,k}(w^{[k]})-D V_k(\boldsymbol\lambda,\mathbf h).
 \quad}
 \tag{JD4}
\]
For $k=1$ the same equality holds with $V_1=0$.  The assertion does
not require a fixed moment presentation for the $h_i$, does not
restrict the original background, and does not assert that $V_k$
is separately closed.
\end{theorem}

\noindent\textbf{Status: \texttt{NEW\_PROVED}} relative to the frozen
R11 modules.  This is a finite current identity over every allowed
background, not a theorem on products of unrelated currents.

\begin{proof}
All summands in (JD3) exist by the bounded energy calculus.  Their
backgrounds are positive separately closed bounded mixed products,
and the positive energy matrix applied to the real difference $f_i$
makes every summand middle-positive.  Thus their finite positive
sum has the asserted degree and coefficient-Radon property.

Fix any original compact complementary test and a common ordinary
neighborhood on which the finitely many potentials are bounded.
Multilinearity makes
\[
 P_i(t)=T_{S,k}(((1-t)w+th_i)^{[k]})
\]
a polynomial with values in original currents.  Termwise
differentiation gives
\[
 \begin{split}
 P_i'(0)&=k\,Df_i\wedge(Dw)^{k-1}\wedge S,\\
 P_i''(t)&=k(k-1)(Df_i)^2\wedge
                         ((1-t)Dw+tDh_i)^{k-2}\wedge S.
 \end{split}
 \tag{JD5}
\]
Taylor's formula with integral remainder is an identity of finite
polynomials, so it applies on that test without any topological
integration assumption.  The weighted first derivatives cancel
because $\sum_i\lambda_i f_i=0$.  Apply (JD2) to the remainder.
The coefficient of
$\Gamma_{(Dh_i)^a(Dw)^{k-2-a}S}(f_i)$ in its positive primitive is
\[
 k(k-1)\binom{k-2}{a}
       \int_0^1t^a(1-t)^{k-1-a}\,dt=k-1-a.
 \tag{JD6}
\]
Indeed the integral is $a!(k-1-a)!/k!$.  This proves precisely
(JD3)--(JD4).  In degree one, bounded recursion is linear in the
potential, giving the assertion directly.

The chosen test was arbitrary.  A finite partition handles its
ordinary compact support.  The frozen bounded identities hold after
every actual moment refinement, so the same polynomial calculation
proves (JD4) on each such refinement.  Only genuine equality of the
finite affine combination of the actual functions was used; no
analytic PS/PL relation quotient or GQ statement entered the proof.
\end{proof}

\begin{corollary}[The quadratic variance current]
\label{cor:r12-quadratic-variance-current}
For $q\geq2$ one has
\[
 \begin{split}
 V_2&=\sum_i\lambda_i\Gamma_S(h_i-w)
       =\sum_i\lambda_i\Gamma_S(h_i)-\Gamma_S(w)\\
    &=\tfrac12\sum_{i,j}\lambda_i\lambda_j
                              \Gamma_S(h_i-h_j),\\
 \sum_i\lambda_i(Dh_i)^2\wedge S-(Dw)^2\wedge S&=-D V_2.
 \end{split}
 \tag{JD7}
\]
\end{corollary}

\noindent\textbf{Status: \texttt{DIRECT\_COROLLARY}}.

\begin{proof}
Expand the symmetric bilinear energy pairing and use
$w=\sum_i\lambda_i h_i$.  The last equality is (JD4).
\end{proof}

\subsection{The global first-energy domain and quantitative coefficient bounds}

For the rest of this section impose exactly the strict analytic box
and background-extension hypotheses (TC1)--(TC2) of R11.  Thus $V$
has compact closure in the ordinary interior of one original
$K$-affinoid domain, $\rho$ is the fixed negative box barrier, and
$S$ is defined on a neighborhood of the associated closed box.
Write
\[
 \mathcal C=\widehat{\mathcal E}^1_{\rm str}(V),\qquad
 d=q+1,\qquad e(u)=\widehat B(u^{[d]}),\qquad I(u,v)=\widehat I(u,v).
 \tag{JD8}
\]
The actual classical first-energy class embeds in $\mathcal C$.
The distinct Borel completion, all its original critical currents,
and its strong topology are \texttt{INTERNAL\_INPUT} from R11
\path{strong_first_energy_completion.tex}, (IC10)--(IC20).
Differences are taken off the joint pole set; their values on that
locally outer-null set may be assigned arbitrarily, because the
frozen completed coefficient currents annihilate it.

\begin{proposition}[The identity in the completed first-energy cone]
\label{prop:r12-completed-jensen-current}
For every finite tuple in $\mathcal C$, formulas (JD3)--(JD4) remain
valid with the completed mixed and energy currents.  In particular
all the actual coefficient multipliers occurring in (JD1)--(JD2)
are integrable before differentiation.
\end{proposition}

\noindent\textbf{Status: \texttt{DIRECT\_COROLLARY}} of
Theorem~\ref{thm:r12-positive-jensen-current} and the frozen critical
completed calculus.

\begin{proof}
First take actual first-energy potentials, put
$h_{i,j}=\max(h_i,-j)$, and set
$w_j=\sum_i\lambda_i h_{i,j}$.  Then $w_j\downarrow w$; one must
use this affine combination, not the generally different
$\max(w,-j)$, in the finite identity.  Upward stability and the
first-energy cone estimate bound the energies of all these rows.
The products of degree $k$ and the energy currents with
$k-2$ background slots have degree budget $k\leq q<d$.
The two multipliers $f^2Q$ and $fD(fQ)$ in (JD1) have the same
budget $k$.  The decreasing graded continuity theorem (GS3), with
its proved decreasing endpoint extension, therefore passes the
bounded identities (JD1)--(JD4), including all signed multilinear
expansions, to the actual unbounded tuple.  Its currents and
multipliers were independently constructed by (PG3), (PG7), and
(PG10)--(PG12), before this passage to the limit.

For completed tuples choose actual strongly Cauchy defining rows
for the $h_i$ and take their finite affine combinations as the row
for $w$.  The completion is a cone and these affine combinations
converge strongly by (IC18)--(IC20).  Every curvature, multiplier
and energy term just listed converges on all original tests by
(IC11).  Pass the actual identities distributionally.  Positivity
of the energy currents and their coefficient integrability are
already part of (IC10)--(IC11); they are not inferred from this
identity.  Representative independence is the frozen completed
independence modulo local ordinary-open outer capacity.
\end{proof}

\begin{lemma}[An exact energy Jensen gap]
\label{lem:r12-energy-jensen-gap}
For a finite tuple $h_i\in\mathcal C$ and its barycenter $w$, put
\[
                       J=\sum_i\lambda_i e(h_i)-e(w).
\]
Then
\[
 \begin{split}
 J&=\sum_i\lambda_i\sum_{a=0}^{q-1}(q-a)
       \int_V\Gamma_{\widehat N(h_i^{[a]},w^{[q-1-a]})}(h_i-w),\\
 \sum_i\lambda_i I(h_i,w)&\leq J
             \leq q\sum_i\lambda_i I(h_i,w).
 \end{split}
 \tag{JD9}
\]
\end{lemma}

\noindent\textbf{Status: \texttt{DIRECT\_COROLLARY}} of the global
first-energy variation (GE16), discrepancy identity (GE18), and
completed mixed-energy continuity.  The bounded common-collar
Bregman coefficient argument was already frozen in R3
\path{energy_convex_blocks.tex}; the finite coefficient
calculation is not claimed anew.

\begin{proof}
The global mixed energy is a symmetric polynomial of degree $d$.
For $u_t=(1-t)w+th_i$, the global variation gives
\[
 \frac{d^2}{dt^2}e(u_t)=q(q+1)
             \int_V\Gamma_{\widehat N(u_t^{[q-1]})}(h_i-w).
\]
The weighted first derivatives at $t=0$ cancel.  In the integral
remainder, the coefficient of the background
$h_i^{[a]},w^{[q-1-a]}$ is
\[
 q(q+1)\binom{q-1}{a}\int_0^1t^a(1-t)^{q-a}\,dt=q-a.
\]
All global integrals are finite by (GE7) and (IC13).  The positive
terms with coefficient one sum to $I(h_i,w)$ by (GE18), proving
both inequalities in (JD9).
\end{proof}

\begin{theorem}[A quantitative bound on every original defect coefficient]
\label{thm:r12-defect-coefficient-bound}
Fix $M<\infty$, an integer $2\leq k\leq q$, and an arbitrary
original compact smooth coefficient test $\eta$ of complementary
degree $(q-k+1,q-k+1)$.  For every finite tuple $h_i\in\mathcal C$
with $e(h_i)\leq M$, and every convex barycenter $w$, one has
\[
 \boxed{\quad
 \|\mu_{V_k,\eta}\|_{\rm TV}
          \leq C_{\eta,q,M} J^{\theta_k},\qquad
       \theta_k=2^{-(q-k+1)}.
 \quad}
 \tag{JD10}
\]
Here $J$ is (JD9), and the constant is independent of the number
of potentials and of the convex weights.  For every additional
actual moment refinement, the same assertion holds with the
control constructed from that actual refined test.  It is a bound
for the whole coefficient measure, not merely its total signed mass.
\end{theorem}

\noindent\textbf{Status: \texttt{NEW\_PROVED}}.

\begin{proof}
Put $m=q-k+1\geq1$.  The original coefficient compensation theorem
(DG2) gives a normalized original smooth psh function $0\leq a\leq1$
on $V$ and a finite $A_\eta$ such that every positive current $T$
of the degree of $V_k$ satisfies
\[
 |\mu_{T,\eta}|
   \leq A_\eta\mathbf1_{\operatorname{supp}\eta}(Da)^m\wedge T.
 \tag{JD11}
\]
Choose $A>0$ with $A\rho<-2$ near $\operatorname{supp}\eta$, and set
$\psi=\max(a-2,A\rho)$.  By the finite-mass tester core theorem,
$\psi$ is an actual first-energy member, it equals $a-2$ near this
support, and
\[
               e(\psi)\leq 2A^q\mu_\rho(V)=:B_\eta<\infty.
 \tag{JD12}
\]
This is an actual finite-energy patch on the original domain, not
an assertion that a local tester belongs globally to the boundary class.

For fixed $i$ put $h=h_i$, $f=h-w$, $z=h+w$, and
$M_* =\max\{2^dM,B_\eta\}$.  Energy convexity gives $e(w)\leq M$;
Minkowski gives $e(z)\leq2^dM$.  The transfer recurrence (EV2)--(EV4),
which is valid in the completion by (IC13), gives for $0\leq s\leq q-1$
\[
 \begin{split}
 L_s&:=\int_V\Gamma_{\widehat N(\psi^{[s]},z^{[q-1-s]})}(f),\\
 L_0&\leq2^{q-1}I(h,w),\qquad
 L_{s+1}\leq(L_sM_*)^{1/2},\\
 L_m&\leq (2^{q-1}I(h,w))^{\theta_k}M_*^{1-\theta_k}.
 \end{split}
 \tag{JD13}
\]
The index $m\leq q-1$ because $k\geq2$.

For each $0\leq a_0\leq k-2$, smooth-slot energy compatibility
and ordinary locality identify its trace in (JD11), on that support,
with the trace having $m$ copies of $\psi$ in the background.
The positive expansion of $(Dh+Dw)^{k-2}$ then gives
\[
 \begin{split}
 &\int_{\operatorname{supp}\eta}(Da)^m\wedge
        \Gamma_{\widehat N(h^{[a_0]},w^{[k-2-a_0]})}(f)\\
 &\hspace{35mm}\leq
       \int_V\Gamma_{\widehat N(\psi^{[m]},z^{[k-2]})}(f)=L_m.
 \end{split}
 \tag{JD14}
\]
This is positivity of energies with the explicitly expanded mixed
backgrounds; it is not an order assertion for arbitrary current products.
Sum (JD14) with the positive coefficients in (JD3), whose sum over
$a_0$ is $k(k-1)/2$, and use (JD13).  Concavity of
$t\mapsto t^{\theta_k}$ and (JD9) give
\[
 \sum_i\lambda_i I(h_i,w)^{\theta_k}
       \leq\left(\sum_i\lambda_i I(h_i,w)\right)^{\theta_k}
       \leq J^{\theta_k}.
\]
Consequently one may take
\[
 C_{\eta,q,M}=A_\eta\frac{k(k-1)}2
                2^{(q-1)\theta_k}M_*^{1-\theta_k}.
 \tag{JD15}
\]
If $M_*=0$, every relevant energy trace vanishes by the
constant-one mixed Holder inequality, and the assertion is interpreted
as zero.  All coefficient domination and smooth-slot identities used
above hold for the original refined test chosen at the start.
\end{proof}

\begin{corollary}[Distributional control of the curvature Jensen error]
\label{cor:r12-curvature-jensen-error}
For an original compact smooth test $\xi$ of degree $(q-k,q-k)$,
the hypotheses of Theorem~\ref{thm:r12-defect-coefficient-bound} give
\[
 \left|\sum_i\lambda_i\langle\widehat N_k(h_i),\xi\rangle
                  -\langle\widehat N_k(w),\xi\rangle\right|
       \leq C_{D\xi,q,M}J^{\theta_k}.
 \tag{JD16}
\]
In particular, vanishing energy Jensen gaps force all these finite
curvature Jensen errors to vanish on all original compact tests.
\end{corollary}

\noindent\textbf{Status: \texttt{DIRECT\_COROLLARY}}.

\begin{proof}
Use (JD4), distributional duality with $D\xi$, and (JD10).
\end{proof}

\subsection{Two false inferences excluded by smooth original examples}

\begin{proposition}[Neither closedness nor an order sign follows]
\label{prop:r12-jensen-defect-sign-counterexamples}
The implications ``$V_k$ is positive, hence separately closed'' and
``$V_k$ is positive, hence $-DV_k$ has a fixed positivity sign'' are
\texttt{FALSE}, already for $k=2$.  The counterexamples exist in every
ambient dimension $n\geq q\geq2$ on original ordinary torus domains.
\end{proposition}

\noindent\textbf{Status: \texttt{NEW\_PROVED}} for the explicit
original-current counterexamples; the two proposed inferences are
\texttt{FALSE}.

\begin{proof}
Take $X=(\mathbf G_{m,K}^n)^{\rm an}$, let
$t_i=\log|T_i|$, restrict to an ordinary inverse image of a bounded
open box in $\mathbf R^n$, and put
\[
 \omega_i=d't_i\wedge d''t_i,\qquad
 \beta=\sum_{i=1}^n\omega_i,\qquad S=\beta^{n-q}.
\]
Here $S$ denotes the original smooth form regarded as its integration
current.  It is a permitted strongly positive, and therefore
middle-positive, separately closed example.  All potentials used
below are bounded original smooth convex functions on this domain;
a sufficiently large constant $C$ makes them negative if desired.

First take $h_1=(t_1^2+t_2^2)/2-C$, $h_2=0$, with weights $1/2$.
Then $V_2=\Gamma_S(h_1)/4$ and
\[
             -DV_2=\tfrac14(\omega_1+\omega_2)^2\wedge S
                         =\tfrac12\omega_1\wedge\omega_2\wedge S.
 \tag{JD17}
\]
It is nonzero and positive, as detected by a nonzero nonnegative
compact scalar cutoff times $\beta^{q-2}$.  Thus $DV_2\ne0$, so
$V_2$ cannot be separately closed.

Next take $h_1=t_1^2/2-C$ and $h_2=t_2^2/2-C$, with the same weights.
Each individual square $(Dh_i)^2$ vanishes, whereas
$(Dw)^2=\omega_1\wedge\omega_2/2$.  Consequently
\[
                -DV_2=-\tfrac12\omega_1\wedge\omega_2\wedge S
 \tag{JD18}
\]
is nonzero and negative on that same positive test.  The positivity
and nonvanishing statements are the ordinary torus top-form integral
on its real box, so they use no delta-test enlargement.  These two
examples refute both order signs.  They do not refute any theorem
with the global zero-boundary hypotheses imposed separately above.
\end{proof}

\paragraph{Precise dependency and novelty boundary.}
The finite formulas use only the frozen bounded original calculus.
The unbounded and completed extensions use R11
\path{energy/power_energy_all_degree_calculus.tex},
\path{comparison/global_first_energy_calculus.tex}, and
\path{comparison/strong_first_energy_completion.tex}.
The quantitative bound uses exactly (DG2), (TC7), and (EV2)--(EV4)
from the corresponding R11 coefficient, tester-core, and energy
topology files.  The old R3 convex-block and Bregman arguments are
not claimed again.  The new mechanism is an explicit positive
original-current primitive for nonlinear convex averaging, together
with its all-original coefficient bound.  It supplies neither
raw/outer capacity equality nor the analytic PS/PL relation quotient.

% END REVIEWED MODULE research_r12/capacity/positive_jensen_defect_currents.tex


% BEGIN REVIEWED MODULE research_r12/capacity/weak_cluster_defect_representation.tex
\section{Strong convex tails and positive-current weak-cluster defects}
\label{sec:r12-weak-cluster-defects}

Retain exactly the strict-box, original-background-extension and
first-energy hypotheses of Section~\ref{sec:r12-positive-jensen-defects}.
Write $\mathcal C=\widehat{\mathcal E}^1_{\rm str}(V)$,
$d=q+1$, $e=\widehat e$, and $I=\widehat I$.  Elements of
$\mathcal C$ are actual Borel classes modulo local ordinary-open
outer $S$-capacity; they are not declared to have ambient classical
psh representatives.  Every current below is an original CLD
current on all compact smooth complementary tests and all actual
moment refinements.

\subsection{Convex tail selection with no prescribed lower barrier}

\begin{theorem}[The minimum-energy center of the convex tails]
\label{thm:r12-strong-convex-tails}
Let $(h_j)_{j\geq1}\subset\mathcal C$ satisfy $e(h_j)\leq M<\infty$.
Set
\[
 \mathcal H_j=\operatorname{conv}_{\rm fin}\{h_l:l\geq j\},\qquad
 L_j=\inf_{h\in\mathcal H_j}e(h),\qquad L=\lim_j L_j.
 \tag{WC1}
\]
Then $0\leq L_j\uparrow L\leq M$.  For every positive sequence
$\epsilon_j\to0$ choose $w_j\in\mathcal H_j$ with
$e(w_j)\leq L_j+\epsilon_j$.  The row $(w_j)$ converges strongly
to a unique $u\in\mathcal C$, independent of these choices, and
$e(u)=L$.  More precisely, for $l\geq j$,
\[
 I(w_j,w_l)\leq
       \frac{L_l-L_j+\epsilon_j+\epsilon_l}{1-2^{-q}}.
 \tag{WC2}
\]
If the original row is actual classical, each $w_j$ may be taken
actual classical.  No actual representative is asserted for $u$.
\end{theorem}

\noindent\textbf{Status: \texttt{DIRECT\_COROLLARY}} of the R11
sharp midpoint inequality and strong completion.  R3 already proved
all-degree bounded local convex-block selection in
\path{energy_convex_blocks.tex}; the present statement records
its global first-energy version without a common barrier and with
the R11 sharp constant.  It is not counted as a new bounded
convex-block theorem.

\begin{proof}
The first energy is convex: its second derivative along any segment
is $q(q+1)$ times a nonnegative global energy trace, by (GE16) and
(IC13).  Thus $0\leq e(h)\leq M$ for every finite tail convex
combination, proving the assertions concerning $L_j$.

For $l\geq j$, both $w_j,w_l$ and their midpoint belong to
$\mathcal H_j$.  The sharp completed midpoint inequality gives
\[
 (1-2^{-q})I(w_j,w_l)
   \leq e(w_j)+e(w_l)-2e((w_j+w_l)/2)
   \leq L_l-L_j+\epsilon_j+\epsilon_l.
\]
This is (WC2).  Its right side tends uniformly to zero for
$j,l$ large.  Strong completeness (IC20) gives $u$, and energy
continuity gives $e(u)=L$.  Applying the same estimate to two
different choices of approximate minimizers proves that their
limits have discrepancy zero.  Hence the relative class is unique.
Finite convex combinations preserve actual classical membership,
which proves the last assertion without constructing any new
classical representative of the limit.
\end{proof}

\begin{corollary}[Zero convexification loss is a strong convergence criterion]
\label{cor:r12-no-convexification-loss}
Suppose in addition that $e(h_j)\to E$.  Then $E\geq L$.  If $E=L$,
the original row itself converges strongly to the center $u$ of
Theorem~\ref{thm:r12-strong-convex-tails}.  In particular all its
original curvature, first-order and energy currents have the strong
completed limits furnished by R11, without taking convex combinations.
\end{corollary}

\noindent\textbf{Status: \texttt{DIRECT\_COROLLARY}}.

\begin{proof}
Since each $h_l$ with $l\geq N$ belongs to $\mathcal H_N$,
$L_N\leq\inf_{l\geq N}e(h_l)$, so $L\leq E$.  For $i,j\geq N$
the same midpoint argument gives
\[
 (1-2^{-q})I(h_i,h_j)\leq e(h_i)+e(h_j)-2L_N.
\]
If $E=L$, its right side tends uniformly to zero as $N\to\infty$.
Thus the row is strongly Cauchy.  Replacing one of $h_i,h_j$ by
an approximate tail minimizer proves that its strong limit is $u$.
The stated current conclusions are the frozen strong continuity
theorem, not a consequence of weak convergence of the measures alone.
\end{proof}

\subsection{Noncircular original-current compactness}

\begin{lemma}[Sequential compactness from original coefficient bounds]
\label{lem:r12-original-current-compactness}
Let $(P_j)$ be symmetric middle-positive original coefficient-Radon
currents of a fixed degree on $V$.  Suppose that for every original
compact smooth coefficient form $\eta$ the total variations
$\|\mu_{P_j,\eta}\|_{\rm TV}$ are uniformly bounded.  Then a
subsequence converges to a symmetric middle-positive original
coefficient-Radon current on every original compact smooth test and
every actual moment refinement.  The coefficient measures for every
fixed such $\eta$ converge weakly against continuous scalar weights.
Separate closedness passes to the limit if it holds for every $P_j$;
closedness is not otherwise asserted.
\end{lemma}

\noindent\textbf{Status: \texttt{DIRECT\_COROLLARY}} of (AE2), the
frozen coefficient-Radon calculus, and weak compactness of bounded
Radon measures on compact metrizable spaces.  The proof is included
to exclude a hidden compactness or current-representation assumption.

\begin{proof}
The containing affinoid is compact metrizable.  Indeed its Banach
algebra is separable over the fixed finite extension of $\mathbf Q_p$,
and evaluations of a countable dense subring separate its points.
The exact norm replacement (AE2) supplies countably many ordinary
moment charts, closed under finite common refinements, on which
every finite actual unit moment is locally identical to one of
these moments.  Choose countable compact exhaustions and original
cutoffs in these charts.  Standard coordinate wedge forms times
these cutoffs give a countable set of compact coefficient frames
$\eta_l$ with the following property: every original compact test
is a finite sum $\sum_l f_l\eta_l$, where $f_l$ are continuous
scalar functions on the relevant compact supports.  One first
makes the exact local replacement, then a finite original smooth
partition on its compact support; there is no approximation of
the differential or moment coordinates.

For each $l$, the signed measures $\mu_{P_j,\eta_l}$ are supported
on one compact metrizable set and have uniformly bounded total
variation.  Choose a countable uniformly dense set of continuous
functions on that compact set.  The diagonal scalar subsequence
argument, together with the uniform variation bound, gives a
weakly convergent subsequence of these signed measures.  A second
diagonal argument makes this true for every $l$ simultaneously.

Define the value of the limiting functional on
$\sum_l f_l\eta_l$ by the corresponding finite sum of limiting
integrals.  This is independent of the representation: every exact
relation has zero value for every $P_j$, so its finite sum of
integrals still has zero value after passage to the limit.  On a
fixed compact original moment stage the coefficient variation
bounds give continuity for the order-zero coefficient seminorms,
hence in particular for the original smooth test topology.
The functional is therefore an original coefficient-Radon current.
Positivity and symmetry pass by testing each of their defining
original cones and identities; no cone is replaced by a different
strong or weak cone.

For an arbitrary fixed coefficient form $\eta=\sum_l f_l\eta_l$,
the same finite argument with an additional continuous scalar
weight proves weak convergence of $\mu_{P_j,\eta}$.  Consequently
its variation is bounded by the lower limit of the corresponding
variation bounds.  Every additional actual refinement has the
same exact local replacement and finite representation, so it is
included.  If the currents are separately closed, evaluate them
on $d'$ or $d''$ of an arbitrary compact test and pass to the
limit.  No such assertion follows without that assumption.
\end{proof}

\begin{lemma}[Uniform first-energy bounds control every curvature coefficient]
\label{lem:r12-energy-ball-curvature-compactness}
For each $0\leq k\leq q$, every compact original coefficient form
$\eta$ of complementary degree $(q-k,q-k)$, and every finite
$M$, the family
\[
           \{\widehat N_k(h):h\in\mathcal C,\ e(h)\leq M\}
 \tag{WC3}
\]
has uniformly bounded $\eta$-coefficient variation.  It is therefore
sequentially weakly relatively compact in the sense of
Lemma~\ref{lem:r12-original-current-compactness}.  The assertion is
local on compact supports and makes no claim of tightness at the
ordinary boundary for unweighted top measures.
\end{lemma}

\noindent\textbf{Status: \texttt{DIRECT\_COROLLARY}} of the global
mixed-energy Holder inequality and the analytic coefficient patch.

\begin{proof}
Choose $a,A_\eta$ by (DG2), and patch $a-2$ to an actual
$\psi=\max(a-2,A\rho)\in\mathcal E^1$ equal to $a-2$ near
$K=\operatorname{supp}\eta$, exactly as in (JD12).
Let $c_K=\min_K(-\rho)>0$; the case of empty support is trivial.
Smooth-slot compatibility and positivity give
\[
 \begin{split}
 \|\mu_{\widehat N_k(h),\eta}\|_{\rm TV}
 &\leq A_\eta\int_K (Da)^{q-k}\wedge\widehat N_k(h)\\
 &\leq\frac{A_\eta}{c_K}
       \int_V(-\rho)\widehat N(h^{[k]},\psi^{[q-k]})\\
 &\leq\frac{A_\eta}{c_K}
              e(\rho)^{1/d}M^{k/d}e(\psi)^{(q-k)/d}.
 \end{split}
 \tag{WC4}
\]
There are $q$ differentiated slots in the middle integral and
one undifferentiated slot, so the final inequality is precisely
the frozen constant-one completed Holder inequality (IC13).
For $k=0$ interpret $M^0=1$; the current is the fixed $S$.
For $k=q$ no differentiated $\psi$ slot is needed.  All constants
depend on the actual test and its support, not on $h$.
Apply Lemma~\ref{lem:r12-original-current-compactness}.
\end{proof}

\subsection{The nonlinear weak cluster, with its positive defect retained}

\begin{theorem}[Positive-current representation of every energy-bounded weak cluster]
\label{thm:r12-energy-weak-cluster-defect}
Let $(h_j)\subset\mathcal C$ have $e(h_j)\leq M$.  After passage
to a subsequence there are an $E\in[0,M]$, original symmetric
middle-positive separately closed currents $\nu_k$ for $1\leq k\leq q$,
a completed class $u\in\mathcal C$, and symmetric middle-positive
coefficient-Radon currents $V_k$ for $2\leq k\leq q$, such that
\[
 \begin{split}
 e(h_j)&\longrightarrow E,\qquad
           \widehat N_k(h_j)\longrightarrow\nu_k,\\
 e(u)&=L\leq E,\qquad
           \nu_1=\widehat N_1(u),\\
 \boxed{\quad\nu_k=\widehat N_k(u)-DV_k\quad}
                 &\qquad(2\leq k\leq q).
 \end{split}
 \tag{WC5}
\]
Here $L$ and $u$ are the convex-tail minimum and center for that
subsequence.  Every prescribed weak cluster of the top currents
arising from a subsequence is included, by applying the theorem
to that subsequence.  The lower-degree clusters may require a
further subsequence.

For every original compact coefficient test $\eta$ of $V_k$,
\[
       \|\mu_{V_k,\eta}\|_{\rm TV}
             \leq C_{\eta,q,M}(E-L)^{\theta_k},
       \qquad\theta_k=2^{-(q-k+1)}.
 \tag{WC6}
\]
The currents $V_k$ have degree $(p+k-1,p+k-1)$; they are not
asserted to be closed.  Their $D$-derivatives are the signed
coefficient-Radon currents specified in (WC5).  For $E=L$ one
has $V_k=0$, and the original subsequence converges strongly to
$u$.  These assertions include every original test and every
actual moment refinement.
\end{theorem}

\noindent\textbf{Status: \texttt{NEW\_PROVED}} for the all-degree
positive-current defect representation and its quantitative bound.
The bounded $q=1$ completed weak-cluster representation was already
frozen in R3; no new linear closure claim is made.

\begin{proof}
First select a subsequence on which the bounded energies converge
to $E$.  Apply Lemma~\ref{lem:r12-energy-ball-curvature-compactness}
and a finite diagonal selection to obtain all $\nu_k$.  Their
positivity, symmetry and separate closedness follow from
Lemma~\ref{lem:r12-original-current-compactness}.

For this subsequence choose finite convex tail combinations
\[
                 w_j=\sum_{l\geq j}\lambda_{j,l}h_l,
\]
with finite support in $l$ and energies at most $L_j+\epsilon_j$.
Theorem~\ref{thm:r12-strong-convex-tails} gives strong convergence
to $u$ and $e(u)=L$.  In particular
$\widehat N_k(w_j)\to\widehat N_k(u)$ on every original test,
by the already proved R11 strong calculus.  This strong continuity
is used only for the convex blocks, not for the unaveraged row.

Since $\widehat N_k(h_l)\to\nu_k$, every finite convex tail average
of their values on a fixed test also tends to that value of $\nu_k$.
Let $V_{k,j}$ be (JD3) for this block.  The exact identity gives
\[
 \sum_l\lambda_{j,l}\widehat N_k(h_l)
                   =\widehat N_k(w_j)-DV_{k,j}.
 \tag{WC7}
\]
Its energy Jensen gap satisfies
\[
 J_j=\sum_l\lambda_{j,l}e(h_l)-e(w_j)
                       \longrightarrow E-L,
 \tag{WC8}
\]
because the first term is a convex tail average of a convergent
scalar sequence and the second tends to $L$.

The full coefficient bound (JD10) makes each positive row
$V_{k,j}$ locally coefficient bounded, independently of the number
of terms in its block.  Apply
Lemma~\ref{lem:r12-original-current-compactness} and a finite
diagonal selection in $k$ to obtain $V_k$.  Passing (WC7)
distributionally gives (WC5); for $k=1$ use exact linearity with
no defect.  Variation lower semicontinuity for each fixed original
coefficient measure, (JD10), and (WC8) give (WC6).

If $E=L$, (WC6) kills every coefficient of $V_k$.  The stronger
statement that the unaveraged row converges strongly follows from
Corollary~\ref{cor:r12-no-convexification-loss}, not from the
vanishing of its $D$-defect alone.  Every limiting identity is an
original-current identity by the compactness lemma; the proof has
not selected a privileged finite family of actual refinements.
\end{proof}

\begin{corollary}[Weak clusters of the full normalized tester family]
\label{cor:r12-full-tester-positive-defect}
Fix $K\Subset V$ and an ordinary neighborhood $W$ of $K$ with
$\overline W\Subset V$.  Let $a_j$ be any sequence of actual
bounded classical psh functions on $V$ with $0\leq a_j\leq1$.
After shrinking $W$ while retaining $K$, every weak cluster of
their original top currents on $W$ has the form
\[
                  \nu_q=\widehat N_q(u)-DV_q\quad\hbox{on }W,
 \tag{WC9}
\]
where $u\in\mathcal C$ has an actual Borel representative in
$[-1,0]$, and $V_q$ is a positive original $(n-1,n-1)$-current.
There is the simultaneous all-degree version (WC5) on $W$.
For $q=1$, the defect is zero.  The completed class $u$ is not
asserted to be an actual classical normalized tester.
\end{corollary}

\noindent\textbf{Status: \texttt{NEW\_PROVED}} for the nonlinear
positive-defect representation of the full original tester family;
the linear completed representation is the frozen predecessor.

\begin{proof}
Choose $A>0$ such that $A\rho<-2$ on an ordinary neighborhood of
$\overline W$, and put $\delta_j=1/j$ and
\[
                   h_j=\max(a_j-1-\delta_j,A\rho).
 \tag{WC10}
\]
Every $h_j$ is actual bounded classical, has compact negative
sublevels, and belongs to the finite-mass tester core up to its
printed range $[-1-\delta_j,0]$.  The tester-core proof, or its
positive rescaling, gives
\[
 \mu_{h_j}(V)\leq A^q\mu_\rho(V),\qquad
                   e(h_j)\leq2A^q\mu_\rho(V).
 \tag{WC11}
\]
On $W$ the strict branch is $h_j=a_j-1-\delta_j$, so ordinary
locality gives $\widehat N_k(h_j)=T_{S,k}(a_j^{[k]})$ there,
for every original degree.  Constants have no effect on these
bounded products.

Apply Theorem~\ref{thm:r12-energy-weak-cluster-defect} to $(h_j)$,
retaining a subsequence that realizes the prescribed top weak
cluster on $W$.  A convex combination of its $j$th tail satisfies
$-1-\sup_{l\geq j}\delta_l\leq w_j\leq0$.  Its strong capacity
limit therefore lies in $[-1,0]$ outside a local outer-null set.
Choose the clipped Borel representative there and on the null
set.  The completed currents are unchanged.  Restrict (WC5) to
$W$ to prove (WC9) and all lower-degree assertions.

The collars on which $h_j=A\rho$ may shrink with $j$.  Uniform
energy in (WC11), not an unproved common-collar hypothesis, is
what permits the strong convex-tail construction.  All patches
use the original field, norm, and arbitrary extending background.
\end{proof}

\paragraph{What this does and does not prove about raw capacity.}
The raw compact capacity is the supremum over actual normalized
classical testers.  The weak-cluster inequality (FG1) from R11
therefore is not a consequence of (WC9): first, the completed
$u$ need not itself be an actual tester; second, a positive
$V_q$ gives no fixed sign to the signed measure $-DV_q$.
The explicit original smooth examples (JD17)--(JD18) exclude
discarding the latter term by a cone argument.  Even in degree
one, where that term is absent, the completed-to-actual no-gain
statement on a compact set remains an independent requirement.
Thus raw/outer equality for arbitrary compact or Borel sets,
full capacity continuity, unrestricted relative extremals, and
the remaining general Dirichlet statement retain status
\texttt{OPEN}.  No identification of BT and energy null ideals
is made.  The proved advance is a constrained, positive-current,
all-degree description of the nonlinear weak defect, rather
than its unproved removal.

% END REVIEWED MODULE research_r12/capacity/weak_cluster_defect_representation.tex


% BEGIN REVIEWED MODULE research_r12/capacity/universal_jensen_probe.tex
\section{One geometric finite-mass probe detects all Jensen defects}
\label{sec:r12-universal-jensen-probe}

Retain the strict original analytic box $V$, its fixed barrier
$-1\leq\rho<0$, and the containing original affinoid from
(TC1)--(TC3).  The geometric construction below is made before
choosing a background.  For every original symmetric middle-positive
separately closed $(p,p)$-current $S$ extending to the specified
closed box, put $q=n-p\geq1$ and use its already defined bounded
and completed first-energy calculus.  The underlying field and
absolute value never change.  A single construction works for
all these $S$ and all their remaining degrees.

\begin{lemma}[A countable family of geometric coefficient compensators]
\label{lem:r12-countable-geometric-compensators}
There are compact sets $K_l\Subset V$ and original smooth classical
psh functions $a_l$ on $V$, with $0\leq a_l\leq1$, such that the
following holds simultaneously for every $1\leq m\leq n$.
For each original compact smooth form $\eta$ of degree $(m,m)$,
there is a finite set $F_\eta\subset\mathbf N$ and finite constants
$A_{\eta,l}>0$ such that every original middle-positive current
$T$ of degree $(n-m,n-m)$ satisfies
\[
       |\mu_{T,\eta}|
         \leq\sum_{l\in F_\eta}A_{\eta,l}\mathbf1_{K_l}
                                      (Da_l)^m\wedge T.
 \tag{UP1}
\]
Every additional actual moment refinement is included.  The sets,
functions, and constants are independent of $T$ and of the
background $S$.
\end{lemma}

\noindent\textbf{Status: \texttt{DIRECT\_COROLLARY}} of the frozen
exact countable original moment replacement (AE2), the analytic
compensator (AE5), and the original middle-cone coefficient theorem.

\begin{proof}
Choose the countable exact unit-moment charts from (AE2), close
them under finite common refinements, and choose countable compact
cores cofinal among compact subsets of those charts.  On each core
apply (AE5) to its finite moment list.  This gives a normalized
global smooth psh $a_l$ whose curvature dominates a positive
multiple of the coordinate identity form
$\beta_l=\sum_i d't_{li}\wedge d''t_{li}$ near that core.
The choices involve only analytic functions on the fixed affinoid
and its geometry.

For an arbitrary original $\eta$, exact replacement of its finite
moment presentations and a finite original partition reduce its
compact support to finitely many of these cores.  The finite
middle-coefficient matrix bound controls its variation there by
a constant times $\beta_l^m\wedge T$.  The difference between
the corresponding powers of $Da_l$ and a positive multiple of
$\beta_l$ is strongly positive, so pairing with the original
middle-positive $T$ is legitimate.  Absorb the finite constants
to obtain (UP1).  This proof applies to any chosen additional
actual refinement, because its finitely many unit norms have the
same exact countable replacement.  It neither changes cones nor
assumes that a positive energy current is closed.
\end{proof}

\begin{theorem}[A background-independent normalized actual probe]
\label{thm:r12-geometric-finite-mass-probe}
There is a bounded continuous actual classical psh function
$\psi_0$ on $V$, determined by the geometric data alone, with
\[
                         \rho\leq\psi_0<0.
 \tag{UP2}
\]
For every allowed extending original background $S$,
\[
 \psi_0\in\mathscr T_0(S,V),\qquad
 \mu_{\psi_0}(V)\leq\mu_\rho(V),\qquad
 e_\sigma(\psi_0)\leq\mu_\rho(V)\quad(\sigma>0).
 \tag{UP3}
\]
For each $l$ there are $c_l>0$ and actual bounded classical
$\psi_l,r_l\in\mathcal E^1_{S,\partial}(V)$ such that
\[
 \psi_0=c_l\psi_l+r_l,\qquad
            \psi_l=a_l-2\text{ near }K_l.
 \tag{UP4}
\]
The same $\psi_0,c_l,\psi_l,r_l$ work for every $S$.
The probe need not be original smooth: its construction uses a
countable uniformly convergent family of actual moment functions.
\end{theorem}

\noindent\textbf{Status: \texttt{NEW\_PROVED}} for this single
geometric finite-mass probe with the simultaneous decomposition
property.  The finite patch and energy comparison used in the
construction are frozen inputs.

\begin{proof}
Choose $A_l\geq1$ with $A_l\rho<-2$ on a neighborhood of $K_l$,
and define
\[
 \psi_l=\max(a_l-2,A_l\rho),\qquad
 c_l=\frac{2^{-l}}{1+A_l},\qquad
                     \psi_0=\sum_{l=1}^\infty c_l\psi_l.
 \tag{UP5}
\]
Then $-2\leq\psi_l<0$, $\psi_l\geq A_l\rho$, and
$\psi_l=a_l-2$ near $K_l$.  The series converges uniformly on
$V$ to a continuous function.  Its partial sums decrease and
are actual classical psh; their finite continuous limit is
therefore actual classical psh by the frozen bounded decreasing
closure.  No local algebraization or current-positivity converse
is needed for this construction.

Since $\sum_lc_lA_l\leq\sum_l2^{-l}=1$ and $\rho<0$,
\[
           \psi_0\geq\left(\sum_lc_lA_l\right)\rho\geq\rho.
\]
At least one negative summand occurs, so $\psi_0<0$.
Its strict negative sublevels have compact closure because of
this lower barrier.  For any $c>1$, the strict comparison sets
$\{c\rho<\psi_0-\epsilon\}$ have compact closure and exhaust
$V$ as $\epsilon\downarrow0$.  The frozen bounded comparison
therefore gives $\mu_{\psi_0}(V)\leq c^q\mu_\rho(V)$, first
by increasing these sets and then by letting $c\downarrow1$.
The right side is finite by the strict-box barrier theorem.
Because $-1\leq\psi_0\leq0$, all the power-energy bounds in
(UP3) follow.  This proof uses neither a common extension collar
for the infinite series nor a differentiated Borel restriction.

Finally set $r_l=\sum_{j\ne l}c_j\psi_j$.  Exactly the same
uniform convergence and lower-barrier argument gives
$\rho\leq r_l\leq0$, the compact negative-sublevel condition,
and finite total mass.  The finite patch $\psi_l$ has the
same properties with its own finite bound and barrier $A_l\rho$.
Thus both are actual first-energy members, and (UP4) is a literal
identity of actual functions, for every background in question.
\end{proof}

\begin{corollary}[One-sided control of the original background coefficients]
\label{cor:r12-universal-probe-background-control}
For every original compact smooth complementary form $\eta$ of
the background $S$, there is a finite geometric constant $C_\eta$
such that
\[
                         |\mu_{S,\eta}|\leq C_\eta\mu_{\psi_0}.
 \tag{UP6}
\]
This is only one-sided domination.  In particular it does not
identify the coefficient-null ideal of $S$ with the null ideal of
$\mu_{\psi_0}$, or with the relative BT-capacity-null ideal.
\end{corollary}

\noindent\textbf{Status: \texttt{DIRECT\_COROLLARY}}.

\begin{proof}
Apply (UP1) with $m=q$ and $T=S$.  On $K_l$, ordinary locality
gives $(Da_l)^q\wedge S=\mu_{\psi_l}$.  The positive bounded
mixed expansion of $\psi_0=c_l\psi_l+r_l$ gives
$\mu_{\psi_0}\geq c_l^q\mu_{\psi_l}$ on all Borel sets.
Thus $C_\eta=\sum_{l\in F_\eta}A_{\eta,l}c_l^{-q}$ works.
The reverse domination has not been proved and is not implied:
curvature products of a nonsmooth psh probe may charge sets
which all smooth coefficient traces of the original background
ignore.  No null-ideal equality is used below.
\end{proof}

\subsection{Linear control by a mixed Jensen loss}

Fix $S$ and $q\geq2$, and let $\mathcal C$ be its full strong
first-energy completion.  For $2\leq k\leq q$ put $m=q+1-k$ and
define the nonnegative polynomial on that cone
\[
             F_{\psi_0,k}(u)=B(\psi_0^{[m]},u^{[k]}).
 \tag{UP7}
\]
For a finite convex tuple $w=\sum_i\lambda_i h_i$ in $\mathcal C$,
write
\[
 J_{\psi_0,k}(\boldsymbol\lambda,\mathbf h)
       =\sum_i\lambda_iF_{\psi_0,k}(h_i)-F_{\psi_0,k}(w).
 \tag{UP8}
\]

\begin{theorem}[A linear detector for every original Jensen coefficient]
\label{thm:r12-universal-jensen-detector}
The polynomial (UP7) is convex.  Its Jensen loss is
\[
 \begin{split}
 J_{\psi_0,k}
  =\sum_i\lambda_i\sum_{a=0}^{k-2}(k-1-a)
       \int_V\Gamma_{\widehat N(\psi_0^{[m]},h_i^{[a]},
                                     w^{[k-2-a]})}(h_i-w)
       \geq0.
 \end{split}
 \tag{UP9}
\]
For every original compact coefficient test $\eta$ of the positive
Jensen current $V_k$ in (JD3), one has the linear bound
\[
 \boxed{\qquad
              \|\mu_{V_k,\eta}\|_{\rm TV}
                         \leq C_\eta J_{\psi_0,k}.
 \qquad}
 \tag{UP10}
\]
The constant is independent of the number and sizes of the potentials
and the convex weights.  It is geometric, with dependence on the
printed degree and the chosen actual test.  In particular the same
probe controls all original coefficients after every actual moment
refinement.  No product of $D\psi_0$ with an arbitrary nonclosed
current is defined or used in this assertion.
\end{theorem}

\noindent\textbf{Status: \texttt{NEW\_PROVED}}.

\begin{proof}
For $h,w\in\mathcal C$, put $f=h-w$ and $u_t=(1-t)w+th$.
Differentiation of the finite symmetric mixed-energy polynomial
gives
\[
 \frac{d^2}{dt^2}F_{\psi_0,k}(u_t)
  =k(k-1)B(\psi_0^{[m]},f,f,u_t^{[k-2]})
  =k(k-1)\int_V
        \Gamma_{\widehat N(\psi_0^{[m]},u_t^{[k-2]})}(f)\geq0.
 \tag{UP11}
\]
There are $m+k-2=q-1$ differentiated background slots in this
global energy trace.  It is therefore precisely a finite positive
critical energy from (GE7), passed to the completion by (IC13).
This proves convexity.  Taylor expansion and cancellation of
the affine first derivatives prove (UP9); the scalar coefficient
is the same elementary beta integral (JD6).

Apply (UP1) to $V_k$, with its complementary degree $m$.
On $K_l$, the function $\psi_l$ is original smooth and equals
$a_l-2$.  Smooth-slot energy compatibility gives, for each
summand of (JD3),
\[
 \begin{split}
 &\int_{K_l}(Da_l)^m\wedge
     \Gamma_{\widehat N(h_i^{[a]},w^{[k-2-a]})}(h_i-w)\\
 &\qquad\leq
     \int_V\Gamma_{\widehat N(\psi_l^{[m]},h_i^{[a]},
                                         w^{[k-2-a]})}(h_i-w).
 \end{split}
 \tag{UP12}
\]
The inequality merely enlarges the domain of a positive trace;
it does not claim that a Borel restriction is closed.
Using $\psi_0=c_l\psi_l+r_l$ from (UP4), the positive mixed
background expansion gives
\[
 \begin{split}
 &\int_V\Gamma_{\widehat N(\psi_0^{[m]},h_i^{[a]},
                                       w^{[k-2-a]})}(h_i-w)\\
 &\qquad\geq c_l^m
     \int_V\Gamma_{\widehat N(\psi_l^{[m]},h_i^{[a]},
                                         w^{[k-2-a]})}(h_i-w).
 \end{split}
 \tag{UP13}
\]
Every omitted term is an energy with an explicitly positive
closed mixed background containing some copies of $r_l$;
its positivity is the frozen middle-cone energy statement.

Sum (UP12)--(UP13) with the coefficients in (JD3) and compare
with (UP9).  Each $l$-trace is at most $c_l^{-m}J_{\psi_0,k}$.
Equation (UP1) proves (UP10) with the explicit constant
\[
                       C_\eta=
                   \sum_{l\in F_\eta}A_{\eta,l}c_l^{-m}.
 \tag{UP14}
\]
All manipulations are finite mixed-energy expansions.  The
infinite construction of the fixed probe was completed before
any such expansion, and its residual $r_l$ is an actual cone
member.  The proof thus neither invokes GQ nor introduces a
nonsmooth wedge against a nonclosed $V_k$.
\end{proof}

\subsection{One literal weighted mass detects a nonlinear weak-cluster defect}

\begin{lemma}[Tightness for a fixed actual finite-energy weight]
\label{lem:r12-fixed-probe-weight-tightness}
Suppose $h_j\in\mathcal C$ have $e(h_j)\leq M$, and
$\widehat\mu_{h_j}\to\nu$ locally on all original compact smooth
scalar tests.  Then
\[
       \int_V(-\psi_0)\,d\widehat\mu_{h_j}
                      \longrightarrow\int_V(-\psi_0)\,d\nu<\infty.
 \tag{UP15}
\]
The positive weighted measures are uniformly tight on $V$.
No tightness of the unweighted top measures at the ordinary
boundary is asserted.
\end{lemma}

\noindent\textbf{Status: \texttt{DIRECT\_COROLLARY}} of the frozen
global mixed Holder bound and the small-value clipping theorem,
with the actual continuity of the constructed probe used explicitly.

\begin{proof}
Let $\psi_0^{(\epsilon)}=\max(\psi_0,-\epsilon)$.
The frozen small-value energy bound gives
$r(\psi_0^{(\epsilon)})\to0$ as $\epsilon\downarrow0$.
The set $K_\epsilon=\overline{\{\psi_0<-\epsilon\}}$ is compact,
and $\psi_0^{(\epsilon)}=\psi_0$ off $K_\epsilon$.  Hence
\[
 \sup_j\int_{V\setminus K_\epsilon}(-\psi_0)
                                     \,d\widehat\mu_{h_j}
 \leq\sup_j B(\psi_0^{(\epsilon)},h_j^{[q]})
 \leq r(\psi_0^{(\epsilon)})M^{q/d}\longrightarrow0.
 \tag{UP16}
\]
The full weighted masses are bounded by
$r(\psi_0)M^{q/d}$.  Original smooth scalars are uniformly dense
in continuous functions on compact stages, and the local masses
are uniformly bounded by (WC4).  Thus the given original-test
convergence passes continuous compactly supported scalar weights,
including their product with the bounded continuous $-\psi_0$.

The same compact-cutoff passage, followed by increasing compact
exhaustion, gives
\[
        \int_V(-\psi_0^{(\epsilon)})\,d\nu
                  \leq r(\psi_0^{(\epsilon)})M^{q/d}.
\]
It bounds the limiting weighted tail outside $K_\epsilon$.
Now pass local continuous weighted tests first, remove the two
tails using (UP16), and let $\epsilon\downarrow0$.  This proves
(UP15) and uniform weighted tightness.  The argument did not
restrict a weak limit to a moving Borel set.
\end{proof}

\begin{theorem}[A single geometric moment is a full-original no-defect criterion]
\label{thm:r12-single-moment-no-defect}
In the top-degree case of (WC5), put
\[
 \Delta_{\psi_0}(\nu,u)
      =\int_V(-\psi_0)\,d\nu-B(\psi_0,u^{[q]}).
 \tag{UP17}
\]
Then this number is finite and nonnegative.  Every positive
Jensen-current cluster $V_q$ constructed from the corresponding
minimum-energy convex tails satisfies
\[
                    \|\mu_{V_q,\eta}\|_{\rm TV}
                           \leq C_\eta\Delta_{\psi_0}(\nu,u)
 \tag{UP18}
\]
for every original compact complementary coefficient test $\eta$.
For $q\geq2$ the following are equivalent:
\[
 \begin{split}
 \Delta_{\psi_0}(\nu,u)=0
 &\quad\Longleftrightarrow\quad \nu=\widehat\mu_u
                  \text{ on every original test}\\
 &\quad\Longleftrightarrow\quad V_q=0
                  \text{ for every such Jensen-current cluster}.
 \end{split}
 \tag{UP19}
\]
For $q=1$, the first equality and current equality already hold
by linearity; there is no nonlinear Jensen current.

The probe in this criterion was chosen from the analytic geometry,
independently of $S$, the sequence, its weak cluster, and the center.
The criterion concerns one actual continuous weighted mass, not a
chosen finite set of moment tests and not equality on extra GK tests.
\end{theorem}

\noindent\textbf{Status: \texttt{NEW\_PROVED}}.

\begin{proof}
For the convex blocks $w_j=\sum_l\lambda_{j,l}h_l$ used in
(WC5), (UP15) gives
\[
 \sum_l\lambda_{j,l}B(\psi_0,h_l^{[q]})
                      \longrightarrow\int_V(-\psi_0)\,d\nu.
\]
The strong convergence $w_j\to u$ gives
$B(\psi_0,w_j^{[q]})\to B(\psi_0,u^{[q]})$ by (IC13).
Thus their nonnegative scalar Jensen losses (UP8), with $k=q$,
converge to (UP17).  This proves its finiteness and nonnegativity.
Pass the linear coefficient estimate (UP10) through the weak
coefficient convergence of $V_{q,j}$ to obtain (UP18).

If the scalar loss is zero, every original coefficient of $V_q$
vanishes, and (WC5) gives $\nu=\widehat\mu_u$.  If these original
Radon measures agree, their two finite integrals against the
actual weight $-\psi_0$ agree, making (UP17) zero.  Vanishing of
any one extracted Jensen current already gives current equality
by (WC5), and then (UP18) kills all extracted clusters.  This
proves (UP19).  In degree one the same weighted-mass passage
combined with the frozen linear representation proves the final
assertion.
\end{proof}

\paragraph{Relation with the final capacity goal.}
This is a proved reduction of the nonlinear \emph{convex-tail}
defect to one fixed geometric weighted mass, with linear bounds
for every original coefficient.  It does not prove that this
scalar loss vanishes for arbitrary normalized tester sequences:
the frozen moving-atom mismatch (NL1)--(NL3) makes it strictly
positive in every nonlinear degree.  Even when it vanishes, the
normalized center is initially a completed Borel tester.  The
actual normalized representative or compact no-gain gate described
after (NL11) is still separate.  Consequently general compact or
Borel raw/outer capacity equality remains \texttt{OPEN}.

% END REVIEWED MODULE research_r12/capacity/universal_jensen_probe.tex


% BEGIN REVIEWED MODULE research_r12/capacity/no_loss_gate_and_positive_transport.tex
\section{The no-loss gate, its frozen counterexample, and a positive transport primitive}
\label{sec:r12-no-loss-gate}

This section does not claim a new moving-atom counterexample.
Its exact frozen source is R3
\path{source/sections/moving_atomic_dirichlet.tex}, with SHA256
\begin{quote}\small\ttfamily
\nolinkurl{c768ae67e256963a1730d9febc7bc6e9077e34d5ddf0bf867ece2f038922fcf4}
\end{quote}
That file records an earlier reviewed pre-reconstruction hash in
its header.  The frozen corollary
\path{cor:r3-moving-convex-block-mismatch} already proves that
nonlinear unaveraged weak limits need not equal the curvature of
uniformly convergent convex blocks.  R3
\path{capacity_growth_instability.tex}, with SHA256
\begin{quote}\small\ttfamily
\nolinkurl{90f8e2b15c0bc62543f182ca6652c8e945cd7df67999149542dae6c529574337}
\end{quote}
also supplies every fixed power capacity-growth estimate for that
same family.  Both results are \texttt{INTERNAL\_INPUT}.

\subsection{The exact energy loss in the frozen moving-atom family}

\begin{proposition}[A common collar and all power bounds do not imply no loss]
\label{prop:r12-frozen-moving-energy-loss}
For every $q\geq1$, the frozen R3 family, after its original
uniformizer coordinate change to the unit $q$-polydisc $V$, gives
bounded actual classical first-energy potentials $h_j,h$ over
$S=[V]$ and a number $a>0$ with
\[
 \begin{gathered}
 -2a\leq h_j\leq h\leq0,\qquad
 \mu_{h_j}=\delta_{z_j},\qquad \mu_h=\delta_y,\qquad
 z_j\longrightarrow z,\\
 e(h_j)=2a,\qquad e(h)=a,\\
 y:\text{ depth vector }(a,\ldots,a),\qquad
 z:\text{ depth vector }(a,2a,\ldots,2a).
 \end{gathered}
 \tag{NL1}
\]
The functions have one common actual finite-energy lower barrier,
one common ordinary boundary collar, and uniformly bounded energies
of every positive power.  The measures are all carried by one compact
subset of $V$.  The convex-tail infimum and center in (WC1) are
\[
                         L=a,\qquad u=h.
 \tag{NL2}
\]
Thus the convexification loss is exactly $E-L=a>0$.
For $q\geq2$, (WC5) necessarily has a nonzero derivative defect:
\[
                      -DV_q=\delta_z-\delta_y\ne0.
 \tag{NL3}
\]
After division by $2a$, this is a sequence in the normalized
range $[-1,0]$, so it also tests the full normalized tester question.
\end{proposition}

\noindent\textbf{Status: \texttt{DIRECT\_COROLLARY}} of the frozen
example and the new defect representation.  The inference that the
printed controls imply $E=L$ or allow omission of $DV_q$ is
\texttt{FALSE}.  Neither the moving-atom construction nor its
weak-capacity instability is counted as new.

\begin{proof}
In the notation of the frozen source take
$a=-\log|\pi|$, and replace its coordinates $T_l$ by $\pi T_l$.
This is an isomorphism over the original field with its unchanged
absolute value.  The frozen all-original calibrations give (NL1).
The value at each source atom is $-2a$, giving $e(h_j)=2a$;
the value at $y$ is $-a$, giving $e(h)=a$.  More generally,
$e_\sigma(h_j)=(2a)^\sigma$ for every $\sigma>0$.

All $h_j$ agree with $\log\max_l|T_l|$ on the same collar
$\{\max_l|T_l|>e^{-a}\}$ of the unit-polydisc boundary.
The fixed strict-box barrier $\rho$ therefore has a positive
multiple $A\rho$ below all $h_j$: on the compact inner core
choose $A$ large enough to lie below $-2a$, and on the collar
use the displayed common formula and $A\geq1$.
The barrier is actual, bounded, finite-mass and belongs to every
positive power-energy class.

The residue-branch sets on which $h_j\ne h$ are pairwise disjoint,
and $0\leq h-h_j\leq a$.  For any $N$ distinct indices,
\[
          0\leq h-\frac1N\sum h_j\leq\frac aN
                                      \quad\hbox{on all of }V.
 \tag{NL4}
\]
Every finite convex combination is at most $h$, so energy
monotonicity gives its energy at least $e(h)=a$.
The equal-weight combinations in (NL4) have energies tending to
$a$.  Here is a direct global check.  Any mixed top product of
members of this family and $h$ has total mass one: replacing
one slot by $h$ changes that product by $D$ of a compactly
supported bounded multiplier of a closed mixed current, whose
total mass is zero by compact Stokes.  The fixed common collar
is what makes this multiplier compactly supported.  For an average
$w_N$ in (NL4), the global first variation consequently gives
\[
 0\leq e(w_N)-e(h)
   =(q+1)\int_0^1\!\int_V(h-w_N)
              \mu_{(1-t)h+tw_N}\,dt
   \leq(q+1)a/N.
 \tag{NL5}
\]
Every tail admits such averages.  Hence its energy infimum is
$a$, proving $L=a$.  Their uniform convergence and energy
convergence make their strong limit $h$ by (IC20), so the
unique center is $h$.  Formula (WC5) now gives (NL3).

For $q=1$, $y=z$ and the curvature map is linear, so the
derivative defect is zero despite the positive energy loss.
For $q\geq2$ the two original Gauss points are distinct, making
the signed measure in (NL3) nonzero on an original compact scalar
test.  Finally positive scaling divides the top currents by
$(2a)^q$ and the first energies by $(2a)^{q+1}$, while producing
the normalized range.  All conclusions retain the original test
category from the frozen measure calibration.
\end{proof}

\subsection{An explicit positive primitive for ordered Gauss transport}

\begin{proposition}[Coordinatewise depth transport by a positive current]
\label{prop:r12-positive-gauss-transport}
On the original unit $q$-polydisc, let
$0<c_i\leq d_i<\infty$, $1\leq i\leq q$, and let
$\zeta_{\mathbf c},\zeta_{\mathbf d}$ be the original joint
Gauss points of radii $e^{-c_i},e^{-d_i}$.  There is an explicit
positive original $(q-1,q-1)$-current $V_{\mathbf c,\mathbf d}$
such that
\[
                 D V_{\mathbf c,\mathbf d}
                        =\delta_{\zeta_{\mathbf c}}
                              -\delta_{\zeta_{\mathbf d}}.
 \tag{NL6}
\]
It is defined across every coordinate zero locus and satisfies
this equality on all original compact smooth scalar tests and
all actual moment refinements.
\end{proposition}

\noindent\textbf{Status: \texttt{DIRECT\_COROLLARY}} of bounded
recursion and the frozen all-original coordinate Gauss calibration;
it is a new explicit interpretation of the defect, not a new
moving-atom counterexample or a general transport theorem.

\begin{proof}
For a finite positive $s$ write
$b_{i,s}=\max(\log|T_i|,-s)$ on the whole polydisc.
These are bounded continuous classical functions.  Set
\[
 \begin{split}
 g_i&=b_{i,c_i}-b_{i,d_i},\\
 R_j&=\left(\bigwedge_{i<j}Db_{i,d_i}\right)
          \wedge\left(\bigwedge_{i>j}Db_{i,c_i}\right)\wedge[V],\\
 V_{\mathbf c,\mathbf d}&=\sum_{j=1}^q g_jR_j.
 \end{split}
 \tag{NL7}
\]
Every wedge in $R_j$ is an explicitly permitted bounded mixed
product over $[V]$; it is positive and separately closed.
Also $0\leq g_i\leq d_i-c_i$.  On the torus, with
$x_i=-\log|T_i|$, the scalar is
\[
              g_i=(x_i-c_i)_+-(x_i-d_i)_+.
 \tag{NL8}
\]
Near $T_i=0$ it is the constant $d_i-c_i$, so (NL7) is an
actual bounded coefficient multiplier on the entire ordinary
domain.  Its nonnegativity makes $V_{\mathbf c,\mathbf d}$
positive; no closedness of that multiplier current is assumed.

Bounded recursion, applied separately to the two psh terms in
$g_j$, gives
\[
 D(g_jR_j)=(Db_{j,c_j}-Db_{j,d_j})\wedge R_j.
\]
The frozen product Gauss calibration identifies the right side as
\[
 \delta_{\zeta_{(d_1,\ldots,d_{j-1},c_j,\ldots,c_q)}}
 -\delta_{\zeta_{(d_1,\ldots,d_j,c_{j+1},\ldots,c_q)}}.
 \tag{NL9}
\]
Its mass normalization is one in both terms.  Summing (NL9)
telescopes to (NL6).  These are identities of actual original
measures obtained from bounded mixed products, so arbitrary
refinements of a compact scalar test are already included.
\end{proof}

For the frozen nonlinear mismatch take
$\mathbf c=(a,\ldots,a)$ and
$\mathbf d=(a,2a,\ldots,2a)$.  Then $g_1=0$ and the current
(NL7) gives $DV=\delta_y-\delta_z$, with exactly the sign
required by (NL3).  It is not asserted to equal a particular
subsequential Jensen current extracted in (WC5).  Their difference
is $D$-closed; neither separate $d'$- nor separate $d''$-closedness
of that difference follows from this equality alone.

\subsection{Why zero nonlinear defect still does not settle raw capacity}

\paragraph{The compact restriction shortcut is false.}
Uniform or strong potential convergence does not license restriction
of weak curvature limits to a fixed compact set.  On the original
unit $q$-polydisc put
\[
 b_t=\max\{\log\max_i|T_i|,-t\},\qquad
                  \mu_{b_t}=\delta_{\zeta_{(t,\ldots,t)}}.
 \tag{NL10}
\]
For distinct $t_j\to t_0>0$, the functions converge uniformly
to $b_{t_0}$, have a common ordinary collar and convergent first
energies $e(b_{t_j})=t_j$, and therefore converge strongly.
Nevertheless on the compact singleton
$K=\{\zeta_{(t_0,\ldots,t_0)}\}$,
\[
                    \mu_{b_{t_j}}(K)=0,
                  \qquad\mu_{b_{t_0}}(K)=1.
 \tag{NL11}
\]
This is a \texttt{DIRECT\_COROLLARY} of the frozen radial
calibration and a counterexample to the restriction inference,
not a counterexample to raw/outer capacity regularity.

\paragraph{The two remaining statements are distinct.}
The proved condition $E=L$ in (WC5) makes the weak cluster the
curvature of a normalized \emph{completed Borel} tester.  The
statement that this curvature has no extra mass on a compact set,
compared with the supremum over actual classical normalized
testers, is a further \texttt{OPEN} statement.  It follows
immediately if a normalized actual classical representative of
that same relative class is supplied independently.  Such a
representative is not supplied by the strong completion, by a
zero Jensen defect, or by (NL11).  A general representative or
compact no-gain theorem has not been proved here.

\paragraph{A failed two-background repair.}
A proposed repair was to make the convex blocks strongly Cauchy
both for $S$ and for the ambient background $[V]$, then select
the ambient classical representative.  This is a
\texttt{FAILED\_ATTEMPT}: the two relative completions need not
select compatible representatives.  The exact frozen R11
power-core example (PO1)--(PO6) has $S=[Y]$, $u_j=a$ on a
neighborhood of $Y$, and $u_j\uparrow H$ with $H=a$ on $Y$
but $H=0$ off $Y$.  Its ambient upper regularization is zero,
while its $S$-curvature remains the nonzero fixed measure of
$a$.  Replacing the relative class by its ambient representative
therefore loses its prescribed $S$-action.  Adding an auxiliary
background, or identifying its polar ideal with the original
relative one, cannot be used without a separate compatibility
proof.  This observation does not assert that the particular
$S$-class lacks every classical representative: $a$ itself is
one.  It refutes only the proposed automatic selection mechanism.

\paragraph{Current status.}
The all-degree positive Jensen current and its energy-loss bound
are proved.  Uniform boundedness, common barriers, common collars,
weak source convergence, and power-capacity growth do not remove
that loss in general.  Zero loss is a sufficient strong-convergence
criterion, but it does not close the actual compact raw-capacity
gate.  General compact/Borel raw equality with ordinary-open outer
capacity therefore remains \texttt{OPEN}; the frozen q=1
completed-capacity statement is not promoted to an actual-test
statement or recounted as new.

% END REVIEWED MODULE research_r12/capacity/no_loss_gate_and_positive_transport.tex


\part{Source compatibility and the appropriate Dirichlet problem}

% BEGIN REVIEWED MODULE research_r12/nonlinear/tropical_curvature_lifts.tex
\section{Closed positive moment lifts and nonlinear source compatibility}
\label{sec:r12-positive-moment-lifts}

This section proves a necessary source condition for every remaining
degree strictly smaller than a full real moment rank.  It applies to
the weak closure of bounded original Bedford--Taylor sources, hence
to the complete first-energy class.  The special factorization of
the background in \eqref{eq:r12-lift-factorization} is an explicit
hypothesis; it is not asserted for an arbitrary background.  The
result supplies counterexamples to universal source existence, not
a replacement of the arbitrary-background problem by this subclass.

The fixed field is the original finite extension $K/\mathbf Q_p$,
with its original absolute value.  Let $X$ be good, Hausdorff,
boundaryless, quasi-smooth and pure of dimension $n$.  All analytic
currents, forms, positivity notions and tests below are original CLD
objects.  ``Positive'' for a current means symmetric
\emph{middle-positive}; strong and weak positivity are not substituted.
Put $D=d'd''$.

\subsection{A real trace lemma with its normalization}

Let $Q\subset\mathbf R^r$ be open, let $0\le q<r$, put $k=r-q$,
and write
\[
 \beta_0=\sum_{i=1}^r d'x_i\wedge d''x_i,
 \qquad h(x)=\tfrac12|x|^2,
 \qquad D h=\beta_0.
 \tag{ML1}
\]
The real supercurrent convention is the one compatible with the
original CLD pullback of a real superform.  In particular
$\beta_0^k$ has no factorial in its definition.

\begin{lemma}[Positive trace controls all coefficients]
\label{lem:r12-positive-trace-coefficients}
Let $T$ be a symmetric middle-positive real supercurrent of
bidegree $(q,q)$ on $Q$.  Then $T$ has order zero and
\[
                         \nu_T=T\wedge\beta_0^k
 \tag{ML2}
\]
is a positive Radon measure.  Every coefficient measure of $T$ is
absolutely continuous with respect to $\nu_T$, with a bound depending
only on the fixed coordinate coefficient convention.  In particular,
\[
                    \operatorname{supp}T=\operatorname{supp}\nu_T.
 \tag{ML3}
\]
If $(T_j)$ is a sequence of such currents and the measures
$\nu_{T_j}$ have uniformly bounded mass on every compact subset of
$Q$, then a subsequence converges on every compact smooth real test
to a symmetric middle-positive order-zero current $T$.  If each
$T_j$ is separately $d'$- and $d''$-closed, then so is $T$.
\end{lemma}

\noindent\textbf{Status: \texttt{DIRECT\_COROLLARY}} of
finite-dimensional positive-semidefinite coefficient duality and
local weak compactness of Radon measures; a proof is included to
fix the cone and factorial conventions.  No tropical fan or
Gauss-fiber identification is involved.

\begin{proof}
Index increasing $k$-subsets by $I,J$ and use the middle-square
convention
\[
 \Theta_k(A)=(-1)^{k(k-1)/2}
        \sum_{I,J}A_{IJ}\,d'x_I\wedge d''x_J.
\]
For a scalar compact test $\varphi$, let $H_{IJ}(\varphi)$ be
the symmetric matrix of the functional
$A\mapsto\langle T,\varphi\Theta_k(A)\rangle$.
For every $z\in\mathbf R^{\binom rk}$ and every nonnegative
$\varphi$, the test $\varphi\Theta_k(zz^{\mathsf T})$ is
middle-positive.  Thus $z^{\mathsf T}H(\varphi)z\ge0$.
The diagonal entries are positive scalar distributions, hence
Radon measures.  Polarization with $e_I\pm e_J$ shows that each
off-diagonal entry is a signed Radon measure.  Positivity on these
two vectors gives the two measure inequalities
$-\tfrac12(H_{II}+H_{JJ})\le H_{IJ}\le
\tfrac12(H_{II}+H_{JJ})$, and therefore
\[
 |H_{IJ}|\le\tfrac12(H_{II}+H_{JJ}).
 \tag{ML4}
\]
After this domination has been established, Radon--Nikodym
densities with respect to $\sum_IH_{II}$ exist for all entries.
Testing first on rational vectors $z$ shows that their density
matrix is positive semidefinite outside one common null set.

The exact exterior expansion is
$\beta_0^k=k!\sum_I\Theta_k(E_{II})$.  Consequently
\[
                         \nu_T=k!\sum_IH_{II}.
 \tag{ML5}
\]
This proves order zero and domination of every coefficient in the
complementary test basis.  That basis is equivalent, by fixed signs
and complementary indices, to the usual coefficient basis of $T$.
The support equality follows from this domination and the definition
of the trace.

For compactness, choose a countable relatively compact exhaustion of
$Q$ and compact cutoffs equal to one on its successive members.
The hypothesis and \textup{(ML4)--(ML5)} bound the variations of all
coefficient measures on each compact set.  Weak compactness of
finite Radon measures on a compact metric space, followed by a
diagonal extraction over the exhaustion and finitely many entries,
gives simultaneous local weak limits of all coefficients.  These
limits define $T$.  Nonnegativity on every compact smooth
middle-positive test and symmetry pass to the limit.  Finally
$d'$ and $d''$ are continuous on currents by duality, so their
separate vanishing also passes to the limit.
\end{proof}

\begin{definition}[Lift-admissible scalar sources]
\label{def:r12-lift-admissible}
For $0\le q<r$ define the cone
\[
 \mathscr L_{r,q}(Q)=
 \left\{T\wedge\beta_0^{r-q}:
 \begin{array}{l}
 T\text{ is a symmetric middle-positive real }(q,q)\text{-current},\\
 d'T=d''T=0
 \end{array}\right\}.
 \tag{ML6}
\]
A member is called \emph{lift-admissible}.  This is a condition on
a positive scalar measure and a linear closed positive-current
problem.  It does not ask for a potential or a Bedford--Taylor
factorization of the lift.
\end{definition}

\begin{proposition}[Closed cone and a scalar collar inequality]
\label{prop:r12-lift-cone-collar}
The cone $\mathscr L_{r,q}(Q)$ is convex and sequentially closed
for local vague convergence of positive Radon measures.  For every
$\nu\in\mathscr L_{r,q}(Q)$ and every
$\chi\in C_c^\infty(Q)$ with $0\le\chi\le1$, set
\[
 \begin{split}
 A_\chi&=\operatorname{supp}(d\chi),\\
 C_\chi&=\sup_{x\in Q}
 \left\|h(x)\operatorname{Hess}\chi(x)
       +\nabla\chi(x)\otimes x+x\otimes\nabla\chi(x)\right\|_{\rm op}.
 \end{split}
 \tag{ML7}
\]
Then
\[
              \int_Q\chi\,d\nu\le C_\chi\nu(A_\chi),
 \qquad
              \nu\{\chi=1\}\le C_\chi\nu(A_\chi).
 \tag{ML8}
\]
In particular, a lift-admissible measure carried by a compact
subset of $Q$ is zero.
\end{proposition}

\noindent\textbf{Status: \texttt{NEW\_PROVED}} for the
explicit scalar necessary condition used below.  Its compact-source
consequence is the elementary exact-positive-form Stokes obstruction,
proved here without an appeal to convex potentials.

\begin{proof}
Convexity follows by adding positive closed lifts.  If
$\nu_j\to\nu$ locally vaguely, choose any corresponding lifts
$T_j$.  Positivity and a compact cutoff show that their traces
have uniformly bounded local masses.  Lemma
\ref{lem:r12-positive-trace-coefficients} supplies a locally weak
subsequential limit $T$ that is symmetric middle-positive and
separately closed.  Wedge with the fixed smooth form
$\beta_0^{r-q}$ is continuous on currents.  Thus its trace is
$\nu$, proving sequential closedness.

Let $T$ be a lift of $\nu$ and retain $k=r-q\ge1$.
The compact test $\chi h\beta_0^{k-1}$ has bidegree
$(k-1,k-1)$.  Separate closedness gives
\[
 0=\langle T,D(\chi h\beta_0^{k-1})\rangle
   =\int\chi\,d\nu+\langle T,G_\chi\rangle,
 \tag{ML9}
\]
where
\[
 G_\chi=
 \bigl(hD\chi+d'\chi\wedge d''h+d'h\wedge d''\chi\bigr)
                         \wedge\beta_0^{k-1}.
 \tag{ML10}
\]
The $(1,1)$ factor has the symmetric matrix in \textup{(ML7)}.
Its support is contained in $A_\chi$: off that closed set, $\chi$
is locally constant and both its first and second derivatives
vanish.  The operator norm bound therefore gives, pointwise,
\[
            -C_\chi\beta_0^k\le G_\chi
                            \le C_\chi\beta_0^k.
 \tag{ML11}
\]
Indeed both differences are a positive $(1,1)$ form wedged with
the strongly positive form $\beta_0^{k-1}$, and hence are
middle-positive.  The order-zero coefficient representation
allows restriction to the Borel set $A_\chi$; it does \emph{not}
assert that such a restriction is closed.  Pairing \textup{(ML11)}
with the positive current on this set yields
$|\langle T,G_\chi\rangle|\le C_\chi\nu(A_\chi)$.
Equation \textup{(ML9)} proves \textup{(ML8)}.

If $\nu$ is carried by a compact $K\Subset Q$, choose $\chi=1$
on a neighborhood of $K$.  Then $A_\chi\cap K=\varnothing$,
so \textup{(ML8)} gives $\nu(Q)=0$.  Equivalently,
\textup{(ML3)} makes the lift compactly supported, and Stokes
against $\chi h\beta_0^{k-1}$ gives the same conclusion.
\end{proof}

\subsection{Lifts of original bounded sources and their weak limits}

Let $V\subset X$ be ordinary open and let
\[
 t=(\log|f_1|,\ldots,\log|f_r|):V\longrightarrow Q\subset\mathbf R^r
 \tag{ML12}
\]
be an actual original moment, where the $f_i$ are analytic units.
Assume $n\ge r$ and that $t$ is proper over $Q$, in the precise
sense that $t^{-1}(H)$ is compact in $V$ for every compact
$H\subset Q$.  The image is not required to be all of $Q$.
Let $1\le q<r$, and let $R$ be an arbitrary original symmetric
middle-positive separately closed $(n-r,n-r)$-current on $V$.
Put $\beta=t^*\beta_0$ and impose the explicit factorization
\begin{equation}
                 S=\beta^{r-q}\wedge R.
 \label{eq:r12-lift-factorization}
\end{equation}
This $S$ is an allowed original middle-positive separately closed
$(n-q,n-q)$-current.  The wedge is with a strongly positive smooth
form; no wedge of two arbitrary currents is being used.

\begin{theorem}[A closed positive lift survives weak source limits]
\label{thm:r12-positive-moment-lift}
Under \textup{(ML12)} and
\eqref{eq:r12-lift-factorization}, let $w_j$ be original locally
bounded classical psh functions on $V$.  Suppose their positive
original top measures
\[
                 \mu_j=(Dw_j)^q\wedge S
 \tag{ML13}
\]
converge to a positive Radon measure $\mu$ on all original compact
smooth scalar tests.  Then
\[
                            t_*\mu\in\mathscr L_{r,q}(Q).
 \tag{ML14}
\]
Consequently $t_*\mu$ satisfies every collar inequality
\textup{(ML8)}.  If $t_*\mu$ is carried by a compact subset of
$Q$, then $\mu=0$.

The hypotheses use equality and convergence on the full original
test sheaf, including every actual moment refinement.  The one
specified moment is a necessary-condition detector, not a
replacement test category and not a definition of $\mu$.
\end{theorem}

\noindent\textbf{Status: \texttt{NEW\_PROVED}} as a
source-compatibility theorem for all $1\le q<r$.  The bounded
arbitrary-background product, its symmetry, separate closedness
and smooth-factor compatibility are \texttt{INTERNAL\_INPUT}
from the frozen R4 bounded BT calculus.

\begin{proof}
The bounded BT product over the arbitrary background $R$ gives
the original current
\[
                        C_j=(Dw_j)^q\wedge R.
 \tag{ML15}
\]
It is symmetric middle-positive and separately closed, of
bidegree $(n-r+q,n-r+q)$.  Define a real $(q,q)$-current by
\[
 \langle T_j,\eta\rangle=\langle C_j,t^*\eta\rangle,
        \qquad \eta\in A_c^{r-q,r-q}(Q).
 \tag{ML16}
\]
The pullback is an actual original smooth form.  Its support
is contained in $t^{-1}(\operatorname{supp}\eta)$, compact by
properness.  Thus every pairing is legitimate.  Actual pullback
preserves symmetry and the middle form cone and commutes with
$d'$ and $d''$.  These facts, and the corresponding assertions
for $C_j$, prove that $T_j$ is symmetric middle-positive and
separately closed.

For every $\varphi\in C_c^\infty(Q)$, the frozen
smooth-factor identity gives the exact equality
\[
 \begin{split}
 \langle T_j\wedge\beta_0^{r-q},\varphi\rangle
 &=\langle C_j,(\varphi\circ t)\beta^{r-q}\rangle\\
 &=\langle\mu_j,\varphi\circ t\rangle
  =\langle t_*\mu_j,\varphi\rangle.
 \end{split}
 \tag{ML17}
\]
There is no omitted factorial: both the background and trace
use the unnormalized exterior power $\beta^{r-q}$.

By properness, $\varphi\circ t$ is an original compact smooth
scalar test.  The given convergence therefore implies
$t_*\mu_j\to t_*\mu$ on every compact smooth real scalar test.
On each compact $H\Subset Q$, choose a nonnegative smooth
cutoff at least one on $H$.  Its convergent pairings bound
$t_*\mu_j(H)$ uniformly.  The trace lemma now supplies a
subsequence $T_j\to T$ on all compact real smooth tests.
The limit is positive, symmetric and separately closed.
Passing \textup{(ML17)} to the limit yields
\[
                         T\wedge\beta_0^{r-q}=t_*\mu.
 \tag{ML18}
\]
This is \textup{(ML14)}.  The scalar collar inequalities follow
from Proposition~\ref{prop:r12-lift-cone-collar}.

If $t_*\mu$ has compact carrier, that proposition makes it zero.
For every compact original set $L\subset V$, the continuous
image $t(L)$ is compact in $Q$, and positivity gives
\[
              \mu(L)\le\mu(t^{-1}(t(L)))=(t_*\mu)(t(L))=0.
\]
Radon inner regularity then gives $\mu=0$ on $V$.
No converse to \textup{(ML14)} and no uniqueness of $T$ is claimed.
\end{proof}

\paragraph{What the lift does not assert.}
The trace of a positive closed current need not determine that
current.  Accordingly \textup{(ML18)} constructs a nonempty set
of lifts, not a canonical cap-current of the completed potential.
These moment-pushforward lifts are also distinct from the frozen
R9--R10 intrinsic full-Gauss-fiber hinge caps.  Refining the original
moment of an analytic test never changes the original current
\textup{(ML15)} or the original measures in the proof.  It does
not, however, turn \textup{(ML18)} into an unproved canonical
lift on a larger real target: the pullback of $\beta_0$ may be
degenerate on that target.  No refinement cofinality, analytic
PS/PL relation quotient, or ordinary-to-all-delta-test equality
is inferred from this argument.

\subsection{The full strong completion and nonlinear counterexamples}

For the next assertion impose exactly R11 \textup{(TC1)--(TC2)}:
$V$ is a nonempty strict Laurent box with compact closed Laurent
box $C\Subset\operatorname{Int}_X(B)$ in a compact original
affinoid $B$, and $S$ extends as an allowed original current to an
ordinary neighborhood of $C$.  Retain the factorization and proper
moment above on $V$.  Write
\[
 \mathscr C_S=\widehat{\mathcal E}^{1}_{\rm str}(V),
 \qquad \widehat\mu_u=\widehat N_q(u)
 \tag{ML19}
\]
for the precise strong completion and its original top measure
constructed in R11 \textup{(IC10)--(IC13)}.  Its elements need
not admit a common actual lower barrier or an ambient classical
representative.

\begin{corollary}[Nonlinear compact-source obstruction in the full completion]
\label{cor:r12-completed-nonlinear-source-obstruction}
For every $u\in\mathscr C_S$,
\[
                      t_*\widehat\mu_u\in\mathscr L_{r,q}(Q).
 \tag{ML20}
\]
In particular, $t_*\widehat\mu_u$ satisfies every
\textup{(ML8)} and cannot have a nonzero compact carrier in $Q$.
If it is carried by such a compact set, then
$\widehat\mu_u=0$ on all original tests.
\end{corollary}

\noindent\textbf{Status: \texttt{NEW\_PROVED}} for every
nonlinear remaining degree $2\le q<r$, and as a different,
positive-current proof of the included degree-one case.  The
strong approximation and clipping assertions used below are
\texttt{INTERNAL\_INPUT} from R11.  The old R11 degree-one
Gauss-convexity obstruction is not recounted as a new result.

\begin{proof}
Choose an actual strongly Cauchy defining row
$a_j\in\mathcal E^1_{S,\partial}(V)$ for $u$.
By R11 \textup{(CC3)--(CC4)}, one can choose $L_j>0$ so that
\[
 w_j=\max(a_j,-L_j),\qquad I(a_j,w_j)<2^{-j}.
 \tag{ML21}
\]
Each $w_j$ is an actual bounded classical member of the same
first-energy class.  Its energy is no greater than that of
$a_j$, so all these energies are bounded.  The proved uniform
discrepancy modulus on bounded-energy sets, R11
\textup{(IC18)}, implies $w_j\to u$ strongly.  This step
uses that modulus, not a nonexistent triangle inequality for
$I$.  The strong completed current convergence of
\textup{(IC10)--(IC13)} gives
$\mu_{w_j}\to\widehat\mu_u$ on every original compact smooth
scalar test.  Theorem~\ref{thm:r12-positive-moment-lift}
applies to this bounded row and proves the assertion.
\end{proof}

\begin{theorem}[Dominated positive data are not universally solvable for any nonlinear degree]
\label{thm:r12-nonlinear-dominated-source-counterexample}
Fix integers $r\ge3$ and $2\le q<r$.  On the original split
$r$-torus take
\[
 \begin{gathered}
 Q=(-1,1)^r,\qquad V=t^{-1}(Q),\qquad t_i=\log|T_i|,\\
 R=[V],\qquad \beta=\sum_{i=1}^r d't_i\wedge d''t_i,
 \qquad S=\beta^{r-q}\wedge[V].
 \end{gathered}
 \tag{ML22}
\]
This is an allowed original background on a strict analytic
Laurent box, with an extension to a neighborhood of its compact
closure.  There exist a bounded actual finite-mass
$\psi\in\mathcal E^1_{S,\partial}(V)$ and a nonzero finite
positive Radon measure $\mu$ such that
\[
                       0\le\mu\le\mu_\psi,
 \qquad \operatorname{supp}(t_*\mu)\Subset Q,
 \tag{ML23}
\]
but
\[
       \widehat\mu_u\ne\mu\quad\hbox{for every }u\in\mathscr C_S.
 \tag{ML24}
\]
Nevertheless R11 \textup{(VF1)--(VF7)} supplies a unique
variational maximizer for this datum in $\mathscr C_S$.
Thus neither domination by one actual source nor finite
energy-dual size makes the unrestricted nonlinear source
equation true in the classical-approximation strong completion.
\end{theorem}

\noindent\textbf{Status: \texttt{NEW\_PROVED}} as a
counterexample in every nonlinear degree.  The universal existence
claim contradicted by \textup{(ML24)} is \texttt{FALSE}.
The variational existence statement itself is an
\texttt{INTERNAL\_INPUT}, not a new equation theorem.

\begin{proof}
The inverse image of a compact subset of $Q$ is a closed subset
of a compact original polyannulus, so $t:V\to Q$ is proper.
All geometric extension hypotheses follow by embedding the
closed box in a slightly larger compact polyannulus.
Let $\rho$ be the R11 strict-box barrier \textup{(TC3)}.
Choose a compact real box $H\Subset Q$ with nonempty interior,
put
\[
                    g(t)=\tfrac12\sum_i t_i^2,
\]
and choose $C>\sup_Qg+1$.  Taking $A>0$ sufficiently large,
the actual bounded classical psh function
\[
                         \psi=\max(g-C,A\rho)
 \tag{ML25}
\]
equals $g-C$ on a neighborhood of $t^{-1}(H)$.
Near the ordinary boundary it equals $A\rho$, because
$g-C<-1$ on all of $V$.  Thus it has compact negative
sublevels and extends as a bounded classical psh function
across the boundary collar.  The frozen original relative CLN
bound gives finite total curvature mass.  It follows that
$\psi$ belongs to every positive power-energy class, in
particular the asserted first-energy class.  This is exactly
the checked core patching mechanism of R11
\textup{(TC7)--(TC8)}, without its optional normalization.

Choose $0\le\chi\le1$ in $C_c^\infty(\operatorname{Int}H)$,
not identically zero, and set
\[
                           \mu=\chi(t)\mu_\psi.
 \tag{ML26}
\]
This is a finite positive measure with \textup{(ML23)}.
On a neighborhood of $t^{-1}(H)$, original BT locality gives
\[
                  \mu_\psi=(Dg)^q\wedge S
                             =\beta^r\wedge[V].
 \tag{ML27}
\]
The frozen all-original torus calibration is
\[
             \frac1{r!}\beta^r\wedge[V]=\sigma_*(dx),
 \tag{ML28}
\]
where $\sigma$ is the original joint Gauss section.  Consequently
$t_*\mu=r!\chi(x)\,dx$, which is nonzero and compactly
supported.  This calculation fixes the normalization explicitly;
all products in \textup{(ML22)} are unnormalized.

If $\widehat\mu_u=\mu$, Corollary
\ref{cor:r12-completed-nonlinear-source-obstruction} would force
$\mu=0$, a contradiction.  Hence \textup{(ML24)} holds.
The domination $\mu\le\mu_\psi$ is precisely the finite
dominated-data hypothesis of R11 \textup{(VF1)}.  Its variational
theorem therefore gives a unique maximizer, but not a source
solution.  No special-field extension, ordinary test restriction,
or analytic PS/PL quotient was used.
\end{proof}

\paragraph{Sharp boundary of this argument.}
The strict inequality $q<r$ is essential: the test
$\chi h\beta_0^{r-q-1}$ ceases to exist when $q=r$.
Top-degree positive currents are measures, whose separate
closedness is automatic, so the lift-admissible cone in that
degree would impose no restriction at all.  The proof also
needs full strict positivity of the real trace form, properness
of the specified moment on compact tests, and the displayed
factorization of $S$.  Dropping any of these requirements is
not justified by a formal change of notation.

\paragraph{Exact difference from the frozen archive.}
R10 \path{full_support_rigidity_of_compact_sources.tex}
proved a remaining-degree-one obstruction from full-support
Radon coordinate marginals and convex Gauss traces.  R11
\path{strong_completion_source_obstruction.tex} extended
that precise degree-one example to the unrestricted strong
completion using energy compactness of those convex traces.
The present proof does not repeat either convexity argument.
It supplies a closed linear positive-current constraint on
\emph{all} weak bounded-source limits, the quantitative scalar
collar inequalities \textup{(ML8)}, and completed counterexamples
in every $2\le q<r$.  R10's separate
\path{all_rank_original_cap_order.tex} concerns intrinsic
hinge-cap order reflection and does not contain this
moment-pushforward lift mechanism.

\paragraph{Remaining minimal gate.}
Lift admissibility is necessary, not proved sufficient for a
nonlinear source equation.  Even its satisfaction at every
available proper factorizable moment does not reconstruct a
global completed potential.  The nonlinear source image and
its variational normal residual therefore remain distinct.
No unrestricted Dirichlet theorem, no general analytic relation
quotient, and no low-real-degree signed constant-order transfer
is claimed here.

% END REVIEWED MODULE research_r12/nonlinear/tropical_curvature_lifts.tex


% BEGIN REVIEWED MODULE research_r12/nonlinear/moment_source_separation.tex
\section{A compact original test separates the nonlinear source image}
\label{sec:r12-moment-source-separation}

Retain the exact strict-box hypotheses \textup{(ED1)}, the
original arbitrary background in the preceding energy-obstacle
section, and, for this section only, the proper moment and
factorization hypotheses \textup{(ML12)} and
\eqref{eq:r12-lift-factorization}.  Thus $1\le q<r$ and
$S=t^*\beta_0^{r-q}\wedge R$, with $R$ an allowed arbitrary
original positive closed current of the printed degree.
The following separator is a concrete original smooth test,
not an abstract Hahn--Banach functional.

\begin{theorem}[A smooth scalar certificate for compact-source nonapproximation]
\label{thm:r12-smooth-source-separator}
Let $\mu$ be a nonzero positive Radon measure on $V$ with
$t_*\mu$ carried by a compact set $K\Subset Q$.  Suppose its
literal integration defines a continuous functional
$\ell_\mu\in\mathbb D_S^*$.
There is a real original compact smooth scalar
$a\in A_c^{0,0}(V)\subset\mathbb D_S$ such that
\[
 \int_Va\,d\mu=\mu(V)>0,
 \qquad
 \int_Va\,d\widehat\mu_v\le0
                         \quad\text{for every }v\in\mathscr C.
 \tag{SC1}
\]
In particular $\|a\|_{\rm ed}>0$ and
\[
 \boxed{\quad
 \operatorname{dist}_*(\ell_\mu,J(\mathscr C))
                       \ge\frac{\mu(V)}{\|a\|_{\rm ed}}>0.
 \quad}
 \tag{SC2}
\]
The datum has finite total mass because its tropical carrier
is compact and the specified moment is proper.
\end{theorem}

\noindent\textbf{Status: \texttt{NEW\_PROVED}} as an
explicit original-test separation and quantitative obstruction
to energy-dual source approximation.  It strengthens merely
knowing that the source equation has no solution.

\begin{proof}
Choose $\chi\in C_c^\infty(Q)$, $0\le\chi\le1$, equal
to one on an open neighborhood of $K$.  The compact set
$A_\chi=\operatorname{supp}(d\chi)$ is disjoint from $K$.
Choose $\kappa\in C_c^\infty(Q\setminus K)$, nonnegative
and at least one on $A_\chi$.  Let $C_\chi$ be the exact
finite constant of \textup{(ML7)}, and put
\[
                    b=\chi-C_\chi\kappa,\qquad a=b\circ t.
 \tag{SC3}
\]
The function $a$ is an actual original smooth scalar.  Its
support lies in the inverse image of the compact support
of $b$, hence is compact by properness.  The R11 exact
finite-mass common-collar difference theorem
\textup{(HC14)--(HC15)} places $a$ in $\mathbb D_S$.

On $K$ one has $\chi=1$ and $\kappa=0$.  Therefore
$\int a\,d\mu=\mu(V)$.  For every $v\in\mathscr C$,
the positive measure $\nu_v=t_*\widehat\mu_v$ is
lift-admissible by \textup{(ML20)}.  Its collar inequality
gives
\[
 \int b\,d\nu_v
    =\int\chi\,d\nu_v-C_\chi\int\kappa\,d\nu_v
    \le C_\chi\nu_v(A_\chi)
                -C_\chi\int\kappa\,d\nu_v\le0.
 \tag{SC4}
\]
Every displayed integral has compact support and finite
mass.  This proves \textup{(SC1)} on the original measures.

Now $\ell_\mu(a)=-\mu(V)\ne0$, so continuity and
well-definedness on the Banach quotient give
$\|a\|_{\rm ed}>0$.  For any $v\in\mathscr C$,
\[
 (\ell_\mu-J(v))(a)
       =-\mu(V)+\int a\,d\widehat\mu_v\le-\mu(V).
\]
Duality yields
$\|\ell_\mu-J(v)\|_*\ge\mu(V)/\|a\|_{\rm ed}$.
Taking the infimum proves \textup{(SC2)}.  The estimate is
in the signed operator dual norm $\delta_*$, not the
absolute-variation distance of R11 \textup{(VF11)}.
\end{proof}

\begin{corollary}[An explicit quantitative failure of obstacle orthogonality]
\label{cor:r12-quantitative-contact-counterexample}
Assume in addition that
$\ell_\mu\in\mathbb D_{S,\mathrm s}^*$ and put
$u=\Pi(\ell_\mu)$.  With the original compact smooth
$a$ from \textup{(SC3)}, the direction $h=-a$ satisfies
\[
 \begin{split}
 (\ell_\mu-J(u))(h)&\ge\mu(V)>0,\\
 \mathfrak d_u(h)&\ge\mu(V),\\
 \liminf_{t\downarrow0}\frac{\mathfrak Q(u+th)}t
                                      &\ge\mu(V).
 \end{split}
 \tag{SC5}
\]
Thus for all sufficiently small $t>0$ the relative obstacle
$u+th$ has strictly positive contact defect.  Its defect is
at least $\tfrac12t\mu(V)$ once $t$ is sufficiently small.
The right and left slopes of its obstacle energy at zero
are separated by at least $\mu(V)$.

For every $2\le q<r$, the explicit dominated data
\textup{(ML22)--(ML28)} satisfy these hypotheses.  Therefore
universal relative contact orthogonality fails in each
nonlinear degree, already along an actual original smooth
direction through a well-defined variational maximizer.
\end{corollary}

\noindent\textbf{Status: \texttt{NEW\_PROVED}} as an
all-degree contact counterexample with a quantitative original
test certificate.  The unqualified orthogonality assertion
being contradicted is \texttt{FALSE}.

\begin{proof}
The first line is \textup{(SC1)} applied to the particular
source $\widehat\mu_u$, with the sign of $a$ changed.
The other two lines and the slope gap follow from
\textup{(OC18)--(OC19)}.  The definition of a positive
lower limit gives the stated eventual lower bound.

For the explicit nonlinear examples,
$0\le\mu\le\mu_\psi$ with $\psi$ an actual first-energy
potential.  The frozen mixed inequality makes its literal
integration a continuous functional on $\mathbb D_S$.
The completed R11 tester estimate \textup{(EV1)} gives
\[
 |\ell_\mu(v_j)-\ell_\mu(v)|
        \le\int|v_j-v|\,d\mu_\psi\longrightarrow0
\]
along every strongly convergent cone row.  Thus it belongs
to the specified strong-continuous dual.  Variational
existence and uniqueness follow from R12 \textup{(NI6)--(NI10)}
or the frozen R11 dominated-data theorem.  All conclusions
then apply with this actual datum.
\end{proof}

\paragraph{Scope of the counterexample.}
The bad obstacle is the sum of a strong completed potential
and an original compact smooth function.  It is not asserted
to be an ordinary smooth obstacle, nor to have an ambient
classical representative.  The counterexample therefore
targets exactly the proposed complete strong-energy
obstacle mechanism.  The source nonexistence and separation
are nevertheless on every original scalar test.  The
special factorization in this section is sufficient to
refute a universal claim; the general arbitrary-$S$ obstacle
theorems of the preceding section do not require it.

% END REVIEWED MODULE research_r12/nonlinear/moment_source_separation.tex


% BEGIN REVIEWED MODULE research_r12/energy/null_fibers_and_nonliteral_duals.tex
\section{Null moment fibers and positive duals which are not literal data}
\label{sec:r12-null-fibers-nonliteral-duals}

The first part uses precisely the proper original moment and
background factorization (ML12), with $1\leq q<r$:
\[
 t:V\longrightarrow Q\subset\mathbf R^r,\qquad
 S=(t^*\beta_0)^{r-q}\wedge R,\qquad
 \beta_0=\sum_{i=1}^r d'x_i\wedge d''x_i.
 \tag{NF1}
\]
Here $R$ is an arbitrary allowed original symmetric middle-positive
separately closed $(n-r,n-r)$-current. Properness means that the
inverse image of every compact real set is compact in the ordinary
analytic open $V$. All capacities use every normalized original
bounded classical tester on $V$, without restricting its actual
moment refinements. The factorization is a printed hypothesis for
these geometric null-set conclusions, not a property of every $S$.

\begin{lemma}[A closed order-zero supercurrent has no sub-top atom]
\label{lem:r12-closed-current-no-atom}
Let $T$ be an order-zero real supercurrent of bidegree $(q,q)$
on $Q$, with $q<r$ and $d'T=0$. Every coefficient measure of
$T$ has zero mass at each point. In particular, if $T$ is
middle-positive, the positive measure
$T\wedge\beta_0^{r-q}$ has no atoms.
\end{lemma}

\noindent\textbf{Status: \texttt{DIRECT\_COROLLARY}} of
distributional closedness and the coefficient-measure definition;
the elementary proof is included. Separate $d''$-closedness and
positivity are not needed for the first assertion.

\begin{proof}
Fix $x\in Q$ and write the vector of coefficient measures near
$x$ as its atom at $x$ plus a vector of measures having no atom
there. Let $k=r-q\geq1$, let $\eta$ be any constant form of
degree $(k-1,k)$, and choose $\varphi\in C_c^\infty(\mathbf R^r)$.
For small $\epsilon>0$ the compact test
\[
               \varphi_\epsilon(y)\eta
           =\epsilon\varphi((y-x)/\epsilon)\eta
\]
is supported in $Q$. Closedness gives
$\langle T,d'\varphi_\epsilon\wedge\eta\rangle=0$.
The coefficients of $d'\varphi_\epsilon$ are uniformly bounded.
Thus the contribution from the atom-free part tends to zero,
because its variations on these shrinking balls tend to zero.
The contribution of the atom is its pairing with
$d'\varphi(0)\wedge\eta$. Every covector can occur as
$d'\varphi(0)$, and forms $\xi'\wedge\eta$ of these degrees
span the entire complementary coefficient space $(k,k)$.
The atomic tensor therefore pairs to zero with every coefficient
test and is zero. This works at every point.
The trace is a fixed finite linear combination of coefficient
measures, so it has no atom either.
\end{proof}

\begin{theorem}[Full original outer capacity vanishes on each proper moment fiber]
\label{thm:r12-null-proper-moment-fibers}
Under (NF1), let $H\Subset Q$ be compact. For sufficiently
small $\epsilon>0$ put
\[
 \omega_H(\epsilon)=
   \sup_{x\in H}\ \sup_{\substack{v\text{ original bounded classical psh on }V\\
                                  0\leq v\leq1}}
     \bigl((Dv)^q\wedge S\bigr)
                           \bigl(t^{-1}(B(x,\epsilon))\bigr).
 \tag{NF2}
\]
This is finite and
\[
                  \lim_{\epsilon\downarrow0}\omega_H(\epsilon)=0.
 \tag{NF3}
\]
Consequently, for every $x\in Q$,
\[
                    C_S^*(t^{-1}\{x\};V)=0.
 \tag{NF4}
\]
The same holds for the inverse image of any countable subset
of $Q$. The conclusions concern the full original capacity
and ordinary open neighborhoods, not a toric-only tester capacity.
No numerical power rate is asserted in (NF3).
\end{theorem}

\noindent\textbf{Status: \texttt{NEW\_PROVED}} as a
uniform full-tester null-fiber statement in every sub-top
degree. It uses the new positive moment lifts and the frozen
arbitrary-background relative CLN estimate.

\begin{proof}
For every tester in (NF2), its original bounded lift is
\[
             T_v=t_*\bigl((Dv)^q\wedge R\bigr),\qquad
             \nu_v=T_v\wedge\beta_0^{r-q}
                       =t_*\bigl((Dv)^q\wedge S\bigr).
 \tag{NF5}
\]
These are the actual currents and exact unnormalized trace
identity of (ML15)--(ML17). On each compact $H'\Subset Q$,
properness makes $t^{-1}(H')$ compact. The original relative
CLN estimate on a slightly larger compact set, with
$\|v\|_\infty\leq1$, bounds $\nu_v(H')$ uniformly over
the full tester class. This proves finiteness in (NF2).
The trace coefficient lemma (ML4)--(ML5) makes the family
of lifts locally sequentially compact. Every limiting lift
is again positive, symmetric and separately closed.

If (NF3) failed, there would be $\delta>0$,
$\epsilon_j\downarrow0$, $x_j\in H$, and normalized
testers $v_j$ with
$\nu_{v_j}(B(x_j,\epsilon_j))\geq\delta$.
After subsequences, $x_j\to x\in H$ and $T_{v_j}\to T$
locally on all real tests. Its trace $\nu$ is the local vague
limit of $\nu_{v_j}$. For any sufficiently small fixed
$\rho>0$, eventually
$B(x_j,\epsilon_j)\subset\overline B(x,\rho)\Subset Q$.
Portmanteau on this compact set gives
\[
      \delta\leq\limsup_j\nu_{v_j}(\overline B(x,\rho))
                              \leq\nu(\overline B(x,\rho)).
\]
Letting $\rho\downarrow0$ gives $\nu\{x\}\geq\delta$,
contradicting Lemma~\ref{lem:r12-closed-current-no-atom}.
This proves the uniform assertion.

The sets $t^{-1}(B(x,\epsilon))$ are ordinary open
neighborhoods of the compact fiber $t^{-1}\{x\}$.
Their raw capacities are exactly the inner supremum in
(NF2). Taking their decreasing numerical bound in the
definition of outer capacity proves (NF4). Countable
subadditivity of the frozen outer capacity gives the last
assertion. Neither a Radon disintegration over residue
fields nor a change of the original field has occurred.
\end{proof}

\subsection{Canonical convex traces of the complete energy class}

For the remainder take the allowed original example
\[
 \begin{gathered}
 X=(\mathbf G_{m,K}^{r})^{\rm an},\quad
 Q=(-1,1)^r,\quad V=t^{-1}(Q),\quad t_i=\log|T_i|,\\
 1\leq q<r,\quad S=\beta^{r-q}\wedge[V],\quad
 \beta=\sum_i d't_i\wedge d''t_i,\quad \sigma:Q\longrightarrow V,\\
 \text{where }\sigma\text{ is the original joint Gauss section}.
 \end{gathered}
 \tag{NF6}
\]
These data satisfy (ED1) and (NF1). The frozen R6
\path{lem:r6-joint-gauss-convexity}, used in R11
\path{comparison/power_class_source_obstruction.tex}
before the choice of its background, says that a bounded
original classical psh potential has a finite convex
restriction to $\sigma(Q)$. That assertion is independent
of the remaining degree $q$. We use it as an exact
\texttt{INTERNAL\_INPUT}; no nonlinear identity between
a Gauss trace Hessian and the original top measure is used.

\begin{proposition}[All-degree completed Gauss traces and point evaluation]
\label{prop:r12-completed-gauss-evaluation}
Every $u\in\mathscr C$ for (NF6) has a canonical finite
nonpositive continuous convex function $g_u$ on $Q$.
It is the locally uniform limit of the Gauss traces of any
actual bounded strong defining row. It agrees with the
chosen relative Borel representative $u\circ\sigma$ only
Lebesgue-almost everywhere, not necessarily at every point.

The assignment is additive and positively homogeneous on
$\mathscr C$, and extends linearly to $\mathbb D_S$.
For each $x\in Q$, the formula
\[
                         \ell_x(u)=-g_u(x)
 \tag{NF7}
\]
defines an order-reversing functional in
$\mathbb D_{S,\mathrm s}^*$. For every original compact
smooth scalar $\phi$ on $V$,
\[
                              \ell_x(\phi)=-\phi(\sigma(x)).
 \tag{NF8}
\]
Nevertheless this functional is not literal integration
against a Radon measure on all relative Borel representatives
in $\mathbb D_S$.
\end{proposition}

\noindent\textbf{Status: \texttt{NEW\_PROVED}} for the
full-degree strong class and the nonliteral positive dual
phenomenon. The old bounded Gauss convexity is not recounted
as a new result.

\begin{proof}
Fix real compact sets $H\Subset\operatorname{Int}H'\Subset Q$.
Patch the quadratic $|t|^2/2-C$ to a large multiple of the
strict-box barrier, as in (ML25), to obtain an actual
bounded first-energy potential $\psi$ equal to that quadratic
on $t^{-1}(H')$. All factors are unnormalized, so locality
and the original torus calibration give
\[
           \mathbf1_{t^{-1}(H')}\mu_\psi
                   =\mathbf1_{t^{-1}(H')}\beta^r\wedge[V]
                   =r!\,\sigma_*(\mathbf1_{H'}\,dx).
 \tag{NF9}
\]
For any actual bounded first-energy potential $a$,
the mixed energy inequality implies
\[
              \int_{H'}(-a\circ\sigma)\,dx
                          \leq\frac{r(\psi)^q}{r!}\,r(a).
 \tag{NF10}
\]
For two such potentials, the completed tester estimate
(EV1) controls their trace difference in $L^1(H')$ by
a quantity tending to zero under strong convergence on
an energy-bounded set.

We recall explicitly the elementary convex compactness step.
If $g\leq0$ is finite convex on a ball $B(x,2\rho)$,
then $-g$ is nonnegative and concave. For
$y=(x+z)/2\in B(x,\rho)$, concavity gives
$-g(y)\geq(-g(x))/2$. Hence
$-g(x)\leq2\operatorname{vol}(B(x,\rho))^{-1}
\int_{B(x,\rho)}(-g)$. An $L^1$ bound on a larger
compact set therefore bounds $|g|$ on an interior compact
set. Secant-slope estimates for convex functions then give
a uniform Lipschitz bound on a still smaller compact set.
Thus an $L^1_{\rm loc}$-Cauchy row of these convex functions
converges locally uniformly to one finite continuous convex
function: compactness gives subsequential limits, and its
$L^1$ limit makes them identical.

Choose an actual bounded strongly defining row for $u$;
such a row is supplied by the R11 clipping diagonal.
The traces are convex by the frozen input, are locally
$L^1$-bounded by (NF10), and are $L^1$-Cauchy by (EV1).
The preceding argument defines $g_u$ and proves independence
of the row. Their $L^1(\mu_\psi)$ convergence to the
relative Borel class, using (NF9), proves the almost-everywhere
identification in the statement. Exact cone addition and
homogeneity, followed by local $L^1$ limits, give the stated
linearity. Equal Borel difference presentations give equal
traces almost everywhere, hence everywhere by continuity.

The point bound from (NF10) yields, for each fixed $x$,
\[
                     0\leq-g_u(x)\leq C_x r(u),
 \tag{NF11}
\]
with a finite constant independent of $u$. Extending by
differences and infimizing proves boundedness of $\ell_x$
on $\mathbb D_S$. If $h=a-b\succeq0$, its relative
exceptional set is $\mu_\psi$-null. Thus
$g_a-g_b\geq0$ almost everywhere on each core, and continuity
makes the inequality hold at every point. This proves order
reversal. For a strongly convergent cone row, (EV1) and
(NF9) again give local $L^1$ convergence of its convex
traces; the uniform convex bounds give local uniform
convergence. Consequently $\ell_x$ is strongly continuous.

The all-degree original test decomposition (HC14)--(HC15)
writes $\phi=a-b$ as an exact difference of actual bounded
finite-mass common-collar potentials. Their actual Gauss
traces give (NF8), including every refined non-toric original
test. In particular any Radon measure giving the same
original test restriction must be $\delta_{\sigma(x)}$.
But (NF4) makes $\{\sigma(x)\}$ outer-capacity null.
Changing a Borel representative at this point does not
change its class, whereas literal integration against this
Dirac measure changes. For example, change the value of
the strictly negative patched $\psi$ there to zero.
Thus the functional cannot be literal Radon integration
on all representatives of $\mathbb D_S$. This does not
conflict with its canonical-trace definition.
\end{proof}

\begin{theorem}[Literal positive data are not closed in the energy dual norm]
\label{thm:r12-literal-dual-data-not-closed}
In (NF6), for every $x\in Q$ there are finite positive Radon
measures $\mu_\epsilon$, each dominated by a finite constant
times one actual bounded first-energy source, such that
their literal integration functionals satisfy
\[
               \|\ell_{\mu_\epsilon}-\ell_x\|_*
                                      \leq C_x'\epsilon\longrightarrow0.
 \tag{NF12}
\]
All these functionals and their limit are order reversing
and strongly continuous. The limit is not a literal Radon
datum on the relative Borel quotient. The same measures
converge on every original compact smooth scalar test to
$\delta_{\sigma(x)}$, which charges a full outer-capacity
null set.

This does not contradict closedness of $J(\mathscr C)$:
the approximating data here are not source functionals.
Their nonzero compact tropical carriers already exclude
that possibility by (ML20).
\end{theorem}

\noindent\textbf{Status: \texttt{NEW\_PROVED}} as a
counterexample to the blanket identification of the positive
strong energy dual with literal finite-energy Radon data.
That identification, and norm-closedness of the literal
positive datum class without a uniform diffuseness gate,
are \texttt{FALSE}.

\begin{proof}
Choose a nonnegative smooth probability density
$\theta_\epsilon$ supported in $B(x,\epsilon)\Subset Q$
and set $\mu_\epsilon=\sigma_*(\theta_\epsilon\,dx)$.
On one fixed larger compact core, (NF9) gives
$\mu_\epsilon\leq(\|\theta_\epsilon\|_\infty/r!)\mu_\psi$.
The mixed energy inequality and (EV1) therefore make
each a literal member of $\mathbb D_{S,\mathrm s}^*$.
For $u\in\mathscr C$ with $r(u)\leq1$, the proof of
(NF11) also gives a common Lipschitz constant for $g_u$
on this core. The almost-everywhere identification of the
representative gives
\[
 \left|\ell_{\mu_\epsilon}(u)-\ell_x(u)\right|
       =\left|\int(g_u(x)-g_u(y))\theta_\epsilon(y)\,dy\right|
       \leq C_x'\epsilon.
\]
The exact cone formula for the operator dual norm (ED10)
proves (NF12). Evaluation on continuous original compact
tests proves the ordinary measure convergence. The previous
proposition proves that its limit is nonliteral on the
relative quotient. Finally the carrier of each tropical
pushforward is a nonempty compact subset of $Q$, so the
source obstruction applies to every approximating measure.
All constructions remain on the original field and norm.
\end{proof}

\paragraph{The exact surviving distinctions.}
There are now three separate notions: a bounded Banach-dual
functional, a strong-continuous order-reversing functional,
and a literal positive Radon datum whose integrals ignore
the relative null ideal. The first need not be strongly
continuous by (CG11)--(CG14), and the second need not be
literal by (NF7)--(NF12). The capacity-modulus condition
(CD1)--(CD2) supplies the third with all required continuity.
The nonlinear source image is a further, closed proper subset,
with the normal and contact obstructions already proved.
None of these distinctions closes OPEN-GQ or identifies
relative null sets with all pluripolar or negligible sets.

% END REVIEWED MODULE research_r12/energy/null_fibers_and_nonliteral_duals.tex


% BEGIN REVIEWED MODULE research_r12/energy/perron_variational_energy_gap.tex
\section{Least-energy source majorants and the Perron--variational gap}
\label{sec:r12-perron-variational-gap}

Retain the full hypotheses (ED1), the original arbitrary background
$S$, and the unchanged strong cone $\mathscr C$. Put $d=q+1$.
For dual functionals write
\[
       \ell\preceq_{\rm psh}\lambda
       \quad\Longleftrightarrow\quad
       \ell(v)\leq\lambda(v)\quad\text{for every }v\in\mathscr C.
 \tag{PV1}
\]
This is the order tested by negative energy potentials. It is
neither ordinary measure order nor middle positivity of currents.

\begin{lemma}[A one-sided Bregman bound]
\label{lem:r12-one-sided-Bregman}
For $a,b\in\mathscr C$,
\[
 \begin{split}
 e(a)-e(b)-dJ(b)(a-b)
  &=\sum_{j=0}^{q-1}(q-j)
       \int_V\Gamma_{\widehat N(a^{[j]},b^{[q-1-j]})}(a-b),\\
 I(a,b)&\leq e(a)-e(b)-dJ(b)(a-b)\leq qI(a,b).
 \end{split}
 \tag{PV2}
\]
Every displayed energy is finite and nonnegative.
\end{lemma}

\noindent\textbf{Status: \texttt{DIRECT\_COROLLARY}} of
the frozen mixed variation and the coefficient computation (JD9).
The bounded Bregman calculation was already present in R3.

\begin{proof}
Apply polynomial Taylor expansion to $e((1-t)b+ta)$ and
retain its first derivative rather than canceling it by
averaging. The same beta integral as in (JD9) gives the
coefficients $q-j$. Their unweighted sum is $I(a,b)$ by
the frozen discrepancy identity, proving both inequalities.
The completed mixed energy calculus justifies all global
integrals before this algebraic computation is made.
\end{proof}

\begin{theorem}[The unique least-energy curvature majorant]
\label{thm:r12-least-energy-curvature-majorant}
Let $\ell\in\mathbb D_{S,\mathrm s}^*$ be nonnegative on
$\mathscr C$, and put $u=\Pi(\ell)$. Then
\[
       \ell\preceq_{\rm psh}J(u),\qquad
       \|J(u)\|_*=\|\ell\|_*.
 \tag{PV3}
\]
Among all $v\in\mathscr C$ satisfying
$\ell\preceq_{\rm psh}J(v)$, the potential $u$ is the
unique one of least energy. Every such $v$ satisfies
\[
               0\leq I(u,v)\leq q\bigl(e(v)-e(u)\bigr).
 \tag{PV4}
\]
Consequently any row of these majorants with
$e(v_j)\to e(u)$ converges strongly to $u$ on the full
original current calculus.

The map $\ell\mapsto J(\Pi(\ell))$ is an idempotent,
positively homogeneous, dual-norm-preserving retraction on
this positive dual cone. Its norm continuity is (NI11).
The phrase \emph{least energy} does not assert a least element
for the order (PV1), and norm preservation does not mean
preservation of the total mass of a measure.
\end{theorem}

\noindent\textbf{Status: \texttt{NEW\_PROVED}} for the
all-degree majorant characterization and its quantitative
energy gap. The normal decomposition is the preceding
proved R12 input.

\begin{proof}
The normal residual in (NI7) is nonpositive on $\mathscr C$,
which proves the order assertion in (PV3). Positivity of
$\ell$ on this cone gives $p(\ell)=\|\ell\|_*$, so (NI9)
and (ED10) prove its norm equality.

For a competitor $v$, the majorant condition tested at $u$
gives $B(u,v^{[q]})=J(v)(u)\geq\ell(u)=e(u)$.
Use (PV2) with $(a,b)=(u,v)$ to obtain
\[
 \begin{split}
 I(u,v)&\leq e(u)-e(v)-d\bigl(B(u,v^{[q]})-e(v)\bigr)\\
       &=e(u)+q e(v)-d B(u,v^{[q]})
                         \leq q(e(v)-e(u)).
 \end{split}
\]
This proves (PV4), minimality, and uniqueness by relative
discrepancy separation. It also proves the approximation
assertion. The retraction, homogeneity and continuity follow
from the corresponding exact properties of (NI7)--(NI11).
No positivity of the signed residual measure is inferred
from its potential-test order.
\end{proof}

For a literal positive Radon datum $\mu$, (PV3) says
\[
      \int_V(-v)\,d\mu\leq\int_V(-v)\,d\widehat\mu_u
               \quad(v\in\mathscr C),\qquad
      \int_V(-u)\,d\mu=e(u).
 \tag{PV5}
\]
One may call this a nonlinear energy sweeping of the datum,
provided (PV1), the specified source class, and the absence
of a total-mass claim remain part of the definition.

\begin{lemma}[The variational source cannot have a positive excess]
\label{lem:r12-variational-positive-excess}
Let $\mu$ be a positive original Radon measure whose literal
functional belongs to $\mathbb D_{S,\mathrm s}^*$, and
put $u=\Pi(\ell_\mu)$. If $\widehat\mu_u\geq\mu$
as original measures, then $\widehat\mu_u=\mu$.
\end{lemma}

\noindent\textbf{Status: \texttt{DIRECT\_COROLLARY}} of
radial variational orthogonality and the frozen completed
weak maximum inequality.

\begin{proof}
For every $a\in\mathscr C$, apply the completed weak
maximum inequality to $a$ and the zero potential, assigning
the equality set to $a$. Since $\max(a,0)=0$, it gives
\[
                   \mathbf1_{\{a=0\}}\widehat\mu_a=0.
 \tag{PV6}
\]
All involved measures also ignore the relative pole set.
Now let $\tau=\widehat\mu_u-\mu\geq0$. It is dominated
by $\widehat\mu_u$, so it ignores both sets just mentioned.
Radial stationarity gives the finite equality
\[
                       \int_V(-u)\,d\tau
                                  =e(u)-\ell_\mu(u)=0.
\]
The integrand is strictly positive and finite on the
remaining full carrier $\{-\infty<u<0\}$. Its increasing
level exhaustion therefore makes $\tau$ vanish there too.
Thus $\tau=0$ as a Radon measure. No subtraction of
infinite unweighted total masses has occurred.
\end{proof}

\begin{theorem}[An exact scalar gap between Perron and variational potentials]
\label{thm:r12-perron-variational-gap}
Let $\mu$ be a positive original Radon measure with a
literal strong-continuous energy functional. Suppose the
completed subsolution family
\[
             \mathscr S_\mu
               =\{v\in\mathscr C:\widehat\mu_v\geq\mu\}
 \tag{PV7}
\]
is nonempty. Let $H_\mu$ be its greatest relative member,
constructed by the frozen R11 full Perron theorem (SE1)--(SE8),
and let $u_\mu=\Pi(\ell_\mu)$. Define
\[
                  \Delta_S(\mu)=e(H_\mu)-e(u_\mu).
 \tag{PV8}
\]
Then
\[
              \Delta_S(\mu)\geq0,\qquad
              I(H_\mu,u_\mu)\leq q\Delta_S(\mu).
 \tag{PV9}
\]
Moreover the following conditions are equivalent:
\begin{enumerate}
\item $\Delta_S(\mu)=0$;
\item $H_\mu=u_\mu$ as relative completed classes;
\item $\widehat\mu_{u_\mu}=\mu$ on every original test;
\item there exists $v\in\mathscr C$ with $\widehat\mu_v=\mu$;
\item there exists a row $v_j\in\mathscr S_\mu$ with
      $e(v_j)\to e(u_\mu)$.
\end{enumerate}
If they hold, this source potential is unique. The family
is nonempty, in particular, whenever $\mu\leq A\mu_\psi$
for an actual first-energy $\psi$ and finite $A\geq0$.

Thus the general subsolution Perron construction and the
general variational construction are linked by a single
nonnegative scalar gap. Its universal vanishing is not
claimed: the compact-source examples (ML22)--(ML28) have
$\Delta_S(\mu)>0$ in every $2\leq q<r$.
\end{theorem}

\noindent\textbf{Status: \texttt{NEW\_PROVED}} for this
all-degree criterion and quantitative comparison. The two
constructions themselves are explicitly frozen or previously
proved inputs. Unqualified vanishing of the gap is
\texttt{FALSE}, not a theorem hidden in the word Perron.

\begin{proof}
The R11 theorem applies to the nonempty family in (PV7),
keeps its measure lower bound, and makes its relative
envelope an actual member of $\mathscr C$. Thus $H_\mu$
is the greatest member of that family. Every measure
majorant is also a majorant for the potential-test order
(PV1). Applying (PV4) proves (PV9).

If the gap is zero, (PV9) gives equality of the two
relative potentials. The common potential is a subsolution,
so Lemma~\ref{lem:r12-variational-positive-excess} makes
its source exactly $\mu$. Conversely, a source potential
$v$ has zero normal residual and hence is the unique
variational potential $u_\mu$. It belongs to (PV7), so
$H_\mu\succeq u_\mu$. Energy monotonicity gives
$e(H_\mu)\leq e(u_\mu)$, and (PV9) gives equality.
This proves the first four equivalences and uniqueness.

Every $v\in\mathscr S_\mu$ is below $H_\mu$, so
$e(v)\geq e(H_\mu)$. A row in condition (5) therefore
forces the gap to be zero. If the gap is zero, the
constant row $v_j=u_\mu$ supplies condition (5).
Alternatively (PV4) directly shows that any such row
converges strongly to $u_\mu$ and its source inequalities
pass to all original compact tests.

For the printed dominated datum, $A^{1/q}\psi$ is a
subsolution (the zero case is immediate), and its literal
datum functional satisfies the strong-dual hypothesis by
mixed Holder and (EV1). The nonlinear compact-source
examples satisfy these same assumptions but have no source
in $\mathscr C$. The equivalence therefore makes their
gap strictly positive. None of these arguments needs a
common actual lower barrier for the whole subsolution family.
\end{proof}

\begin{proposition}[A quantitative positive-test lower bound for the gap]
\label{prop:r12-positive-test-gap-bound}
Under the preceding theorem put $u=u_\mu$, $H=H_\mu$,
$R=r(H)$, and define
\[
 \varepsilon_+(\mu,u)=
   \sup_{\substack{0\leq\phi\in A_c^{0,0}(V)\\
                   \|\phi\|_{\rm ed}\leq1}}
                   \left(\int_V\phi\,d(\mu-\widehat\mu_u)\right)_+.
 \tag{PV10}
\]
This is finite. If $R>0$, then
\[
       \Delta_S(\mu)\geq
          \frac{\varepsilon_+(\mu,u)^2}{q^2R^{q-1}}.
 \tag{PV11}
\]
If $R=0$, the datum and all quantities are zero. Also
$\varepsilon_+(\mu,u)=0$ if and only if the source equation
holds. Every positive test in (PV10) is from the full
original sheaf, not a selected moment family.
\end{proposition}

\noindent\textbf{Status: \texttt{DIRECT\_COROLLARY}} of
the new gap theorem and the proved source-dual estimate.

\begin{proof}
All original compact smooth scalars lie in $\mathbb D_S$
by (HC14)--(HC15). Finiteness follows by the dual norm of
$\ell_\mu-J(u)$. Since $\widehat\mu_H\geq\mu$, each
nonnegative test in (PV10) gives
\[
 \int\phi\,d(\mu-\widehat\mu_u)
       \leq\int\phi\,d(\widehat\mu_H-\widehat\mu_u)
       =(J(u)-J(H))(\phi).
\]
Consequently $\varepsilon_+\leq\|J(H)-J(u)\|_*$.
The energy minimum gives $r(u)\leq R$. Equations (NI2)
and (PV9) now give
$\varepsilon_+\leq qR^{(q-1)/2}\Delta_S(\mu)^{1/2}$,
which is (PV11). When $R=0$, the subsolution $H=0$
already forces $\mu=0$.

Vanishing in (PV10), together with original positive
scalar test detection, says $\widehat\mu_u\geq\mu$.
Lemma~\ref{lem:r12-variational-positive-excess} upgrades
this to equality. The reverse implication is immediate.
This is a full original-measure statement, not an
ordinary-to-delta-test promotion.
\end{proof}


% END REVIEWED MODULE research_r12/energy/perron_variational_energy_gap.tex


% BEGIN REVIEWED MODULE research_r12/nonlinear/single_probe_normal_residual.tex
\section{A geometric probe detects every literal variational source defect}
\label{sec:r12-single-probe-normal-residual}

Retain the strict original analytic box and background-extension
hypotheses \textup{(ED1)}.  The background $S$ is arbitrary:
it is the original symmetric middle-positive separately closed
$(p,p)$-current, with $q=n-p\ge1$.  Set $d=q+1$ and keep
the unchanged strong classical-approximation cone $\mathscr C$.
Let $\psi_0$ be the single geometry-dependent bounded actual
probe of \textup{(UP2)--(UP5)}.  Its construction precedes
the choice of $S$, and it is a finite-mass, strictly negative,
continuous zero-boundary actual potential for every allowed $S$.

The result below concerns compact original smooth directions.
It does not claim that $\mathscr C$ has nonempty Banach
interior, or that a pointwise lower barrier permits subtraction
of that barrier inside the psh cone.

\begin{lemma}[A simultaneous two-sided original test probe]
\label{lem:r12-two-sided-geometric-probe}
For every real original compact smooth scalar $\phi$ on $V$
there is $\epsilon_\phi>0$, depending only on the geometric
probe and the chosen actual test, such that
\[
             \psi_0+t\phi\in\mathscr T_0(S,V)
                  \qquad(|t|\le\epsilon_\phi)
 \tag{NP1}
\]
for every allowed background $S$ and every remaining degree.
The actual functions in \textup{(NP1)} are bounded continuous
classical psh, strictly negative, of ordinary zero boundary
value and finite total top mass.  They agree with $\psi_0$
outside one compact set.  Every actual moment refinement of
$\phi$ is permitted.

On one fixed finite original moment presentation, compact
support and smooth partition, one may choose a constant
$c_{\mathfrak a}>0$ such that
\[
       \epsilon_\phi\ge
           \frac{c_{\mathfrak a}}{\max(1,\|\phi\|_{2,\mathfrak a})}.
 \tag{NP2}
\]
Here $\|\cdot\|_{2,\mathfrak a}$ is the sum of the
usual real $C^2$ norms of the smooth representatives in
that fixed presentation.  No presentation-independent
global norm on all original tests is asserted.
\end{lemma}

\noindent\textbf{Status: \texttt{NEW\_PROVED}} for the
simultaneous probe.  The finite smooth perturbation argument
is the frozen explicit classical compensator construction in
R11 \textup{(AE5)} and the proof of \textup{(HC15)}, rather
than an inference from positivity against one current.

\begin{proof}
Take a finite original smooth partition of the support and
write $\phi=\sum_{l\in F}\phi_l$, with each $\phi_l$
supported in the interior of a compact moment core used
in \textup{(UP1)}.  Include the finitely many moment
functions of the partition itself in those presentations.
The countable exact moment replacement and its closure under
finite refinements allow these cores and compensators to
be selected from the family defining $\psi_0$.
If an index repeats, first combine its finitely many
localized test functions.  The construction then gives
\[
 \psi_0=\sum_{l\in F}c_l\psi_l+r_F,
 \qquad \psi_l=a_l-2\text{ near }\operatorname{supp}\phi_l,
 \tag{NP3}
\]
where $r_F$ is the uniformly convergent sum of the remaining
nonpositive actual psh terms.  It is itself actual bounded
classical psh and belongs to every printed finite power class;
the zero residual is harmless.

On the relevant smaller original moment stage, the explicit
convex-flat numerator and denominator summands in $a_l$
have a strictly positive Hessian in the finite selected
logarithmic directions.  The real Hessian coefficients of
$\phi_l$ are bounded on their fixed compact support.
Consequently there is $\delta_l>0$ such that
\[
                     a_l+\tau\phi_l
             \text{ is actual classical psh near that support}
                  \qquad(|\tau|\le\delta_l).
 \tag{NP4}
\]
This is exactly the frozen explicit smooth perturbation
argument: the selected convex-flat terms dominate both
signs of the Hessian of the smooth real representative,
and the remaining compensator terms are already classical
psh.  Convex composition is used in those actual finite
unit charts, followed by ordinary-open gluing.  No claim
that a current inequality for one background implies
classical psh is made.

Since $\phi_l$ vanishes off its compact support,
$\psi_l+\tau\phi_l$ is therefore classical on all of
$V$ for these small $\tau$.  Decrease $\delta_l$ if
necessary so that it is strictly negative: the continuous
$-\psi_l$ has a positive minimum on that support.
Choose $\epsilon_\phi>0$ with
$\epsilon_\phi/c_l\le\delta_l$ for every $l\in F$.
Then
\[
     \psi_0+t\phi
       =r_F+\sum_{l\in F}c_l
                    \left(\psi_l+\frac{t}{c_l}\phi_l\right)
 \tag{NP5}
\]
is an actual bounded continuous negative classical psh
function.  The normalization required by $\mathscr T_0$
also survives after decreasing $\epsilon_\phi$.
Indeed \textup{(UP5)} gives
\[
 T:=\sum_l c_lA_l=1-\sum_lc_l<1,
       \qquad\psi_0\ge T\rho\ge-T>-1.
\]
Choose in addition
$\epsilon_\phi\|\phi\|_\infty\le(1-T)/2$
(with no restriction if $\phi=0$).  Then
$-1<\psi_0+t\phi<0$.  It equals $\psi_0$ off a fixed compact set,
so its ordinary boundary value and compact strict negative
sublevels are unchanged.  Its top source equals that of
$\psi_0$ off the compact perturbation support by ordinary
bounded locality.  On a slightly larger compact set,
relative CLN gives finite mass.  Since $\psi_0$ has
finite total mass, so does the perturbed function.  All
finite positive power energies follow from boundedness.

For a fixed finite presentation, the Hessian bounds in
\textup{(NP4)}, the function-value bounds used for
strict negativity, and the finitely many $c_l$ are
uniform when $\|\phi\|_{2,\mathfrak a}\le1$.
The lower normalization uses the same fixed positive
margin $1-T$ and is uniform on that unit ball as well.
Scaling this fixed estimate proves \textup{(NP2)}.
The choices depend on the actual test presentation but
not on $S$ or its remaining degree.  This proves all
claims, including those for every further actual refinement.
\end{proof}

\begin{definition}[The normal probe defect]
Suppose $\ell\in\mathbb D_S^*$ and $u\in\mathscr C$
satisfy the variational normal relations
\[
 n=\ell-J(u),\qquad n(a)\le0\ (a\in\mathscr C),
                       \qquad n(u)=0.
 \tag{NP6}
\]
For example these hold for the uniquely constructed
$u=\Pi(\ell)$ whenever $\ell\in\mathbb D_{S,\mathrm s}^*$.
Define the finite scalar
\[
        \Xi_{\psi_0}(\ell,u)
                   =J(u)(\psi_0)-\ell(\psi_0)=-n(\psi_0)\ge0.
 \tag{NP7}
\]
The definition requires an existing normal pair.  It
does not assert existence of a variational maximizer for
every continuous functional on $\mathbb D_S$.
\end{definition}

\begin{theorem}[One scalar controls every original normal test]
\label{thm:r12-one-probe-normal-detector}
For every pair \textup{(NP6)} and every real original
compact smooth scalar $\phi$,
\[
            |n(\phi)|
                 \le\epsilon_\phi^{-1}\Xi_{\psi_0}(\ell,u).
 \tag{NP8}
\]
On a fixed finite presentation the right side may be
replaced by
$c_{\mathfrak a}^{-1}\max(1,\|\phi\|_{2,\mathfrak a})
\Xi_{\psi_0}(\ell,u)$; by homogeneity one obtains the
corresponding linear $C^2$-seminorm estimate for nonzero
tests.  This is a bound on original distributional test
pairings, not a total-variation norm assertion.

In particular $\Xi_{\psi_0}=0$ forces $n$ to vanish on
all original compact smooth scalar tests and on their
order-norm closed core $\mathbb L_{S,0}$ whenever $n$
extends continuously to that lattice.  It need not force
$n=0$ on $\mathbb D_S$ for a nonliteral energy dual.

Now let $\mu$ be a positive original Radon measure whose
\emph{literal} integration functional belongs to
$\mathbb D_{S,\mathrm s}^*$, and put $u=\Pi(\ell_\mu)$.
Then the following are equivalent:
\[
 \begin{split}
 \Xi_{\psi_0}(\mu)
   &:=\int_V(-\psi_0)\,d\widehat\mu_u
                      -\int_V(-\psi_0)\,d\mu=0;\\
 \widehat\mu_u&=\mu
        \quad\text{on all original compact smooth scalar tests};\\
 \text{there exists }v&\in\mathscr C
                         \text{ with }\widehat\mu_v=\mu.
 \end{split}
 \tag{NP9}
\]
If they hold, $v=u$ is unique.  For all these literal data,
\[
 \left|\int_V\phi\,d(\mu-\widehat\mu_u)\right|
                        \le\epsilon_\phi^{-1}\Xi_{\psi_0}(\mu).
 \tag{NP10}
\]
The same single geometry-only $\psi_0$ works for every
allowed arbitrary $S$, every remaining degree, every
datum in the printed class, and all actual refinements.
\end{theorem}

\noindent\textbf{Status: \texttt{NEW\_PROVED}} as the
single-probe reduction and its all-original quantitative
bound.  This is distinct from the Jensen-cluster probe
theorem: the defect here is a variational normal, not
a positive Jensen current.

\begin{proof}
Apply the polar inequality in \textup{(NP6)} to the
two actual cone members
$\psi_0\pm\epsilon_\phi\phi$ supplied by \textup{(NP1)}.
It gives
\[
               0\ge n(\psi_0)\pm\epsilon_\phi n(\phi),
\]
which is \textup{(NP8)}.  The presentation-wise bound
follows from \textup{(NP2)}.  If $n$ is continuous on
$\mathbb L_S$, vanishing on compact original smooth
tests extends to their defined closed core.  No density
of that core in the full lattice or difference space
is used.

For a literal datum, both integrals in \textup{(NP9)}
are finite by the energy-dual bound and mixed Holder.
The normal is literally
$n(\phi)=-\int\phi\,d(\mu-\widehat\mu_u)$,
so \textup{(NP8)} proves \textup{(NP10)}.
If the scalar defect is zero, every original smooth
test integral of the two Radon measures agrees.
Original scalar density and Radon uniqueness give
equality as original measures.  Conversely equality
of these measures gives equality of their finite
integrals against the actual continuous weight
$-\psi_0$, so the scalar defect is zero.

Any source potential $v$ has $J(v)=\ell_\mu$ by
the literal energy extension.  It is the unique
variational maximizer by the frozen convex energy
support inequality and strict separation.  Thus it
equals $u$.  This proves the third equivalence and
uniqueness.  Neither equality on additional GK test
forms nor an actual ambient classical representative
of every completed potential is inferred.
\end{proof}

\begin{corollary}[An additive strictness criterion for the source equation]
\label{cor:r12-additive-probe-source-criterion}
Under the literal datum hypotheses of \textup{(NP9)},
suppose its variational maximizer admits a decomposition
\[
                  u=v+\varepsilon\psi_0,
                  \qquad v\in\mathscr C,\quad\varepsilon>0.
 \tag{NP11}
\]
Then $\widehat\mu_u=\mu$ on all original tests.

The hypothesis is an \emph{additive cone decomposition}.
The pointwise inequalities $u\ge\varepsilon\psi_0$ or
$u\le\varepsilon\psi_0$ do not imply it.  Nor is it
enough that $u$ merely has a bounded classical lower
barrier.  Those replacements are contradicted by the
explicit incompatible bounded-source examples.
\end{corollary}

\noindent\textbf{Status: \texttt{NEW\_PROVED}} conditional
on the explicit geometric variation hypothesis, not on
an assertion restating the source equation.

\begin{proof}
The normal relation gives
$0=n(u)=n(v)+\varepsilon n(\psi_0)$.
Both terms are nonpositive, so $n(\psi_0)=0$.
The one-probe theorem applies.  Equivalently,
\textup{(NP1)} shows directly that every original
compact smooth direction can be varied with both
signs at $u$ within the cone; the normal must
annihilate it.  This is the step unavailable at
an arbitrary maximizer.
\end{proof}

\begin{proposition}[Density of the test-interior source family]
\label{prop:r12-dense-test-interior-sources}
For every $a\in\mathscr C$ and $\varepsilon>0$,
$a+\varepsilon\psi_0$ admits both signs of every
original compact smooth direction within $\mathscr C$.
These potentials approach $a$ in the energy-difference
norm as $\varepsilon\downarrow0$.  Their curvature
functionals satisfy the explicit estimate
\[
 \|J(a+\varepsilon\psi_0)-J(a)\|_*
       \le (r(a)+\varepsilon r(\psi_0))^q-r(a)^q.
 \tag{NP12}
\]
Thus this test-interior source family is dense in
$J(\mathscr C)$ for the signed operator dual norm.
It is not asserted dense among all positive data,
nor is an arbitrary prescribed datum automatically
the source of a test-interior potential.
\end{proposition}

\noindent\textbf{Status: \texttt{DIRECT\_COROLLARY}}
of the two-sided probe and the bounded mixed energy
form on $\mathbb D_S$.

\begin{proof}
For a fixed test $\phi$, use
$a+\varepsilon(\psi_0+t\phi)\in\mathscr C$
with $|t|\le\epsilon_\phi$.  The difference from
$a$ has norm $\varepsilon r(\psi_0)$.
Expand the $q$ slots of $J(a+\varepsilon\psi_0)$.
Each mixed functional has operator norm at most
the product of the roots of its differentiated
slots, by the proved mixed energy dual bound.
The binomial sum is precisely the right side of
\textup{(NP12)}.  This also proves the density
claim without approximating the original background.
\end{proof}

\begin{proposition}[Stability and a scalar lower bound for image distance]
\label{prop:r12-normal-probe-stability}
For $\ell,m\in\mathbb D_{S,\mathrm s}^*$ let
$\Xi(\ell)=\Xi_{\psi_0}(\ell,\Pi(\ell))$,
$\delta=\|\ell-m\|_*$ and
$R=\max(\|\ell\|_*,\|m\|_*)$.  Then
\[
            |\Xi(\ell)-\Xi(m)|
            \le r(\psi_0)\bigl(\delta+\sqrt{2qR\delta}\bigr).
 \tag{NP13}
\]
If $r(\psi_0)>0$, the normal satisfies
$\|\ell-J(\Pi(\ell))\|_*\ge\Xi(\ell)/r(\psi_0)$.
For $\|\ell\|_*>0$ this gives
\[
 \operatorname{dist}_*(\ell,J(\mathscr C))
 \ge
 \left(\sqrt{q\|\ell\|_*+\Xi(\ell)/r(\psi_0)}
                        -\sqrt{q\|\ell\|_*}\right)^2.
 \tag{NP14}
\]
Thus a strictly positive scalar normal defect certifies
a quantitative failure of source-image approximation.
The converse zero-defect assertion uses literal data
as in \textup{(NP9)}, not an arbitrary dual shadow.
\end{proposition}

\noindent\textbf{Status: \texttt{DIRECT\_COROLLARY}}
of the scalar definition and the proved nonlinear
retraction estimate \textup{(NI11)--(NI12)}.

\begin{proof}
Apply the dual norm to the difference of the two
values at $\psi_0$, then use \textup{(NI11)}.
This proves \textup{(NP13)}.  Testing the residual
norm at $\psi_0$ gives its displayed lower bound.
The right side of \textup{(NI12)}'s lower estimate
is increasing in that residual norm, yielding
\textup{(NP14)}.  No absolute-variation norm
$\delta_1$ is substituted for the signed dual norm.
\end{proof}

\paragraph{Why literal data cannot be omitted from the equivalence.}
In the radial example of Section~\ref{sec:r12-radial-pole-order-dual},
the functional $L_\infty$ vanishes on every bounded
lattice element and is nonnegative on $\mathscr C$.
Use the geometric bounded probe for that same background
and set $\ell=J(\psi_0)-L_\infty$, $u=\psi_0$.
Then $n=-L_\infty$ is nonpositive on the potential
cone and $n(u)=0$.  The convex support inequality
shows that $u$ is its unique variational maximizer.
The scalar defect is zero, but $n\ne0$ on the
energy-difference space.  Its restriction to every
original compact smooth test is zero, as the theorem
predicts.  This is a \texttt{NEW\_PROVED} example
of the nonliteral residual left undetected by a
current-test criterion.  It does not contradict
\textup{(NP9)} and does not assert strong continuity
of this $\ell$.

\paragraph{Remaining minimal issue.}
For literal finite-energy data the full original source
equation is now reduced to a single nonnegative
geometry-weighted variational defect.  Its vanishing
is still a genuine source-compatibility restriction:
the positive-capacity atom in \textup{(RF18)} has
strictly positive defect by \textup{(NP9)}.  A useful
general source theorem must prove vanishing from
additional prescribed-data or geometric hypotheses.
The probe theorem neither claims such unrestricted
vanishing nor changes OPEN-GQ, OPEN-CTO, or the
ordinary/extra-delta-test distinction.

% END REVIEWED MODULE research_r12/nonlinear/single_probe_normal_residual.tex


% BEGIN REVIEWED MODULE research_r12/capacity/weak_cluster_dual_reduction.tex
\section{Every localized tester cluster is a literal strong-dual datum}
\label{sec:r12-capacity-cluster-strong-dual}

Retain exactly the strict original box and background-extension
hypotheses (M2), with arbitrary original symmetric middle-positive
separately closed $S$ of degree $(p,p)$. Put $q=n-p\geq1$,
$d=q+1$, and $C_0=d2^d$. All capacities in this section use the
full original normalized classical psh tester class on the same
whole domain $V$. Let
\[
 \mathcal A(V)=\{a: a\text{ is original classical psh on }V,
                                      \ 0\leq a\leq1\}.
 \tag{CW1}
\]
Fix an original compact smooth scalar $\chi$ with $0\leq\chi\leq M$,
where $M<\infty$. For any row $a_j\in\mathcal A(V)$ put
\[
       \nu_j=\chi\,(Da_j)^q\wedge S.
 \tag{CW2}
\]
These are positive Radon measures, not necessarily global source
measures of potentials in the zero-boundary energy cone. The latter
distinction is essential below.

\begin{proposition}[An unconditional full-tester cluster datum theorem]
\label{prop:r12-capacity-cluster-literal-dual}
Every subsequential weak original-current limit $\nu$ of (CW2)
is a finite positive Radon measure supported in
$H=\operatorname{supp}\chi$, and for every Borel $E\subset V$,
\[
                         \nu(E)\leq M C_S^*(E;V).
 \tag{CW3}
\]
In particular, with $m=\nu(V)$, literal integration on the complete
original Borel energy-difference space defines
\[
 \begin{split}
 \ell_\nu(h)&=-\int_Vh\,d\nu\in\mathbb D_{S,\mathrm s}^*,\\
 \|\ell_\nu\|_*
       &\leq \frac{d}{d-1}
                 m^{1-1/d}(M C_0)^{1/d}=:B_{m,M}.
 \end{split}
 \tag{CW4}
\]
The right side is zero if $m=0$. For every strongly convergent
row $v_j\to v$ in the unchanged full cone $\mathscr C$,
\[
                         \int_V|v_j-v|\,d\nu\longrightarrow0.
 \tag{CW5}
\]
There is no no-loss, common-barrier, actual-representability or
source-compatibility hypothesis in this proposition.
\end{proposition}

\noindent\textbf{Status: \texttt{DIRECT\_COROLLARY}} of
positive weak compactness, the definition of the full tester
capacity and the proved capacity-modulus strong-dual theorem
(CD1)--(CD10). The full arbitrary-$S$ weak-cluster application
is recorded here to isolate its unconditional content.

\begin{proof}
Relative CLN gives one finite bound for the masses of (CW2),
independent of every $a_j$ and of its actual moment presentations.
All measures are supported in the fixed compact $H$. Original
smooth scalar density and positivity identify weak convergence
on original tests with weak Radon convergence on this compact
stage. For instance, approximate a continuous compact test
uniformly by original smooth tests and use the common mass bound
to pass the approximation error. Positivity and the fixed support
therefore pass to the limit, and $m<\infty$.

For every ordinary open $G\subset V$, the open-set Portmanteau
inequality and the original tester definition give
\[
    \nu(G)\leq\liminf_j\nu_j(G)
        \leq M C_S(G;V).
 \tag{CW6}
\]
For a Borel $E\subset G$, $\nu(E)\leq\nu(G)$.
Taking the infimum over all ordinary open neighborhoods $G$
of $E$ proves (CW3). This does not replace $C_S^*$ by raw
$C_S$ on a compact or Borel set. In particular it uses no
instance of the missing compact no-gain assertion (FG1).

The modulus $F(s)=Ms$ satisfies (CD1), because
$d>1$. Formula (CD10) with $\gamma=1$ gives exactly (CW4),
and (CD5) gives (CW5). The cited theorem proves integrability
and representative independence on the actual Borel difference
space, not just an abstract positive extension from compact
tests. Its strict-box hypotheses are precisely those retained
above. Every actual refinement remains in the original test
class used in (CW6) and in the capacity definition.
\end{proof}

\begin{corollary}[A canonical normal pair for each compact tester cluster]
\label{cor:r12-cluster-normal-pair}
For every $\nu$ as above there is a uniquely determined
variational potential $u_\nu=\Pi(\ell_\nu)\in\mathscr C$ and
normal $n_\nu=\ell_\nu-J(u_\nu)$ satisfying
\[
 \begin{gathered}
 n_\nu(w)\leq0\quad(w\in\mathscr C),\qquad n_\nu(u_\nu)=0,\\
                  e(u_\nu)\leq B_{m,M}^{d/q}.
 \end{gathered}
 \tag{CW7}
\]
The geometry-only actual probe $\psi_0$ defines a finite scalar
\[
 \Xi_\nu=\int_V(-\psi_0)\,d\widehat\mu_{u_\nu}
                       -\int_V(-\psi_0)\,d\nu\geq0,
 \tag{CW8}
\]
and, for every original compact smooth scalar $\phi$,
\[
 \left|\int_V\phi\,d(\nu-\widehat\mu_{u_\nu})\right|
                   \leq\epsilon_\phi^{-1}\Xi_\nu.
 \tag{CW9}
\]
The condition $\Xi_\nu=0$ is equivalent to existence of a
\emph{completed global zero-boundary source} for $\nu$.
It is not equivalent to compact raw-capacity no-gain.
\end{corollary}

\noindent\textbf{Status: \texttt{DIRECT\_COROLLARY}} of
(CW4), the variational normal decomposition (NI6)--(NI10),
and the original-test normal detector (NP6)--(NP10).

\begin{proof}
The preceding proposition places $\ell_\nu$ in the exact
strong-continuous dual where the variational construction is
unconditional. Since it is nonnegative on the nonpositive
potential cone, (NI9) gives
$r(u_\nu)^q=\|\ell_\nu\|_*\leq B_{m,M}$.
Raising to the power $d/q$ proves the energy bound in (CW7).
The normal relations are (NI7), and the remaining assertions
are the literal-data part of the one-probe theorem. No actual
classical representative, normalized tester bound for $u_\nu$,
or zero normal residual is inferred from variational existence.
\end{proof}

\begin{proposition}[Localization can give positive normal defect with no compact gain]
\label{prop:r12-localized-normal-defect-not-capacity-gate}
For every $q\geq1$ and every $n>q$ there are allowed original
data, one normalized actual tester $a$ and one $0\leq\chi\leq1$
as above such that the constant row $a_j=a$ has localized limit
$\nu=\chi(Da)^q\wedge S$ satisfying
\[
 \begin{gathered}
       \nu(E)\leq C_S(E;V)\quad\text{for every Borel }E\subset V,\\
       \Xi_\nu>0,
       \qquad \widehat\mu_v\ne\nu\quad\text{for every }v\in\mathscr C.
 \end{gathered}
 \tag{CW10}
\]
In particular compact no-gain holds for this row, although its
localized datum has no completed global source.
\end{proposition}

\noindent\textbf{Status: \texttt{DIRECT\_COROLLARY}}
of the proved proper-moment compact-source obstruction and
normal detector, with the constant normalized tester written
explicitly. The attempted assertion that every localized tester
cluster has $\Xi_\nu=0$ is \texttt{FALSE}.

\begin{proof}
Let $X=(\mathbf G_{m,K}^n)^{\rm an}$ over the original $K$,
$t_i=\log|T_i|$, $Q=(-1,1)^n$ and $V=t^{-1}(Q)$. Put
\[
       \beta=\sum_{i=1}^n d't_i\wedge d''t_i,
       \qquad S=\beta^{n-q}\wedge[V],
       \qquad a=\frac{1}{2n}\sum_{i=1}^n t_i^2.
 \tag{CW11}
\]
The background extends across the closed box, is original
middle-positive and separately closed, and has the required
degree. The actual smooth convex tester satisfies $0\leq a<1/2$
on the whole $V$. Its full original measure is
\[
                      (Da)^q\wedge S=n^{-q}\beta^n\wedge[V],
 \tag{CW12}
\]
which is nonzero and has positive density on the original
Gauss section by the frozen torus integration calibration.

Choose a nonzero nonnegative real smooth $\chi_0\leq1$
compactly supported in $Q$, and put $\chi=\chi_0\circ t$.
The moment $t:V\to Q$ is proper over compact subsets,
so $\chi$ is an original compact smooth scalar.
Then $\nu=\chi(Da)^q\wedge S$ is nonzero, positive,
finite and has compact tropical carrier.
The first line of (CW10) follows immediately by using the
single normalized tester $a$ in the raw capacity supremum.
No limiting or outer-regularity argument is needed.

On the other hand the proper-moment lift theorem applies
with real rank $r=n$, lift background $[V]$, and $q<n$.
It excludes every nonzero source with compact tropical
carrier from the full completed source image. Equivalently,
the smooth separator theorem supplies an original compact
smooth $b$ with
\[
       \int b\,d\nu=\nu(V)>0,\qquad
       \int b\,d\widehat\mu_v\leq0
                                      \quad(v\in\mathscr C).
 \tag{CW13}
\]
The datum is in the literal strong dual by (CW4).
The one-probe estimate applied to $b$ gives the quantitative
strictness $\Xi_\nu\geq\epsilon_b\nu(V)>0$.
This proves (CW10) on all original tests and refinements.
The special background is used solely to disprove the
universal zero-defect shortcut; the positive cluster theorem
assumes no such factorization of its arbitrary $S$.
\end{proof}

\begin{theorem}[A quantitative weakly compact family inside the strong dual]
\label{thm:r12-compact-capacity-dual-equicontinuity}
Fix a compact $H\Subset V$ and finite $m_0,M\geq0$. Let
$\mathcal G(H,m_0,M)$ be the positive Radon measures $\lambda$
with
\[
 \operatorname{supp}\lambda\subset H,\qquad
 \lambda(V)\leq m_0,\qquad
 \lambda(E)\leq M C_S^*(E;V)\quad(E\text{ Borel}).
 \tag{CW14}
\]
This is a weakly compact family, and its literal images
$\ell_\lambda$ form a weak-$*$ compact subset of
$\mathbb D_S^*$ contained entirely in
$\mathbb D_{S,\mathrm s}^*$. Here weak-$*$ means
$\sigma(\mathbb D_S^*,\mathbb D_S)$, not operator-norm
convergence.

More quantitatively, fix $R<\infty$ and let
$r(u),r(v)\leq R$, with $u,v\in\mathscr C$. Put
$\alpha=2^{-q}$ and $\varepsilon=I(u,v)^\alpha$.
There is a finite $A=A_{H,R,M,S}$, independent of $u,v$
and of the measure, such that for $0<\varepsilon\leq1$,
\[
 \sup_{\lambda\in\mathcal G(H,m_0,M)}
       \int_V|u-v|\,d\lambda
 \leq\varepsilon\left[
        m_0+\frac{2MC_0R^d}{d-1}
             +\frac{dA}{d-1}\log\frac1\varepsilon\right].
 \tag{CW15}
\]
For $I(u,v)=0$ the left side is zero. In particular this
family is uniformly continuous on every bounded-energy
strong cone, with the explicit logarithmic modulus in (CW15).

If $\lambda_j\to\lambda$ weakly in this family and
$v_j\to v$ strongly in $\mathscr C$, then
\[
 \begin{split}
   \int_V|v_j-v|\,d\lambda_j&\longrightarrow0,\\
   \int_V v_j\,d\lambda_j&\longrightarrow\int_Vv\,d\lambda.
 \end{split}
 \tag{CW16}
\]
For fixed $h\in\mathbb D_S$ one consequently has
$\ell_{\lambda_j}(h)\to\ell_\lambda(h)$.
These conclusions apply in particular to every weakly
convergent localized normalized tester row (CW2), with
$m_0$ supplied by its uniform CLN mass bound.
\end{theorem}

\noindent\textbf{Status: \texttt{NEW\_PROVED}} for the
uniform logarithmic strong modulus, the full weak-$*$
compact datum family and the jointly varying test conclusion.
Its inputs are the frozen first-energy sublevel estimate,
the all-original compact capacity estimate (EV8), completed
outer quasicontinuity and the proved datum criterion (CD).
No Banach-norm compactness or nonlinear source closure under
weak measure convergence is asserted.

\begin{proof}
The frozen compact gap estimate, passed to the strong
completion, gives a finite $C_{H,R,S}$ such that
\[
 C_S^*(H\cap\{|u-v|>t\};V)
       \leq C_{H,R,S} I(u,v)^\alpha/t\qquad(t>0).
 \tag{CW17}
\]
Here is a direct justification of this use of completion.
Choose actual defining rows for $u,v$ with uniformly
bounded energy roots, tending strongly to them. For a
fixed $t$, the limiting gap is contained in the actual
gap at $t/2$ together with the two approximation gaps
at $t/4$, outside their common null set. The latter
two capacities tend to zero. Apply (EV8) to the actual
gap and use continuity of the mixed first energies in
the strong completion. Absorb its factor two in
$C_{H,R,S}$. This uses capacity subadditivity and
capacity convergence, not continuity from above on
arbitrary compact sets or raw/outer equality.

For each $\lambda$ in (CW14), the distribution of
$|u-v|$ therefore has the simultaneous bounds
\[
 \lambda\{|u-v|>t\}
       \leq\min\left\{m_0,\frac{A\varepsilon}{t},
                            \frac{B}{t^d}\right\},
 \qquad A=MC_{H,R,S},\quad B=2MC_0R^d.
 \tag{CW18}
\]
The third bound follows from
$\{|u-v|>t\}\subset\{u<-t\}\cup\{v<-t\}$ for
nonpositive representatives, and the complete global
sublevel estimate $C_S^*\{w<-t\}\leq C_0e(w)t^{-d}$.
All pole and representative exceptional sets are
annihilated by (CW14).

Integrate (CW18) by layer cake. Use $m_0$ on
$(0,\varepsilon)$, $A\varepsilon/t$ on
$(\varepsilon,\varepsilon^{-1/(d-1)})$, and $B/t^d$
on the remaining interval. The three integrals are
respectively
\[
       m_0\varepsilon,\qquad
       \frac{dA}{d-1}\varepsilon\log(1/\varepsilon),
       \qquad \frac{B}{d-1}\varepsilon.
\]
This is (CW15). If $I=0$, the relative classes agree
outside an outer-null set, so the asserted zero follows
directly. A strongly convergent row has bounded energy
roots; hence (CW15) proves the first line of (CW16).

It remains to pass a fixed completed weight through
ordinary weak measure convergence. Put
$v^{(L)}=\max(v,-L)$. Its restriction to $H$ is bounded
and outer-quasicontinuous. For any $\delta>0$ choose
an ordinary open $G$ with $C_S^*(G;V)<\delta$ so
that this restriction is continuous on $H\setminus G$.
The ordinary continuous extension and compact cutoff
construction gives $f\in C_c(V)$ with $|f|\leq L$
and $f=v^{(L)}$ on $H\setminus G$. Thus both the
approximating measures and the limit satisfy
\[
       \int|v^{(L)}-f|\,d\lambda_j\leq2LM\delta,
       \qquad
       \int|v^{(L)}-f|\,d\lambda\leq2LM\delta.
\]
Weak convergence passes the $f$ integral; one may
equivalently use original smooth scalar density and
the common mass bound. Letting $\delta\downarrow0$
passes the fixed clipped weight. The global sublevel
bound also gives, uniformly over the entire family,
\[
 \int\bigl((-v)-L\bigr)_+\,d\lambda_j
       \leq\frac{MC_0e(v)}{d-1}L^{1-d},
\]
and the same bound for $\lambda$. Therefore
$\int v\,d\lambda_j\to\int v\,d\lambda$.
Combine this with the first line of (CW16) to obtain
its second line. Applying the fixed-weight assertion
to each of two cone members whose difference is $h$
proves convergence on all $\mathbb D_S$.

Finally, positive measures supported in the compact $H$
with mass at most $m_0$ form a weakly compact set.
The subclass (CW14) is weakly closed: on ordinary
opens use Portmanteau and
$C_S^*(G;V)=C_S(G;V)$, then take ordinary open
neighborhood infima for each Borel set, just as in
(CW6). Every member is a literal strong-dual datum
by (CD). The fixed-weight argument just proved shows
that the map $\lambda\mapsto\ell_\lambda$ is
continuous into the weak-$*$ topology. The same proof
works for nets, since its two error bounds are uniform
and its middle continuous-test limit is the definition
of weak convergence. Its image is consequently
weak-$*$ compact. This proves all the claims without
claiming that the entire strong dual is weak-$*$ closed.
\end{proof}

\paragraph{The attempted compact no-gain proof and its exact stopping point.}
The new Jensen theorem represents a nonlinear weak cluster,
before an arbitrary compact cutoff, by
$\nu=\widehat\mu_u-DV_q$ with $V_q\geq0$.
Its universal scalar loss controls every original coefficient
of $V_q$. A compact-set extremizing sequence, however, need
not minimize this Jensen loss. The finite smooth examples
in the Jensen section show that $-DV_q$ has neither sign;
convex averaging cannot be asserted to preserve, increase,
or decrease the mass on the target compact set.

Even vanishing Jensen loss yields the source of a normalized
\emph{completed} Borel potential. It does not construct a
normalized actual classical tester with at least the same
mass on that compact. Replacing a sharp compact restriction
by weak convergence is the already recorded moving-atom
failure. These are distinct obstacles, not a counterexample
to (FG1).

The variational alternative always produces the well-defined
normal pair (CW7), but (CW10) proves that its global source
defect is too strong a condition for localized compact no-gain.
Consequently the route from either scalar defect directly to
(FG1) remains a \texttt{FAILED\_ATTEMPT}. The missing
assertion is still an actual compact mass comparison for the
full tester family, not existence of a finite-energy dual
functional for each localized weak limit. No unrestricted
comparison, contact orthogonality, Radon-probe null-ideal
identification or actual-representability theorem is assumed.

\paragraph{Exact new dependencies.}
The positive dual criterion is
\path{energy/capacity_controlled_dual_data.tex}, SHA256:
\begin{quote}\small
\nolinkurl{415cab8b5265bd34d298db187317ac28cd075f4265939e0d362a0b30c42da7e0}.
\end{quote}
The variational construction is
\path{energy/nonlinear_dual_source_geometry.tex}, SHA256:
\begin{quote}\small
\nolinkurl{29a083c0f65d93f777cae9bbe0c02a36385510e7b1a6c9542552322439168743}.
\end{quote}
The normal test detector is
\path{nonlinear/single_probe_normal_residual.tex}, SHA256:
\begin{quote}\small
\nolinkurl{f09fb0f75efdce5fd637e088a70262091ed650c13cb216b5611541b6e0de28d4}.
\end{quote}
The compact-source separator is
\path{nonlinear/moment_source_separation.tex}, SHA256:
\begin{quote}\small
\nolinkurl{42c807dff24816a0e7466a089af094276fcab1a48c800821e02d28b0c098084c}.
\end{quote}
These are proofs in the present round; the frozen bounded calculus,
relative CLN and exact original torus calibration retain their
\texttt{INTERNAL\_INPUT} status.
The frozen outer gap estimate is (EV8) in
\path{R11/research_r11/comparison/first_energy_capacity_topology.tex},
SHA256:
\begin{quote}\small
\nolinkurl{5e1c0eba5cfeef5ddc15a25d4c7cb5f1c4972fc325a4568d0f01339ee203da94}.
\end{quote}
Finite local outer quasicontinuity is precisely
\path{prop:r11-strong-cauchy-borel-limit} in
\path{R11/research_r11/comparison/strong_first_energy_completion.tex},
SHA256:
\begin{quote}\small
\nolinkurl{11392009f63f5de3b21dd2d90d8ded48682ede8cb79665ab74e3da8ebf6534df}.
\end{quote}

% END REVIEWED MODULE research_r12/capacity/weak_cluster_dual_reduction.tex


% BEGIN REVIEWED MODULE research_r12/capacity/energy_ball_weak_compactness_obstruction.tex
\section{Energy balls need not be weakly compact in either Banach realization}
\label{sec:r12-energy-ball-weak-compactness-obstruction}

The weak-$*$ compact datum family in (CW14) does not by itself
supply a weakly compact family of candidate potentials. The
following consequence of the proved pole example makes that
obstruction explicit. It is a statement about the two new
Banach norms, not about failure of the complete strong topology.

Fix $n>q\geq1$ and the original radial-mixture background
$S$ of (PO3), on its original unit polydisc $V$. Retain the
original field and absolute value. Put $d=q+1$ and use the
already proved actual potentials
\[
 \psi_j=\max\!\left(\log\max_{1\leq i\leq n}|T_i|,-j\right),
 \qquad
 u=\log\max_{1\leq i\leq n}|T_i|,
                    \qquad j=1,2,\ldots .
 \tag{WB1}
\]
The value of $u$ at the origin is $-\infty$. It represents
a member of the same no-loss first-energy cone, and (PO7)
gives
\[
       e(u)=1,\qquad e(\psi_j)=1-e^{-j},\qquad
                               I(\psi_j,u)=e^{-j}.
 \tag{WB2}
\]
The pole functional is continuous on both $\mathbb D_S$
and $\mathbb L_S$ and satisfies
\[
                   L_\infty(u)=1,\qquad L_\infty(\psi_j)=0.
 \tag{WB3}
\]
All these statements are exactly the earlier original-current
pole construction, including its actual Borel representatives.

\begin{proposition}[No weak cluster and nonreflexivity in every remaining degree]
\label{prop:r12-energy-ball-nonreflexivity}
For the allowed data just specified, the sequence $(\psi_j)$
has no weakly convergent subnet, and in particular no weakly
convergent subsequence, in either $\mathbb D_S$ or
$\mathbb L_S$. Consequently
\[
      \mathcal B_{\mathscr C}
            =\{v\in\mathscr C:r(v)\leq1\}
 \tag{WB4}
\]
is not weakly compact in either Banach space, and neither
$\mathbb D_S$ nor $\mathbb L_S$ is reflexive.

This holds in every remaining degree $q\geq1$ by the
specified example in every $n>q$. It does not say that
these spaces are nonreflexive for every possible background
$S$, including the zero or other degenerate backgrounds.
\end{proposition}

\noindent\textbf{Status: \texttt{DIRECT\_COROLLARY}} of
the present round's exact Banach realizations, source-functional
continuity and pole functional. The proposed universal weak
compactness of bounded potential energy balls is
\texttt{FALSE}. The weak closedness of norm-closed convex
sets and weak compactness of bounded closed balls in a
reflexive Banach space are \texttt{EXTERNAL\_KNOWN}
functional-analysis facts.

\begin{proof}
Let $E$ denote either $\mathbb D_S$ or $\mathbb L_S$.
The proved realization theorems say that the same cone
$\mathscr C$ is norm closed and convex in $E$, and that
\[
                         \|v\|_E=r(v)\qquad(v\in\mathscr C).
 \tag{WB5}
\]
It is therefore weakly closed, by the Hahn--Banach
separation theorem. Every source functional $J(a)$,
$a\in\mathscr C$, is continuous on $E$. In the lattice
case this uses its proved literal extension (OL6),
not a newly assumed extension of an arbitrary dual.
On the cone these source functionals are also continuous
for the strong discrepancy topology, by the frozen
tester estimate and (NI4).

Suppose that a subnet $(\psi_{j_\gamma})$ converges
weakly in $E$ to $v$. The subnet indices are cofinal
in the integer row, so (WB2) still makes it converge
strongly to $u$. Weak closedness gives $v\in\mathscr C$.
For every fixed $a\in\mathscr C$, the weak and strong
continuities just stated give, respectively,
\[
      J(a)(v)=\lim_\gamma J(a)(\psi_{j_\gamma})=J(a)(u).
 \tag{WB6}
\]
Choose $a=v$ and $a=u$. The exact energy-discrepancy
identity then yields
\[
              I(v,u)=(J(v)-J(u))(v-u)=0.
 \tag{WB7}
\]
The strong cone is separated, so $v=u$ as its original
relative Borel class. On the other hand, continuity of
$L_\infty\in E^*$ and (WB3) would give
$L_\infty(v)=0$, whereas $L_\infty(u)=1$. This is a
contradiction. It proves the absence of a weakly
convergent subnet in either space.

The sequence lies in (WB4), by (WB2) and (WB5).
Every net in a compact space has a convergent subnet;
thus (WB4) is not weakly compact. Finally (WB4) is
the intersection of the norm-closed convex cone with
the closed unit ball, so it is weakly closed in that
ball. If $E$ were reflexive, its closed unit ball
would be weakly compact, making (WB4) weakly compact
as well. This contradiction proves nonreflexivity.
\end{proof}

\paragraph{Consequence for the compact-capacity attempt.}
The mechanism ``bounded first energies imply weakly compact
potential candidates; then use weak-$*$ compactness of the
localized source data'' is invalid in these Banach realizations.
The first assertion is disproved by (WB1)--(WB7), even though
the displayed row actually converges strongly. This does not
contradict strong completeness of $\mathscr C$, completeness
of the curvature metric $d_J$, or the weak-$*$ compactness
of the specified datum family in the \emph{dual} space.
It also gives no counterexample to (FG1), to compact ordinary
outer regularity, or to the analytic relation quotient.

\paragraph{Exact dependency snapshots.}
The energy-difference closed-cone theorem is
\path{energy/energy_difference_banach.tex}, SHA256:
\begin{quote}\small
\nolinkurl{12bde9976ef8d8dcb3abbfc23f71f8e1a1ae37a6195de0dcdd13f09c5c0f5f45}.
\end{quote}
The lattice realization and source extension are
\path{nonlinear/order_energy_lattice_completion.tex}, SHA256:
\begin{quote}\small
\nolinkurl{754db6baac08f682e94d5342aa3e757914f3442592bd10cb288f8f2c7ce1bb69}.
\end{quote}
The exact energies and pole dual are
\path{nonlinear/radial_pole_order_dual_obstruction.tex}, SHA256:
\begin{quote}\small
\nolinkurl{f57f4ccac82fc36cbfb4c6663091bdc738d800ebaf022eb9c602176854fa4365}.
\end{quote}
The source strong-continuity and discrepancy identity are
\path{energy/nonlinear_dual_source_geometry.tex}, SHA256:
\begin{quote}\small
\nolinkurl{29a083c0f65d93f777cae9bbe0c02a36385510e7b1a6c9542552322439168743}.
\end{quote}

% END REVIEWED MODULE research_r12/capacity/energy_ball_weak_compactness_obstruction.tex


% BEGIN REVIEWED MODULE research_r12/capacity/weak_dual_variational_energy_stability.tex
\section{Weak dual convergence with scalar variational energy no-loss}
\label{sec:r12-weak-dual-variational-energy-stability}

Retain the complete strict-box and extending arbitrary-background
hypotheses (M2), with $q=n-p\geq1$, $d=q+1$, and the unchanged
full strong cone $\mathscr C$. The following criterion supplies
a positive compactness mechanism without asserting that all
bounded potential energy balls are weakly compact.

\begin{proposition}[An exact scalar criterion for variational strong stability]
\label{prop:r12-weak-dual-energy-no-loss}
Let $\ell_j,\ell\in\mathbb D_{S,\mathrm s}^*$ satisfy
\[
            \ell_j(h)\longrightarrow\ell(h)
                                 \quad(h\in\mathbb D_S).
 \tag{VE1}
\]
They may be signed. In particular (VE1) is weak-$*$
convergence, not assumed operator dual-norm convergence.
Put $u_j=\Pi(\ell_j)$, $u=\Pi(\ell)$ and
\[
 \begin{split}
 \delta_j
   &=F_{\ell_j}(u_j)-F_{\ell_j}(u)\\
   &=\frac qd\bigl(e(u_j)-e(u)\bigr)-(\ell_j-\ell)(u).
 \end{split}
 \tag{VE2}
\]
Then
\[
                     0\leq I(u_j,u)\leq d\delta_j,
               \qquad \liminf_j e(u_j)\geq e(u).
 \tag{VE3}
\]
Moreover the following are equivalent:
\[
 \begin{gathered}
             e(u_j)\longrightarrow e(u),\qquad
             \delta_j\longrightarrow0,\qquad
             I(u_j,u)\longrightarrow0.
 \end{gathered}
 \tag{VE4}
\]
If $r(u_j),r(u)\leq R$, the source functionals satisfy
the quantitative estimate
\[
       \|J(u_j)-J(u)\|_*
              \leq\sqrt{qd}\,R^{(q-1)/2}\delta_j^{1/2}.
 \tag{VE5}
\]
When (VE4) holds, all original mixed-current and weighted
strong limits apply to the potentials, and the normal
residuals $\ell_j-J(u_j)$ converge weak-$*$ to
$\ell-J(u)$. Residual operator-norm convergence is not
asserted from (VE1).
\end{proposition}

\noindent\textbf{Status: \texttt{DIRECT\_COROLLARY}} of
the proved variational normal decomposition, the one-sided
all-degree Bregman bound (PV2), and the curvature-dual
strong estimate (NI2). The criterion is a scalar energy
no-loss condition for already constructed canonical
variational potentials; it does not assume that the data
have source solutions.

\begin{proof}
The normal decomposition gives
\[
      \ell_j=J(u_j)+n_j,\qquad
      n_j(w)\leq0\ (w\in\mathscr C),\qquad n_j(u_j)=0.
 \tag{VE6}
\]
Radial stationarity gives
$\ell_j(u_j)=e(u_j)$ and $\ell(u)=e(u)$, including
the zero-potential case. Substituting these equalities
in the definition of $F$ proves the second line of
(VE2). Since $u_j$ maximizes $F_{\ell_j}$, $\delta_j\geq0$.

There is also the exact identity
\[
 d\delta_j
    =e(u)-e(u_j)-dJ(u_j)(u-u_j)-d n_j(u).
 \tag{VE7}
\]
Apply (PV2) with $a=u$ and $b=u_j$ to the first three
terms. Their sum is at least $I(u,u_j)$. The last
term is nonnegative by (VE6), so (VE7) proves the
first assertion of (VE3). Since (VE1) passes the
fixed test $u\in\mathbb D_S$, the second line of
(VE2) and $\delta_j\geq0$ imply
$\liminf e(u_j)\geq e(u)$.

The same fixed-test convergence makes the first two
conditions in (VE4) equivalent. The bound just proved
makes the second imply the third. Conversely the
frozen energy-root inequality
$|r(u_j)-r(u)|^d\leq I(u_j,u)$ shows that strong
convergence implies energy convergence. This proves
the equivalence without an assumption of weak
compactness of the potential energy ball.

Apply the second estimate in (NI2), followed by
$I(u_j,u)\leq d\delta_j$, to obtain (VE5).
Under (VE4) the energy roots are bounded, so this
estimate gives norm convergence of the source
functionals. Subtracting it from (VE1) proves
weak-$*$ convergence of the residuals. The frozen
complete strong calculus supplies the stated original
mixed-current and weighted limits after every actual
moment refinement. No additional GK delta-test
equality is inferred.
\end{proof}

\begin{corollary}[A source-closure criterion for weakly convergent Radon data]
\label{cor:r12-weak-measure-energy-source-closure}
Let $\nu_j\to\nu$ weakly in one family
$\mathcal G(H,m_0,M)$ of (CW14), and let
$u_j=\Pi(\ell_{\nu_j})$, $u=\Pi(\ell_\nu)$.
Then (VE2)--(VE5) apply. Thus scalar convergence of
their canonical variational energies is equivalent
to strong convergence of these potentials.

If, additionally, every datum is a completed source
$\nu_j=\widehat\mu_{v_j}$ with $v_j\in\mathscr C$
and $e(v_j)\to e(u)$, then
\[
             v_j\longrightarrow u\text{ strongly},
                         \qquad\widehat\mu_u=\nu
 \tag{VE8}
\]
on all original compact scalar tests and all actual
refinements. An actual normalized classical
representative of $u$ is not concluded.
\end{corollary}

\noindent\textbf{Status: \texttt{DIRECT\_COROLLARY}}.

\begin{proof}
The compact capacity-controlled family theorem puts
all literal functionals in the strong dual and
proves (VE1) from ordinary weak measure convergence.
The proposition therefore applies. If the additional
source hypothesis holds, uniqueness in the variational
normal decomposition gives $u_j=v_j$. The no-loss
criterion makes them converge strongly. Formula
(VE5) then gives $J(v_j)\to J(u)$ in dual norm.
Since their weak-$*$ limit is already $\ell_\nu$,
one has $J(u)=\ell_\nu$. Every original compact
smooth scalar is in $\mathbb D_S$, so this equality
gives (VE8) on the full original current test sheaf.
\end{proof}

\paragraph{What remains missing in the unrestricted capacity argument.}
There is no proof here that weakly convergent full
tester clusters automatically satisfy variational
energy no-loss. Moreover a compactly localized tester
measure $\chi\mu_a$ is not automatically a global
completed source: (CW10) explicitly disproves that
shortcut. Even the successful source conclusion
(VE8) gives a completed potential, not the normalized
actual compact mass comparison required by (FG1).
The new criterion therefore supplies a correct
conditional passage through weaker data convergence,
while keeping these further gates separate.

\paragraph{Exact dependency snapshots.}
The variational normal decomposition and source estimate are
\path{energy/nonlinear_dual_source_geometry.tex}, SHA256:
\begin{quote}\small
\nolinkurl{29a083c0f65d93f777cae9bbe0c02a36385510e7b1a6c9542552322439168743}.
\end{quote}
The one-sided coefficient identity is (PV2) in
\path{energy/perron_variational_energy_gap.tex}, SHA256:
\begin{quote}\small
\nolinkurl{76769b5072af4d886ea97abd1b56e63084d3cdc858dd0635ed3cafd5bd75a702}.
\end{quote}
The weak Radon to weak-$*$ datum passage is
\path{capacity/weak_cluster_dual_reduction.tex}, SHA256:
\begin{quote}\small
\nolinkurl{aabd324949c7faf87352ce2c39689aafb05c08dab11a5e11fffb1eb89c120268}.
\end{quote}

% END REVIEWED MODULE research_r12/capacity/weak_dual_variational_energy_stability.tex


% BEGIN REVIEWED MODULE research_r12/nonlinear/radial_flux_source_inversion.tex
\section{An exact nonlinear flux law and compatible radial source equations}
\label{sec:r12-radial-flux-source-inversion}

Keep the original $n$-polydisc, $n>q\ge1$, and the single
fixed original background $S$ of \textup{(PO3)}.  This is
an explicit family of allowed backgrounds in every remaining
degree, not a replacement of the arbitrary background in the
general theorems.  Write $d=q+1$ and retain the original
radial parameter $t$ and diagonal Gauss map $\sigma$.

For a positive Radon measure $\nu$ on $(0,\infty)$ with
finite tails, use the \emph{open-tail} convention
\[
 H_\nu(t)=\nu((t,\infty)),\qquad t>0.
 \tag{RF1}
\]
Thus $H_\nu$ is right continuous and $-dH_\nu=\nu$
as distributions on $(0,\infty)$.  An atom at $t$ is
excluded from $H_\nu(t)$, but included in its left limit.
This convention is important in the clipping formulas below.

\begin{definition}[Admissible radial slopes and source fluxes]
\label{def:r12-radial-admissible-flux}
Let $\mathscr A_q$ be the cone of finite-valued,
nonnegative, nonincreasing right-continuous functions
$A:(0,\infty)\to[0,\infty)$ such that
\[
                 \int_0^\infty e^{-t}A(t)^d\,dt<\infty.
 \tag{RF2}
\]
Set
\[
 F_A(t)=\int_0^t A(s)\,ds,\qquad
 u_A(x)=-F_A(t(x))\quad(x\ne o),\qquad
 u_A(o)=-F_A(\infty).
 \tag{RF3}
\]
The value at $o$ may be $-\infty$.  A source $\nu$ is
\emph{radially flux-compatible} if it has finite tails,
the function $e^tH_\nu(t)$ is nonincreasing, and
\[
              \int_0^\infty e^{t/q}H_\nu(t)^{d/q}\,dt<\infty.
 \tag{RF4}
\]
These are conditions on the ordinary one-dimensional
Radon datum.  They do not invoke an unknown potential
or an unspecified solution of a source equation.
\end{definition}

\begin{theorem}[Exact radial flux and finite-energy inversion]
\label{thm:r12-exact-radial-flux-inverse}
Every $A\in\mathscr A_q$ defines an actual extended
classical psh potential $u_A$ in the original no-loss
zero-boundary first-energy class.  Its original top source is
\[
 \mu_{u_A}=\sigma_*\nu_A,\qquad
 H_{\nu_A}(t)=e^{-t}A(t)^q.
 \tag{RF5}
\]
Its energy and every mixed first energy have the exact formulas
\[
 \begin{split}
 e(u_A)&=\int_0^\infty e^{-t}A(t)^d\,dt,\\
 B(u_{A_0},\ldots,u_{A_q})
     &=\int_0^\infty e^{-t}\prod_{i=0}^q A_i(t)\,dt.
 \end{split}
 \tag{RF6}
\]
All source identities hold on every original compact smooth
scalar test and every actual moment refinement.

The source map from these radial potentials is a bijection
onto the flux-compatible Radon data.  The inverse is explicit:
\[
 A_\nu(t)=\bigl(e^tH_\nu(t)\bigr)^{1/q},\qquad
 u_\nu=-\int_0^{t(\cdot)}
                     \bigl(e^sH_\nu(s)\bigr)^{1/q}ds.
 \tag{RF7}
\]
For such a datum $u_\nu$ is the unique solution of
$\widehat\mu_v=\sigma_*\nu$ in the \emph{full} strong
classical-approximation cone $\mathscr C$, not just among
radial candidates.  The flux condition is asserted necessary
only within the displayed radial class; no classification of
all arbitrary nonradial source measures is claimed.
\end{theorem}

\noindent\textbf{Status: \texttt{NEW\_PROVED}} for the
explicit fixed allowed background.  The arbitrary-background
bounded calculus, no-loss extension, first-energy mixed
calculus, and original radial-shell calibration are frozen
\texttt{INTERNAL\_INPUT}.

\begin{proof}
On every interval $(0,R)$ the exponential weight is bounded
below, so Holder's inequality makes $F_A(R)$ finite.
The function $F_A$ is nonnegative, increasing and concave,
with $F_A(0)=0$.  The scalar function
$z\mapsto-F_A(-z)$ on $z<0$ is convex and nondecreasing.
The frozen radial/convex composition calculus therefore gives
the actual extended classical psh function in \textup{(RF3)}.
One can also construct it by decreasing the bounded classical
radial potentials with slopes
\[
                  A_{N,R}(t)=\min(A(t),N)\mathbf1_{(0,R)}(t).
 \tag{RF8}
\]
Each of these extends across the boundary by a finite linear
continuation of its radial scalar function.  The limits are
taken on the original space, with no extension of $K$.
Their finite values at every finite positive depth show
that the limit is not identically $-\infty$.  The boundary value of $u_A$ is zero
because $F_A(t)\to0$ as $t\downarrow0$, and its strict
negative sublevels are compactly contained smaller polydiscs.

For the moment take a bounded radial slope, clipped to zero
after a finite depth.  Polarization of the same original
radial-shell calibration used in \textup{(PO6)} gives
\[
 H_{T_{[V],n}(u_A^{[q]},\psi_s^{[n-q]})}(t)
                      =A(t)^q\mathbf1_{(0,s)}(t)
 \tag{RF9}
\]
at all continuity points.  Here the tail is taken along the
original diagonal Gauss ray; the formula asserts the full
original measure identity.  Indeed the top measure of
$a u_A+b\psi_s$ has tail
$(aA(t)+b\mathbf1_{(0,s)}(t))^n$ away from jump points.
The coefficient of $a^qb^{n-q}$ has the same binomial
factor on both sides of mixed multilinearity, so there is
no additional factorial.  Right continuity determines the
measure and all jump masses uniquely.

Insert the fixed background \textup{(PO3)}.  The coefficient
Fubini argument of \textup{(PO5)}, applied on each fixed
compact original stage, gives
\[
          H_{\mu_{u_A}}(t)
                =A(t)^q\int_t^\infty e^{-s}ds
                =e^{-t}A(t)^q.
 \tag{RF10}
\]
For a slope not bounded near zero, first fix $R$.  One has
\[
 \sup_V\left|u_{A_{N,R}}-u_{A\mathbf1_{(0,R)}}\right|
        \le\int_0^R(A-\min(A,N))\,dt\longrightarrow0.
\]
Frozen bounded uniform stability therefore gives the
clipped original source, including on compact sets
containing the origin.  Away from clipping depth $R$
its tail is the stated one.  Local boundedness on
every $(\varepsilon,R)$ gives the same formula there.
Only after this fixed-$R$ passage do we let $R$ tend
to infinity, using \textup{(RF11)} below to control
the entire collapsed tail and the value-truncation loss.
A compact subset of $V$ stays a positive radial depth
away from the outer boundary, so the zero-depth boundary
is not silently added to $V$.

Let $H_A(t)=e^{-t}A(t)^q$.  It is right continuous,
nonincreasing, finite for $t>0$, and tends to zero at
infinity.  It therefore defines a unique positive Radon
measure $\nu_A=-dH_A$ with the displayed open tail.
Clipping the slope at depth $R$ collapses its entire
remaining source to that depth:
\[
 \nu_{A\mathbf1_{(0,R)}}
       =\mathbf1_{(0,R)}\nu_A+\nu_A([R,\infty))\delta_R,
 \qquad
 \nu_A([R,\infty))=e^{-R}A(R-)^q.
 \tag{RF11}
\]
There is no ambiguity when $A$ jumps at $R$: the left
limit includes the pre-existing atom.  For $R\ge1$,
this collapsed mass is bounded by $e^{-R}A(1-)^q$
after choosing any fixed earlier positive depth instead
of $1$ if necessary.  In particular it tends to zero.

If $F_A$ is unbounded, its value truncation at $-j$
corresponds to the unique depth $R_j$ with $F_A(R_j)=j$;
on any possible plateau one uses the first such depth,
which gives the same clipped potential.  When $F_A$
is unbounded its slope cannot vanish at a finite depth,
so these depths tend to infinity.  Formula \textup{(RF11)}
proves absence of truncation loss.  If $F_A$ is bounded,
the potential is already bounded and there is no pole
loss.  Thus the frozen no-loss class and original source
in \textup{(RF5)} are established without inferring
membership from a formally finite nonpluripolar energy.

Integration by parts now gives
\[
              \int_0^\infty F_A\,d\nu_A
                   =\int_0^\infty A(t)H_A(t)\,dt.
 \tag{RF12}
\]
For completeness, both endpoint terms vanish.  At infinity
$F_A$ grows at most linearly and $H_A$ is exponentially
bounded.  At zero, monotonicity and Holder imply
\[
 F_A(t)A(t)^q
 \le \left(t^{q/d}\Bigl(\int_0^t A^d\Bigr)^{1/d}\right)
       \left(t^{-q/d}\Bigl(\int_0^t A^d\Bigr)^{q/d}\right)
 =\int_0^t A^d\longrightarrow0.
\]
The exponential factor does not affect this limit.
Substitution in \textup{(RF12)} proves the first identity
in \textup{(RF6)}, hence actual first-energy membership.
The second follows by polarizing that identity in
$A_0,\ldots,A_q$, using the frozen symmetric mixed first
energy.  All integrals are finite by ordinary Holder.

Conversely, the flux hypotheses make $A_\nu$ a member
of $\mathscr A_q$: the integral in \textup{(RF4)} is
exactly \textup{(RF2)} after substitution.  Formula
\textup{(RF5)} constructs $u_\nu$ with the specified
original Radon source.  The right-continuous tail recovers
$A$, and integration with boundary value zero recovers
the radial potential, proving the radial bijection.

Finally suppose any $v\in\mathscr C$ has the same
original source measure as $u_\nu$.  The literal mixed
energy extension identifies their curvature functionals
on every cone member, hence on $\mathbb D_S$.
The strong discrepancy is zero and frozen separation
gives $v=u_\nu$ as relative classes.  This proves
uniqueness in the full cone without asserting radiality
of arbitrary other source data.
\end{proof}

\begin{corollary}[A linear compatibility condition for finite radial data]
\label{cor:r12-finite-radial-data-compatibility}
Let $\nu$ be a finite positive Radon measure on $(0,\infty)$,
of total mass $M$.  It is flux-compatible if and only if
\[
                         \nu\ge H_\nu(t)\,dt.
 \tag{RF13}
\]
Equivalently, $e^tH_\nu(t)$ is nonincreasing; no further
finite-energy condition is needed in this finite-mass case.
The solution \textup{(RF7)} has bounded initial slope
$A_\nu(0+)=M^{1/q}$, belongs to every finite positive
power-energy class, and satisfies
\[
             e_\alpha(u_\nu)
                 \le M^{1+\alpha/q}\Gamma(\alpha+1)
                 \qquad(\alpha>0).
 \tag{RF14}
\]
The same bound holds for its bounded radial truncations.
\end{corollary}

\noindent\textbf{Status: \texttt{DIRECT\_COROLLARY}} of
the exact source theorem, with an explicit datum criterion.

\begin{proof}
Distributional differentiation gives
\[
                 d(e^tH_\nu)=e^t(H_\nu\,dt-\nu).
\]
Thus \textup{(RF13)} is exactly monotonicity of the
right-continuous flux.  Since $H_\nu(0+)=M$,
compatibility gives $H_\nu(t)\le Me^{-t}$ and
$A_\nu(t)\le M^{1/q}$.  This proves the energy
integrability and the initial-slope assertion.  The
bound $F_{A_\nu}(t)\le M^{1/q}t$ and layer cake yield
\[
 \int F_{A_\nu}(t)^\alpha\,d\nu(t)
 \le M^{\alpha/q}\alpha\int_0^\infty
                              t^{\alpha-1}H_\nu(t)dt
 \le M^{1+\alpha/q}\Gamma(\alpha+1).
\]
The no-loss proof is already part of the theorem.
In \textup{(RF11)} each source tail is moved to the
clipping depth and its potential value becomes the
corresponding smaller clipped value.  Its power
energy is therefore no larger than the original
integral, proving the last assertion.
\end{proof}

\begin{proposition}[Exact discrepancy and stability of the inverse]
\label{prop:r12-radial-flux-stability}
For $A,B\in\mathscr A_q$ one has
\[
 I(u_A,u_B)=\int_0^\infty
               e^{-t}(A-B)(A^q-B^q)\,dt.
 \tag{RF15}
\]
With all norms below taken in $L^d(e^{-t}dt)$,
\[
 \|A-B\|_d^d\le I(u_A,u_B)
   \le q(\|A\|_d+\|B\|_d)^{q-1}\|A-B\|_d^2.
 \tag{RF16}
\]
Consequently strong convergence in this radial class
is exactly convergence of slopes in this weighted
$L^d$ space.  For compatible sources it is equivalently
\[
 \int_0^\infty e^{t/q}
       |H_{\nu_j}(t)^{1/q}-H_\nu(t)^{1/q}|^d\,dt
                                                        \longrightarrow0.
 \tag{RF17}
\]
In particular, vaguely convergent compatible finite
data of uniformly bounded total mass have strongly
convergent solutions, provided the vague limit is
taken on $(0,\infty)$; the conclusion includes loss
of source mass through the excluded outer boundary.
\end{proposition}

\noindent\textbf{Status: \texttt{NEW\_PROVED}} for
the exact radial nonlinear inverse; no general data
continuity outside this fixed-background family is inferred.

\begin{proof}
Apply the original discrepancy identity and then
integrate by parts, with the vanishing endpoint
terms justified as in \textup{(RF12)}.  This gives
\[
 I(u_A,u_B)
  =\int(F_A-F_B)d(\nu_A-\nu_B)
  =\int(A-B)(H_{\nu_A}-H_{\nu_B})dt,
\]
which is \textup{(RF15)}.  For nonnegative scalars
$a,b$ and integer $q\ge1$,
\[
 |a-b|^{q+1}\le(a-b)(a^q-b^q)
                      \le q(a+b)^{q-1}|a-b|^2.
\]
Integrate the first inequality.  Holder with exponents
$d/(q-1)$ and $d/2$, followed by Minkowski, proves
the second; for $q=1$ it is simply an equality.
The energy identity in \textup{(RF6)} makes the
topological equivalence immediate.  Substitution
of \textup{(RF7)} gives \textup{(RF17)}.

For the final statement let the total masses be
bounded by $M$.  Compatibility gives the uniform
tail bound $H_{\nu_j}(t)\le Me^{-t}$, so no mass
can escape to infinity.  Vague convergence therefore
implies tail convergence at every continuity point
of the limit tail, by first restricting to a compact
interval and then using the exponential bound.
The limit is finite and compatible: the distributional
inequality \textup{(RF13)} passes locally to the limit,
and its tails obey the same bound.  Almost-everywhere
convergence of $A_{\nu_j}$, their uniform bound
$M^{1/q}$, and dominated convergence in $L^d(e^{-t}dt)$
prove strong convergence.  Source mass may escape to
zero, which is outside $V$'s ordinary radial domain;
the proof does not assert preservation of total mass
under this vague convergence.
\end{proof}

\subsection{An explicit nonpolar datum with no full-cone source solution}

\begin{theorem}[A positive-capacity atom and its exact variational replacement]
\label{thm:r12-positive-capacity-atom-source-obstruction}
Fix $R>0$, write $E_R=1-e^{-R}$, and let
\[
 \lambda_R=\delta_{\sigma(R)},\qquad
 \ell_R(h)=-h(\sigma(R)),\qquad
 c_R=\left(\frac R{E_R}\right)^{1/q},\qquad
 v_R=c_R\psi_R.
 \tag{RF18}
\]
The datum is a literal positive Radon measure supported
at a point of positive full outer capacity.  It is
dominated by the actual bounded source
$e^R\mu_{\psi_R}$, so $\ell_R\in\mathbb D_{S,\mathrm s}^*$.
Its exact norm and its unique variational maximizer
over the entire completed cone are
\[
       \|\ell_R\|_* =\frac R{E_R^{1/d}},\qquad
       \mathop{\rm argmax}_{a\in\mathscr C}
                  \{\ell_R(a)-e(a)/d\}=v_R.
 \tag{RF19}
\]
Nevertheless there is no $a\in\mathscr C$ satisfying
$\widehat\mu_a=\lambda_R$ on all original tests.
Indeed the maximizer has the different exact source
\[
 \widehat\mu_{v_R}
     =\frac R{E_R}\,
           \sigma_*\bigl(\mathbf1_{(0,R)}e^{-s}ds+e^{-R}\delta_R\bigr).
 \tag{RF20}
\]
Thus positive-capacity support, domination by an actual
bounded source, and all finite power-energy bounds on
the variational solution do not suffice for the
unrestricted source equation.
\end{theorem}

\noindent\textbf{Status: \texttt{NEW\_PROVED}}.  The
unrestricted source-existence claim under only the
listed conditions is \texttt{FALSE}, in every $q\ge1$.
The background remains the single explicit allowed
current \textup{(PO3)}.

\begin{proof}
Positivity of capacity and the literal atom bound
follow from \textup{(PO4)}.  Frozen strong tester
continuity against the fixed actual source proves
strong continuity of $\ell_R$; this is literal
evaluation at a non-null point, not the nonliteral
null-fiber evaluation considered elsewhere.

For arbitrary $a\in\mathscr C$ put
$F_a(t)=-a(\sigma(t))$.  The completed trace theorem
makes $F_a$ nonnegative and concave with $F_a(0)=0$.
Its chord bound on $[0,R]$ gives
\[
 \begin{split}
 B(a,\psi_R^{[q]})
   &=\int_0^R F_a(s)e^{-s}ds+e^{-R}F_a(R)\\
   &\ge\frac{F_a(R)}R
             \left(\int_0^R se^{-s}ds+Re^{-R}\right)
     =\frac{E_R}{R}F_a(R).
 \end{split}
 \tag{RF21}
\]
Mixed Holder bounds the left side by $r(a)E_R^{q/d}$.
Consequently $\ell_R(a)\le R E_R^{-1/d}r(a)$,
with equality at a positive multiple of $\psi_R$.
The cone/difference dual norm identity proves the
first assertion of \textup{(RF19)}.  Optimizing
$cr\mapsto\ell_R(c\psi_R)-e(c\psi_R)/d
=cR-c^dE_R/d$ gives $c=c_R$, and the same norm
bound proves this is a global maximizer on the full
cone.  Frozen strict midpoint energy separation
proves uniqueness.  No radiality assumption is
imposed on the competitors in this argument.

Homogeneity of the original source gives
\textup{(RF20)}.  Its diffuse positive measure on
every compact radial subinterval of $(0,R)$ makes
it different from $\lambda_R$.  This difference is
visible on original compact smooth scalar tests:
choose a nonnegative smooth compact real function
$\chi$ supported in $(-R,0)$ and put, on the
coordinate-unit torus,
$\eta=\prod_{i=1}^n\chi(\log|T_i|)$, extended
by zero near all coordinate zero loci.  Its support
is an ordinary compact polyannulus in $V$,
$\eta(\sigma(R))=0$, and its integral against
\textup{(RF20)} is strictly positive when $\chi$
is not identically zero.  Thus the discrepancy is
not merely a failure on auxiliary or delta tests.

If an $a\in\mathscr C$ solved the source equation,
then $J(a)=\ell_R$ by literal integration on the
energy-difference space.  The differentiable energy
support inequality makes such an $a$ a variational
maximizer.  Uniqueness would force $a=v_R$, in
contradiction with \textup{(RF20)} and the preceding
original test.  This proves full-cone nonexistence.
\end{proof}

\begin{corollary}[The normal residual is exactly a concavity defect]
\label{cor:r12-radial-atom-normal-concavity}
For the datum and maximizer in \textup{(RF18)}, the
normal functional $n_R=\ell_R-J(v_R)$ satisfies,
for every $a\in\mathscr C$,
\[
 n_R(a)=-\frac R{E_R}\int_0^R
          \left(F_a(s)-\frac{s}{R}F_a(R)\right)e^{-s}ds\le0.
 \tag{RF22}
\]
It vanishes precisely on those cone members whose
actual radial trace is linear on $[0,R]$.
In particular $n_R(v_R)=0$, but $n_R$ is nonzero
as an original-test functional.
\end{corollary}

\noindent\textbf{Status: \texttt{NEW\_PROVED}} exact
geometric description of one full-cone normal face.

\begin{proof}
Use \textup{(PO4)} and $c_R^q=R/E_R$ to expand
$J(v_R)(a)$.  Since
$\int_0^Rse^{-s}ds=E_R-Re^{-R}$, the result is
\textup{(RF22)}.  The integrand is nonnegative by
concavity and continuous, so its weighted integral
vanishes exactly when it vanishes identically.
That is precisely the asserted linearity.  The
original test constructed above gives a nonzero
value of the residual, so its negativity on the
potential cone is not an equality on the original
smooth test sheaf.
\end{proof}

\paragraph{What this advances, and what it does not.}
The flux condition gives an explicit compatible
nonlinear source class, a literal actual solution,
full-cone uniqueness, and a quantified strong inverse
in every remaining degree.  The atom computation
identifies a nonzero normal residual as a concrete
concavity defect, rather than assuming it vanishes
because the datum is positive or capacity controlled.
It is a full-dimensional stress test for the general
source problem, not a universal classification of
arbitrary-$S$ sources and not a proof of OPEN-GQ or
OPEN-CTO.  No auxiliary analytic field, change of
absolute value, or restriction to toric test forms
is used in any displayed current equality.

% END REVIEWED MODULE research_r12/nonlinear/radial_flux_source_inversion.tex


% BEGIN REVIEWED MODULE research_r12/nonlinear/smooth_density_source_instability.tex
\section{A one-sided source threshold under smooth positive density perturbations}
\label{sec:r12-smooth-density-source-instability}

Keep the explicit original background \textup{(PO3)} in
dimension $n>q\ge1$, and put $d=q+1$.  Fix $R>0$ and abbreviate
\[
 \begin{gathered}
 E=1-e^{-R},\\
 \nu_0=\mathbf1_{(0,R)}e^{-t}dt+e^{-R}\delta_R,
       \qquad \mu_0=\sigma_*\nu_0=\mu_{\psi_R},\\
             \pi(dt)=\frac{t}{E}\nu_0(dt).
 \end{gathered}
 \tag{SD1}
\]
The last measure is a probability, since
$\int t\,d\nu_0=E$.  Both its diffuse part and its atom
are retained.  All solution and maximization statements
below use the entire strong classical-approximation cone
$\mathscr C$ for this fixed allowed background.

\begin{corollary}[An elementary curvature background already suffices]
\label{cor:r12-elementary-codimension-one-background}
In every remaining degree $q\ge1$, one may take
$n=q+1$ in all the results of this section.  The
background \textup{(PO3)} then has the elementary
original-current description
\[
                  S=Dv,\qquad
                  v=\max_i|T_i|-1\quad\text{on }V.
 \tag{SD1a}
\]
It extends as the curvature of a bounded continuous
actual classical psh function on the larger original
open $W$ used in \textup{(PO1)}.  Thus the source
obstructions already occur over the curvature of
one explicitly given continuous potential.
\end{corollary}

\noindent\textbf{Status: \texttt{DIRECT\_COROLLARY}}
of the defining radial integral, with its full
original-current identity made explicit.

\begin{proof}
Put $z=\log\max_i|T_i|$ and define on $W$
\[
 \widetilde v=\Theta(z),\qquad
 \Theta(z)=
 \begin{cases}
 e^z-1,&z\le0,\\
 z,&z\ge0,
 \end{cases}
 \qquad\Theta(-\infty)=-1.
\]
The scalar function is continuous, convex and
nondecreasing.  The original convex radial calculus
therefore makes $\widetilde v$ actual bounded
continuous classical psh, including at the rigid
origin.  Direct integration gives
\[
              \Theta(z)=\int_0^\infty
                               e^{-s}\max(z,-s)\,ds.
\]
On compact subsets of $W$, the absolute integrands
are bounded by $e^{-s}(s+C)$, so the integral
converges uniformly.  Distributional $D$ commutes
with it on every original compact complementary
test by coefficient Fubini.  For $n-q=1$ the
result is exactly \textup{(PO3)}.  On $V$ one
has $z<0$, so $\widetilde v=e^z-1=\max_i|T_i|-1$.
This proves \textup{(SD1a)} and its required
neighborhood extension without a new field or norm.
\end{proof}

\begin{theorem}[Nondecreasing density tilts have an exact full-cone normal]
\label{thm:r12-monotone-density-normal}
Let $f:(0,R]\to[0,\infty)$ be bounded, Borel and
nondecreasing.  Define the literal Radon datum and scalar
\[
 \mu_f=\sigma_*(f\nu_0),\qquad
       K_f=\int f\,d\pi=\frac1E\int_{(0,R]} t f(t)\,d\nu_0(t).
 \tag{SD2}
\]
Then $\ell_f(h)=-\int h\,d\mu_f$ belongs to
$\mathbb D_{S,\mathrm s}^*$.  Its unique variational
maximizer over the full cone is
\[
              u_f=K_f^{1/q}\psi_R,
           \qquad\widehat\mu_{u_f}=K_f\mu_0,
           \qquad\|\ell_f\|_*=K_f E^{q/d}.
 \tag{SD3}
\]
For every $a\in\mathscr C$, write
$F_a(t)=-a(\sigma(t))$ and $G_a(t)=F_a(t)/t$.
The normal residual has the exact covariance formula
\[
 \begin{split}
 n_f(a)&:=\ell_f(a)-J(u_f)(a)\\
   &=E\left(\int fG_a\,d\pi
                   -\int f\,d\pi\int G_a\,d\pi\right)\le0.
 \end{split}
 \tag{SD4}
\]
Its value on $u_f$ is zero.  If $f$ is strictly increasing,
then $n_f(a)=0$ exactly when $F_a$ is linear on $[0,R]$.

The full original source equation
$\widehat\mu_a=\mu_f$ has a solution $a\in\mathscr C$
if and only if $f$ is constant $\nu_0$-almost everywhere.
Thus every nonconstant nondecreasing tilt is incompatible,
even though its datum is dominated by the actual bounded
source $\|f\|_\infty\mu_0$.
\end{theorem}

\noindent\textbf{Status: \texttt{NEW\_PROVED}} for this
full-cone fixed-background family.  The source obstruction
holds in every remaining degree, not only in degree one.

\begin{proof}
Domination by $\|f\|_\infty\mu_0$, mixed Holder, and
the frozen strong tester estimate give the literal finite
energy norm and strong continuity.  The completed radial
trace theorem gives finite nonnegative concave $F_a$ with
$F_a(0)=0$, so $G_a$ is nonnegative and nonincreasing.
It is integrable for $\pi$, because
\[
                E\int G_a\,d\pi
                       =\int F_a\,d\nu_0=B(a,\psi_R^{[q]})<\infty.
\]
The boundedness of $f$ makes every covariance integral
finite.  The elementary double-integral identity is
\[
 \operatorname{Cov}_\pi(f,G_a)
   =\frac12\iint(f(t)-f(s))(G_a(t)-G_a(s))\,d\pi(t)d\pi(s).
 \tag{SD5}
\]
Its integrand is nonpositive by the opposite monotonicities.
This proves \textup{(SD4)}.  At $a=\psi_R$ one has
$G_a\equiv1$ on $(0,R]$, so equality holds.

In particular
$\ell_f(a)\le K_f B(a,\psi_R^{[q]})
\le K_f E^{q/d}r(a)$, with equality on the ray
of $\psi_R$.  This proves the exact norm in
\textup{(SD3)} by the cone/difference dual norm identity.
The functional $n_f=\ell_f-K_fJ(\psi_R)$ is nonpositive
on $\mathscr C$ and annihilates $K_f^{1/q}\psi_R$.
The convex energy support inequality therefore proves
that this potential is a global maximizer.  Frozen
strict midpoint separation gives uniqueness.  Its
source follows from homogeneous original BT scaling.
If $K_f=0$, nonnegativity makes $f=0$ almost everywhere
for $\pi$, hence also for $\nu_0$; all assertions then
reduce to the zero potential and zero datum.

If $f$ is strictly increasing, equality in \textup{(SD5)}
forces $G_a$ to be constant $\pi$-almost everywhere.
Indeed any two positive-measure sets on which its values
are separated would contribute a strictly negative
double integral.  The measure $\pi$ has positive density
on $(0,R)$, and $G_a$ is continuous there and at $R$.
Hence it is constant throughout $(0,R]$, which is
exactly linearity of $F_a$ on $[0,R]$.  The converse
is immediate.

Any source solution is the variational maximizer by
literal integration and the energy support inequality.
It must therefore be $u_f$.  The source in
\textup{(SD3)} equals the prescribed datum exactly
when $f=K_f$ almost everywhere for $\nu_0$.
This proves the equivalence, including its converse.
All competing potentials in this argument were arbitrary
members of the full cone; no reduction of the competitors
to radial potentials has occurred.
\end{proof}

\begin{proposition}[The opposite monotone tilts are compatible]
\label{prop:r12-decreasing-density-solutions}
If $f:(0,R]\to[0,\infty)$ is bounded, Borel and
nonincreasing, then $\mu_f=\sigma_*(f\nu_0)$ has a
unique source solution in the full cone.  It is the
actual bounded radial potential of \textup{(RF7)} with
\[
 \begin{aligned}
 H_f(t)&=\int_{(t,R)}f(s)e^{-s}ds+f(R)e^{-R}
                                  &&(0<t<R),\\
 H_f(t)&=0 &&(t\ge R),\\
                  A_f(t)&=(e^tH_f(t))^{1/q}.
 \end{aligned}
 \tag{SD6}
\]
It has ordinary zero boundary value and every finite
positive power energy.
\end{proposition}

\noindent\textbf{Status: \texttt{DIRECT\_COROLLARY}}
of the explicit flux criterion, not an unrestricted
positive-density existence theorem.

\begin{proof}
For $0<t<R$, monotonicity gives
$H_f(t)\le f(t)e^{-t}$, apart from immaterial
Lebesgue-null changes in $f(t)$ at its jump points.
Thus the diffuse part of $f\nu_0$ dominates
$H_f(t)dt$; its atom at $R$ is nonnegative.
The compatibility condition \textup{(RF13)} holds.
The exact inverse gives the actual solution and its
full-cone uniqueness.  Its slope is zero after $R$
and bounded by the $q$th root of the finite total
source mass, so the potential is bounded.  The
all-power and boundary assertions follow from the
finite-data part of the inverse theorem.
\end{proof}

\begin{theorem}[A sharp one-sided threshold for original smooth densities]
\label{thm:r12-smooth-density-one-sided-threshold}
There is a bounded nonnegative original smooth scalar
$G$ on $V$ with
\[
                 G(\sigma(t))=t\quad(0<t\le R),
                        \qquad 0\le G\le C_R<\infty,
 \tag{SD7}
\]
such that the family of positive original Radon data
\[
                       \mu_\epsilon=(1+\epsilon G)\mu_0,
                   \qquad -C_R^{-1}<\epsilon<\infty,
 \tag{SD8}
\]
has the following exact source behavior:
\begin{enumerate}
\item If $\epsilon>0$, there is no source potential in
the full completed cone, although $1\le1+\epsilon G
\le1+\epsilon C_R$ and its variational maximizer is
the bounded actual potential
\[
 u_\epsilon=(1+\epsilon m_R)^{1/q}\psi_R,
         \qquad m_R=2-\frac{2R}{e^R-1}.
 \tag{SD9}
\]
\item If $\epsilon=0$, the source solution is $\psi_R$.
\item If $\epsilon=-a<0$, there is a unique source
solution in the full cone.  It is the actual bounded
radial potential with slope
\[
 A_{-a}(t)=
 \begin{cases}
 \bigl[1-a(t+1-e^{t-R})\bigr]^{1/q},&0<t<R,\\
 0,&t\ge R.
 \end{cases}
 \tag{SD10}
\]
\end{enumerate}
Every density in \textup{(SD8)} is bounded above and
bounded below by a positive constant.  Its datum has
exactly the same Borel null ideal as the actual bounded
source $\mu_0$.  The densities approach $1$ in every
original compact smooth test seminorm as
$\epsilon\to0$, and their literal dual functionals
approach $J(\psi_R)$ in the signed operator dual norm.
Consequently source solvability is not open, even
along a one-parameter family of arbitrarily small
strictly positive original smooth density perturbations.
\end{theorem}

\noindent\textbf{Status: \texttt{NEW\_PROVED}} exact
threshold and \texttt{FALSE} universal smooth-density
openness assertion, for this explicit allowed background
in every $q\ge1$.  No assertion about an unrelated
special smooth background is substituted for it.

\begin{proof}
Choose a smooth nonnegative function on the real line
which equals $1$ on $[0,R]$, vanishes on $[R+1,\infty)$,
and lies in $[0,1]$ on $[R,R+1]$.  Integrating it
from zero produces a smooth function $F$ on $(0,\infty)$
which equals $t$ on $(0,R]$ and is a nonnegative
constant for $t\ge R+1$, bounded by $R+1$.
On $T_1\ne0$ put
\[
                         G=F(-\log|T_1|).
 \tag{SD11}
\]
Extend it by the same constant on $T_1=0$.  It is
original smooth there, because it is already constant
on the ordinary neighborhood $|T_1|<e^{-(R+1)}$.
Elsewhere it is a smooth function of one original
unit moment.  This proves \textup{(SD7)} with,
for example, $C_R=R+1$.  The densities are genuine
original smooth functions, not merely smooth
functions on the diagonal Gauss ray.

On the support of $\mu_0$ they are exactly
$f_\epsilon(t)=1+\epsilon t$.  For $\epsilon>0$
the first theorem applies, and
\[
 K_{f_\epsilon}=1+\epsilon\int t\,d\pi
 =1+\frac{\epsilon}{E}
       \left(\int_0^R t^2e^{-t}dt+R^2e^{-R}\right)
 =1+\epsilon m_R.
\]
The value $m_R$ lies strictly between $0$ and $R$,
as the mean of the nonconstant coordinate for the
probability $\pi$.  The nonconstant density is
therefore incompatible by \textup{(SD3)--(SD4)}.
The difference from its variational source is visible
on original compact tests supported in a diagonal
polyannular interval not containing $m_R$:
take $\eta=\prod_i\chi(\log|T_i|)$ as in the
proof following \textup{(RF20)}.  Its pairing with
$\mu_\epsilon-(1+\epsilon m_R)\mu_0$ has a strict
sign.  Thus the nonexistence statement is on the
full original test sheaf, not just on formal tails.

For $\epsilon=-a<0$, the density is nonincreasing
and strictly positive on the support.  Its tail is
\[
 H_{-a}(t)=e^{-t}
               \bigl[1-a(t+1-e^{t-R})\bigr]\quad(0<t<R),
\]
and zero thereafter.  This proves \textup{(SD10)}.
The bracket is positive since $aC_R<1$ and $t\le R$;
its derivative is $-a(1-e^{t-R})\le0$.
Thus the exact admissible-slope construction applies,
as also follows from the preceding proposition.

The uniform positive density bounds give equality
of the Borel null ideals.  Since $G$ is fixed
original smooth, multiplication by $\epsilon$
makes every compact smooth seminorm of
$\epsilon G$ tend to zero.  The signed dual norm
estimate follows directly from domination by
$|\epsilon|C_R\mu_0$ and mixed Holder.  In fact
the exact norm, by \textup{(SD3)} applied to $f=t$,
is
\[
       \|\ell_{\mu_\epsilon}-J(\psi_R)\|_*
                         =|\epsilon|m_R E^{q/d}.
 \tag{SD12}
\]
The entire family retains the same original field,
norm, space, background, domain and test category.
Only the prescribed positive density changes.
\end{proof}

\paragraph{A quantitative nonzero residual.}
For $0<s<R$ put $E_s=1-e^{-s}$ and
\[
 \kappa_{R,s}
   =m_R E_s-\bigl(2-(s+2)e^{-s}-s e^{-R}\bigr)>0.
 \tag{SD13}
\]
The strict inequality follows either from
\textup{(SD5)} with $f(t)=t$ and
$G_{\psi_s}(t)=\min(t,s)/t$, or from the exact
integral.  In fact
\[
 n_\epsilon(\psi_s)=-\epsilon\kappa_{R,s},\qquad
 \|n_\epsilon\|_*
             \ge\frac{\epsilon\kappa_{R,s}}{E_s^{1/d}}
                       \quad(\epsilon>0).
 \tag{SD14}
\]
The signed source-image distance bound \textup{(NI12)}
therefore gives a strictly positive, explicit distance
for each positive $\epsilon$, even though that distance
can tend to zero with the perturbation.  These formulas
are exact auxiliary checks of the normal mechanism;
they are not a numerical replacement for its proof.

\paragraph{Interpretation for the general problem.}
The failure is not caused by charging a relative-null
set, losing finite energy, changing the boundary, or
removing positivity of the density.  The actual source
cone is a proper constrained image even under all of
those controls.  The opposite sign of the same smooth
perturbation satisfies an explicit flux condition and
does have a unique actual source.  Any general
arbitrary-background Dirichlet theory must therefore
retain a genuine source-compatibility condition or
change its explicitly defined admissible potential
class.  This example does neither silently, and it
does not address the unresolved analytic PS/PL quotient.

% END REVIEWED MODULE research_r12/nonlinear/smooth_density_source_instability.tex


% BEGIN REVIEWED MODULE research_r12/energy/parameter_finite_energy_dirichlet.tex
\section{A complete parameter module of finite-energy mixed-slot solutions}
\label{sec:r12-parameter-energy-Dirichlet}

The purpose of this section is to give a positive source theorem
with a genuinely explicit admissibility condition. It is a theorem
for a new relative parameter module, not a solution of the
unrestricted ambient classical psh Dirichlet problem.

Assume the strict-box geometry (M2). Let $s=(\log|f_1|,\ldots,
\log|f_c|)$ be original unit parameters defined on an ordinary
neighborhood of the closed box, and let $R$ be a symmetric
middle-positive separately closed original current on that
neighborhood, of degree $(P,P)$ with $k=n-P\geq1$, satisfying
\[
                         d's_i\wedge R=d''s_i\wedge R=0.
 \tag{PD1}
\]
In particular, for \emph{every} original allowed $S$ with
$n-p>c$, one may take $R=\Omega_s\wedge S$. The results then
concern the mixed-slot operator with the fixed factor $\Omega_s$;
they do not replace the given $S$ by a special assumed background.

Let $\psi$ be the geometry-dependent potential of the universal
probe construction. We use its actual bounded continuous
representative, normalized so that
\[
 \rho\leq\psi<0,\qquad
 \omega=(D\psi)^k\wedge R,\qquad \omega(V)<\infty.
 \tag{PD2}
\]
Put $d=k+1$, $Q_j=T_{R,j}(\psi^{[j]})$, $Q_0=R$, and define
finite positive parameter measures
\[
                      \lambda=s_*\omega,\qquad
                      \theta=s_*\bigl((-\psi)\omega\bigr).
 \tag{PD3}
\]
They have the same Borel null sets because $-\psi$ is strictly
positive everywhere on $V$. The measure $\theta$, not an
unproved analytic structure on a fiber, supplies the parameter
energy domain.

\begin{lemma}[The probe dominates every fixed tail coefficient]
\label{lem:r12-parameter-tail-domination}
For every $0\leq j\leq k$ and every original compact smooth
complementary coefficient test $\eta$, there is a finite
$C_\eta$ such that
\[
                         |\mu_{Q_j,\eta}|\leq C_\eta\,\omega.
 \tag{PD4}
\]
The inequality is an inequality of Radon measures on the
coefficient support. Every actual refinement of $\eta$ is
allowed. If $\theta(E)=0$ for a Borel parameter set $E$, then
$R_E=\mathbf1_{s^{-1}(E)}R=0$ and, on every ordinary
$W\Subset V$,
\[
                     C_R^*(W\cap s^{-1}(E);V)=0.
 \tag{PD5}
\]
\end{lemma}

\noindent\textbf{Status: \texttt{DIRECT\_COROLLARY}} of the
geometry-only probe construction, positivity of bounded mixed
expansions, the closed-band theorem and (CB7).

\begin{proof}
On every core in the probe construction there is a decomposition
$\psi=c_l\psi_l+r_l$, with $c_l>0$, both $\psi_l,r_l$ actual
classical psh, and $\psi_l$ equal near that core to a smooth
analytic compensator whose curvature dominates the relevant
original moment trace. For the positive current $Q_j$, finite
middle-coefficient domination gives
\[
 |\mu_{Q_j,\eta}|
       \leq C(D\psi_l)^{k-j}\wedge Q_j
       \leq Cc_l^{-(k-j)}(D\psi)^{k-j}\wedge Q_j
       = Cc_l^{-(k-j)}\omega
\]
there. The second inequality is the positive bounded mixed
expansion of $(c_lD\psi_l+Dr_l)^{k-j}$ over $Q_j$.
A finite cover of the support proves (PD4). At $j=k$ it
is simply the scalar bound for a top measure. No product
with a nonclosed nonsmooth current occurs: every $Q_j$ is
already a positive closed bounded BT tail.

If $\theta(E)=0$, strict negativity of $\psi$ implies
$\omega(s^{-1}(E))=0$. Apply (PD4) at $j=0$ to every
original compact coefficient. This gives $R_E=0$.
For every original normalized classical tester $a$, the
closed-band identity then gives
$\mu_a^R(W\cap s^{-1}(E))=T_{R_E,k}(a^{[k]})(W)=0$.
The full-tester raw capacity of this band is zero.
The proved band raw/outer equality gives (PD5).
\end{proof}

\begin{corollary}[Exact null detection on the parameter sigma algebra]
\label{cor:r12-parameter-exact-null-ideal}
For every Borel $E\subset J$, the following conditions are equivalent:
\[
 \begin{gathered}
 \theta(E)=0,\qquad R_E=0,\\
 C_R^*(W\cap s^{-1}(E);V)=0
       \quad\text{for every ordinary }W\Subset V.
 \end{gathered}
 \tag{PD5a}
\]
This identifies null ideals only on parameter-saturated sets.
It does not identify coefficient-null, source-null and capacity-null
ideals on arbitrary analytic Borel sets.
\end{corollary}

\noindent\textbf{Status: \texttt{DIRECT\_COROLLARY}}.

\begin{proof}
The forward implications are (PD5). If $R_E=0$, the closed-band
identity gives $\mathbf1_{s^{-1}(E)}\omega=T_{R_E,k}(\psi^{[k]})=0$,
hence $\theta(E)=0$. Finally, if all the displayed local outer
capacities vanish, the bounded actual tester $\psi+1\in[0,1]$
shows that $\omega(W\cap s^{-1}(E))=0$ on each such window.
A countable compact exhaustion of the fixed strict box gives
$\omega(s^{-1}(E))=0$, and therefore $\theta(E)=0$.
\end{proof}

\begin{definition}[The parameter first-energy cone]
For a nonnegative Borel function $c$ with
$c\in L^d(\theta)$, choose a finite version outside a
$\theta$-null set, give it the value zero on that null set,
and put
\[
                           h_c=\psi\,c(s).
 \tag{PD6}
\]
Let $\mathscr P^1_{s,\psi,R}$ be these actual Borel functions
modulo local ordinary-open outer $R$-capacity zero.
A \emph{bounded parameter potential} means that $c$ is bounded.
Such a function has ordinary boundary value zero. For a general
$c\in L^d(\theta)$ the boundary condition means approximation
by the bounded parameter potentials $h_{\min(c,N)}$ in
the norm defined below; no ordinary boundary value is inferred.
\end{definition}

\begin{theorem}[All-degree products and a complete relative energy module]
\label{thm:r12-parameter-Lp-BT}
The map $c\mapsto h_c$ identifies $L^d_+(\theta)$ injectively
with $\mathscr P^1_{s,\psi,R}$; it respects addition and
nonnegative scaling. For $c_1,\ldots,c_j\in L^d_+(\theta)$
and $0\leq j\leq k$, the intrinsic original recursion is
well-defined and has the exact formula
\[
 T^{s,\psi}_{R,j}(h_{c_1},\ldots,h_{c_j})
       =\left(\prod_{i=1}^j c_i(s)\right)Q_j.
 \tag{PD7}
\]
It is middle-positive, separately closed, mixed symmetric,
and conormal. Every coefficient is Radon and every equality
holds on all original compact complementary tests and
all actual moment refinements.

The global mixed energy and energy root are
\[
 \begin{split}
 B_{s,\psi}(h_{c_0},\ldots,h_{c_k})
       &=\int_V(-h_{c_0})\,
                 T^{s,\psi}_{R,k}(h_{c_1},\ldots,h_{c_k})\\
       &=\int_J c_0\cdots c_k\,d\theta,\\
 e_{s,\psi}(h_c)&=\|c\|_{L^d(\theta)}^d,\qquad
 r_{s,\psi}(h_c)=\|c\|_{L^d(\theta)} .
 \end{split}
 \tag{PD8}
\]
In particular the parameter first-energy cone is complete
for the induced $L^d$ difference metric and has the
constant-one mixed energy inequality.

For $c,a\geq0$ its strong discrepancy is
\[
 I_{s,\psi}(h_c,h_a)
       =\int_J(c-a)(c^k-a^k)\,d\theta.
 \tag{PD9}
\]
Writing $R_0=\|c\|_d+\|a\|_d$, one has
\[
 \|c-a\|_d^d\leq I_{s,\psi}(h_c,h_a)
       \leq kR_0^{k-1}\|c-a\|_d^2.
 \tag{PD10}
\]
The right side is interpreted as $\|c-a\|_2^2$ when $k=1$.
Thus the strong and $L^d$ topologies coincide on this new
cone. This does not assert that they coincide on the
unchanged ambient classical strong cone $\mathscr C$.
\end{theorem}

\noindent\textbf{Status: \texttt{NEW\_PROVED}} for the
explicit relative parameter module, with its complete
all-degree calculus and energy. The bounded inputs are
the closed-band theorem and the frozen original BT
recursion. This is not an appeal to OPEN-GQ.

\begin{proof}
First let the coefficients be bounded nonnegative Borel
functions. Every $h_c$ is an ordinary psh level trace:
on $s=a$ it equals the classical potential $c(a)\psi$.
Alternatively it is the explicitly Borel slice-psh family
$H(a,x)=c(a)\psi(x)$ with reference map the identity.
The closed-band identity applied at every insertion gives
(PD7), either directly by induction or through the proved
Borel-family calculus. Separate closedness holds because
the product of a bounded Borel parameter coefficient with
a closed conormal current is closed.

For unbounded $c_i\in L^d(\theta)$, truncate them at $N$.
On a fixed compact coefficient support $H\Subset V$,
continuity and strict negativity of $\psi$ give
$-\psi\geq\epsilon_H>0$. Hence (PD4) implies, for every
nonnegative Borel parameter function $b$,
\[
 \int_H b(s)\,d|\mu_{Q_j,\eta}|
       \leq \frac{C_\eta}{\epsilon_H}\int_J b\,d\theta.
 \tag{PD11}
\]
Products of $j\leq k<d$ members of $L^d(\theta)$ belong
to $L^{d/j}(\theta)\subset L^1(\theta)$, since $\theta$
is finite. The same holds for differences of products
along $L^d$-convergent rows, by the telescoping product
identity and Holder's inequality. Therefore the currents
in (PD7) for the bounded truncations converge in each
original local coefficient total variation to the
displayed right side.

Separate closedness, symmetry and conormality pass to
these original-current limits. To justify the intrinsic
recursion, its multiplier before the $j$th derivative
is exactly
\[
             \psi\left(\prod_{i=1}^j c_i(s)\right)Q_{j-1}.
\]
It has locally integrable coefficients by (PD11) and
boundedness of $\psi$. The same products for truncated
coefficients converge in local coefficient variation.
Distributional $D$ passes to the limit and gives (PD7).
Thus no product of arbitrary unbounded distributions has
been defined without its integrability check.

A change on a $\theta$-null parameter set changes $h_c$
only on a locally outer-null band by (PD5), and changes
none of these coefficient currents. Conversely, suppose
$h_c=h_a$ outside a locally outer-null set. The bounded
actual source $\omega$ ignores such sets by the original
capacity estimate. Since $\psi<0$ everywhere, it follows
that $\lambda(\{c\ne a\})=0$, hence
$\theta(\{c\ne a\})=0$. A countable compact exhaustion
of the fixed box justifies passing from local nullity
to this global measure statement. This proves injectivity.

Formula (PD7) at top degree, followed by (PD3), gives
the literal mixed integral (PD8); Holder proves its
finiteness, symmetry and constant-one inequality.
Completeness is that of $L^d_+(\theta)$, and the bounded
truncations converge in this norm. For bounded $c$,
$|h_c|\leq\|c\|_\infty(-\psi)$ shows its ordinary
zero boundary value and compact negative sublevels.
No such pointwise boundary conclusion has been used
for unbounded $c$.

For nonnegative numbers $x,y$ one has
\[
 (x-y)(x^k-y^k)\geq |x-y|^{k+1},
\]
and
\[
 (x-y)(x^k-y^k)\leq
                  k(x+y)^{k-1}|x-y|^2.
\]
Integrate the first inequality. For $k>1$, apply Holder
to the second with exponents $d/(k-1)$ and $d/2$, whose
reciprocals sum to one. For $k=1$ it is equality.
This gives (PD10). Formula (PD9) itself is the expansion
of (PD8) into the four discrepancy terms. All these
calculations occur on the proved parameter module;
they do not identify its members with ambient classical
psh representatives.
\end{proof}

\begin{proposition}[Comparison, domination, and explicit parameter obstacles]
\label{prop:r12-parameter-comparison-domination}
For any $c,a\in L^d_+(\theta)$ the module satisfies comparison:
\[
 \int_{\{h_c<h_a\}} T^{s,\psi}_{R,k}(h_a^{[k]})
       \leq
 \int_{\{h_c<h_a\}} T^{s,\psi}_{R,k}(h_c^{[k]}).
 \tag{PD10a}
\]
The integrals are allowed to be infinite. If
$h_c\geq h_a$ almost everywhere for
$T^{s,\psi}_{R,k}(h_c^{[k]})$, then $h_c\geq h_a$
outside a locally outer $R$-capacity-null set.

For a signed $a\in L^d(\theta)$, the greatest member of this
module below the actual Borel obstacle $\psi a(s)$, in the
relative outer-capacity order, is
$h_{a_+}$. For a Borel parameter set $E$, the greatest module
potential satisfying $\psi\leq h\leq0$ and $h\leq\psi$ on
$s^{-1}(E)$, with inequalities in the relative outer-capacity
order, is
\[
             h_E=\psi\mathbf1_E(s),\qquad
             T^{s,\psi}_{R,k}(h_E^{[k]})=\mathbf1_E(s)\omega,
             \qquad e_{s,\psi}(h_E)=\theta(E).
 \tag{PD10b}
\]
These are comparison, domination and obstacle statements only
within the specified parameter module. They are not statements
about arbitrary ambient psh envelopes or all level traces.
\end{proposition}

\noindent\textbf{Status: \texttt{NEW\_PROVED}} for the explicit
module, as consequences of its intrinsic source formula and
the exact saturated null ideal.

\begin{proof}
Strict negativity of $\psi$ identifies
$\{h_c<h_a\}$ with $s^{-1}(\{c>a\})$. On this set
$a(s)^k\leq c(s)^k$, so (PD7) proves (PD10a) directly
as an inequality of nonnegative measure integrals.
For domination put $E=\{c>a\}$. The hypothesis says
$\int_{s^{-1}(E)}c(s)^k\,d\omega=0$.
Since $c>0$ on $E$, taking the countable sets $E\cap\{c\geq1/m\}$
gives $\lambda(E)=0$, hence $\theta(E)=0$. Formula (PD5a)
then removes this band in the exact local outer-capacity order.

The order inequality $h_c\leq\psi a(s)$ is equivalent to
$c\geq a$ $\theta$-almost everywhere. Among nonnegative
$c\in L^d(\theta)$ the smallest such coefficient is $a_+$,
which gives the greatest potential. Likewise the constraints
for $h_E$ are exactly $0\leq c\leq1$ and $c\geq1$ on $E$,
in the same measure-algebra order. Their smallest solution is
$\mathbf1_E$. The source and energy formulas now follow from
(PD7)--(PD8). Existence and uniqueness in the relative class
use (PD5a), not an identification with ambient polar sets.
\end{proof}

\begin{theorem}[Explicit compatible data and their unique parameter solution]
\label{thm:r12-parameter-Dirichlet}
Let $b\geq0$ be Borel with
\[
                             b\in L^{d/k}(\theta),
 \qquad \mu_b=b(s)\,\omega.
 \tag{PD12}
\]
Then $\mu_b$ is a locally finite positive original Radon
measure, independent of the chosen finite Borel version.
Its unique source solution within
$\mathscr P^1_{s,\psi,R}$ is
\[
                         h_b=\psi\,b(s)^{1/k},\qquad
                  T^{s,\psi}_{R,k}(h_b^{[k]})=\mu_b.
 \tag{PD13}
\]
Its energy is exactly $\int b^{d/k}\,d\theta$.
For every original compact scalar test $\varphi$,
the second identity in (PD13) is the literal identity
of $\varphi$-pairings; it is not limited to parameter
tests. When $b$ is bounded, $h_b$ has ordinary boundary
value zero. For general (PD12), retain the bounded-data
completion boundary condition.

Equivalently the positive dual of the linear parameter
space $L^d(\theta)$ is precisely $L^{d/k}_+(\theta)$,
with the curvature gradient $c\mapsto c^k$ onto that
cone. For these data the variational functional
\[
             \mathcal F_b(c)=\int cb\,d\theta
                                  -\frac{1}{d}\int c^d\,d\theta,
                         \qquad c\in L^d_+(\theta),
 \tag{PD14}
\]
has its unique maximum at $c=b^{1/k}$, and
\[
                    \max\mathcal F_b=\frac{k}{d}\int b^{d/k}\,d\theta.
 \tag{PD15}
\]
Here there is no source normal residual, because the
specified parameter module has the full positive
$L^d$ cone as its domain.
\end{theorem}

\noindent\textbf{Status: \texttt{NEW\_PROVED}} as an
explicitly compatible mixed-slot Dirichlet family in
every remaining degree and for every arbitrary $S$
through $R=\Omega_s\wedge S$. The nonlinear solution
space is the one just defined. Neither uniqueness
among all level traces nor ambient classical existence
is asserted.

\begin{proof}
Condition (PD12) is exactly
$b^{1/k}\in L^d(\theta)$. On a compact $H$,
$\omega|_H\leq\epsilon_H^{-1}(-\psi)\omega|_H$,
so Holder gives $\mu_b(H)<\infty$. The same domination
gives local Radon regularity, or it follows from the
locally integrable density with respect to $\omega$.
Formula (PD7) proves (PD13) and (PD8) proves its energy.

If another $h_c$ in this module has the same original
top measure, then
$c(s)^k\omega=b(s)\omega$. Their parameter pushforwards
give $c^k=b$ $\lambda$-almost everywhere, hence
$\theta$-almost everywhere. Nonnegativity gives
$c=b^{1/k}$, proving the stated uniqueness.
For (PD14), the scalar Young inequality
$cb-c^d/d\leq(k/d)b^{d/k}$ has equality precisely
at $c=b^{1/k}$. Its integrable bound and equality
case prove (PD15) and uniqueness. This also gives
the displayed dual description by the usual
$L^d$ duality; no energy-dual shadow identification
from a different Banach space is being used.
\end{proof}

\begin{proposition}[All-original stability and local outer quasicontinuity]
\label{prop:r12-parameter-source-stability}
If $b_\nu,b\geq0$ and
$b_\nu\to b$ in $L^{d/k}(\theta)$, their parameter
solutions converge in the complete energy topology
of (PD10). Every mixed original coefficient current
converges locally in total variation when all its
parameter slots converge in $L^d(\theta)$.

Every member of $\mathscr P^1_{s,\psi,R}$ has a finite
representative which is locally ordinary-open outer
$R$-quasicontinuous. These statements neither identify
raw and outer capacity on arbitrary sets nor provide
an ambient classical psh representative.
\end{proposition}

\noindent\textbf{Status: \texttt{DIRECT\_COROLLARY}}
for stability of the explicit solution family, with
a separate elementary quasicontinuity proof.

\begin{proof}
For nonnegative numbers,
$|x^{1/k}-y^{1/k}|^k\leq|x-y|$. Therefore
\[
 \|b_\nu^{1/k}-b^{1/k}\|_{L^d(\theta)}^d
      \leq\|b_\nu-b\|_{L^{d/k}(\theta)}^{d/k}\longrightarrow0.
\]
The discrepancy convergence follows from (PD10).
The coefficient variation statement is the telescoping
product estimate and (PD11) already proved, and includes
each actual refinement of each compact original test.

For quasicontinuity fix $W\Subset V$ ordinary open.
Its uniform parameter CLN control $\nu_W$ from (CB4)
is the pushforward of finitely many original compact
coefficient traces of $R$. By (PD4), each such trace
is dominated by a finite constant times $\omega$.
Consequently $\nu_W\ll\lambda$, and hence
$\nu_W\ll\theta$. The parameter coefficient $c$ is
finite $\nu_W$-almost everywhere. Give it a finite
Borel value on its exceptional null set, which changes
$h_c$ only on a band of local outer capacity zero.

For $\epsilon>0$, Lusin's theorem for the finite Radon
measure $\nu_W$ gives a compact real set $F\subset J$
with $\nu_W(J\setminus F)<\epsilon$ such that
$c|_F$ is finite and continuous. The exceptional set
$G=W\cap s^{-1}(J\setminus F)$ is ordinary open in
$W$, and (CB4) gives
$C_R(G;V)\leq\nu_W(J\setminus F)<\epsilon$.
On $W\setminus G$, the function
$\psi\,c(s)$ is finite and continuous in the relative
ordinary topology. An exhaustion by such windows
gives the claimed local outer quasicontinuity.
This uses the full original normalized tester family
through (CB4), not just the explicit parameter
solutions themselves.
\end{proof}

\paragraph{Exact relation to the final general problem.}
The construction adds a background-independent parameter
module over every conormal mixed-slot current, including
one derived from every allowed arbitrary $S$. Its source
condition (PD12) is a measurable integrability condition
on prescribed data, not an assertion that a source
solution already exists. Its products are intrinsic
original currents, and its energy, boundary completion,
uniqueness class and exceptional-set convention are
specified explicitly. Its members may fail ambient
classical psh even when bounded, so none of these
positive results closes the unrestricted classical
Dirichlet gate or the analytic PS/PL relation quotient.


% END REVIEWED MODULE research_r12/energy/parameter_finite_energy_dirichlet.tex


% BEGIN REVIEWED MODULE research_r12/gq/parameter_energy_capacity.tex
\section{An attained finite-energy capacity for the parameter module}
\label{sec:r12-parameter-energy-capacity}

Keep all hypotheses and notation of
Section~\ref{sec:r12-parameter-energy-Dirichlet}. In particular
$R$ is a conormal original background on the strict original box,
$k=n-P\ge1$, $d=k+1$, and $\psi$ is the bounded continuous
actual geometry-dependent probe with $-1\le\psi<0$.
Write
\[
        \omega=T_{R,k}(\psi^{[k]}),\qquad
        \theta=s_*((-\psi)\omega),\qquad
        h_c=\psi c(s),\qquad c\in L^d_+(\theta).
                                                                    \tag{PC1}
\]
The preceding module theorem proves the literal all-original-test
identities
$\mu_{h_c}=c(s)^k\omega$ and
$e_{s,\psi}(h_c)=\|c\|_d^d$. The present admissible family is
normalized by that energy. It is neither the full classical
$[0,1]$ tester family nor the full band-pasting family.

\begin{definition}[Energy-normalized parameter capacity]
\label{def:r12-parameter-energy-capacity}
For Borel $A\subset V$, set
\[
       \mathfrak c_\psi(A)
          =\sup_{\substack{c\in L^d_+(\theta)\\\|c\|_d\le1}}
                                \mu_{h_c}(A).
                                                                    \tag{PC2}
\]
The value may be infinite. Its ordinary-open outer regularization is
\[
       \mathfrak c_\psi^{\mathrm o}(A)
          =\inf_{\substack{G\supset A\\G\text{ ordinary open in }V}}
                                      \mathfrak c_\psi(G),
                                                                    \tag{PC3}
\]
where (PC3) is defined for arbitrary subsets. The original field,
absolute value, underlying potentials and current test sheaf have
not changed. Only the explicitly displayed scalar normalization
and the explicitly specified parameter module define this capacity.
\end{definition}

For Borel $A$, the finite parameter measure
$\lambda_A=s_*(\mathbf1_A\omega)$ is absolutely continuous with
respect to $\theta$, because the parameter measures
$s_*\omega$ and $\theta$ have the same null sets. Choose a
nonnegative Borel Radon--Nikodym density
\[
                            g_A=\frac{d\lambda_A}{d\theta}.
                                                                    \tag{PC4}
\]
It is specified only as a $\theta$-class. This is a conditional
\emph{measure} density, not a pointwise capacity on an intrinsic
analytic fiber. It retains the mass of the complete original
measure $\omega$ on $A$.

\begin{theorem}[Exact capacity and an attained bounded optimizer]
\label{thm:r12-parameter-energy-capacity-formula}
For every Borel $A\subset V$,
\[
                         \mathfrak c_\psi(A)=\|g_A\|_{L^d(\theta)},
                                                                    \tag{PC5}
\]
with the value $+\infty$ allowed on either side. Suppose
$A\subset W\Subset V$, where $W$ is ordinary open, and put
$\varepsilon_W=\min_{\overline W}(-\psi)>0$. Then
\[
              0\le g_A\le\varepsilon_W^{-1},\qquad
              \mathfrak c_\psi(A)
                    \le\varepsilon_W^{-k/d}\omega(A)^{1/d}.
                                                                    \tag{PC6}
\]
If $\omega(A)>0$, the supremum in (PC2) has the unique
maximizing parameter class
\[
                    c_A=\frac{g_A}{\|g_A\|_d},\qquad
                    h_A=\psi\,c_A(s),\qquad
                    e_{s,\psi}(h_A)=1.
                                                                    \tag{PC7}
\]
This coefficient has a bounded Borel representative, so $h_A$
has ordinary zero boundary value and is one of the already
constructed bounded original-current parameter potentials.
Its entire original source, not just its parameter pushforward,
is
\[
                        \mu_{h_A}
                             =\frac{g_A(s)^k}{\|g_A\|_d^k}\omega.
                                                                    \tag{PC8}
\]
For $\omega(A)=0$ the capacity is zero; uniqueness of a maximizing
coefficient is not asserted in that case.
\end{theorem}

\noindent\textbf{Status: \texttt{NEW\_PROVED}} for this
explicit finite-energy capacity and its original-current optimizer.
Radon--Nikodym and Holder duality are \texttt{EXTERNAL\_KNOWN};
the required actual products and energy are already proved in
the parameter-module theorem.

\begin{proof}
For each admissible coefficient, the exact source formula gives
\[
                       \mu_{h_c}(A)=\int_J c^k g_A\,d\theta.
                                                                    \tag{PC9}
\]
The map $c\mapsto c^k$ sends the nonnegative unit ball of
$L^d(\theta)$ onto that of $L^{d/k}(\theta)$. Since the
conjugate exponent of $d/k$ is $d$, Holder's inequality gives
the upper bound in (PC5) whenever $g_A\in L^d$.
For a nonzero such $g_A$, the coefficient in (PC7) has norm one
and gives equality in (PC9): its value is
$\|g_A\|_d^{-k}\int g_A^{k+1}d\theta=\|g_A\|_d$.
The equality case in Holder, including vanishing on $\{g_A=0\}$,
gives the uniqueness of this nonnegative parameter class.
If $g_A$ is not in $L^d$, let $g_{A,N}=\min(g_A,N)$ and use
$c_N=g_{A,N}/\|g_{A,N}\|_d$ whenever the denominator is nonzero.
Then (PC9) is at least $\|g_{A,N}\|_d\to\infty$.
If $g_A=0$, all values in (PC9) are zero. These cases prove (PC5).

For the local bound, on $W$ one has
$\omega\le\varepsilon_W^{-1}(-\psi)\omega$. Hence
$\lambda_A\le\varepsilon_W^{-1}\theta$, which gives the
pointwise density bound in (PC6). Therefore
\[
             \int g_A^d\,d\theta
                    \le\varepsilon_W^{-(d-1)}\int g_A\,d\theta
                    =\varepsilon_W^{-k}\omega(A).
\]
This proves the remaining bound and the boundedness of the
optimizer. The parameter-module theorem supplies (PC8), its
original-test meaning, and ordinary zero boundary value for a
bounded coefficient. No maximizer in the unrestricted classical
psh tester family is inferred.
\end{proof}

\begin{theorem}[Local full Borel regularity and the exact null ideal]
\label{thm:r12-parameter-energy-capacity-regularity}
Fix any ordinary $W\Subset V$. Then for every Borel $A\subset W$,
\[
                  \mathfrak c_\psi(A)=\mathfrak c_\psi^{\mathrm o}(A).
                                                                    \tag{PC10}
\]
The outer infimum may equivalently be taken over ordinary opens
contained in $W$. On Borel subsets of $W$, the capacity is inner
regular by compact subsets of the original analytic space and
outer regular by ordinary
open subsets. It is continuous on increasing and on decreasing
Borel sequences in $W$. More quantitatively, for Borel $A,B\subset W$,
\[
       |\mathfrak c_\psi(A)-\mathfrak c_\psi(B)|
                 \le\varepsilon_W^{-k/d}\omega(A\mathbin\triangle B)^{1/d}.
                                                                    \tag{PC11}
\]
For every such Borel $A$,
\[
       \mathfrak c_\psi(A)=0
          \quad\Longleftrightarrow\quad\omega(A)=0
          \quad\Longleftrightarrow\quad
                            \mathfrak c_\psi^{\mathrm o}(A)=0.
                                                                    \tag{PC12}
\]
The finite ordinary-open outer extension on $W$ is a Choquet
capacity on all subsets of $W$.
\end{theorem}

\noindent\textbf{Status: \texttt{NEW\_PROVED}} for the
energy-normalized parameter module. This is a full Borel
regularity theorem for the capacity (PC2), not for unrestricted
classical Bedford--Taylor capacity.

\begin{proof}
As a supremum of positive Radon measures, (PC2) is monotone and
countably subadditive. It is finite on $W$ by (PC6). For any
Borel $E\subset W$, finite Radon regularity of $\omega|_W$
provides an ordinary open $G$ in $W$ containing $E$ with
$\omega(G\setminus E)$ arbitrarily small. The set $G$ is also
ordinary open in $V$ because $W$ is open. Subadditivity and
(PC6) give
\[
             \mathfrak c_\psi(G)
               \le\mathfrak c_\psi(E)
                      +\varepsilon_W^{-k/d}\omega(G\setminus E)^{1/d}.
\]
Taking the error to zero proves (PC10). Intersecting an ambient
open neighborhood with $W$ shows equality of the two outer
infima. This argument uses an open window, not an unjustified
compact-window neighborhood replacement.

Subadditivity bounds the difference of the capacities of $A,B$
by the larger of their capacities on the two set differences;
(PC6) gives (PC11). Inner regularity follows by choosing compact
$K\subset A$ with $\omega(A\setminus K)$ small. Increasing
or decreasing Borel sequences in $W$ have symmetric differences
of $\omega$-mass tending to zero, since $\omega(W)<\infty$.
Thus (PC11) proves both continuity assertions. Formula (PC5),
together with $\int g_A\,d\theta=\omega(A)$, proves the first
equivalence in (PC12); (PC10) gives the second.

For completeness, every subset of $W$ has a Borel hull of the
same outer capacity: intersect a sequence of ordinary open
neighborhoods whose capacities approach the infimum and use
the already proved Borel regularity. For increasing arbitrary
sets $A_j$, choose such Borel hulls $B_j$ and set
$H_j=\bigcap_{l\ge j}B_l$. Then $H_j$ increases, contains
$A_j$, and has capacity $\mathfrak c_\psi^{\mathrm o}(A_j)$.
Continuity below on Borel sets proves continuity below of the
outer extension on arbitrary sets. Continuity above on compact
sets follows from (PC11). These, monotonicity and the zero value
on the empty set are the claimed finite Choquet properties.
\end{proof}

\begin{corollary}[Parameter bands and comparison with original capacity]
\label{cor:r12-parameter-energy-capacity-bands}
For an ordinary $W\Subset V$, let $g_W$ denote (PC4) with
$A=W$, and set $\Theta_W=g_W^d\theta$. This is a finite positive
Radon measure on the real parameter space, and for every Borel $E$,
\[
         \mathfrak c_\psi(W\cap s^{-1}(E))=\Theta_W(E)^{1/d}.
                                                                    \tag{PC13}
\]
For every Borel $A\subset W$ one also has
\[
          \mathfrak c_\psi(A)
               \le\varepsilon_W^{-k/d} C_R(A;V)^{1/d}.
                                                                    \tag{PC14}
\]
The same inequality holds for the two ordinary-open outer
regularizations on arbitrary subsets of $W$. Only this direction
is asserted. In particular (PC12) does not identify the full
classical $R$-capacity-null ideal with the $\omega$-null ideal
on arbitrary analytic sets.
\end{corollary}

\noindent\textbf{Status: \texttt{DIRECT\_COROLLARY}}.

\begin{proof}
The parameter density of $\mathbf1_{W\cap s^{-1}(E)}\omega$
is $\mathbf1_Eg_W$. Formula (PC5) gives (PC13), and (PC6)
makes its measure finite. The actual potential $\psi+1$ is
an original normalized classical tester on the same $V$, so
$\omega(A)\le C_R(A;V)$. Combine this with (PC6) to obtain
(PC14). For outer capacities, intersect any proposed ordinary
open neighborhood with $W$, apply the same bound there, and
take the infimum. No equivalence between the two tester
families or their null ideals has been used.
\end{proof}

\paragraph{Background, boundary and test scope.}
Every arbitrary original allowed $S$ yields this mechanism
through the proved conormal mixed-slot current
$R=\Omega_s\wedge S$. The coefficients are the exact
parameter $L^d$ classes of the preceding energy module.
The bounded capacitary optimizer (PC7) has ordinary zero
boundary value; general energy-admissible potentials retain
only the printed bounded-parameter completion boundary convention.
Their products are literal original currents on all original
tests, not currents on a new field or a parameter space alone.

The density $g_A$ in (PC4) averages the fixed original source
measure along parameters. There is no maximization over an
uncountable family of fiberwise psh functions, and no measurable
choice of intrinsic analytic-fiber solutions. This is why the
exact formula and attainment are available in this explicitly
specified module. They do not resolve OPEN-GQ, OPEN-CTO, or
the raw/outer problem for the full original classical capacity.

\paragraph{Exact dependency snapshot.}
The current and energy theorem used here is
\path{energy/parameter_finite_energy_dirichlet.tex}, including
(PD4)--(PD11), at SHA256
{\small\nolinkurl{dd3fbcc3753ad9e3df7145688fa5ff6c1d5e39cb4400c18a5c7a077dfcfdba7c}}.
Its subsequent additions of comparison and exact saturated-null
statements have been independently reviewed; (PC5)--(PC14) need
only the already proved intrinsic source and energy formulas,
strict negativity of the actual probe, and its normalized bounded
classicality.

% END REVIEWED MODULE research_r12/gq/parameter_energy_capacity.tex


% BEGIN REVIEWED MODULE research_r12/energy/countable_measure_capacity_obstruction.tex
\section{Why countably many Radon probes cannot recover full original capacity}
\label{sec:r12-countable-measure-capacity-obstruction}

The exact energy-normalized capacity in the preceding section is
controlled by one original Radon measure. Its null ideal is therefore
different in nature from the full original tester capacity. The
following all-dimensional example proves that the distinction cannot
in general be removed by replacing that measure by a countable family.
This is a counterexample within the allowed background class, not a
restriction on the backgrounds of the positive theorems.

Fix integers $n\geq q\geq1$ and a real number $L>0$. Work over the
original finite extension $K/\mathbf Q_p$ with its unchanged absolute
value. On $X=(\mathbf G_{m,K}^n)^{\rm an}$ put
\[
 \begin{gathered}
 t_i=\log|T_i|,\qquad
 V=\{x:-L<t_i(x)<L\text{ for }1\leq i\leq n\},\\
 Y=\{T_{q+1}=\cdots=T_n=1\},\qquad R=[Y]|_V.
 \end{gathered}
 \tag{CM1}
\]
If $q=n$, $Y=X$. The current $R$ is an original symmetric strongly
positive, hence middle-positive, separately closed current of degree
$(n-q,n-q)$. It extends across the closed box. Taking a larger compact
polyannulus supplies the precise strict-box hypothesis. In particular
no original field, absolute value or boundary convention has changed.

For $a\in(-L,L)^q$, denote by $x_a$ the original Gauss point of $Y$
with radii $(e^{a_1},\ldots,e^{a_q})$. These points exist for arbitrary
real $a$, without a value-group or algebraic-closure assumption.
The map $a\mapsto x_a$ is a continuous injection. Define on the
whole original domain $V$
\[
                      h_a=\frac{1}{2L}
                                  \max_{1\leq i\leq q}|t_i-a_i|.
 \tag{CM2}
\]

\begin{lemma}[A full original normalized tester at every Gauss point]
\label{lem:r12-gauss-singleton-full-capacity}
The function $h_a$ is an original continuous classical psh function
with $0\leq h_a\leq1$ on $V$. Its full original top current is
\[
                         (Dh_a)^q\wedge R=L^{-q}\delta_{x_a}.
 \tag{CM3}
\]
Consequently
\[
             C_R(\{x_a\};V)\geq L^{-q},\qquad
             C_R^*(\{x_a\};V)\geq L^{-q}.
 \tag{CM4}
\]
These are assertions for the unrestricted original normalized tester
family on the whole $V$, not for a fixed finite presentation class.
\end{lemma}

\noindent\textbf{Status: \texttt{DIRECT\_COROLLARY}} of the
frozen original bounded BT calculus, its integration-current
projection formula, and the elementary toric convex PL calibration.
The explicit calibration below fixes its normalization. It is not a
new instance of the arbitrary-background analytic relation quotient.
The exact frozen torus locator is R6
\path{sections/capacity_gauss_convex_reduction.tex},
\path{eq:r6-toric-convex-laplacian}, whose source SHA256 is
{\small\nolinkurl{2ab2277e81fdad829815f0fe30dc6e921f1d0d8c1fb70de71b718fe4f9d83e38}}.
It explicitly distinguishes the all-original-test toric identity
from a mere pushforward identity. The original real-radius
calibration is also recorded in R6
\path{sections/capacity_laminar_graphs.tex},
\path{eq:r6-arbitrary-radius-capacity}. The targeted nested
archive and source checks are retained in
\path{checks/targeted_r6_calibration.json}.

\begin{proof}
The function in (CM2) is the pullback of a convex real PL function
under an actual original unit moment. Convex smoothing and the frozen
bounded uniform convergence theorem make it original classical psh.
For all $x\in V$, $|t_i(x)-a_i|<2L$, which proves its normalization.

On the coordinate subtorus $Y$, the subgradient at $a$ is the
cross-polytope
\[
                  P_a=\operatorname{conv}
                   \left\{\frac{e_i}{2L},-\frac{e_i}{2L}:1\leq i\leq q\right\},
             \qquad \operatorname{vol}(P_a)=\frac{1}{q!L^q}.
 \tag{CM5}
\]
Here is a direct smooth calculation that also checks support on
the full original test sheaf. Put
\[
 F(z)=\log\!\left(\sum_{i=1}^q
              \bigl(e^{z_i/(2L)}+e^{-z_i/(2L)}\bigr)\right),
       \qquad g_\epsilon(t)=\epsilon F((t-a)/\epsilon).
 \tag{CM5a}
\]
Then $0\leq g_\epsilon-h_a\leq\epsilon\log(2q)$ on the entire
real moment cube. The Hessian of $F$ is the positive definite
covariance matrix of the affinely spanning vectors
$\{\pm e_i/(2L)\}$ with strictly positive exponential weights.
The map $\nabla F$ is a diffeomorphism onto the interior of $P_a$:
for any interior $v$, the strictly convex function $F(z)-v\cdot z$
is coercive, since the support function of $P_a$ exceeds $v\cdot z$
by a positive multiple of $|z|$. It has one critical point;
the positive Hessian gives the smooth inverse. Change of variables
therefore gives
\[
                   \int_{\mathbf R^q}\det(\nabla^2F(z))\,dz
                          =\operatorname{vol}(P_a).
 \tag{CM5b}
\]
Since $\det\nabla^2g_\epsilon(t)
=\epsilon^{-q}\det\nabla^2F((t-a)/\epsilon)$, its finite
measure converges weakly to $\operatorname{vol}(P_a)\delta_a$.
This follows directly by dominated convergence after the displayed
change of scale; no general Monge--Ampere continuity theorem for
unrestricted unbounded potentials is involved.

For completeness the elementary original torus integration formula
has its full-test meaning here. At any point $x$ of the $q$-torus
outside its Gauss skeleton, some Laurent polynomial $f=\sum_m c_mT^m$
satisfies
$|f(x)|<\max_m |c_m|e^{m\cdot t(x)}$. This strict inequality persists
on an ordinary neighborhood. At each point of that neighborhood at
least two monomials attain the displayed maximum, by the ultrametric
inequality. Its moment image thus lies in the tropical hypersurface
of $f$, of real dimension at most $q-1$. Every smooth moment-pulled
$(q,q)$ form vanishes there by the definition of original CLD pullback.
The resulting smooth top measure is therefore carried by the original
Gauss skeleton. On that skeleton canonical torus tropical integration
has weight one and gives, for every compact original smooth scalar
test $\eta$,
\[
  \int_Y \eta\,(Dg_\epsilon)^q
       =q!\int_{(-L,L)^q}\eta(x_t)
                          \det\nabla^2g_\epsilon(t)\,dt.
 \tag{CM5c}
\]
Skeleton support is what makes this an identity on all original
tests, not merely a statement about the moment pushforward.
The normalization agrees with the unit simplex radial shell (PO6).
Now (CM5b), continuity of $t\mapsto\eta(x_t)$ and the uniform
bounded-potential convergence theorem give
$q!\operatorname{vol}(P_a)\eta(x_a)=L^{-q}\eta(x_a)$.
For $q<n$, the frozen integration-current projection formula applies
this same calculation on $Y$ and pushes the Dirac measure into the
original ambient $V$. Thus all original tests and all actual
moment refinements are retained in (CM3).
Using $h_a$ as a permitted normalized tester gives the first
inequality in (CM4); monotonicity of the ordinary-open outer
regularization gives the second.
\end{proof}

\begin{theorem}[No countable Radon family controls the full capacity-null ideal]
\label{thm:r12-no-countable-radon-capacity-control}
For the allowed data (CM1), let $(\nu_j)_{j\geq1}$ be any countable
family of positive Radon measures on $V$. There is one original
point $x_a$, lying in a fixed ordinary relatively compact window
independent of the measures, such that
\[
            \nu_j(\{x_a\})=0\quad\hbox{for every }j,
                  \qquad C_R^*(\{x_a\};V)\geq L^{-q}>0.
 \tag{CM6}
\]
Hence the full original outer-capacity-null ideal cannot be
characterized as the simultaneous null ideal of countably many
positive Radon measures. In particular, there is no single Radon
measure $\nu$ and function $F:[0,\infty)\to[0,\infty)$ with
$F(0)=0$ such that
\[
                         C_R^*(A;V)\leq F(\nu(A))
 \tag{CM7}
\]
for every Borel subset of that window.
\end{theorem}

\noindent\textbf{Status: \texttt{NEW\_PROVED}} as an
all-dimensional obstruction to the proposed countable Radon
control mechanism. The proposed universal equivalence and (CM7)
are \texttt{FALSE}.

\begin{proof}
Choose the fixed window
$W=\{-3L/4<t_i<3L/4\ (1\leq i\leq n)\}\Subset V$
and the compact uncountable Gauss set
\[
                    K_G=\{x_a:a\in[-L/2,L/2]^q\}\subset W.
 \tag{CM8}
\]
Each $\nu_j$ is finite on $K_G$. A finite measure has at most
countably many atoms there: for each positive integer $m$, only
finitely many atoms have mass at least $1/m$. Thus the union,
over all $j$, of the atoms of $\nu_j$ in $K_G$ is countable.
Choose $x_a$ outside it. All the vanishing assertions in (CM6)
follow, while (CM4) gives the strictly positive full capacity.
Applying the proposed (CM7) to this singleton would give
$L^{-q}\leq F(0)=0$, a contradiction.
\end{proof}

\begin{corollary}[Strict null-ideal separation for the parameter energy capacity]
\label{cor:r12-parameter-capacity-strict-null-ideal}
In every $n\geq q\geq1$, choose (CM1) and use the empty parameter
tuple in the parameter finite-energy construction. For any one of
its allowed probes $\psi$, there is a Borel singleton $A\subset W$
such that
\[
                  \mathfrak c_\psi(A)
                    =\mathfrak c_\psi^{\rm o}(A)=0,
                         \qquad C_R^*(A;V)>0.
 \tag{CM9}
\]
The same singleton can be chosen simultaneously for any countable
family of such probes and their energy-normalized capacities.
\end{corollary}

\noindent\textbf{Status: \texttt{NEW\_PROVED}} explicit
strict separation, using the already proved exact null ideal
in (PC12).

\begin{proof}
The empty parameter tuple makes conormality automatic; it does
not change the current $R$. Apply (CM6) to the finite positive
Radon measure $\omega=(D\psi)^q\wedge R$, or to the countable
family of these measures. Formula (PC12) gives precisely the
two zeros in (CM9), whereas (CM4) gives its last inequality.
All resulting capacities still use their explicitly different
admissible families. No comparison of those families has been
suppressed.
\end{proof}

\paragraph{What this does not contradict.}
The geometry-only probe detects specified normal or Jensen defects
by closedness and two-sided original test perturbations; it was
never asserted to dominate every normalized source measure.
Countable cofinality of original smooth test stages is likewise
not countable cofinality of all normalized nonlinear tester
\emph{measures}. The band formulas choose their essential
suprema for each fixed set and parameter control; they do not
claim a single countable list detecting every analytic set.
Thus (CM6) is an obstruction to a particular route to unrestricted
capacity, not a counterexample to its Choquet regularity or to
the analytic PS/PL relation quotient.

% END REVIEWED MODULE research_r12/energy/countable_measure_capacity_obstruction.tex


% BEGIN REVIEWED MODULE research_r12/nonlinear/relative_hessian_cone_obstructions.tex
\section{Two failed enlargements of the nonlinear relative potential cone}
\label{sec:r12-relative-Hessian-cone-obstructions}

The compact-source obstruction shows that a universal relative
source equation cannot be repaired merely by completing the
ambient classical psh class.  This section checks two tempting
larger classes directly, before any regularization or capacity
theorem is attempted.  Both failures occur for explicit smooth
strongly positive backgrounds, which are permitted members of
the fixed arbitrary-middle-positive background class.  They
therefore refute universal constructions; they do not change
the background in the general positive theorems.

All examples use the original finite $K/\mathbf Q_p$, its
original norm, and an ordinary open polyannulus
\[
 U=t^{-1}((-1,1)^n)\subset(\mathbf G_{m,K}^{n})^{\rm an},
 \quad t_i=\log|T_i|,\quad
 \omega_i=d't_i\wedge d''t_i,\quad\beta=\sum_i\omega_i.
 \tag{RH1}
\]
The forms $\omega_i$ are closed and strongly positive.
Their products commute, and $\omega_i^2=0$.  Every scalar
quadratic in the $t_i$ is an actual original smooth function
on the whole domain, with no extension of the ground field.

\begin{theorem}[All lower relative curvatures can be positive while the next is negative]
\label{thm:r12-relative-depth-positivity-failure}
For every integer $q\ge2$, take $n=q+1$ and
$S=\beta\wedge[U]$.  Choose a rational number
\[
                         \frac1q<a<\frac1{q-1},
\]
and put
\[
 u=\tfrac12\left(-a t_1^2+\sum_{i=2}^{q+1}t_i^2\right),
 \qquad\alpha=Du=-a\omega_1+\sum_{i=2}^{q+1}\omega_i.
 \tag{RH2}
\]
Then every smooth current
\[
                        \alpha^j\wedge S,
                             \qquad1\le j<q,
 \tag{RH3}
\]
is strongly positive, hence middle-positive, symmetric and
separately closed on every original test and every actual
moment refinement.  But the next top current is strictly
negative on a nonzero nonnegative original compact scalar test:
\[
       \alpha^q\wedge S
                 =q!(1-qa)\,\omega_1\cdots\omega_{q+1}\wedge[U].
 \tag{RH4}
\]
In particular the condition $D(uS)\ge0$ cannot by itself
define a nonlinear relative psh class on which all iterative
products remain positive.  Even imposing every lower-depth
positivity condition in \textup{(RH3)} does not force the
next one.
\end{theorem}

\noindent\textbf{Status: \texttt{NEW\_PROVED}} as an
explicit all-degree obstruction to the proposed relative
admissibility inference.  That inference is \texttt{FALSE}.
The potential in \textup{(RH2)} is not ambient classical psh,
so this does not contradict the frozen bounded BT theorem.

\begin{proof}
For diagonal coefficients $\lambda_i$, elementary exterior
expansion gives
\[
 \left(\sum_i\lambda_i\omega_i\right)^j\beta
       =j!\sum_{|J|=j+1}\sigma_j(\lambda_J)\omega_J,
 \tag{RH5}
\]
where $\sigma_j$ is the elementary symmetric polynomial of
degree $j$ and $\omega_J=\prod_{i\in J}\omega_i$.
For the coefficients $(-a,1,\ldots,1)$, a subset $J$ not
containing index one has $\sigma_j(\lambda_J)=j+1$.
A subset containing index one has one term equal to one
and $j$ terms equal to $-a$, so its coefficient is $1-ja$.
Both are strictly positive for $1\le j<q$ by the choice
of $a$.  Thus \textup{(RH3)} is a positive linear
combination of decomposable strongly positive original
forms wedged with $[U]$.  Separate closedness follows
from the constant logarithmic differential coefficients.
Strong positivity and these identities are intrinsic
original identities and persist under every actual
refinement.

At $j=q$ there is only the full index set, and its
coefficient is $1-qa<0$.  This proves \textup{(RH4)}.
Choose any nonzero $\chi\in C_c^\infty((-1,1)^n)$ with
$\chi\ge0$.  The original scalar $\chi(t)$ is compactly
supported, and the full original torus calibration gives
\[
 \langle\alpha^q\wedge S,\chi(t)\rangle
                =q!(1-qa)\int\chi(x)\,dx<0.
\]
Hence this is an actual original-current counterexample,
not only a formal Hessian sign computation.  Subtracting
any constant from $u$ changes none of the conclusions;
one may therefore make it negative and bounded if desired.
No zero-boundary energy membership is asserted for $u$.
\end{proof}

\subsection{Strong positivity of the background does not imply hyperbolicity}

For a constant symmetric real $(1,1)$-form
$\alpha_A=\sum A_{ij}d'x_i\wedge d''x_j$ and a fixed
constant strongly positive $(n-q,n-q)$-form $\Theta$,
let $P_\Theta(A)$ be the scalar coefficient of
$\alpha_A^q\wedge\Theta$ relative to
$\omega_1\cdots\omega_n$.
A homogeneous polynomial $P$ of degree $q$ is called
hyperbolic in the identity direction if $P(I)>0$ and,
for every real symmetric $A$, all roots of the degree-$q$
polynomial $t\mapsto P(A+tI)$ are real.

\begin{theorem}[A sum of allowed positive backgrounds destroys the hyperbolicity mechanism]
\label{thm:r12-positive-background-nonhyperbolicity}
For every $q\ge2$, take $n=2q$ and
\[
 \Theta=\omega_1\cdots\omega_q
                  +\omega_{q+1}\cdots\omega_{2q},
 \qquad S=\Theta\wedge[U].
 \tag{RH6}
\]
This is a smooth strongly positive separately closed allowed
background of remaining degree $q$.  Its scalar Hessian
polynomial is
\[
 P_\Theta(A)=q!\left(\det A_{1:q,1:q}
                          +\det A_{q+1:2q,q+1:2q}\right),
 \tag{RH7}
\]
and is not hyperbolic in the identity direction.  Each of
the two summands of the background separately has a
hyperbolic determinant-compression polynomial.

Already for $q=2$, the set containing $I$ defined by the
two scalar inequalities
$P_\Theta(A)>0$ and $D_I P_\Theta(A)>0$ is not convex,
even if one takes its connected component containing $I$.
Thus scalar positivity and selection of that component
cannot be invoked as a universal convex admissibility cone.
\end{theorem}

\noindent\textbf{Status: \texttt{NEW\_PROVED}} as an
exact all-degree failure mechanism for arbitrary-background
Hessian relaxations.  The blanket hyperbolicity and convexity
inferences are \texttt{FALSE}.

\begin{proof}
Wedging with the first coordinate $q$-plane form kills
every coefficient of $\alpha_A^q$ except the complementary
principal block; the coefficient of that block is $q!$
times its determinant.  The second summand gives the
other principal block.  This proves \textup{(RH7)}.
Each block is a real symmetric compression, so its
identity pencil has real eigenvalue roots, proving the
last assertion about individual summands.

Now take $A_0=\operatorname{diag}(I_q,-I_q)$.  Then
\[
 P_\Theta(A_0+tI)=q!\bigl((t+1)^q+(t-1)^q\bigr).
 \tag{RH8}
\]
If $q$ is even, this polynomial is strictly positive
on the real line and has degree $q\ge2$.  If $q$ is
odd, then $q\ge3$; its derivative is strictly positive
on the real line, so its only real root is the simple
root zero.  In either case it has nonreal roots.
Also $P_\Theta(I)=2q!>0$.  Thus hyperbolicity fails.

For the convexity assertion let $q=2$ and take diagonal
matrices
\[
 A=\operatorname{diag}(4,0,-1,-1),\qquad
 B=\operatorname{diag}(1,-3,2,2).
 \tag{RH9}
\]
They satisfy $P_\Theta(A)=P_\Theta(B)=2$ and
$D_I P_\Theta(A)=D_I P_\Theta(B)=4$.  Their average
has diagonal $(5/2,-3/2,1/2,1/2)$, and
\[
                       P_\Theta((A+B)/2)=-7<0.
 \tag{RH10}
\]
Both belong to the indicated component.  To see this
explicitly, write the diagonal as
$(u+a,u-a,v+b,v-b)$.  Then
\[
 P_\Theta/2=u^2+v^2-a^2-b^2,
                 \qquad D_I P_\Theta=4(u+v).
\]
For $A$ the parameters are $(u,v,a,b)=(2,-1,2,0)$;
for $B$ they are $(-1,2,2,0)$.  In either case first
decrease $a$ to zero, keeping both inequalities strict.
Next move $(u,v)$ linearly to $(1,1)$.  Its sum stays
positive, so $(u,v)$ never vanishes and the first
inequality stays strict as well.  This gives paths to
$I$ inside the set, while \textup{(RH10)} excludes
their midpoint.

All these matrices are Hessians of actual original
quadratic moment functions on \textup{(RH1)}.  Thus
the algebra is realized on an allowed original analytic
space.  The scalar inequalities here are deliberately
not confused with middle positivity of all intermediate
original currents; the first counterexample addressed
that distinct proposed condition.
\end{proof}

\begin{proposition}[A restricted scalar repair for a power background]
\label{prop:r12-power-background-hyperbolicity}
For $n\ge q\ge1$ and $\Theta=\beta^{n-q}$, one has
\[
 P_\Theta(A)=q!(n-q)!\,\sigma_q(A),
 \qquad
 P_\Theta(A+tI)=q!\frac{d^{n-q}}{dt^{n-q}}\det(A+tI).
 \tag{RH11}
\]
Consequently this scalar polynomial is hyperbolic in the
identity direction.  This elementary conclusion is
specific to the displayed power background.  It does
not supply a relative regularization theorem, positivity
on all original mixed tests for a new nonclassical class,
or a nonlinear source existence theorem.
\end{proposition}

\noindent\textbf{Status: \texttt{DIRECT\_COROLLARY}} of
the exterior determinant expansion and Rolle's theorem;
it records a qualified structural distinction, not an
unrestricted repair.

\begin{proof}
Orthogonal diagonalization, or direct determinant
expansion, gives the first formula.  Comparing the
elementary symmetric coefficients of the two sides
gives the second.  The polynomial $\det(A+tI)$ has
only real roots for symmetric $A$.  Repeated Rolle's
theorem, including multiplicities, gives the same
property for its derivative of order $n-q$.
The leading coefficient is positive, and
$P_\Theta(I)>0$, which proves the stated hyperbolicity.
\end{proof}

\paragraph{Minimal remaining nonlinear-class gate.}
A larger relative theory must specify a class stable
under its proposed mixed products and the intended
limiting operations.  The linear condition $D(uS)\ge0$
is insufficient in every nonlinear degree, and a scalar
Garding prescription is not available merely from middle
or even strong positivity of $S$.  Background classes
carrying additional Hessian/hyperbolicity structure may
support such a separate construction, but that is a
\texttt{CONDITIONAL} direction, not an assumed property
of the arbitrary background fixed in the main problem.
The construction of an adequate nonlinear relative
potential class for every original $S$ remains
\texttt{OPEN}.

% END REVIEWED MODULE research_r12/nonlinear/relative_hessian_cone_obstructions.tex


\part{The analytic relation quotient and remaining gates}

% BEGIN REVIEWED MODULE research_r12/gq/curvature_depth_conormal_reconstruction.tex
\section{Curvature depth and sharper conormal reconstruction}
\label{sec:r12-curvature-depth}

The original field $K/\mathbf Q_p$ and its absolute value are fixed.
Let $X$ be good, Hausdorff, boundaryless, quasi-smooth and pure of
dimension $n$, and let all opens be ordinary analytic opens.
An allowed background is an arbitrary original symmetric CLD
middle-positive, separately $d'$- and $d''$-closed current.
All current equalities below are tested against every original compact
smooth complementary CLD form, after every actual moment refinement.
No ordinary current equality is promoted to equality on all GK delta
tests.

We use the following exact frozen R10 inputs, not new assertions:
the intrinsic finite PS coefficient module $\mathscr N_{\rm PS}$,
its fixed-PS continuity against convergent positive closed tails,
the corrected insertion operators $\mathfrak B_v$, the every-parameter
corrected-cap/smooth-wedge equivalence (PC1)--(PC3), and the intrinsic
alternating quotient of a Hessian-invisible coefficient-Radon current.
Their sources are
\path{signed_descent/pws_differential_curvature_module.tex},
\path{signed_descent/pws_curvature_generated_reconstruction.tex}, and
\path{signed_descent/alternating_quotient_and_residue.tex} in the exact
R10 archive. R10 retains only one regularizable curvature factor in
its reconstruction threshold. The new invariant below retains their
number.

\subsection{A filtration by the number of curvature factors}

Let $\mathcal P_{\mathrm{sm}\downarrow}$ have its frozen meaning:
locally bounded original classical psh functions which locally admit
decreasing original smooth classical psh approximations, uniformly
bounded on compact subsets. Every actual convex PL moment potential
belongs to this class. We do not assert that every bounded classical
psh function belongs to it.

\begin{definition}[Regularizable curvature depth]
\label{def:r12-curvature-depth}
For $0\le k\le\min(a,b)$, let
$\mathscr G_k^{a,b}(U)$ be the actual PS submodule locally generated
by finite terms
\[
       \theta\wedge T_{S,k}(u_1,\ldots,u_k),
       \qquad u_i\in\mathcal P_{\mathrm{sm}\downarrow},
 \tag{HD1}
\]
where $S$ ranges over all allowed original backgrounds, $\theta$ is
an arbitrary actual PS form of the required degree, and
$T_{S,0}=S$. Finite real signed sums are allowed. Thus
$\mathscr G_0=\mathscr N_{\rm PS}$, and $\mathscr G_1$ is the
frozen R10 PS curvature-generated module. Membership is the existence
of such a local witness; a chosen witness is not part of the current.
\end{definition}

\begin{lemma}[Depth is retained by the intrinsic operations]
\label{lem:r12-depth-operations}
One has $\mathscr G_{k+1}\subseteq\mathscr G_k$. The depth-$k$
module is stable under actual PS coefficients and every permitted
corrected bounded classical insertion. If the inserted function
$v$ belongs to $\mathcal P_{\mathrm{sm}\downarrow}$, then
\[
                    \mathfrak B_v\mathscr G_k
                                         \subseteq\mathscr G_{k+1}.
 \tag{HD2}
\]
For any fixed arbitrary allowed background $S$, a prescribed pure
ambient cycle action of codimension $b$, with any actual PS
coefficient, belongs to $\mathscr G_b$.
\end{lemma}

\noindent\textbf{Status: \texttt{DIRECT\_COROLLARY}} of frozen
mixed symmetry, the intrinsic PS insertion formula, and the finite
convex-PL presentation of a constant cycle.

\begin{proof}
Absorb one positive closed BT factor into $S$ to obtain the first
inclusion. The frozen intrinsic formula gives
\[
 \mathfrak B_v\bigl(\theta\wedge T_{S,k}(\mathbf u)\bigr)
             =\theta\wedge T_{S,k+1}(\mathbf u,v).
\]
If $v$ is only bounded classical, absorb $D(vS)$ into the background
and keep the original $k$ regularizable slots; this preserves depth
$k$. If $v$ is regularizable, keep all $k+1$ slots, giving (HD2).
A codimension-$b$ constant cycle has a finite signed sum of length-$b$
convex-PL presentations. Each convex-PL slot is regularizable by the
frozen real convolution argument. No background is approximated.
\end{proof}

\subsection{The exterior identity with all curvature factors retained}

For a real vector space $L$, write $s_{a,b}$ for the normalized
alternating shuffle from $\Lambda^{a+b}L^*$ to doubled superforms,
with $\operatorname{Alt}\circ s_{a,b}=\mathrm{id}$.
For an original unit $g$ put
$\omega_g=d'\log|g|\wedge d''\log|g|$.

\begin{lemma}[Depth-$k$ conormal shuffle annihilation]
\label{lem:r12-depth-shuffle}
Let $k\ge1$, $j\ge0$, and $\ell\ge1$. Let
$\Omega=\omega_{g_1}\wedge\cdots\wedge\omega_{g_j}$ for original
analytic units on an ordinary open. If
\[
                              n<j+2\ell+k,
 \tag{HD3}
\]
then for all original smooth scalars $v_1,\ldots,v_k$ and all
original prime forms $\gamma$ of degree $2\ell$,
\[
       \Omega\wedge Dv_1\wedge\cdots\wedge Dv_k
                                      \wedge s_{\ell,\ell}\gamma=0.
 \tag{HD4}
\]
Consequently, if $R\in\mathscr G_k^{P,P}$ and
$\ell=n-P-j\ge1$, then under (HD3)
\[
                 \langle R,\Omega\wedge s_{\ell,\ell}\gamma\rangle=0
                         \qquad(\gamma\in\mathscr A_c^{2\ell,0}).
 \tag{HD5}
\]
Neither symmetry nor separate closedness of $R$ is needed.
\end{lemma}

\noindent\textbf{Status: \texttt{NEW\_PROVED}}. The case $k=1$
is the frozen R10 result; the full retained-depth identity is new.

\begin{proof}
Represent all forms at one actual moment stage and restrict to one
actual tangent space $L$, of dimension at most $n$. A symmetric
Hessian tensor is a finite real linear combination of rank-one
tensors $h'\wedge h''$. Expand each of the $k$ Hessians in that
way. One resulting term has $j+k$ repeated prime/double-prime
covectors, including the $j$ unit covectors in $\Omega$.
If those covectors are dependent, their doubled exterior product
is zero. Otherwise let $H\subset L^*$ be their span, of dimension
$j+k$. Their doubled top product annihilates the ideal generated
by $H'$ and $H''$. The shuffle is natural under $L^*\to L^*/H$.
Since
\[
                      \dim(L^*/H)\le n-j-k<2\ell,
\]
the image of $\gamma$ in $\Lambda^{2\ell}(L^*/H)$ is zero.
Its shuffle lift is in the annihilated ideal. This proves (HD4)
term by term. The argument allows arbitrary smooth coefficient
arrays and is unchanged on every actual refinement.

For (HD5), take one witness term in (HD1) and local smooth decreasing
approximations $u_{i,m}\downarrow u_i$. Frozen joint bounded BT
continuity gives
$T_{S,k}(\mathbf u_m)\to T_{S,k}(\mathbf u)$ on every original
test. The tails are positive and separately closed, with the finite
compact coefficient controls supplied by the same bounded theorem.
The frozen fixed-PS continuity theorem therefore also gives the
limit after multiplication by the fixed actual PS coefficient
$\theta$. At every finite $m$, (HD4) and the original smooth rule
make the required pairing zero. Pass to the limit. Finite sums and
ordinary localization prove (HD5). Different approximating terms
may use different finite moments: each finite stage includes the
test moment, and no common moment for the whole sequence is assumed.
No derivative of $\theta$ is taken, and no Borel face restriction is
treated as a closed background.
\end{proof}

\subsection{Sharp cap reconstruction at every depth}

Fix $1\le k\le P<n$ and put $q=n-P$. Let
$\mathscr G_{k,c}^{P,P}$ be the currents in
$\mathscr G_k^{P,P}$ killed by every local $c$-fold smooth unit
wedge, for $1\le c\le q$, and put $\mathscr G_{k,0}^{P,P}=0$.
The frozen corrected-cap criterion identifies these with the
kernels of all corrected intrinsic rank-$c$ caps, at every real
parameter and on every original intrinsic fiber test.

\begin{theorem}[Sharp depth-sensitive kernel jumps and zero detection]
\label{thm:r12-depth-sharp-reconstruction}
The universal inclusion
$\mathscr G_{k,c-1}^{P,P}\subseteq\mathscr G_{k,c}^{P,P}$
is strict on an original bounded $n$-torus precisely when
\[
                              n\ge 2q-c+1+k.
 \tag{HD6}
\]
In the symmetric submodule, strictness additionally requires
$q-c+1$ even. Otherwise the two kernels agree on every allowed
original analytic space.

Define
\[
 \begin{split}
 h_k^{\rm all}(n,P)&=\max\{0,n-2P+k\},\\
 h_k^{\rm sym}(n,P)&=
      \max\left\{q\bmod2,
                  q-2\left\lfloor\frac{P-k}{2}\right\rfloor\right\}.
 \end{split}
 \tag{HD7}
\]
These integers lie between zero and $q$. The sharp universal
reconstruction ranges are
\[
 \begin{array}{ll}
 R\in\mathscr G_{k,c}^{P,P}\Longrightarrow R=0
             &\Longleftrightarrow c\le h_k^{\rm all}(n,P),\\[2pt]
 R\in\mathscr G_{k,c}^{P,P},\ R\text{ symmetric}\Longrightarrow R=0
             &\Longleftrightarrow c\le h_k^{\rm sym}(n,P).
 \end{array}
 \tag{HD8}
\]
Sharpness refers to the universal statement over all allowed spaces,
not to strictness on every particular analytic space.
\end{theorem}

\noindent\textbf{Status: \texttt{NEW\_PROVED}} for $k>1$;
the $k=1$ thresholds are \texttt{INTERNAL\_INPUT} from R10 and are
recovered exactly.

\begin{proof}[The universal vanishing assertion]
Put $j=c-1$ and $\ell=q-j$. For $R\in\mathscr G_{k,c}$ choose
any $j$ unit factors $\Omega$ and set $T=\Omega\wedge R$.
Every additional unit square kills $T$. On a common original moment,
monomial units and polarization imply annihilation of the full
symmetric Hessian ideal. The frozen intrinsic alternating quotient
therefore writes the complete pairing of $T$ as
\[
       \langle T,\eta\rangle
             =\mathcal R_T(\operatorname{Alt}\eta),\qquad
       \mathcal R_T(\gamma)
             =\langle R,\Omega\wedge s_{\ell,\ell}\gamma\rangle.
 \tag{HD9}
\]
If (HD6) fails, then $n<j+2\ell+k$, so (HD5) gives
$\mathcal R_T=0$. Hence every $j$-compression vanishes. This proves
equality of the two consecutive kernels. If $R$ is symmetric and
$\ell$ is odd, the same frozen alternating quotient vanishes by
the prime/double-prime swap rule, independently of dimension.
No exactness or derivative of $R$ is used.
\end{proof}

\begin{proof}[Exact original witnesses and the threshold arithmetic]
Suppose (HD6) holds. With $j=c-1$, $\ell=q-j$, put
$z=n-j-2\ell\ge k$ and split the original torus coordinates into
groups $A,V,Z$ of sizes $j,2\ell,z$. On the bounded original norm
box let $t$ be their logarithms, let $\gamma_V$ be the prime volume
form in the $V$ coordinates, and put
\[
 R=\Phi_V\wedge\bigwedge_{a\in Z}\omega_a\wedge[U],
                 \qquad \Phi_V=s_{\ell,\ell}(\gamma_V).
 \tag{HD10}
\]
Its degree is $P=\ell+z$. Extract any $k$ distinct $Z$ coordinates
as the bounded smooth classical potentials $t_a^2/2$. The remaining
constant form is an allowed smooth PS coefficient, and the background
is the original $[U]$. Thus $R\in\mathscr G_k^{P,P}$.

Every $j+1$ additional rank-one wedges kill (HD10): the $Z$ factors
discard their $Z$ components; an independent $(j+1)$-dimensional
covector span in $A^*\oplus V^*$ has a nonzero vector wholly in
$V^*$; its rank-one square kills the top $V$ shuffle. If the span
was dependent, its exterior product was already zero. Changing
basis among the $j+1$ covectors multiplies their doubled exterior
product by the square of the determinant, so this argument applies
to their original product.

The constant identity holds on every actual refined tangent space.
Indeed for an $n$-dimensional tangent space approximate its possibly
singular projection to the original coordinate space by isomorphisms,
keeping the added covectors fixed, and pass the polynomial identity
to the limit. For a smaller tangent space extend the map and covectors
to dimension $n$ and restrict. Integration therefore gives zero on
every refined original test, not only toric tests.

The $j$ coordinate wedges in $A$, however, leave a nonzero constant
form. A constant complementary exterior test with nonzero pairing,
times a nonzero compact nonnegative toric cutoff, detects it by the
original calibrated Lebesgue integral. Thus
$R\in\mathscr G_{k,c}\setminus\mathscr G_{k,c-1}$.
The swap rule makes (HD10) symmetric exactly when $\ell$ is even;
for odd $\ell$ it is antisymmetric. This proves precisely the
claimed jump locations, including their original-refinement scope.

Without symmetry, the first strict jump is $c=j+1$ for the least
integer $j\ge0$ with $j\ge2q+k-n=n-2P+k$. This is the first line
of (HD7). In the symmetric case $j$ must have parity $q$, so the
least such integer is the second line of (HD7). If this integer is
$q$, there is no jump at a permitted rank. Otherwise its first
nonzero witness belongs to every subsequent kernel. Starting from
the zero rank-zero kernel proves both implications and all sharpness
claims in (HD8).
\end{proof}

\subsection{The actual relation consequence and its unchanged premise}

\begin{corollary}[A cycle-codimension-sensitive arbitrary-background reduction]
\label{cor:r12-depth-analytic-response}
Fix any allowed original background $S$. Suppose a finite compatible
prescribed response has final degree $(P,P)$, $P<n$, and is a sum
\[
                R=\mathcal A_S(e)
                     =\sum_i\theta_i\wedge\mathcal A(C_i)S,
                \qquad\operatorname{codim}C_i\ge k\ge1,
 \tag{HD11}
\]
with actual PS coefficients. If the same presentation acts by zero
on every local original pushed rank-$c$ cap background of $S$, at
every real parameter, then $R=0$ whenever
$c\le h_k^{\rm all}(n,P)$; for symmetric $R$ it suffices that
$c\le h_k^{\rm sym}(n,P)$.

For a pure codimension-$b$ constant-cycle response over a
degree-$(p,p)$ background, the symmetric threshold is
\[
       \max\left\{(n-p-b)\bmod2,
                    n-p-b-2\left\lfloor p/2\right\rfloor\right\}.
 \tag{HD12}
\]
All conclusions hold on every ordinary germ and all original tests
after every actual refinement of the given presentations.
\end{corollary}

\noindent\textbf{Status: \texttt{DIRECT\_COROLLARY}} as a
reconstruction statement; \texttt{CONDITIONAL} as an analytic
relation quotient because its capped-background premise remains
independent of actual analytic equality.

\begin{proof}
Lemma~\ref{lem:r12-depth-operations} puts (HD11) in $\mathscr G_k$.
The frozen PS-linearity and hinge/cycle commutation identify its
corrected cap with the action of the same presentation on the
original capped background. The printed premise makes it zero.
Apply (HD8) and the corrected-cap/smooth-wedge equivalence. In the
pure case take $P=p+b$ and $k=b$ in (HD7), giving (HD12).
The proof nowhere invokes a BT operation over an auxiliary residue
field.
\end{proof}

\paragraph{What remains open.}
Actual equality of two positive cycle pullbacks still does not supply
the cap-response equality assumed in (HD11). If $P=n$, no positive
cap rank is available at all. The new depth filtration reduces the
number of conormal residues which could survive an independently
proved cap relation; it does not turn cap order into middle order.
Indeed the frozen R10 signed-strong-span order counterexamples are
already finite sums of length-$P$ smooth positive curvature products
of $[U]$, so they have the maximal depth $P$. They continue to refute
every excluded universal cap-order reflection even at that depth.
Full OPEN-GQ and low-real-degree OPEN-CTO therefore remain open.

\paragraph{Exact source recovery.}
The R10 archive was read from the unchanged embedded R11 file and
verified against all 226 entries of its \path{SHA256SUMS.txt}, with
ZIP CRC passing. Its SHA256 is
{\small\nolinkurl{4891735a31504e5154d7e027efb2fdf739dbf4645fb963f45a4b4943360b1402}}.
No R/B/A2 input was re-audited, and no earlier source was modified.

% END REVIEWED MODULE research_r12/gq/curvature_depth_conormal_reconstruction.tex


% BEGIN REVIEWED MODULE research_r12/gq/original_fiber_transfer_obstructions.tex
\section{Two explicit obstructions to a fiberwise quotient shortcut}
\label{sec:r12-original-fiber-transfer-obstructions}

This section tests two proposed implications which would bypass
the remaining actual relation gate. Both implications are false
even over the original fixed field. Neither example is a
counterexample to OPEN-GQ: in particular, no pair of distinct
cycle presentations with equal actual analytic pullbacks is
produced here. Their role is to rule out invalid transfers in a
proof of that conjecture.

\subsection{The reference-space pushforward loses original mixed tests}

\begin{proposition}[A nonfaithful reference-space transfer]
\label{prop:r12-nonfaithful-reference-pushforward}
For every $h\ge0$ there is an original smooth $(h+2)$-dimensional
torus domain $U$, an original unit parameter $s$, and two nonzero
positive separately closed $s$-annihilated currents $R_0,R_1$ of
degree $(2,2)$ such that:
\begin{enumerate}
\item both currents are carried by $s^{-1}(0)$;
\item their proper-support pushforwards under one reference
projection $\pi:U\to Y$ agree as original currents on $Y$;
\item they are distinct on an explicit compact original smooth
test on $U$ using a mixed unit.
\end{enumerate}
In particular, an original conditional current supported on a
norm fiber is not determined by its reference-space pushforward,
even within positive closed currents obtained by an ordinary
divisor cap of a smooth analytic graph.
\end{proposition}

\noindent\textbf{Status: \texttt{NEW\_PROVED}} counterexample
to the proposed faithful transfer; that transfer is \texttt{FALSE}.

\begin{proof}
Choose a uniformizer $\varpi$ of $K$, an integer $k\ge1$, and put
$\lambda_0=1$, $\lambda_1=1+\varpi^k$. Both have norm one.
Let $A$ be a bounded open torus annulus with logarithmic interval
$(-L,L)$ containing zero. Let $Z$ be any bounded open
$h$-dimensional torus polyannulus; for $h=0$ take a point. Set
\[
       U=A_{T_1}\times_K A_{T_2}\times_K Z,\qquad
       Y=A_{T_2}\times_K Z,\qquad
       s=\log|T_1|,
                                                                    \tag{FT1}
\]
and let $\pi$ forget $T_1$. For $j=0,1$, the graph
$G_j=\{T_2=\lambda_jT_1\}\subset U$ is a closed smooth analytic
subspace, and $\pi|_{G_j}:G_j\to Y$ is an analytic isomorphism.
Its original integration current $[G_j]$ is positive and
separately closed. Define the actual bounded local hinge caps
\[
                         R_j=D\bigl(s_+[G_j]\bigr),
                  \qquad s_+=\max(s,0).                            \tag{FT2}
\]
Frozen BT and the original annulus calibration show that these
currents are positive, separately closed, and nonzero. Their
coefficient measures are carried by $G_j\cap\{s=0\}$. Since
$sR_j=0$ and each $R_j$ is separately closed, the ordinary
distributional product rule gives
$d's\wedge R_j=d''s\wedge R_j=0$.

The graph isomorphisms identify $s$ with $\log|T_2|$ because
$|\lambda_j|=1$. Consequently
\[
                      \pi_*R_0=\pi_*R_1
                          =D\bigl((\log|T_2|)_+[Y]\bigr).           \tag{FT3}
\]
This is an equality on every original test of $Y$: it is obtained
by an analytic isomorphism of the original graphs, not by limiting
attention to toric tests. The pushforward is legitimate although
$\pi$ is not proper on all of $U$, since its restriction to the
closed graph carriers is proper; equivalently one may define (FT3)
through those isomorphisms with compact cutoffs near each inverse
test support.

Now consider the original analytic function
\[
                          g=T_2/T_1-1+\varpi^{k+1}.                 \tag{FT4}
\]
It is constant $\varpi^{k+1}$ on $G_0$ and constant
$\varpi^k(1+\varpi)$ on $G_1$. Their norms are respectively
$|\varpi|^{k+1}$ and $|\varpi|^k$, which are distinct, including
in residue characteristic two. Choose $F\in C_c^\infty(\mathbf R)$
with
\[
                  F((k+1)\log|\varpi|)=1,
                  \qquad F(k\log|\varpi|)=0.
\]
Define $F(\log|g|)=0$ where $g=0$. This is an original smooth
scalar on $U$: where $g\ne0$ it is a smooth function of an
original unit moment, and near each zero of $g$ it is identically
zero because $F$ has support bounded below.

Choose a nonnegative smooth compact real cutoff
$\chi(\log|T_1|,\log|T_2|)$ equal to one near $(0,0)$.
Choose a compact original strongly positive top test $\eta_Z$
on $Z$ with integral one (use $\eta_Z=1$ for $h=0$). The test
\[
                  \eta=\chi\,F(\log|g|)\,\operatorname{pr}_Z^*\eta_Z
                                                                    \tag{FT5}
\]
is a compact original smooth $(h,h)$-form on $U$. The frozen
degree-one annulus cap normalization on the graph gives
\[
                            \langle R_0,\eta\rangle=1,
                            \qquad\langle R_1,\eta\rangle=0.        \tag{FT6}
\]
Thus $R_0\ne R_1$ on original ambient tests. The separating moment
$g$ mixes both torus variables; it is not a moment pulled back
from the reference space $Y$ alone.
\end{proof}

The proposition specifically refutes the inference
\[
 \begin{gathered}
 \text{same parameter and same current on a reference space}\\
     \Longrightarrow\ \text{same original conditional current}.
 \end{gathered}
                                                                    \tag{FT7}
\]
It does not say that product tests never suffice for a necessary
positivity detector. Such detectors were used in the product
order-transfer theorem with a separate original-current positivity
argument. It also does not refute the conditional-current theorem:
that theorem retains the whole ambient test sheaf and hence
distinguishes (FT2) correctly.

\subsection{Fixed-moment caps do not constitute all-unit reconstruction}

\begin{proposition}[Invisible to all caps in one finite moment tuple]
\label{prop:r12-fixed-moment-cap-blindness}
Fix integers $r\ge2$, $1\le b<r$, and $h\ge1$. There is an original
smooth torus domain of dimension $n=r+h$, an original moment tuple
$t=(t_1,\ldots,t_r)$ of real rank $r$, a positive closed background
$S$ of degree $(r-b,r-b)$, and a positive constant codimension-$b$
real cycle $C$ such that the response
\[
                            R=\mathcal A_t(C)S                     \tag{FT8}
\]
is nonzero, positive, and of degree $(r,r)$, but every elementary
unit hinge cap from the fixed moment tuple annihilates $R$.
In fact
\[
        D(\phi(t)R)=0
        \quad\text{for every locally bounded Borel }\phi,
        \qquad
        (d't_{r+1}\wedge d''t_{r+1})\wedge R\ne0                  \tag{FT9}
\]
for one additional original unit $T_{r+1}$, where
$t_i=\log|T_i|$. Thus fixed-moment cap vanishing cannot be inserted
into a reconstruction theorem requiring all original unit caps.
\end{proposition}

\noindent\textbf{Status: \texttt{NEW\_PROVED}} counterexample
to the proposed fixed-moment reconstruction; that implication is
\texttt{FALSE}. The example is a general-dimensional diagnostic,
not a nonmatroidal analytic relation counterexample.

\begin{proof}
Let $U$ be a bounded open $n$-torus polyannulus whose logarithmic
box contains zero, and put $p=r-b$. Write
$\omega_i=d't_i\wedge d''t_i$, and set
\[
         S=\omega_1\wedge\cdots\wedge\omega_p\wedge[U],
         \qquad
         C=\{t_{p+1}=\cdots=t_r=0\}
              \quad\text{with weight one}.                         \tag{FT10}
\]
This is a positive balanced linear real cycle. Its frozen ambient
action is the complete intersection of the original convex hinges:
\[
     R=\omega_1\wedge\cdots\wedge\omega_p\wedge
               D\bigl((t_{p+1})_+\bigr)\wedge\cdots\wedge
               D\bigl((t_r)_+\bigr)\wedge[U].                      \tag{FT11}
\]
The notation denotes the frozen recursive mixed BT current;
no product of arbitrary distributions is being defined.

For $i\le p$, wedging (FT11) by either $d't_i$ or $d''t_i$ repeats
a factor. For $p<i\le r$, the individual hinge cap is conormal
to $t_i$; this follows from its closedness and its carrier
$t_i=0$, or from the frozen unit-cap lemma. Conormality survives
every remaining curvature insertion by the distributional identity
(PB8). Therefore $R$ is annihilated by all $d't_i,d''t_i$,
$1\le i\le r$. The closed Borel band theorem, now applied with
parameter tuple $t$, gives the first identity in (FT9).

To see both nonvanishing assertions, choose a smooth nonnegative
compact real cutoff $\zeta$ in the logarithmic box, positive near
zero. Test $\omega_{r+1}\wedge R$ against
\[
                        \zeta(t_1,\ldots,t_n)
                                  \omega_{r+2}\wedge\cdots\wedge\omega_n.
\]
The original torus cap calibration evaluates the pairing as
\[
       \int \zeta(t_1,\ldots,t_p,0,\ldots,0,t_{r+1},\ldots,t_n)
                  \,dt_1\cdots dt_p\,dt_{r+1}\cdots dt_n>0.        \tag{FT12}
\]
There is no hidden field normalization: all moments, hinges,
Lebesgue factors, and graph caps use the original fixed norm and
the frozen weight-one original torus convention. The displayed
positive value proves $\omega_{r+1}\wedge R\ne0$, hence $R\ne0$.
\end{proof}

For the requested low real degree $b=2<r=3$, take $p=1$ and any
$n\ge4$. The positive response of the line
$\{t_2=t_3=0\}$ over $S=\omega_1[U]$ is killed by every cap built
from $(t_1,t_2,t_3)$ but not by the fourth unit square. This is not
the required genuinely nonmatroidal pressure test: it only refutes
a step which a general proposed proof would have to justify.

\paragraph{Failed proof route and the corrected remaining statement.}
One may try to multiply a degree-$b$ actual relation by enough
caps in its fixed real moment to reach the frozen top-real-degree
quotient. Even if this proves vanishing after all those caps, it
does not meet the all-original-unit premise of conormal
reconstruction. Proposition~\ref{prop:r12-fixed-moment-cap-blindness}
is a concrete obstruction to that inference. Adding further
actual moments also adds further real moment rank, so repeating
the fixed-rank argument is not a proof of the missing premise.
The route is recorded as \texttt{FAILED\_ATTEMPT}.

The original-current disintegration removes a measure-theoretic
and test-coherence obstacle, but neither example removes the
actual relation gate. A successful next argument must prove
vanishing on the original conditional currents using all their
mixed tests, or produce a background-independent trace-lifting
certificate as in Section~\ref{sec:r12-level-trace-calculus}.
General OPEN-GQ and signed low-degree OPEN-CTO retain their status.

\paragraph{Dependencies.}
Only frozen original analytic graph integration, the degree-one
torus hinge calibration, arbitrary-background bounded BT,
ordinary support locality, and the new closed-band theorem are
used. The calculations are exact symbolic calibrations, not
numerical evidence for an unproved cycle equality.

% END REVIEWED MODULE research_r12/gq/original_fiber_transfer_obstructions.tex


% BEGIN REVIEWED MODULE research_r12/gq/conditional_action_and_division_obstructions.tex
\section{Conditional responses and two invalid fiber-transfer repairs}
\label{sec:r12-conditional-response-gate}

This section returns directly to the analytic relation problem.
The conditional currents constructed in (DI4) retain the original
test sheaf, but their existence alone does not annihilate a signed
actual relation. We first give the exact response-level statement,
then isolate two stronger implications which are false. Neither
counterexample below is a counterexample to OPEN-GQ.

\begin{proposition}[Disintegration of a fixed finite prescribed response]
\label{prop:r12-conditional-PS-response}
Work on a second-countable ordinary neighborhood \(U\) as in (DI1).
Let \(R\) be an allowed conormal background and fix its original
disintegration \(R=\int R_a\,d\nu(a)\). Let \(e\) be one finite
compatible PS/PL presentation, with every prescribed degree permitted.
Then its original coefficient-Radon responses satisfy
\[
 \mathcal A_R(e)=\int_J\mathcal A_{R_a}(e)\,d\nu(a),\qquad
 \mathbf1_{\{s\in E\}}\mathcal A_R(e)
       =\mathcal A_{R_E}(e)
       =\int_E\mathcal A_{R_a}(e)\,d\nu(a)
       \quad(E\text{ Borel}).
 \tag{CQ1}
\]
Every fixed compact coefficient variation of the conditional
response is integrable. Moreover
\[
 \mathcal A_R(e)=0
   \quad\Longleftrightarrow\quad
 \mathcal A_{R_a}(e)=0\text{ on every original test for }\nu
       \text{-almost every }a.
 \tag{CQ2}
\]
The exceptional set in (CQ2) is common to every original test and
every actual refinement of this fixed finite presentation. It may
depend on \(R,e\); no simultaneous assertion for uncountably many
new presentations is made.
\end{proposition}

\noindent\textbf{Status: \texttt{DIRECT\_COROLLARY}} of the new
original disintegration and the frozen finite ambient cycle action,
with signed response cancellation verified explicitly.

\begin{proof}
On each protected compact original stage, choose a common finite
actual refinement for the given presentation. Each constant cycle
occurring in the prescribed compatible presentation has its specified
finite signed convex-PL presentation. Thus its response is a
finite signed sum of positive bounded classical mixed products,
followed by multiplication by the specified finite PS coefficient.
Apply the measurable original BT-family construction with slices
independent of \(a\). Its recursive test formula and relative CLN
give integrable compact coefficient variations. Multiplication by
each fixed PS coefficient is finite facewise Borel multiplication
with bounded coefficients on this stage. The coefficient kernels
therefore remain measurable and integrable. The frozen actual PS
multiplication and its level-tangency identifications retain all
original test refinements. This proves the first equality in (CQ1).

The closed-band identity applies to every positive mixed tail before
PS multiplication. Multiplying its coefficient-Radon identity by the
fixed PS coefficient gives the middle equality in (CQ1); the same
conditional construction for \(R_E\) gives the last equality. No
PS-weighted or Borel-restricted response is declared closed.

If the conditional responses vanish almost everywhere, their integral
is zero. Conversely suppose the left side of (CQ2) is zero. For any
fixed original compact test \(\eta\), (CQ1) gives
\[
        \int_E\langle\mathcal A_{R_a}(e),\eta\rangle\,d\nu(a)=0
                    \quad\text{for every Borel }E.
 \tag{CQ3}
\]
The integrand is integrable, so it is zero almost everywhere. Use
the countable cofinal original moment-and-cutoff stages from (DI1)
and countable dense real smooth coefficient arrays in each stage.
The union of their exceptional sets is still null. Outside it,
continuity of the conditional original current on each fixed stage
extends zero to all arrays, and exact cofinality includes every
actual refinement. This proves (CQ2). In particular the proof uses
all Borel bands and coefficient-current identities, not cancellation
of a single signed total integral.
\end{proof}

\subsection{Closed conormal currents need not admit closed division}

\begin{proposition}[Failure of closed division in every dimension]
\label{prop:r12-no-closed-conormal-division}
For every \(n\ge1\), there is an original open polyannulus \(U\),
an original unit parameter \(s\), and a nonzero positive separately
closed \((1,1)\)-current \(R\), with
\[
       d's\wedge R=d''s\wedge R=0,\qquad
       \operatorname{supp}R\subset\{s=a\},
 \tag{CQ4}
\]
for which no positive separately closed original \((0,0)\)-current
\(Q\) satisfies
\[
                           R=(d's\wedge d''s)\wedge Q.
 \tag{CQ5}
\]
The failure can be arranged as an ordinary-germ failure at a joint
Gauss point. This does not deny the frozen coarea identity expressing
the smooth wedge as an integral of hinge caps; it denies reversing
that operation into closed division of an individual cap.
\end{proposition}

\noindent\textbf{Status: \texttt{FALSE}} for the proposed universal
division assertion, with the following explicit all-dimensional proof.

\begin{proof}
Write \(U=A_I\times_K A_{J_2}\times_K\cdots\times_K A_{J_n}\)
inside the original split torus, with unit coordinates
\(T_1,\ldots,T_n\), and let \(s=\log|T_1|\), \(a\in I\).
Set
\[
                             R=D((s-a)_+[U]).
 \tag{CQ6}
\]
The frozen bounded calculus makes it positive and separately closed.
Its carrier and both conormal identities are the original hinge
locality and level-tangency identities.

Choose a nonnegative original compact top form \(\eta_Y\) on the
last \(n-1\) annuli with integral one. Explicitly it may be a
nonnegative smooth compact function of their logarithms times
their standard doubled top form, normalized by its real integral.
For \(n=1\), this means the scalar one on the zero-dimensional factor.
The top degree gives \(d'\eta_Y=d''\eta_Y=0\).
Suppose (CQ5) held. Define the positive Radon distribution on \(I\)
by
\[
 \langle m_Q,\chi\rangle
   =\langle(d's\wedge d''s)\wedge Q,\chi(s)\eta_Y\rangle
                      \quad(\chi\in C_c^\infty(I)).
 \tag{CQ7}
\]
These are genuine original compact tests: inverse images of compact
logarithmic intervals are compact annuli, and \(\eta_Y\) is compact.
Testing \(d'Q=0\) against
\(\chi(s)d''s\wedge\eta_Y\) gives
\(\langle m_Q,\chi'\rangle=0\). A one-dimensional distribution with
zero derivative on the connected interval \(I\) is a constant
multiple of Lebesgue measure. On the other hand the exact original
hinge calibration gives
\[
           \langle R,\chi(s)\eta_Y\rangle=\chi(a).
 \tag{CQ8}
\]
Thus (CQ5) would identify a constant Lebesgue measure with
\(\delta_a\), which is impossible. Product tests here are only a
necessary detector of an alleged original-current equality; they
are not claimed to span the full test sheaf.

For the germ assertion choose a vector
\(\mathbf a=(a_1,\ldots,a_n)\) whose coordinates are rationally
independent modulo \(\log|K^\times|\), and take \(a=a_1\).
The complete tropical fiber over \(\mathbf a\) is the single joint
Gauss point \(\sigma(\mathbf a)\): each nonzero Laurent polynomial
has a unique monomial of maximal norm at those radii, so every
multiplicative seminorm with those coordinate radii equals its
Gauss seminorm. Laurent polynomials determine the analytic point.
Tropicalization of an open polyannulus is proper, since compact
real boxes have compact affinoid inverse images. A proper map is
closed; hence every ordinary neighborhood of this singleton fiber
contains the inverse image of a smaller real neighborhood, and
in particular a smaller product box about \(\mathbf a\).
Any proposed local \(Q\) would restrict to such a box. Repeating
(CQ7)--(CQ8) there gives the same contradiction. No degree-zero
rigidity theorem on arbitrary ordinary domains is used.
\end{proof}

\subsection{A norm level is too large for the absolute-cap kernel}

\begin{proposition}[An overstrong fiber-transfer principle fails in every degree]
\label{prop:r12-norm-level-absolute-cap-counterexample}
For every \(b\ge1\) and \(n\ge b+1\), there exist an original
ordinary open \(U\) of dimension \(n\), a unit \(f\), a positive
finite PL cycle presentation \(e\) of degree \(b\), and positive
separately closed degree-one currents
\[
       C=D((\log|f|)_+[U]),\qquad Q,
\]
such that \(Q\) is conormal to \(\log|f|\) and carried by its
zero norm level, but
\[
                       \mathcal A_C(e)=0,\qquad
                       \mathcal A_Q(e)\ne0.
 \tag{CQ9}
\]
The presentation \(e\) is not actual-zero on \(U\).
Moreover this \(Q\) is not a relative hinge cap for \(f\):
it lies in an ordinary region where \(\log|f|\) is constant.
Thus (CQ9) disproves transfer to an \emph{arbitrary conormal
norm-level current}, not the restricted cap-derived gate in GQ.
\end{proposition}

\noindent\textbf{Status: \texttt{FALSE}} for that overstrong
absolute-to-relative norm-level principle. The actual relation
quotient itself retains status \texttt{OPEN}.

\begin{proof}
Choose an original uniformizer \(\varpi\) and put \(r=|\varpi|<1\).
Take a bounded annulus in the coordinate \(f\) containing the
unit-radius skeleton, remove the rigid point \(f=1\), and multiply
it by an annulus in \(y\) about radius \(r\) whose upper radius
is less than one. Add \(b-1\) annuli in coordinates
\(z_1,\ldots,z_{b-1}\), each about radius one, and \(n-b-1\)
additional annuli. All of \(f,f-1,y,z_j\) are units on this
ordinary open. Define
\[
 h=\max(\log|f-1|,\log|y|),\qquad
 k_j=(\log|z_j|)_+,\qquad
 e=D_\delta h\prod_{j=1}^{b-1}D_\delta k_j,\qquad
 Q=[\,f=1+\varpi\,].
 \tag{CQ10}
\]
The displayed potentials are continuous locally bounded actual
classical psh functions and form a finite convex-PL presentation.
The original analytic hypersurface current \(Q\) is positive and
separately closed. On its carrier, \(\log|f|=0\), giving the
conormal assertions without changing the ground field.

The support of \(C\) is contained in
\(\{|f|=1,\ |f-1|=1\}\). Indeed the hinge is locally either zero
or pluriharmonic away from \(|f|=1\); and on the ordinary open
\(|f-1|<1\), the ultrametric inequality gives \(|f|=1\)
identically, so its curvature is zero there as well.
On this support \(h=0\), because \(|y|<1\). Thus \(hC=0\)
as an actual coefficient current and \(D(hC)=0\). Mixed symmetry
proves the first equality in (CQ9), including every remaining
hinge slot and every original test.

On the original hypersurface \(f=1+\varpi\), the function \(h\)
is \(\max(\log r,\log|y|)\). Testing the defining integration
current and commuting restriction with \(D\) gives its nonzero
unit hinge at \(|y|=r\). The independent \(z_j\) hinges give
unit point factors at their unit radii. A compact original cutoff
equal to one near \(f=1+\varpi\), a nonnegative scalar test equal
to one near these hinge levels, and a positive compact top form
in the remaining coordinates give a strictly positive pairing.
This is the frozen original toric calibration, with coefficient
one in each independent factor. It proves the second assertion.

Finally \(e\) is independently seen to be nonzero as an actual
form. On an ordinary punctured annulus
\(0<|f-(1+\varpi)|<r\), one has \(|f-1|=r\), so its absolute
action is exactly the nonzero product of the \(y\)-hinge and
the \(z_j\)-hinges, times the untouched integration factors.
Compact positive top tests in the \(f-(1+\varpi)\) and remaining
directions give a positive absolute pairing. The actual
delta-current injection therefore excludes actual zero; no
assumption about GQ is used for this exclusion.

In contrast, a neighborhood \(|f-(1+\varpi)|<r\) has
\(\log|f|=0\) identically. Every relative \(f\)-hinge cap vanishes
there by the original bounded recursion. Hence this \(Q\) cannot
be substituted for a cap-derived current in the surviving gate.
\end{proof}

\paragraph{The surviving signed gate, with exact frozen provenance.}
The frozen R8 signed affine coarea theorem already supplies
continuous every-parameter original hinge families and the
smooth-wedge/cap equivalence. Frozen R9
\path{analytic_mechanism/gauss_carried_factorization.tex},
SHA256
{\small\nolinkurl{8d6cd0fe6120003a8d4b6072e980016fb6341d16a4f0ce1e1724d6383a8c343d}},
already supplies the full joint Gauss carrier, all-original-test
fiber factorization, induced trace products and cap Fubini.
These are \texttt{INTERNAL\_INPUT}, not discoveries of this section.
They are stronger than mere norm-level support, and (CQ9) must not
be used against them.

For an actual-zero original presentation \(e\), multiplication
by an actual positive cap remains actual-zero. Frozen absolute
compatibility consequently gives zero response on the corresponding
absolute cap integration current. What is still missing is its
zero response on the particular cap of an arbitrary original
\(S\), even though both retain all original tests and the correct
full Gauss carrier. Formula (CQ2) makes the signed cancellation
question exact but does not supply this implication. Closed
division is false by (CQ4)--(CQ8); enlarging the current class to
all norm-level currents loses the necessary carrier by (CQ9).
Neither manipulation resolves the genuine degree-two,
three-dimensional-image nonmatroidal case.

% END REVIEWED MODULE research_r12/gq/conditional_action_and_division_obstructions.tex


% BEGIN REVIEWED MODULE research_r12/gq/carrier_cln_relation_certificates.tex
\section{A background-independent closed certificate for relation descent}
\label{sec:r12-carrier-cln-certificates}

The preceding conditional construction does not itself prove the
analytic relation quotient. Here we give a quantitative sufficient
certificate which involves only actual functions on a closed carrier,
not a choice of positive background. Its completion is taken in a
potential norm controlled by relative CLN. The missing assertion is
whether actual analytic relations supply this certificate.

Fix an allowed original ordinary open \(U\), and let
\(\mathcal D(U)\) be the real vector space of finite continuous
locally bounded classical psh differences. Write
\[
                \mathcal P^b(U)=\operatorname{Sym}^b_{\mathbb R}
                                      \mathcal D(U),\qquad b\ge1.
 \tag{KC1}
\]
A monomial \(h_1\cdots h_b\) here is a \emph{formal curvature
polynomial}; its prescribed action on a positive closed \(Q\) is
\(L_{h_1}\cdots L_{h_b}Q\). It is not ordinary multiplication of
the resulting currents. The frozen signed bounded calculus makes
this action well-defined, symmetric and linear in each factor.
Every finite cycle presentation under consideration, with its
prescribed finite convex-PL representation, has such a polynomial
representative after expanding its specified convex differences.
We do not quotient (KC1) by the unknown actual analytic kernel.

Let \(F\subset U\) be closed. Call \(h\in\mathcal D(U)\)
\(F\)-harmonic if, near each point of \(F\), its values on \(F\)
agree with a finite continuous original function whose original
current curvature is zero. Define \(\mathcal Z_F^b\) to be the
linear span of the following explicitly permitted relations:
\begin{enumerate}
\item monomials having one \(F\)-harmonic factor;
\item monomials \(h_1\cdots h_b\) whose ordinary pointwise product
      is zero on \(F\);
\item finite identities already zero in a common ambient real
      PL cycle algebra, with their original moment pullbacks.
\end{enumerate}
The third item invokes only the frozen ambient functional calculus.
An identity only known after actual analytic pullback is not an
allowed generator. Local versions on ordinary neighborhoods of
the carrier are permitted when applying the resulting assertion
to a correspondingly localized test.

\begin{lemma}[The finite carrier certificate]
\label{lem:r12-finite-carrier-certificate}
If \(Q\) is any original middle-positive separately closed
background carried by \(F\), then
\[
                              \mathcal A_Q(Z)=0
                              \qquad(Z\in\mathcal Z_F^b)
 \tag{KC2}
\]
in every permitted degree, on all original tests and refinements.
The same is true after any additional finite signed mixed tail.
\end{lemma}

\noindent\textbf{Status: \texttt{DIRECT\_COROLLARY}} of the
ordinary harmonic quotient, (LT7), and the frozen ambient action.

\begin{proof}
Every mixed tail remains carried by \(F\), by ordinary support
locality. For an \(F\)-harmonic factor, replace it locally on \(F\)
by its original harmonic witness in multiplication with that tail.
Their coefficient products agree before applying \(D\), and the
frozen original harmonic potential quotient kills the insertion.
Finite signed positive-tail expansion treats signed tails. A
pointwise zero-product monomial is exactly (LT7), whose induction
uses the ordinary open where its first factor is nonzero.
Ambient relations act by zero by the printed frozen input.
Linearity and ordinary locality prove (KC2).
\end{proof}

\subsection{Relative CLN with bounds only on the actual carrier}

For a compact \(L\Subset U\), define
\[
 \|h\|_{F,L}=\sup_{F\cap L}|h|,\qquad
 \operatorname{osc}_{F,L}u
          =\inf_{c\in\mathbb R}\sup_{F\cap L}|u-c|.
 \tag{KC3}
\]
Both quantities are set to zero if \(F\cap L\) is empty.
The latter is a centered oscillation, not the uncentered
supremum minus infimum.

\begin{lemma}[A carrier-only mixed estimate]
\label{lem:r12-carrier-only-CLN}
Fix an original compact complementary test \(\eta\) for a
length-\(b\) insertion and the degree of the background.
There is a fixed larger compact stage \(L\Subset U\), finitely
many compact positive coefficient controls \(\Xi_\ell\) on that
stage, and \(C_\eta<\infty\), independent of \(F,Q\) and of the
potentials, such that
\[
 \left|\left\langle L_hL_{u_2}\cdots L_{u_b}Q,\eta\right\rangle\right|
 \le C_\eta\left(\sum_\ell\langle Q,\Xi_\ell\rangle\right)
       \|h\|_{F,L}\prod_{j=2}^b\operatorname{osc}_{F,L}u_j.
 \tag{KC4}
\]
Here \(Q\) is positive and separately closed, carried by \(F\);
\(h\in\mathcal D(U)\); and \(u_2,\ldots,u_b\) are continuous
locally bounded original classical psh on the same \(U\).
The empty product when \(b=1\) is one. Every additional actual
refinement of \(\eta\) is allowed.

For positive tuples the corresponding usual local coefficient
CLN estimate likewise uses only the centered bounds on \(F\cap L\).
For the signed first factor in (KC4), only the displayed test
pairing estimate is asserted, not coefficient total-variation
control by its small supremum.
\end{lemma}

\noindent\textbf{Status: \texttt{DIRECT\_COROLLARY}} of the
frozen finite-control relative CLN proof, with its carrier
restriction made explicit.

\begin{proof}
Choose the usual \(b\) nested compact stages between the support
of \(\eta\) and \(L\). Every positive mixed tail is carried by
\(F\): neither multiplication by a continuous scalar nor
distributional differentiation enlarges the support.
In the one-step control (LT2), multiplication by the potential
is consequently estimated by its supremum on \(F\cap L\), not
by values outside that carrier. The same fixed positive trace
controls and cutoff derivatives work; their construction does
not use the potentials. Iterating proves the carrier-only
positive CLN bound. Subtracting an arbitrary constant from any
psh slot preserves classicality and changes no insertion.
Taking the infimum over these finitely many constants gives the
centered bounds in (KC3).

For (KC4), put \(T=L_{u_2}\cdots L_{u_b}Q\). It is positive,
separately closed and carried by \(F\), and
\[
               \langle L_hT,\eta\rangle=\langle hT,D\eta\rangle.
\]
Finite original coefficient domination bounds this last pairing
by \(\|h\|_{F,L}\) times the compact coefficient controls of \(T\).
The just proved positive-tail estimate controls these by the
right side of (KC4). No estimate for the variation of \(D(hT)\)
has been inferred from a bounded scalar multiplier.
\end{proof}

\begin{definition}[The carrier-CLN relation seminorm]
\label{def:r12-carrier-cln-seminorm}
For \(P\in\mathcal P^b(U)\) and a compact \(L\Subset U\), put
\[
 \begin{aligned}
 \|P\|^{\rm cln}_{F,L}
  &=\inf
       \sum_{\nu=1}^N |c_\nu|\,\|h_\nu\|_{F,L}
                  \prod_{j=2}^b\operatorname{osc}_{F,L}u_{\nu j},\\
 P&=Z+\sum_{\nu=1}^N c_\nu h_\nu u_{\nu2}\cdots u_{\nu b},
             \qquad Z\in\mathcal Z_F^b .
 \end{aligned}
 \tag{KC5}
\]
The infimum ranges over all displayed finite decompositions.
Their identities are formal polynomial identities,
or identities in the explicitly permitted frozen ambient cycle
quotient. Each \(h_\nu\) lies in \(\mathcal D(U)\), and each
\(u_{\nu j}\) is continuous locally bounded classical psh on \(U\).
The infimum is finite: expand each difference in all but its
first slot. Finite concatenation and scaling show that (KC5)
is a seminorm on the displayed formal quotient. It depends on
the original functions and the closed set \(F\), but not on \(Q\).
It is not declared to be a seminorm on actual analytic delta-forms.
\end{definition}

\begin{theorem}[The closed carrier certificate annihilates every background]
\label{thm:r12-carrier-cln-zero-certificate}
For every \(Q\) in Lemma~\ref{lem:r12-carrier-only-CLN},
\[
          |\langle\mathcal A_Q(P),\eta\rangle|
          \le C_\eta\left(\sum_\ell\langle Q,\Xi_\ell\rangle\right)
                                      \|P\|^{\rm cln}_{F,L}.
 \tag{KC6}
\]
In particular, if the seminorm is zero on every required enlarged
compact stage, then \(\mathcal A_Q(P)=0\) on all original tests,
simultaneously for every allowed background carried by \(F\).
Finite certificates may be replaced by arbitrarily small
remainders in (KC5); no background regularization is needed.
\end{theorem}

\noindent\textbf{Status: \texttt{NEW\_PROVED}} for this
background-independent closed-certificate principle. It is a
sufficient relation test, not identification of the analytic kernel.

\begin{proof}
For any decomposition in (KC5), the \(Z\)-term vanishes by (KC2).
Apply (KC4) to each remaining term, sum, and take the infimum.
The controls and \(C_\eta\) are fixed before that decomposition,
so they do not grow with its number of terms or its real
subdivision complexity. This gives (KC6). Zero seminorm kills
each arbitrary original compact pairing. The stage may depend
on that test, as in ordinary current equality; there is no
restriction to a fixed finite moment test family.
\end{proof}

\subsection{A conditional all-unit route back to the pure-cycle quotient}

For a tuple of original units \(f=(f_1,\ldots,f_c)\) on an
ordinary open \(O\subset U\), put \(s_i=\log|f_i|\) and define
\[
 Q_{f,\mathbf a}(S)
       =T_{S,c}((s_1-a_1)_+,\ldots,(s_c-a_c)_+),\qquad
 F_{f,\mathbf a}
       =\operatorname{supp}Q_{f,\mathbf a}([O]).
 \tag{KC7}
\]
The second set is the \emph{absolute cap carrier}, not the whole
norm level. Frozen R7 positive-zero descent gives
\[
                    \operatorname{supp}Q_{f,\mathbf a}(S)
                                            \subset F_{f,\mathbf a}.
 \tag{KC8}
\]
Indeed on the ordinary complement the positive point cycle has
zero absolute actual pullback, so its prescribed action on every
allowed \(S\) vanishes. For smooth original parameter maps,
frozen R9 identifies this carrier with the complete joint Gauss
fiber and supplies its all-original-test factorization. We use
only (KC8); no new-field quotient is assumed.

\begin{proposition}[A uniform carrier-certificate criterion]
\label{prop:r12-carrier-certificate-pure-cycle-quotient}
Let \(P\) represent a finite pure degree-\(b\) PL cycle expression
on \(U\), let \(S\) have arbitrary allowed degree \(p\), and put
\(P_0=p+b<n\), \(q=n-P_0\). Choose an integer
\[
       1\le c\le
       \max\{\,q\bmod2,\ q-2\lfloor p/2\rfloor\,\}.
 \tag{KC9}
\]
Suppose that for every ordinary localization and every tuple of
\(c\) original units, the following background-independent
condition holds: for Lebesgue-almost every parameter \(\mathbf a\),
\[
                      \|P\|^{\rm cln}_{F_{f,\mathbf a},L}=0
 \tag{KC10}
\]
on every compact stage required by (KC6). A countable exhaustion
on a local affinoid neighborhood may be used for these stages.
Then \(\mathcal A_S(P)=0\) for every such arbitrary \(S\), on
all original tests and all actual refinements.

This proposition does not assert (KC10) for every actual-zero
presentation. It concerns pure cycle expressions; variable PS
coefficient relations are not silently reduced to this statement.
If \(q=0\), or the right side of (KC9) is zero, it supplies no
positive-rank cap criterion.
\end{proposition}

\noindent\textbf{Status: \texttt{CONDITIONAL}} for the analytic
quotient implication, with the stated reduction completely proved.

\begin{proof}
For fixed \(f,\mathbf a\), (KC8) makes its relative cap an allowed
positive closed background carried by \(F_{f,\mathbf a}\).
The closed-certificate theorem gives
\(\mathcal A_{Q_{f,\mathbf a}(S)}(P)=0\) for almost every parameter.
For each original compact test, these particular pairings are
continuous in \(\mathbf a\): expand the fixed \(P\) into finitely
many bounded convex-PL tuples, adjoin the Lipschitz hinge tuple,
and apply the frozen joint local-uniform bounded BT theorem.
Consequently they vanish at every parameter. This uses the
canonical cap-derived family, not continuity postulated for an
arbitrary Radon disintegration.

The frozen signed coarea identity now gives
\[
       \bigwedge_{i=1}^c(d's_i\wedge d''s_i)
                         \wedge\mathcal A_S(P)=0
                         \quad\text{for every such original tuple}.
 \tag{KC11}
\]
Its commutation with the pure response follows from signed mixed
symmetry; no coefficient derivative is suppressed. The response
is symmetric and has regularizable curvature depth \(b\), because
its original finite PL representatives have \(b\) such slots.
The sharp depth-sensitive all-unit detector (HD7)--(HD8), with
total degree \(P_0=p+b\) and retained depth \(b\), has threshold
exactly (KC9). It therefore gives zero on the entire original
test sheaf. Neither fixed-moment caps alone nor tests only on a
chosen Gauss section have been used.
\end{proof}

\paragraph{The attempted degree-two proof and its exact stopping point.}
For two nonzero positive degree-two cycles with the same actual
pullback on a full ordinary three-dimensional-image germ, write
their difference using a fixed finite polynomial \(P\).
Actual zero implies zero after multiplication by every actual cap,
and hence zero on each corresponding absolute integration cap.
This is valid also for a genuinely nonmatroidal,
non-simplicial-polytopal germ; it is not the missing step.

The attempted next implication was
\[
 \mathcal A_{Q_{f,\mathbf a}([O])}(P)=0
       \quad\Longrightarrow\quad
                       \|P\|^{\rm cln}_{F_{f,\mathbf a},L}=0.
 \tag{KC12}
\]
It has not been proved. The estimate (KC6) goes in the opposite
direction and cannot be reversed. A vanishing signed cap cycle
does not by itself provide a finite pointwise zero-product
decomposition, harmonic ordinary traces, or approximations of
these with a small carrier remainder. Disintegration only
represents the already defined response measures; it does not
construct such potential identities. Replacing the correct
carrier by a norm level invokes the false assertion (CQ9),
while reversing coarea into closed division invokes (CQ5).

Thus (KC12), restricted to presentations already actual-zero on
the original ordinary germ and with the full actual refinements
retained, is a precise additional analytic-to-carrier closure
gate. Its status is \texttt{OPEN}; the proposed automatic use
of (KC6) to prove it is a \texttt{FAILED\_ATTEMPT}.
No new nonmatroidal degree-two quotient, no pair of actual-equal
positive cycles with different arbitrary-\(S\) actions, and no
low-degree signed order transfer has been established here.
The exact frozen \(\Gamma_4\) quotient is not counted anew.

\paragraph{Dependencies and non-circularity.}
The proofs use only frozen original bounded BT, its finite-control
relative CLN estimate, its continuous harmonic quotient and
ambient PL action; the carrier zero-product lemma (LT7);
the positive-zero cap support input of R7; frozen R8/R9 coarea
and parameter continuity; and the independently proved R12
depth detector. They use neither OPEN-GQ nor OPEN-CTO, no
\(\mathcal W_{\rm loc}\) density, and no arbitrary-background
BT theory over a completed residue field. The new seminorm is
defined before any actual analytic relation quotient.

% END REVIEWED MODULE research_r12/gq/carrier_cln_relation_certificates.tex


% BEGIN REVIEWED MODULE research_r12/remaining_gates_and_status.tex
\section{What has advanced, what has failed, and what remains}
\label{sec:r12-final-status}

This record is a proof continuation, not a new audit of the
authorized papers R, B and A2. Statements labelled new are new
relative to the exact frozen R11 chain. The original R8 product-fan
audit and the later nonmatroidal fan calculations remain where
they were frozen in that chain. No new unrestricted fan quotient
is inferred from the analytic constructions here.

\subsection{The analytic quotient remains a separate gate}

For a compatible finite actual PS/PL presentation $e$, the
unrestricted target is still
\[
 e=0\ \hbox{as an actual analytic delta-form}
       \quad\Longrightarrow\quad \mathcal A_S(e)=0
 \tag{GQ}
\]
on every ordinary germ, for every original allowed background $S$,
every compact original complementary smooth test, and every actual
moment refinement. This implication is \texttt{OPEN}. Its
low-real-degree signed constant-tensor order-transfer counterpart
is also \texttt{OPEN}.

The new conditional-current construction solves a measure and
test-coherence problem: its $R_a$ are original ambient currents,
not formal objects on a residue space. The ordinary level-trace
calculus gives a background-independent sufficient lifting
criterion. The missing step is still to extract such an ordinary
carrier/harmonic-trace certificate from an arbitrary actual
analytic relation. The depth-$k$ conormal reconstruction theorem
lowers or identifies a sufficient cap threshold only after the
required cap vanishing has been proved. It does not prove that
premise from (GQ).

The final direct-at-gate attempt defines a background-independent
carrier-CLN seminorm on formal curvature polynomials before taking
the unknown analytic quotient. Its zero class annihilates every
allowed background carried by the specified closed set, with an
all-original-test estimate. Applied to the correct absolute cap
carriers and all unit tuples, it gives a fully proved conditional
pure-cycle criterion. The unproved analytic step is (KC12):
absolute cap vanishing for an already actual-zero presentation
must supply zero carrier seminorm. The proved estimate goes in
the opposite direction and cannot be reversed.

Conditional response disintegration handles signed cancellation
by all Borel bands, but does not establish (KC12). Closed division
of an individual conormal cap is false even as an ordinary germ.
Transfer from an absolute cap to an arbitrary conormal norm-level
current is also false; the explicit latter current is not
cap-derived. The frozen R8/R9 canonical coarea and full Gauss
carrier theorems remain valid and are not recounted as new.

Two exact counterexamples delimit the route. A reference-space
pushforward loses mixed original unit tests even for explicit
positive closed graph caps. Vanishing of all caps in one finite
moment tuple does not imply vanishing after all additional original
unit moments. Both proposed inferences are \texttt{FALSE}; neither
is a counterexample to (GQ). The particular $b=2<r=3$ instance of
the latter calculation is not a genuinely nonmatroidal pressure
test, and is not presented as one.

The \texttt{OPEN} $W_{\rm loc}$ density and Gauss-residue
auxiliary-field quotient questions are unchanged. No argument
here uses them, and no ordinary-current equality is strengthened
to all GK delta-test equality.

\subsection{Capacity, exceptional sets, and envelopes}

The full normalized original tester family is retained wherever
the notation $C_R$ or $C_S$ is used. The distinct band-pasting,
completed-seed and energy-normalized capacities always name their
own admissible families explicitly. The new closed-band theorem proves raw/outer
equality on every Borel parameter-saturated set in an ordinary
relatively compact window, for a parameter-annihilated background.
For every arbitrary $S$, it gives the stated mixed-slot corollary
with the derived current $\Omega_s\wedge S$. It does not replace
$S$ itself or prove unrestricted raw/outer equality.

The continuous-branch band seed now has a bounded capacitary
completion with intrinsic mixed products, relative outer
quasicontinuity, joint decreasing continuity and the stated
maximum carriers. Its seed-capacity completion preserves open
values and hence the outer capacity. The version completed in
the larger full band capacity has only the printed one-sided
comparison; cofinality and inclusion of every original tester
are not asserted. The empty-parameter case recovers a frozen
R11 completion and is not recounted as a new theorem.

The separate parameter finite-energy module has an exact
$L^{k+1}$ model and an attained energy-normalized capacity.
That capacity is raw/ordinary-outer regular on all local Borel
sets, but its null ideal is exactly the one fixed source measure's
null ideal. These facts do not settle unrestricted original
capacity, quasicontinuity or classical polar-set equivalences.
In fact Section~\ref{sec:r12-countable-measure-capacity-obstruction}
proves a strict obstruction in every $n\geq q\geq1$: for one
allowed integration background, every countable family of
positive Radon probes misses some singleton of uniformly positive
full original capacity. Thus even a countable collection of the
energy-normalized probe capacities cannot in general recover the
full original null ideal. This does not disprove Choquet regularity
of the original capacity.

The general compact no-gain statement remains \texttt{OPEN}:
a weak curvature cluster represented by a completed normalized
tester is not automatically the curvature of an actual normalized
classical tester with no gain on a specified compact set.
The Jensen-current formula and the single-probe criterion identify
one obstruction exactly, but the probe deficit can be nonzero.
Even zero deficit does not justify restricting weak convergence
to a fixed compact set.

For the full unrestricted tester family, every compactly localized
weak cluster now belongs unconditionally to a literal strong-dual
datum class. The larger class with a common linear outer-capacity
bound is weakly compact and uniformly continuous on bounded-energy
strong rows, with an explicit $\varepsilon\log(1/\varepsilon)$
modulus. None of this produces an actual normalized tester.
A constant tester row with a cutoff can already have positive
global source normal, despite satisfying raw-capacity no-gain on
every Borel set. Global source compatibility is therefore not a
necessary local no-gain condition.

Nor can the argument be completed by asserting weak compactness
of potential energy balls in either new Banach norm. The pole
row has no weakly convergent subnet there, and the two Banach
spaces are nonreflexive in that explicit allowed all-degree
example. The strong cone itself remains complete.

There is a positive stability replacement for that false shortcut.
For weak-$*$ convergent strong-dual data, the energies of the
canonical variational potentials are lower semicontinuous;
convergence of this energy to that of the limit datum's canonical
potential is equivalent to strong potential convergence.
The exact identity (VE7) proves the
quantitative bound $I(u_j,u)\leq d\delta_j$.
Together with the compact capacity-controlled datum theorem,
this preserves completed source solutions under weak measure
convergence and the stated scalar energy no-loss hypothesis.
That hypothesis is not automatic, and the completed conclusion
is still not the normalized actual tester required for (FG1).

Relative outer quasicontinuity and relative exceptional-set
statements already proved in R11 are \texttt{INTERNAL\_INPUT}.
The new order-energy lattice inherits relative outer
quasicontinuity. The complete general equivalence between all
classical pluripolar, negligible and capacity-null notions is
not asserted. In particular the frozen distinction between
relative and ambient polar ideals is preserved; equating them
is not a harmless change of representative.

The new obstacle map and its convex energy functional are on the
unchanged completed relative cone. Their envelopes are actual
Borel classes, not automatically ambient classical psh envelopes.
The contact functional is nonnegative, and its vanishing is a
real additional condition, not a consequence of formal variation.

\subsection{Comparison, domination, and the correct source problem}

The frozen bounded and specified no-loss/finite-energy relative
comparison and domination results retain every printed domain,
boundary and null-quotient hypothesis. They are not recounted as
R12 discoveries. Unrestricted unbounded comparison without loss
control already has frozen counterexamples.

The source map $J$ has a closed image in the signed energy dual,
and the new variational retraction is defined for the printed
strong-continuous dual. It gives a maximizer with an explicit
normal residual. A general source equation requires that residual
to vanish on all original tests. With literal positive data and
a nonempty subsolution family, this is equivalent to zero
Perron--variational energy gap.

The new single-probe theorem removes the subsolution hypothesis
from the exact source test: for literal strong-continuous data,
one nonnegative number $\Xi_{S,\psi_0}(\mu)$ vanishes if and
only if the full original source equation is solved by the
variational maximizer. The probe is geometry-only, and its
deficit quantitatively controls every compact original smooth
test pairing. Proving automatic zero for an appropriate new
data class remains an additional problem.

The assertion that every positive dominated datum has a classical
or completed source is \texttt{FALSE} in the unrestricted
background class. The new positive closed moment lift extends
the compact-carrier obstruction to every nonlinear sub-top
degree. Its special factorization hypothesis is part of the
obstruction construction; it is not imposed on the arbitrary
background of the positive results. An appropriate data class,
a verified compatibility condition and a boundary/nondegeneracy
hypothesis must still be specified before claiming a general
Monge--Ampere Dirichlet existence theorem.

The radial flux theorem gives an explicit compatible family
and actual source inverse in every remaining degree. Its
smooth-density threshold proves a sharper negative statement:
source solvability is not open under arbitrarily small strictly
positive original smooth density perturbations, even with the
same Borel null ideal and domination by a bounded actual source.
The opposite sign of the same perturbation is exactly solvable.
This is a full-cone obstruction, not just failure of radial
competitors.

The exact order-dual topology counterexample shows why positivity
alone cannot remove the strong-continuity hypothesis. It retains
all finite power energies and the original base field, yet
strong convergence fails in the order norm and for a norm-one
order-reversing functional. A positive compact-test shadow also
does not by itself mean literal Radon integration on the whole
relative Borel energy space. These are different obstructions:
the shadow can charge a relative null set, or an energy dual can
have zero shadow while remaining nonzero.

The narrower assertion that every \emph{literal} positive Radon
energy-bounded functional is strongly continuous remains
\texttt{OPEN}. Section~\ref{sec:r12-literal-dual-rooftop-gate}
reduces it to a precise ordered two-entry lower-rooftop
continuity statement, proves the common-lower-bound consequence,
and verifies that mechanism on the all-degree radial energy
family. Neither the nonliteral pole functional nor the smooth
density source obstruction refutes this narrower gate.

Fixed Radon domination does repair literalness: every measure
order interval $0\leq\mu\leq M\widehat\mu_v$ is weak-$*$ compact
and norm closed inside the literal strong dual. Mere finiteness
of the integrals of every cone member also forces a bounded
literal energy functional by a summable-cone-series argument,
but that last assertion alone does not prove strong continuity.

\subsection{Dependency and non-circularity check}

The closed-band proof starts with coefficient measures and the
smooth Leibniz rule, then proves Borel-band closedness before
using bounded BT on those bands. Disintegration is constructed
from countably cofinal original test stages and the Radon
coefficient control. Positivity and closedness of the conditional
currents are proved before inserting any conditional potentials.

The two energy spaces are constructed from the already frozen
strong potential cone. Their completeness never serves to
define that cone. The order norm uses the frozen relative
Perron theorem with an explicitly nonempty minorant family.
Literal current multipliers are justified by a finite-energy
trace-completing potential, not by differentiating arbitrary
Borel functions as forms. The variational problem uses the
specified strong-dual continuity, which the counterexamples
show cannot be dropped without another argument.

The positive Jensen current is first an explicit bounded
formula. Its extension uses the frozen weighted strong
continuity, and its weak compactness uses the original
coefficient controls. The universal probe is constructed
independently of the background before it is used to detect
defects; no product with a nonsmooth positive nonclosed
defect current is introduced.

Finally, every negative example is realized over the original
finite field and norm and separated on actual original smooth
tests, unless the example explicitly concerns an energy dual
rather than a current. No numerical calculation is used as a
substitute for the displayed symbolic or distributional proof.

\paragraph{Minimum next lemmas.}
The following are concrete targets rather than completed
consequences:
\begin{enumerate}
\item An ordinary ambient certificate for arbitrary actual
analytic PS/PL relations after conditional-current
disintegration, with all original mixed unit tests retained.
\item A compact no-gain theorem for the full original normalized
tester class, or an explicit counterexample to its proposed
strongest formulation.
\item A usable intrinsic characterization of literal
strong-continuous Radon data which enforces zero source normal,
equivalently zero Perron--variational gap when a subsolution exists.
\item A nonlinear enlarged relative potential class with proved
mixed positivity and closure, if one is intended to solve data
excluded by the ambient classical cone. Linear relative
subharmonicity and a general scalar hyperbolicity prescription
have been disproved as automatic substitutes.
\end{enumerate}
These gates remain logically distinct. The final general theory
has advanced through new exact mechanisms and excluded shortcuts;
it has not become an unconditional complete classical
Bedford--Taylor theory on every allowed analytic domain.


% END REVIEWED MODULE research_r12/remaining_gates_and_status.tex

\end{document}

